\documentclass[11pt,a4paper]{article}

\usepackage[utf8]{inputenc}
\usepackage[T1]{fontenc}
\usepackage{textcomp}
\usepackage{lmodern}

\usepackage[a4paper,margin=1in]{geometry}
\usepackage{amsmath,amssymb,amsthm}

\usepackage{xcolor}
\usepackage{hyperref}
\hypersetup{
  colorlinks=true,
  linkcolor=blue!60!black,
  citecolor=blue!60!black,
  urlcolor=blue!60!black,
  breaklinks=true,
}

\usepackage{tabularx}
\usepackage{array}
\newcolumntype{Y}{>{\raggedright\arraybackslash}X}
\usepackage{booktabs}

\usepackage{enumitem}
\setlist[itemize]{leftmargin=*,topsep=2pt,partopsep=0pt,parsep=2pt,itemsep=2pt}
\setlist[enumerate]{leftmargin=*,topsep=2pt,partopsep=0pt,parsep=2pt,itemsep=2pt}

\usepackage{etoolbox}
\setlength{\emergencystretch}{3em}
\tolerance=1000

\usepackage{mdframed}
\definecolor{quotebg}{gray}{0.965}
\definecolor{quoterule}{gray}{0.55}
\mdfdefinestyle{quotestyle}{%
  backgroundcolor=quotebg,
  linecolor=quoterule,
  linewidth=0pt,
  leftline=true,
  innerleftmargin=10pt,
  innerrightmargin=6pt,
  innertopmargin=4pt,
  innerbottommargin=4pt,
  skipabove=4pt,
  skipbelow=4pt
}
\newenvironment{qiqtquote}{\begin{mdframed}[style=quotestyle]}{\end{mdframed}}

\definecolor{codebg}{gray}{0.96}
\mdfdefinestyle{codestyle}{%
  backgroundcolor=codebg,
  linecolor=quoterule,
  linewidth=0pt,
  leftline=true,
  innerleftmargin=8pt,
  innerrightmargin=4pt,
  innertopmargin=4pt,
  innerbottommargin=4pt,
  skipabove=4pt,
  skipbelow=4pt
}
\newenvironment{qiqtcode}{\begingroup\footnotesize\ttfamily\begin{mdframed}[style=codestyle]}{\end{mdframed}\endgroup}

\usepackage{newunicodechar}
\newunicodechar{−}{$-$}
\newunicodechar{×}{$\times$}
\newunicodechar{÷}{$\div$}
\newunicodechar{±}{$\pm$}
\newunicodechar{≈}{$\approx$}
\newunicodechar{≠}{$\neq$}
\newunicodechar{≤}{$\leq$}
\newunicodechar{≥}{$\geq$}
\newunicodechar{∈}{$\in$}
\newunicodechar{∉}{$\notin$}
\newunicodechar{∞}{$\infty$}
\newunicodechar{→}{$\to$}
\newunicodechar{↦}{$\mapsto$}
\newunicodechar{⊗}{$\otimes$}
\newunicodechar{⊕}{$\oplus$}
\newunicodechar{⊂}{$\subset$}
\newunicodechar{⊆}{$\subseteq$}
\newunicodechar{∩}{$\cap$}
\newunicodechar{∪}{$\cup$}
\newunicodechar{∂}{$\partial$}
\newunicodechar{∇}{$\nabla$}
\newunicodechar{∑}{$\sum$}
\newunicodechar{∫}{$\int$}
\newunicodechar{√}{$\sqrt{}$}
\newunicodechar{α}{$\alpha$}
\newunicodechar{β}{\ensuremath{\beta}}
\newunicodechar{γ}{$\gamma$}
\newunicodechar{δ}{$\delta$}
\newunicodechar{ε}{$\varepsilon$}
\newunicodechar{ζ}{$\zeta$}
\newunicodechar{η}{$\eta$}
\newunicodechar{θ}{$\theta$}
\newunicodechar{λ}{$\lambda$}
\newunicodechar{μ}{$\mu$}
\newunicodechar{ν}{$\nu$}
\newunicodechar{ξ}{$\xi$}
\newunicodechar{π}{$\pi$}
\newunicodechar{ρ}{$\rho$}
\newunicodechar{σ}{$\sigma$}
\newunicodechar{τ}{$\tau$}
\newunicodechar{φ}{$\varphi$}
\newunicodechar{χ}{$\chi$}
\newunicodechar{ψ}{$\psi$}
\newunicodechar{ω}{\ensuremath{\omega}}
\newunicodechar{Γ}{$\Gamma$}
\newunicodechar{Δ}{$\Delta$}
\newunicodechar{Θ}{$\Theta$}
\newunicodechar{Λ}{$\Lambda$}
\newunicodechar{Π}{$\Pi$}
\newunicodechar{Σ}{$\Sigma$}
\newunicodechar{Φ}{$\Phi$}
\newunicodechar{Ψ}{$\Psi$}
\newunicodechar{Ω}{$\Omega$}
\newunicodechar{ℏ}{$\hbar$}
\newunicodechar{ℓ}{$\ell$}
\newunicodechar{ℝ}{$\mathbb{R}$}
\newunicodechar{ℂ}{$\mathbb{C}$}
\newunicodechar{ℕ}{$\mathbb{N}$}
\newunicodechar{ℤ}{$\mathbb{Z}$}
\newunicodechar{ℍ}{$\mathbb{H}$}
\newunicodechar{ℰ}{$\mathcal{E}$}
\newunicodechar{✓}{\checkmark}
\newunicodechar{✗}{$\times$}
\newunicodechar{…}{\ldots}
\newunicodechar{—}{---}
\newunicodechar{–}{--}
\newunicodechar{‘}{`}
\newunicodechar{’}{'}
\newunicodechar{“}{``}
\newunicodechar{”}{''}
\newunicodechar{·}{\ensuremath{\cdot}}
\newunicodechar{∼}{$\sim$}
\newunicodechar{‐}{-}
\newunicodechar{‑}{-}

% --- QIQTH extended symbol coverage (arrows/relations/logic/brackets) ---
\newunicodechar{↔}{$\leftrightarrow$}
\newunicodechar{⟷}{$\longleftrightarrow$}
\newunicodechar{⟶}{$\longrightarrow$}
\newunicodechar{⟵}{$\longleftarrow$}
\newunicodechar{⟹}{$\Longrightarrow$}
\newunicodechar{⟸}{$\Longleftarrow$}
\newunicodechar{⟺}{$\Longleftrightarrow$}
\newunicodechar{←}{$\leftarrow$}
\newunicodechar{⇒}{$\Rightarrow$}
\newunicodechar{⇐}{$\Leftarrow$}
\newunicodechar{⇔}{$\Leftrightarrow$}
\newunicodechar{⇏}{$\nRightarrow$}
\newunicodechar{⇎}{$\nLeftrightarrow$}
\newunicodechar{⇌}{$\rightleftharpoons$}
\newunicodechar{↑}{$\uparrow$}
\newunicodechar{↓}{$\downarrow$}
\newunicodechar{↕}{$\updownarrow$}
\newunicodechar{↪}{$\hookrightarrow$}
\newunicodechar{≡}{$\equiv$}
\newunicodechar{≢}{$\not\equiv$}
\newunicodechar{≅}{$\cong$}
\newunicodechar{≃}{$\simeq$}
\newunicodechar{≲}{$\lesssim$}
\newunicodechar{≳}{$\gtrsim$}
\newunicodechar{≪}{$\ll$}
\newunicodechar{≫}{$\gg$}
\newunicodechar{∝}{$\propto$}
\newunicodechar{≔}{$:=$}
\newunicodechar{∀}{$\forall$}
\newunicodechar{∃}{$\exists$}
\newunicodechar{∄}{$\nexists$}
\newunicodechar{∅}{$\emptyset$}
\newunicodechar{¬}{$\neg$}
\newunicodechar{∧}{$\wedge$}
\newunicodechar{∨}{$\vee$}
\newunicodechar{⊙}{$\odot$}
\newunicodechar{⊥}{$\perp$}
\newunicodechar{∥}{$\parallel$}
\newunicodechar{‖}{$\Vert$}
\newunicodechar{∘}{$\circ$}
\newunicodechar{∙}{$\bullet$}
\newunicodechar{•}{$\bullet$}
\newunicodechar{∗}{$\ast$}
\newunicodechar{⊇}{$\supseteq$}
\newunicodechar{⊃}{$\supset$}
\newunicodechar{∖}{$\setminus$}
\newunicodechar{⟨}{$\langle$}
\newunicodechar{⟩}{$\rangle$}
\newunicodechar{†}{$\dagger$}
\newunicodechar{‡}{$\ddagger$}
\newunicodechar{′}{$'$}
\newunicodechar{″}{$''$}
\newunicodechar{⟂}{$\perp$}
\newunicodechar{⊤}{$\top$}
\newunicodechar{⊨}{$\models$}
\newunicodechar{⊢}{$\vdash$}
\newunicodechar{∎}{$\blacksquare$}
\newunicodechar{□}{$\square$}
\newunicodechar{△}{$\triangle$}
\newunicodechar{Ξ}{$\Xi$}
\newunicodechar{Υ}{$\Upsilon$}
\newunicodechar{ϕ}{$\phi$}
\newunicodechar{ϵ}{$\epsilon$}
\newunicodechar{ϑ}{$\vartheta$}
\newunicodechar{∓}{$\mp$}
\newunicodechar{°}{\textdegree}
\newunicodechar{µ}{$\mu$}
\newunicodechar{²}{\textsuperscript{2}}
\newunicodechar{³}{\textsuperscript{3}}
\newunicodechar{¹}{\textsuperscript{1}}
\newunicodechar{⁰}{\textsuperscript{0}}
\newunicodechar{₀}{\textsubscript{0}}
\newunicodechar{₁}{\textsubscript{1}}
\newunicodechar{₂}{\textsubscript{2}}
\newunicodechar{₃}{\textsubscript{3}}
\newunicodechar{ₙ}{\textsubscript{n}}
\newunicodechar{⋆}{$\star$}
\newunicodechar{⟪}{$\langle\!\langle$}
\newunicodechar{⟫}{$\rangle\!\rangle$}
\newunicodechar{κ}{\ensuremath{\kappa}}
\newunicodechar{ᶜ}{\textsuperscript{c}}
\newunicodechar{ᶠ}{\textsuperscript{f}}

\title{One Wave Function, One World: $\Phi$-Monism, a Holographic Entropy Bound, and a Machine-Verified Path Toward Emergent Gravity}
\author{Pawel Kaplanski}
\date{2026-05-25}

\begin{document}
\maketitle

\begin{abstract}
\noindent We develop \textbf{QIQT-H}, a single-world, holographic formulation of quantum mechanics whose ontology is \textbf{$\Phi$-monism}: there is one substance — the universal wave function $\Phi$, evolving exactly unitarily — of which observers are macroscopic patterns, and a non-dynamical selector $\lambda$ marks exactly one decoherent record actual per run. No collapse term, no branching, no fundamental probability. The framework rests on five postulates — \textbf{(P1)} the $(\Phi,\lambda)$ ontology, \textbf{(P2)} quantum kinematics, \textbf{(P3)} microcausality, \textbf{(P4)} finite holographic capacity, and \textbf{(P5)} quantum equilibrium of the typicality measure — of which (P2)–(P3) are the standard quantum-relativistic arena, so the genuinely irreducible new physics is \textbf{(P4) + (P5), on the (P1) ontology}. The mathematical home of (P4) is the crossed-product / Type II construction of Chandrasekaran-Penington-Witten (2022) and Witten (2022). In ordinary QFT the local algebras of bounded regions are Type III$_1$, admitting no trace, no factorization, and UV-divergent entanglement entropy; gravitational dressing by the modular flow of a reference state produces a Type II algebra $\hat{\mathcal{A}}(R)$ with a semifinite trace and a well-defined renormalized entropy whose differences match the generalized entropy $A/(4\ell_P^2) + S_{\rm matter}$. In this setting (P4) is a \textbf{finite-entropy} axiom: the renormalized (von Neumann) entropy of the state $\omega_\Psi$ induced on $\hat{\mathcal{A}}(R)$ is bounded by $Q_R := A(\partial R)/(4\ell_P^2)$. We stress the reading that survives adversarial audit: \emph{$Q_R$ bounds an entanglement entropy $S(\rho_R)$ of the pre-selection $\Phi$ — its effective modular rank $e^{S(\rho_R)}$ — \textbf{not} a finite matter Hilbert-space dimension, and \textbf{not} a finite specification precision of amplitudes.} Because it is an entanglement entropy of $\Phi$, the inertness of $\lambda$ is \emph{required}, not merely interpretive: conditioning on the actual branch would destroy the coherence the entropy measures, so \textbf{geometry responds to $\Phi$, never to the actualized branch}. Algebra-state equivalence supplies the accompanying basis-independent notion of regional physical indistinguishability — two abstract Hilbert vectors are the same regional state iff they agree on every observable in $\hat{\mathcal{A}}(R)$; the crossed-product construction is borrowed from CPW/Witten, the bound $S_{\rm ren} \le Q_R$ is our postulate, and the foundational application to the measurement problem is our contribution. The machine-checked \textbf{P4-MICRO} refinement (§1.1a) sharpens what is irreducibly postulated: only that the regional capacity is \emph{finite} ($N_R<\infty$); the area \emph{floor} $S_{\rm vN}(\rho_R)\le Q_R$ is then a \emph{derived theorem} (\texttt{area\_floor\_vonNeumann}), the area \emph{form} $\propto A$ comes from the Sakharov bridge (a \emph{re-derivation} of standard induced gravity), and the relation $G = 1/(N\Lambda_s^2)$ is \emph{derived} (the dimensionless $G/a_0^2 = 1/N$ a theorem) — only $G$'s \emph{numerical value} is carried, blocked at the $(1/6-\xi)$ Seeley–DeWitt species coefficient (a Riemannian heat-kernel gap Mathlib lacks). We distinguish the \textbf{formal wave function} of standard QM (an idealized ensemble description) from the \textbf{per-run wave function} (whose regional physical content is given by its values on the regional Type II algebras). The measurement problem dissolves without collapse: decoherence renders the records non-interfering and the selector $\lambda$ (P1) makes exactly one actual, with the global evolution left exactly unitary and no amplitude trimmed. We \textbf{retire}, as a category error, the earlier \emph{Macroscopic Definiteness Conjecture} that the finite capacity \emph{itself} forbids a two-record regional state ($I_\Sigma \approx K\cdot I_0 > Q_R$): a second macroscopic branch adds only $\sim\!\log K$ to $S(\rho_R)$ — far below $Q_R$ — so a two-record state does \emph{not} overflow capacity, and no kinematic entropy bound can \emph{select} an outcome in a unitarily-evolving theory. Single outcomes are the work of $\lambda$ + decoherence, \textbf{never} of the capacity bound; the earlier conditional-on-Macroscopic-Definiteness language in §4, §7.3, and §7.6 is superseded accordingly. The Born rule is likewise not a primitive: it is \textbf{reduced} — provably underivable from unitarity alone, by a battery of machine-checked no-go theorems — to (P5) alone (both earlier Born premises collapse to refinement-equivariance). No Bohmian particles, no MWI branches, no modal value-rule are added. Gravity emerges holographically: for the free field the entanglement first law at every probe is \emph{equivalent} to the linearized Einstein equations, and the Sakharov $1/4$ is a theorem. We claim a fully machine-verified derivation chain from finite information \emph{toward} quantum gravity and single-outcome quantum mechanics — \textbf{not} a completed theory: every derivation is conditional on named, shrinking inputs (the Clausius/area law where not yet discharged, the matching of the trace-defined area to external geometry, the value of $G$, the continuum Type III$_1$ limit, interacting matter), every gap named, checkable, and independently re-runnable (\texttt{bash verify/verify.sh} emits a claim card). We state two propositions and two explicit conjectures, identify four explicit open problems, and frame the construction as a research program rather than a completed interpretation.


\medskip
\noindent\textbf{Keywords:} foundations of quantum mechanics; Bekenstein-Bousso bound; holographic entropy bound; Type II von Neumann algebras; crossed product; generalized entropy; $\Phi$-monism; single-world interpretation; typicality; induced gravity.

\textbf{Formal verification.} The deductive core of this framework is machine-checked in Lean 4 / Mathlib and is now \textbf{axiom-free}: every audited theorem depends only on the three standard Lean axioms (\texttt{propext}, \texttt{Classical.choice}, \texttt{Quot.sound}), with \textbf{no \texttt{sorry}} and \textbf{no project-specific axioms} (CI-guarded by \texttt{scripts/axiom\_budget\_check.sh} at budget 0; the per-theorem \texttt{\#print axioms} audit in \texttt{lean/mathlib/QIQTH/AxiomAudit.lean} carries roughly 2,800 directives, within a full library of over 420 files and more than 4,400 theorems, the whole building green). This is the endpoint of a sustained discharge effort in which the project axiom total fell \textbf{57 → 40 → 37 → 35 → … → 0}, each interface axiom either proved in a concrete finite model or recast as an explicit hypothesis. The machine-checked, axiom-free content includes: Theorems 3, 6, 7 and Lemma 1; Donald's identity and the entire \texttt{ArakiInterface} relative-entropy layer (11 axioms discharged to theorems, including the Holevo bound $\chi \le H(p)$ via operator-monotonicity of the logarithm through a \texttt{CStarMatrix} bridge) together with the \texttt{EntropyBridge}; a complete axiom-free \textbf{operator-convexity → Lieb's concavity → data-processing → strong-subadditivity} tower (\texttt{QIQTH/Entropy/}, \textasciitilde{}19 files — the finite-dimensional Carlen DPI–Lieb machinery that Mathlib lacks, formalized bottom-up: Peierls' inequality, trace-function convexity, the Loewner order, Ando's joint convexity, Lieb concavity, joint convexity of relative entropy, the data-processing inequality, and \texttt{strong\_subadditivity} / \texttt{condMutualInfo\_nonneg} as the capstones); Klein positivity; the Goldstein-Struyve Schur classification (\texttt{GoldsteinStruyveStep1.schur\_classification\_real}, formerly a named interface axiom, now proved outright) and the Tsirelson-attainability bound; the no-signaling chain; the seven structural audits; the \textbf{free-field modular-energy bound} (\texttt{ModularEnergyBound} — the Umegaki identity, the Casini bound $S(\rho)-S(\sigma)\le\langle K_\sigma\rangle_\rho-\langle K_\sigma\rangle_\sigma$, the first law, and its rigidity, with Bisognano–Wichmann giving the Unruh form $\Delta S \le 2\pi\,\Delta\langle B_{\rm boost}\rangle$; formalized modular QFT, \emph{not} a derivation of the holographic $A/4G$ bound, §1.1a); and a \textbf{finite no-collapse Born representation} (\texttt{BornJoin.finite\_noCollapseBornRepresentation}): in the finite model a unique actual record is isolated (with $\lambda$ selecting it, not the capacity bound — §1.1a), effect-Gleason (\texttt{EffectGleason.finite\_effect\_gleason}) forces the outcome weights to be Born \emph{from positivity}, product preparation gives independent trials, and the actual-value histories carry the Born product law and are Chebyshev-typical. The formalization also includes deliberate \textbf{negative audits} proving what does \emph{not} follow: linear unitary decoherence does not concentrate branch weights (\texttt{NoConcentration}); the structural axioms do not single out Born (\texttt{NoBornFromNothing}); support preservation is strictly weaker than Born equivariance (\texttt{EquivarianceGap}); operational click-statistics underdetermine the IC measure (\texttt{OperationalNoGo}). The single-record \emph{mechanism} (the number-bound form of the Macroscopic Definiteness Conjecture, §7.6) is mechanized as a conditional theorem (\texttt{CoreNoCollapse}, \texttt{CapacityModel}, \texttt{SBSBridge}, \texttt{CollisionalGamma}): a redundant Spectrum-Broadcast record forces storage $\ge R\log n$, a finite additive capacity then admits at most one macroscopic record (the threshold \emph{derived} from pointer-state orthonormality, not stipulated), and the per-collision distinguishability $\gamma = |\cos 2\theta| < 1$ is \emph{derived} from the toy Hamiltonian $H_{\rm int}=g\,\sigma_z^S\otimes\sigma_x^E$; the one remaining open input is the \emph{field-theoretic} origin of that scattering premise. \textbf{(Revised 2026.} This conditional theorem concerns an \emph{additive} record cost; since the additive (branch-summed) cost is provably \emph{not} the holographic von Neumann entropy that $Q_R$ bounds — \texttt{BranchLedger.branchSummed\_not\_bounded\_by\_Shannon} — it does \textbf{not} establish that the physical capacity forbids a two-record state. The "Macroscopic Definiteness" reading is accordingly \textbf{retired} as a category error, §1.1a: single outcomes are fixed by the non-dynamical selector $\lambda$, not by the capacity bound, and the formalization's own negative audits confirm capacity alone neither concentrates weights nor selects a basis — \texttt{NoConcentration}, \texttt{RealmSelection.capacity\_underdetermines\_realm}.)\textbf{ What machine verification does and does not establish must be stated plainly:} it certifies that the framework's \emph{conditional and structural} mathematics is correct and rests on no hidden axiom; it does \textbf{not} establish the framework's \emph{physical} postulates — the (FQ) holographic bound (P4), the Canonical IC Measure (Born) Principle (P5), and the dynamical, Lorentz-covariant law of the selector $\lambda$ — which remain the open problems of §11.4. (The Macroscopic Definiteness Conjecture is \emph{not} among them: it is \textbf{retired} as a category error, §1.1a, not a pending verification target.) A continuum Type II / Fock / modular-flow tower (\texttt{QIQTH/Fock}, \texttt{QIQTH/Spectral}, \texttt{QIQTH/Entropy}, \texttt{QIQTH/TowerGNS}) is in active development toward those open problems; its finite-stage inductive-limit \textbf{Tomita–Takesaki modular theory is now complete} — the closed Tomita operator $\bar S$, the self-adjoint modular operator $\Delta$, the modular unitary group $\Delta^{it}$ (identified with the physical flow), the modular conjugation $J$ with the polar decomposition $\bar S = J\Delta^{1/2}$ on the core, Tomita's theorem first half ($\Delta^{it}\mathcal{M}\Delta^{-it}=\mathcal{M}$) and second half in inclusion form ($J\mathcal{M}J\subseteq\mathcal{M}'$), and a genuine non-tracial KMS state ($\Delta\neq 1$), to our knowledge the first complete such in any proof assistant — while the full equality $J\mathcal{M}J=\mathcal{M}'$, the continuum Type III$_1$ limit, and the type classification remain cited frontiers. The complete audit is in \texttt{lean/mathlib/QIQTH/AxiomAudit.lean}; see \texttt{lean/mathlib/QIQTH/} in the project repository.
\end{abstract}

\tableofcontents
\newpage


\section{Introduction}

\subsection{The hidden inconsistency in standard quantum foundations}

Standard foundational discussion of quantum mechanics typically holds four claims simultaneously:

\begin{enumerate}
  \item The wave function is physically real.
  \item Bounded spacetime regions have finite physical information capacity (the holographic principle, supported by Bekenstein, 't Hooft, Susskind, Bousso, Banks, and the recent algebraic-QFT-plus-gravity work of Chandrasekaran-Penington-Witten and Witten).
  \item The amplitudes of the wave function are exact real numbers (specifiable to arbitrary precision in principle).
  \item Macroscopic superpositions, once formed, persist as physically real components forever.
\end{enumerate}

These four claims cannot all be true — but the operative tension is \emph{entropic}, not about amplitude precision. If the wave function is physically real (1) and macroscopic superpositions, once formed, persist as co-actual physically real components forever (4), then the regional reduced state $\rho_R$ of $\Phi$ accumulates ever more mutually-decohered branches, so its von Neumann entropy $S(\rho_R)$ grows without bound — contradicting the finite regional entropy budget (2). The exactness of amplitudes (3) is a \emph{red herring}: continuous amplitudes cost nothing entropically (a pure state has $S=0$), so (3) can be kept. The genuine culprit is the co-actual persistence in (4). The standard interpretations of QM each resolve this tension by denying one of the four claims:

\begin{itemize}
  \item \textbf{Many-Worlds} denies single-world realism (accepts (1)–(4) but at the cost of branching ontology with infinitely many physically real branches; pays in spacetime-information overflow that is rarely confronted).
  \item \textbf{Bohm} denies that the wave function is the complete physical content (denies a strong form of (1); adds primitive particle ontology).
  \item \textbf{GRW / CSL / DP / OR} deny that Schrödinger evolution is exact (modify dynamics so that (4) fails dynamically).
  \item \textbf{QBism / RQM} deny that the wave function is straightforwardly physically real (deny (1) in favor of epistemic or relational status).
  \item \textbf{Modal interpretations / decoherent histories} add a value-rule or history-selection structure to break (4) without modifying dynamics.
\end{itemize}

Each pays a distinct foundational price. Each accepts the tension as a \emph{given} of the theory and pays the price to dissolve it.

This paper proposes a different resolution: \textbf{deny the co-actual persistence in (4), keeping (1), (2), and (3)}. The holographic bound is read as a bound on the regional \emph{entropy} $S(\rho_R)$ of $\Phi$ (its effective modular rank $e^{S(\rho_R)}$), never a finite matter Hilbert-space dimension and never a finite-precision floor on amplitudes; amplitudes (3) remain exact. After decoherence the per-run regional content is a single actual record fixed by the non-dynamical selector $\lambda$ — so the branches do not persist as \emph{co-actual} physically real components, and $S(\rho_R)$ stays within the finite budget (2). With (4)'s co-actual persistence denied, all four survivors reconcile: the wave function is physically real (1), bounded spacetime has finite entropy capacity (2), amplitudes are exact (3), and a single record is actual per run via the $(\Phi,\lambda)$ ontology (P1), the global evolution left exactly unitary. [\emph{Historical note:} an earlier version instead denied (3), positing a finite physical specification precision on amplitudes ("amplitudes within $\epsilon(R)$ of $0$ are physically $0$"), \emph{and} attributed the single record to the finite capacity \emph{itself} forbidding a two-record state ("a structural consequence, not an added postulate"). \textbf{Both moves are retired} (§1.1a): the finite-resolution reading is not load-bearing, and capacity — being the kinematic constraint behind the area law and the field equations — does not, and cannot, perform the outcome selection, which is $\lambda$'s role. The finite-resolution reading is not relied on below.]

This is arguably the \emph{least radical} resolution. It does not multiply worlds, add particles, modify dynamics, or subjectivize the wave function. It accepts one physically natural constraint, that bounded spacetime carries finite entanglement \emph{entropy} (holography and quantum gravity motivate this independently), together with the non-dynamical actuality fact $\lambda$. The measurement problem is resolved by recognizing that the co-actual persistence in (4) was the illegitimate assumption: once decoherence renders the records non-interfering and $\lambda$ makes one actual, the others do not persist as co-actual realities, and no entropy overflow arises.

\subsection{Two facts, one world}

The framework rests on two physical commitments.

\textbf{Fact 1, There is only the wave function, and it carries no probability.} $\Phi$ is the whole of physical reality, evolving by one exactly-unitary law: no collapse, no branches-as-substances, no observer standing outside it, and no probability built into it. What is conventionally called the Born rule is an \emph{across-run frequency}, not a number the world assigns to a single run. The only further fact is $\lambda$, \emph{which} of $\Phi$'s macroscopic realizations is the actual one. $\lambda$ is not a substance, not an observer, and not a probability; it is actuality, and nothing more (§7.6, §11.3).

\textbf{Fact 2, Regional entropy is limited by surface area.} For every bounded region $R$ the regional entanglement (von Neumann) entropy of $\Phi$ is bounded by the surface area, $S(\rho_R) \le Q_R = A(\partial R)/4\ell_P^2$ (Bekenstein–Bousso; (FQ) below). Equivalently, the effective modular rank $e^{S(\rho_R)}$ is finite, so a region carries only \emph{finitely many distinguishable records}. This is a bound on \emph{entanglement entropy}, \textbf{not} a finite matter Hilbert-space dimension and \textbf{not} a finite-precision floor on amplitudes.

From these, with decoherence, the measurement problem dissolves, not by adding a mechanism but by removing a mistake: \textbf{the binary, definite outcomes we observe are a feature of macroscopic record-structures, which are area-limited and einselected, not of the wave function.} We measure in $0/1$ because \emph{we} are macroscopic.

The positive thesis makes the division of labor explicit, and is deliberate about what each ingredient does and does \emph{not} do:

\begin{itemize}
  \item \textbf{Continuous, conserved amplitude.} The weights $|c_k|^2$ are continuous and are \emph{conserved} under unitary evolution; decoherence never drives them to $0$ or $1$ (formalized as \texttt{NoConcentration}; §9). The wave function itself has no binary structure and no probability.
  \item **Decoherence + einselection ⟹ stable, Boolean \emph{records}.** What becomes effectively $0/1$ is a \emph{different} variable from the amplitude: macroscopic record-existence (whether a definite pointer-record is instantiated in this region). Einselection makes such records redundant, robust, and mutually exclusive; this is \emph{why a macroscopic system: apparatus, brain, or AI, can only register a definite result.} This explains the \textbf{stability and classicality} of outcomes, not their \textbf{uniqueness}.
  \item \textbf{The holographic bound ⟹ finitely many records.} $Q_R$ caps \emph{how many} distinguishable records a region can hold (its genuine, non-redundant role). It does \textbf{not}, by itself, forbid a superposition of records (§7.6).
  \item \textbf{$\lambda$ ⟹ which record is actual.} That exactly one record is the lived one is the non-dynamical actuality fact $\lambda$. Born statistics are the across-run tally of $\lambda$ under a typicality measure: emergent bookkeeping, not a per-run probability (§8, Open Problem on the typicality measure).
\end{itemize}

The dissolution is then this: there is \textbf{no external observer to whom a probability is assigned.} "The outcome" is simply which macroscopic realization of $\Phi$ is actual. Collapse and fundamental probability were posited to explain a definiteness that lives in \emph{us}, macroscopic, area-limited record-structures, and was never a feature of $\Phi$. There is only $\Phi$; the binary is the texture of our scale, not of the world.

\subsection{The position}

This paper develops a position in the foundations of quantum mechanics whose structure is:

\begin{enumerate}
  \item The wave function of the universe evolves unitarily under the Schrödinger / Heisenberg evolution of the underlying field theory. Standard QM dynamics is preserved exactly.
  \item The Bekenstein-Bousso holographic bound is treated as a foundational constraint on the \textbf{regional entropy of the wave function in any bounded region $R$}: it bounds the von Neumann (renormalized) entropy $S(\rho_R)$ of $\Phi$'s regional reduced state. It is not a bit-count of "total physical information", not a finite matter Hilbert-space dimension, and not a finite precision on amplitudes.
  \item The mathematical home of this constraint is the algebraic QFT framework with gravitational dressing. The regional algebra $\hat{\mathcal{A}}(R)$, obtained from the standard Type III$_1$ local algebra $\mathcal{A}(R)$ via the crossed product with the modular flow of a reference state, is of \textbf{Type II} and admits a semifinite trace whose finite values realize finite renormalized entropy for regional states. This is the framework of Chandrasekaran-Penington-Witten (2022), Witten (2022), and successors.
  \item The (FQ) axiom: \emph{the physical content of the per-run wave function in any bounded region $R$ is given by its expectation values on $\hat{\mathcal{A}}(R)$, and the renormalized entropy is bounded by $Q_R = A(\partial R)/(4\ell_P^2)$.}
  \item The Type II structure of $\hat{\mathcal{A}}(R)$ gives a basis-independent regional identity of states together with a finite renormalized entropy: states agreeing on all operators in $\hat{\mathcal{A}}(R)$ are physically identical in $R$, and the regional entropy $S(\rho_R)$ is finite. The equivalence is exact algebraic-state equality between regional states, \textbf{not} a finite precision on amplitudes.
  \item The standard textbook formal wave function, a superposition $\sum_i c_i |s_i\rangle|A_i\rangle|E_i\rangle$, is an idealized statistical description of an ensemble of per-run wave functions across many runs. The per-run wave function evolves unitarily from specific microscopic initial conditions; decoherence renders the macroscopic record components non-interfering; and the non-dynamical selector $\lambda$ (P1) marks exactly one decohered record actual (capacity is kinematic and does \emph{not} forbid a $\ge 2$-record state — §1.1a) leaves the per-run regional content single-record. The run's microscopic initial conditions index which record.
  \item Born statistics emerge from typicality of microscopic initial conditions across runs.
\end{enumerate}

\subsection{The irreducible postulates}

After the machine-checked discharge-and-grounding effort (§11.4b), the framework's content reduces to five postulates; the irreducible new physics is \textbf{(P4) + (P5)} on the (P1) ontology, with (P2)–(P3) the standard quantum-relativistic arena:

\begin{itemize}
  \item \textbf{(P1) The $(\Phi,\lambda)$ ontology.} The universal wave function $\Phi$ is the complete ontology (no external observer, no fundamental probability); a non-dynamical selector $\lambda$ makes exactly one decoherent record actual per run — no collapse term, no branching.
  \item \textbf{(P2) Quantum kinematics.} The complex-Hilbert-space / operator-algebraic framework (states as positive normalized functionals, observables as a net of algebras, the $\operatorname{tr}(\rho\,\cdot)$ pairing). The Born \emph{squared modulus is not a separate postulate}: it is the inner-product geometry of this arena ($\operatorname{tr}(|\psi\rangle\langle\psi|\,P_k)=|\langle k|\psi\rangle|^2$), forced into linear/trace form by (P5) and \emph{selected} — exponent $1$ over the $\alpha$-family — by refinement no-signaling.
  \item \textbf{(P3) Microcausality.} Spacelike-separated regional algebras commute. This is the Lorentz-covariance input; it is \emph{not} the source of the Born premises (see below).
  \item \textbf{(P4) The (FQ) holographic capacity bound.} The regional renormalized entropy is bounded by $Q_R = A(\partial R)/(4\ell_P^2)$. Together with (P5) this is the irreducible new physics; (P4) is the sole holographic input to the machine-checked QIQT$\to$GR chain (§11.4b). That chain is now instantiated for an explicit free Klein–Gordon field on a curved \textbf{pp-wave} spacetime (\texttt{qiqt\_gr\_ppwave\_showcase}), in which \emph{every geometric and analytic premise is discharged} — the metric and tetrad, the Raychaudhuri \textbf{area derivative} (an expansion-free congruence has constant cross-sectional area), and the \textbf{entropy bound} $S\le\eta A$ (Shannon's maximum at the holographic capacity) — leaving as hypotheses \emph{exactly} the irreducible floor: the matter equation of motion, (P4) itself, and the \textbf{localization map} (the field-coupled record law whose entropy rate equals the stress flux $2\pi\hbar^{-1}T_{kk}$). The localization map is provably \emph{not} dischargeable by analysis: at the uniform reference the Shannon entropy is stationary ($\sum p'=0$), so the heat rate's value is forced to be the stress flux. Thus the Einstein field equations for the pp-wave spacetime follow, machine-checked and axiom-free, from the matter equation of motion $+$ (P4) $+$ the localization map.
  \item \textbf{(P5) Quantum equilibrium of the typicality measure.} $\lambda$'s measure is refinement-equivariant (the Dürr–Goldstein–Zanghì / Valentini condition).
\end{itemize}

\textbf{Internal-consistency constraint ($\Phi$, not the selected branch, carries the holographic entropy).} Because the holographic value $Q_R = A(\partial R)/4\ell_P^2$ is an \emph{entanglement} (von Neumann) entropy of the regional state — not a Boltzmann log-count of decohered branches — (P4) is necessarily a statement about the \textbf{pre-selection constitution $\Phi$}, whose reduced state $\rho_R$ carries that entropy. The non-dynamical, inert character of $\lambda$ in (P1) is therefore not merely interpretive but \emph{required}: conditioning on the $\lambda$-selected record would destroy the off-diagonal coherence the area measures (the entropy would collapse from $S(\rho_R)$ to the within-branch $S(\rho_{R,i})$ in the decomposition $S(\rho_R) = H(\{p_i\}) + \sum_i p_i\,S(\rho_{R,i})$), making (P4) incompatible with the Bekenstein–Hawking / Ryu–Takayanagi reading. So geometry responds to $\Phi$, never to the actualized branch; equivalently, the regional "record count" $\log|R|$ must be read as the effective modular rank $e^{S(\rho_R)}$ — a quantum edge-mode degeneracy of the pre-selection state — not as a cardinality of classical outcomes. (This is also why a microscopic \emph{derivation} of (P4) — replacing the postulate by a JLMS-type identity $K_{\partial R} = A/4\ell_P^2 + K_{\rm bulk}$ on the modular/edge algebra — is the natural next target rather than any counting of actualized records.)

The refinement of this account is that the two Born-rule premises of earlier formulations — the additivity / non-contextuality \emph{bridge} and the $\mu$-\emph{selection} — both reduce to \textbf{(P5) alone}. In the Lean development: refinement-equivariance implies selector no-signaling (\texttt{BornMuSelection.equivariant\_no\_signaling}); selector no-signaling is \emph{equivalent} to additivity (\texttt{BornActualityConsistency.apc\_iff\_positiveAdditive}), hence to the Born rule; and equivariance grounds the measure-selection directly (\texttt{equivariant\_context\_independent}, \texttt{mu\_selection\_martingale}), with the no-go theorems certifying that a selection principle is unavoidable (\texttt{NoBornFromNothing.any\_anti\_born\_realizable}) and non-vacuous (\texttt{BornRoutes.sqRule\_refinement\_signals}). Crucially, (P5) is \textbf{not} reducible to (P3): the formalization proves that observable microcausality does \emph{not} entail selector no-signaling (\texttt{SelectorRefinement.alphaSq\_selector\_signals}) — the two live at different layers (observable algebra vs.\textbackslash\{\} actuality measure), so they cannot be merged. The genuinely irreducible \emph{physics} therefore reduces to \textbf{(P4) + (P5)}, on the (P1) ontology, with (P2)–(P3) the standard quantum-relativistic arena.

**(P4) refined — the area \emph{floor} is now a derived theorem; only the finite \emph{capacity} is postulated (the P4-MICRO formalization, 2026).** The machine-checked development separates two things earlier bundled in (P4). The \emph{postulate} is only that the regional capacity is \textbf{finite} — a UV-finite record structure (the "Quantized Information" core). This finiteness has \textbf{two machine-checked, provably-distinct layers} (verified by direct read of the hypothesis classes): in the \textbf{finite-dimensional model} it is a literal record \textbf{count}, \texttt{HolographicCapacityBound} $= (\log\operatorname{card}R \le$ areaTerm$)$, i.e. $N_R = \operatorname{card}R < \infty$; in the \textbf{continuum / Type III$_1$} setting it is the corresponding finite \textbf{entropy} bound, \texttt{Phase5Master} $= (S_{\mathrm{vN}} + S_{\mathrm{rel}} \le$ areaTerm$)$ over an infinite-dimensional one-particle space (no cardinality), and \texttt{EntropyNotCardinality} \emph{proves} the entropy bound is \textbf{not} a count ($S_\tau \le Q \nRightarrow \operatorname{card}\le e^{Q}$). So the continuum/physical claim is finite \emph{entropy}; the finite \emph{count} is the tractable finite-dimensional shadow (strictly stronger). The finiteness is on \emph{records/entropy}, \textbf{never} a finite matter Hilbert space (D2/D3). Moreover the count layer is \textbf{not derivable} from the entropy/area bound (the \texttt{EntropyNotCardinality} no-go forbids it): the only sound \emph{operational} count is a \textbf{Holevo capacity} $\log M_\epsilon \le (Q+h_2(\epsilon))/(1-\epsilon)$ under a relative-entropy bound, which is finite as a \emph{number} only given an \emph{imported} energy cutoff — the \textbf{Bekenstein/microcanonical} bound, standard holography, not new physics. QIQT-H is distinctive here only via a $Q_R$ differing from standard generalized entropy $S_{\rm gen}=A/4G+S_{\rm bulk}$ — and such a $Q_R$ \textbf{cannot be derived} from the program's principles (a conditional no-go: area/JLMS use $S_{\rm vN}$, the finite count is independent of $S_{\rm vN}$ via \texttt{EntropyNotCardinality}, and $\lambda$ is inert). It is possible \emph{only} by adding the explicit \textbf{max-entropy bridge postulate} — gravity's capacity is $S_{\max}$ (the finite \emph{count}), not $S_{\rm vN}$ — a new postulate, not a derivation, which predicts the capacity-of-entanglement gap $Q_R-S_{\rm gen}=S_{\max}-S_{\rm vN}\sim\sqrt{V_{\rm gen}}$ (coefficient and the value of $G$ are cited frontiers; the no-go \texttt{svn\_underdetermines\_smax}, the gap \texttt{gap\_nonneg}, the capacity of entanglement \texttt{capEnt\_nonneg}, and the conditional prediction \texttt{distinctive\_gap} under the \texttt{MaxEntropyCapacity} typeclass postulate are \emph{proved} axiom-free in \texttt{QIQTH/MaxEntropyCapacity.lean}). The operational bound itself is \emph{proved} axiom-free in \texttt{QIQTH/OperationalCapacity.lean} (\texttt{record\_capacity}, the $\varepsilon=0$ form \texttt{exact\_distinguishable\_capacity}, the Bekenstein \texttt{gibbs\_entropy\_bound}) — all Holevo/Bekenstein-class, none deriving the count from the area law. Finiteness \emph{alone} gives only $S_{\mathrm{vN}}(\rho_R) \le \log N_R$; it does \textbf{not} by itself fix whether $\log N_R$ scales with \textbf{area} or volume (a generic finite local cutoff has \emph{volume}-scaling maximum entropy — the area law is a property of \emph{vacuum entanglement}, not an automatic property of the capacity). That the bound takes the \textbf{holographic area form} ($\log N_R \le A(\partial R)/4\ell_P^2$) is \textbf{derived}, but in a \emph{conditional} Sakharov / induced-gravity \textbf{bridge} — assuming local relativistic QFT on a smooth background with a covariant UV cutoff identified with the finite microstructure — \textbf{not} from finiteness alone: there the area law $S\propto A$ \emph{emerges} from the conical-deficit geometry (the cone curvature is a $\delta$-function on the boundary surface, whose integral is the area), and the $1/4$ is the universal ratio between the conical replica-entropy coefficient and the induced Einstein–Hilbert coefficient — two quantities sharing one UV coefficient (the \emph{ratio} machine-checked, \texttt{SakharovRatio.sakharov\_ratio}; the emergence is the standard Susskind–Uglum/Solodukhin heat-kernel result, Stage B, \texttt{docs/SAKHAROV\_KG\_STAGE\_B.md}). The carried inputs are the \emph{value} of $G$/$\ell_P$ (the species/cutoff problem) and, for the full effective action, $\Lambda$ and higher-curvature terms. The capacity $Q_R$ is read, per the internal-consistency constraint above, as the effective modular rank $e^{S(\rho_R)}$ of the pre-selection constitution $\Phi$, and is carried in Lean as the \texttt{HolographicCapacityBound} typeclass hypothesis, \emph{not} a Lean \texttt{axiom}. The \textbf{only} carried datum is the \emph{value} of $G$/$\ell_P$ (the species/cutoff problem). \textbf{(2026 — $G$ promoted from carried to derived.)} Even that carried datum can be reduced: positing a fundamental \textbf{record-granularity scale} $\Lambda_s$ ($a_0 = 1/\Lambda_s$, the finite-information "pixel size") \emph{in place of} $\ell_P$ delivers $G = 1/(N\,\Lambda_s^2)$ — the Sakharov/Dvali species bound, machine-checked axiom-free (\texttt{QIQTH/InducedNewtonConstant.lean}: \texttt{inducedG\_delivers}). So $G$ moves from \emph{carried} to \emph{derived}, and P4-MICRO's carried inputs collapse to a \textbf{single scale} $\Lambda_s$, from which the finite capacity \emph{and} $G$ both follow — the capacity exponent re-expresses in primitives as $A/4\ell_P^2 = (A/4)\,N\Lambda_s^2$. What is derived is the \emph{relation} (and the dimensionless $G/a_0^2 = 1/N$, \texttt{inducedG\_ratio\_is\_pure\_number}); the \emph{numerical value} of $G$ still requires the species-coefficient accounting (a cited frontier), and $\Lambda_s$ remains the one carried scale — a length cannot come from a count, exactly as AdS/CFT carries the string scale $\alpha'$. This is a reframing of P4-MICRO's primitive (granularity in place of $\ell_P$), not a computation of $G$'s value. This granularity capacity \textbf{maps onto the standard holographic dictionary}: with the induced $G = 1/(N\Lambda_s^2)$, the boundary-CFT Cardy microstate count of a BTZ horizon \emph{equals} QIQT-H's own bulk capacity exponent $(A/4)\,N\Lambda_s^2$ — a machine-checked correspondence (\texttt{QIQTH/HolographicBridge.lean}: \texttt{btz\_cardy\_eq\_qiqth\_capacity}; the AdS radius cancels). This is a \emph{correspondence} (the two holographic bookkeepings agree under the shared $G$), \textbf{not} an import of a boundary CFT, the Cardy formula, bulk reconstruction, or AdS/CFT's independent cross-check. The \emph{area floor itself}, $S_{\mathrm{vN}}(\rho_R) \le Q_R$, is then a \textbf{theorem} — \texttt{area\_floor\_vonNeumann} (\texttt{QIQTH/FQBoundMicro.lean}): the elementary finite-dimensional maximum-entropy bound $S_{\mathrm{vN}} \le \log\dim$ applied to the spectrum, axiom-free (standard three). Three honesty refinements are \emph{enforced} in the formal statements: the bound is on the \textbf{von Neumann} entropy of the spectrum (not the Shannon entropy of a decohered record law — the two coincide only diagonally, $S_{\mathrm{vN}}\le H(\text{record})$ in general), consistent with the entanglement-entropy reading above; only the inequality $\le$ is needed for the floor, with equality reserved for the maximally-mixed/saturating sector (\texttt{HolographicCapacityBound} vs \texttt{HolographicCapacityExact}); and $N_R$ is the \textbf{regional} capacity — an explicit finite type-I/code cutoff of the genuinely type-III$_1$ local algebra, not a global Hilbert-space dimension. The (P4)$\to$GR chain is recast accordingly: \texttt{gr\_from\_p4micro} (\texttt{QIQTH/GRFromMicro.lean}) feeds this \emph{derived} area floor into the Jacobson capstone \texttt{qiqt\_bekenstein\_gives\_gr}, yielding the Einstein field equations with \textbf{exactly} the residual labelled inputs — the Bisognano–Wichmann / Unruh modular flux \texttt{hFlux} (discharged for the free field via \texttt{Fock.OneParticleBW}), Raychaudhuri focusing, stress-energy conservation, metric regularity, and — the genuine residual — the \textbf{joint reference-state identification}: that the horizon equilibrium reference is \emph{both} the Bisognano–Wichmann modular state (whose modular Hamiltonian gives the Unruh flux, $\beta = 2\pi/\kappa$) \emph{and}, in its record/edge sector, capacity-saturating ($S = \log\dim = \eta A$ for the same geometric area). These are different states in general ($\beta = 2\pi/\kappa$ vs the maximally-mixed $\beta = 0$), so the physical content sits in this identification. The entropy-area \emph{variation} $\delta S = \eta\,\delta A$ that Jacobson needs is itself a \textbf{derived theorem} — \texttt{DifferentialAreaLaw.differential\_area\_law}: the capacity bound + point-saturation (\texttt{area\_floor\_saturates}, $S(0)=\eta A(0)$ at the maximally-mixed record) + the entanglement first law \emph{entail} it, with \textbf{no} hypothesis asserting $S=\eta A$ or $\delta S=\eta\,\delta A$. So there is \textbf{no} separate "Clausius / area-saturation postulate" — saturation is discharged; the open Lean target is to prove the reference of \texttt{relEntropy\_self}/\texttt{hFlux} coincides (or shares its record marginal) with the one \texttt{area\_floor\_saturates} uses, likely via an edge$\otimes$bulk split (the area entropy from a uniform record/edge sector, the BW flux from the bulk). The Lean now \emph{enforces} the honest scope: \textbf{the capacity postulate alone does not give GR} — a microstate \emph{count} cannot supply a \emph{temperature}, so the thermal \texttt{hFlux} slot is irreducibly modular (Route 1 / Type II) and is not derivable from (P4). The $1/4$ coefficient is the separately-derived Sakharov induced-gravity ratio (\texttt{SakharovRatio.sakharov\_ratio}, regulator- and matter-independent); the value of $G$ remains a carried UV datum. Deriving the capacity \emph{law} itself — replacing the postulate by a JLMS-type modular identity $K_{\partial R} = A/4\ell_P^2 + K_{\rm bulk}$ — remains the natural next target (Route 1), as noted above. \textbf{(Reframed, 2026, GPT-5.5-pro expert review.} For a fixed-background \emph{free scalar} the $A/4\ell_P^2$ term is \textbf{not} derivable — the free theory has no Newton constant and no geometric area operator, the wedge entropy's cutoff coefficient is matter/scheme-dependent (not universally $1/4G$), and the $\delta A/4G = 2\pi\!\int\!\delta T_{kk}$ step \emph{uses the Einstein equations}, not pure Bisognano–Wichmann kinematics. So BW supplies the Unruh $2\pi$ but not the $1/4G$; the $A/4\ell_P^2$ identification stays a \emph{gravitational input}, and the continuum Type III$_1\!\to$II crossed-product dual-weight trace is a multi-year cited frontier. What \emph{is} derivable — and is the honest content of Route 1 — is the free-field \textbf{modular-energy bound} $\Delta S \le \Delta\langle K_W\rangle = 2\pi\,\Delta\langle B_{\rm boost}\rangle$ (Casini/first law, from relative-entropy positivity + the machine-checked one-particle BW $K_W = 2\pi B_{\rm boost}$), which upgrades \texttt{Phase5Master}'s modular pieces from a carried hypothesis to theorems — \emph{formalized modular QFT}, not a derivation of the holographic $A/4G$ bound. Built in \texttt{QIQTH/ModularEnergyBound.lean}, \texttt{ROUTE1\_MODULAR\_PLAN.md}.)

\textbf{Independent convergent work — QIQT-H as the machine-verified counterpart (2026).} Dorau and Much (\emph{From Quantum Relative Entropy to the Semiclassical Einstein Equations}, Phys. Rev. Lett. 2026; arXiv:2510.24491) independently give a quantum-field-theoretic extension of Jacobson's argument, deriving the \textbf{semiclassical Einstein equations} from the \textbf{Araki–Uhlmann relative entropy} between the vacuum and a coherent excitation of a Klein–Gordon field on a local Rindler (bifurcate Killing) horizon. Their chain is, step for step, the free-field GR chain machine-checked here: (i) the modular flow acts as the geometric boost — Bisognano–Wichmann / Summers–Verch, $\Delta_\mathcal{R}^{it}=\mathfrak{D}_{2\pi t}$ — our unconditional \texttt{Fock.OneParticleBW}; (ii) the coherent-state relative entropy equals the horizon energy flux $\delta Q$, $S^{\rm rel}(\omega_0\|\omega_\phi)=-2\pi\!\int U\langle{:}T_{ab}{:}\rangle\xi^a\xi^b$ (Casini–Grillo–Pontello) — our \texttt{ModularEnergyBound} (the first law $\Delta S=\Delta\langle K_W\rangle=2\pi\,\Delta\langle B_{\rm boost}\rangle$); (iii) the area variation is fixed by Raychaudhuri focusing on the horizon — our \texttt{DifferentialAreaLaw.differential\_area\_law} together with the Raychaudhuri machinery; and (iv) stress-energy conservation closes it to $R_{ab}-\tfrac12 R\,g_{ab}+\Lambda g_{ab}=\alpha\langle{:}T_{ab}{:}\rangle$ — our capstone \texttt{qiqt\_gr\_freefield\_complete} / \texttt{qiqt\_bekenstein\_gives\_gr}. Where they \emph{assume} the Bekenstein–Hawking relation $S^{\rm rel}=\delta A/4$ to fix $\alpha=8\pi$, QIQT-H additionally \emph{re-derives} the $1/4$ ratio (Sakharov, \texttt{SakharovRatio.sakharov\_ratio}). We regard their Letter as independent, peer-reviewed confirmation that the relative-entropy $\to$ modular-theory $\to$ Jacobson route to the Einstein equations is a live and correct research direction; QIQT-H is its \textbf{machine-verified counterpart} — the \emph{same} derivation reduced to a kernel-checked chain with every physical input in an explicit hypothesis ledger (the claim card of \texttt{qiqt\_gr\_freefield\_complete}). Both meet the \emph{same} honest frontier: their closing caveat — higher-order corrections on a genuinely curved horizon, "\emph{technically demanding, especially regarding the modular data}" — is precisely QIQT-H's cited continuum/curved-correction frontier (the localization map and the Riemannian-heat-kernel / Seeley–DeWitt $(1/6-\xi)$ gap).

**What P4-MICRO implies about the \emph{kind} of theory QIQT-H is: a constraint (boundary) theory, and its "realisations".** The P4-MICRO split has a structural consequence worth stating explicitly. Because the only distinctive \emph{postulate} is the \textbf{finiteness} of the regional capacity — the area \emph{form} and the area \emph{floor} being \emph{derived} (the floor outright via \texttt{area\_floor\_vonNeumann}, the form conditionally via the Sakharov bridge) — and because the framework everywhere \emph{accepts} a field algebra and \emph{bounds} its records rather than constructing them, QIQT-H is most accurately read not as a generative theory but as a \textbf{descriptive constraint (boundary) theory built on a small postulate core}: the (P1) ontology, the (P5) equilibrium measure, and finite regional capacity, together with a family of \emph{necessary conditions} on the regional record structure of any quantum theory. A \textbf{generative theory} — a constructive interacting QFT, the full Standard Model, or a specific quantum-spacetime construction — is then a \textbf{realisation of QIQT-H} when its regional data, encoded into the finite-capacity microstate space, satisfies those conditions. The conditions are exactly the machine-checked content: Born-weighting of the records (\texttt{records\_born\_and\_area\_bounded}), the area floor $S_{\mathrm{vN}}\le Q_R$ (\texttt{area\_floor\_vonNeumann}), corner-faithful representation of the field algebra — the encoding lands in the corner $P\cdot\mathrm{End}(\mathcal{H}_R)\cdot P$, $P=VV^\dagger$, never the ambient identity, with the no-overclaim guards — the necessary \textbf{truncation} of bosonic modes (\texttt{no\_finiteDim\_CCR}, \texttt{finite\_weyl\_qpow\_eq\_one}), the absence of exact finite dilation/boost scaling (\texttt{finiteDim\_scaling\_forces\_zero}), the Ryu–Takayanagi area bounds with min-cut in its (non-metric) \emph{area} role (\texttt{minCut\_area\_not\_metric}, \texttt{entropy\_le\_cut}), and the selector no-go theorems (\texttt{NoBornFromNothing}). P4-MICRO additionally pins what a realisation must \emph{supply} beyond mere compatibility — precisely the inputs the area-\emph{form} derivation assumed: a local relativistic QFT on a smooth background, a covariant UV cutoff identified with the finite microstructure, the value of $G/\ell_P$, and, for the gravitational sector, the reference-state identification (the horizon-equilibrium state being simultaneously Bisognano–Wichmann modular and capacity-saturating). The characterization is \textbf{non-vacuous} — the boundary genuinely \emph{excludes} theories (exact finite bosonic CCR; an anti-Born or volume-saturating selector; a region whose entropy exceeds its area) — and it is the same descriptive/generative relationship that thermodynamics bears to statistical mechanics, or the analytic/conformal bootstrap to specific Lagrangian QFTs: the constraint layer is genuine physics \emph{precisely because} it bounds, and the generative theory is the realisation that makes the bounds hold. The bounds are, of course, \textbf{conditional on the postulate core holding of nature} — the part experiment ultimately adjudicates — and the open \emph{generative} problems (constructive interacting QFT, the background-independent $3{+}1$ manifold) remain cited frontiers, untouched by this reframing. In short: **QIQT-H specifies \emph{what must hold} of any region's record structure; a realisation is \emph{what makes it hold}.**

\textbf{The code–capacity bridge — a formalized, separated refinement (2026).} A small axiom-free Lean module (\texttt{QIQTH/CodeCapacityBridge.lean}) makes the capacity constraint a \emph{typed interface} between the field and the microstructure, rather than silently identifying field states with microstates. It keeps the field's regional \textbf{code space} $C$ separate from the microstate space $\mathcal{H}_R$ and links them only by an explicit isometric encoding $V:C\hookrightarrow\mathcal{H}_R$. Then \texttt{encoded\_field\_entropy\_le\_area}: given the encoding (so $\dim C \le \dim\mathcal{H}_R$) and the \texttt{HolographicCapacityBound} hypothesis $\log\dim\mathcal{H}_R \le A/4\ell_P^2$, every code density obeys $S_{\mathrm{vN}}(\rho) \le \log\dim C \le \log\dim\mathcal{H}_R \le A/4\ell_P^2$, with all encoded record expectations preserved (\texttt{encoded\_record\_expectation}: $\operatorname{Tr}((V\rho V^\dagger)(VOV^\dagger))=\operatorname{Tr}(\rho O)$); instantiated for the actual free-field code dimensions — the CAR Fock $2^n$ and the number-cutoff symmetric Fock $\binom{d+N}{N}$ — and capped by a record-count bound $\log|I|\le A/4\ell_P^2$ for any family of perfectly distinguishable records (\texttt{record\_log\_card\_le\_area}). The one genuine \emph{structural} asymmetry it certifies is the CAR/CCR cutoff distinction: exact finite-dimensional canonical commutation $[a,a^\dagger]=1$ is impossible ($\operatorname{Tr}[a,a^\dagger]=0\ne\dim\mathcal{H}$, \texttt{no\_finiteDim\_CCR}), so the photon's bosonic mode requires a number/energy cutoff to fit a finite microstate sector, whereas a fermionic CAR sector $\bigwedge h$ is finite-dimensional and fits exactly. \textbf{Scope, enforced in the statements:} these are \emph{conditional} bounds — the Lean proofs \emph{transport} an assumed finite microstate capacity through an isometric code encoding; they do \textbf{not} derive the holographic bound, the value of $G$, the Type-II renormalized entropy, or the matter spectrum beyond the CAR-vs-CCR cutoff distinction. Capacity here constrains the field's entropy and record count \emph{through} the fitting inequality; it does not generate the field.

\textbf{The corner construction — the field's records inside the capacity-bounded microstate memory (2026).} A follow-on axiom-free Lean module (\texttt{QIQTH/CornerConstruction.lean}) deepens the code–capacity bridge into a faithful representation of the field's records inside the microstate space, while keeping the honest field/microstate separation rigorous. The load-bearing structural device is the \textbf{corner}: everything transported by the encoding $A \mapsto VAV^\dagger$ lands in $P\cdot\mathrm{End}(\mathcal{H}_R)\cdot P$ with the code projector $P=VV^\dagger$ as its unit — \emph{never} the ambient identity $1_{\mathcal{H}_R}$ unless the code fills the whole microstate space ($P=1$, \texttt{codeProjector\_eq\_one\_iff\_encode\_one}). Within this discipline the development establishes, axiom-free (standard three), that the field's regional records are: \textbf{faithfully read back} — the encoded $n$-point correlator equals the bare one, $\operatorname{Tr}\!\big((V\rho V^\dagger)\,\iota_V(A_1)\cdots\iota_V(A_n)\big) = \operatorname{Tr}(\rho\,A_1\cdots A_n)$ (\texttt{encoded\_npoint}); \textbf{Born-weighted with area-bounded information} — the Shannon entropy of the Born record law obeys the area floor (\texttt{born\_readout\_entropy\_le\_area}); \textbf{algebraically represented} — the electron's CAR transports to the corner, $\{\iota_V(a(f)),\iota_V(a^\dagger(g))\} = \langle f,g\rangle\,P$ (\texttt{encoded\_anticomm}), with a \emph{no-overclaim guard} certifying that ambient-identity CAR would force $P=1$ (\texttt{encoded\_CAR\_ambient\_forces\_full}), while the photon's truncated ladder carries an \emph{explicit, surviving} truncation defect $[\iota_V(a),\iota_V(a)^\dagger]=P-N\,\iota_V(|N{-}1\rangle\langle N{-}1|)$ (\texttt{encoded\_truncated\_ladder\_commutator}) — the bosonic mode is necessarily truncated in finite capacity, the error quantified rather than hidden; \textbf{modularly structured} — a finite (Type-I) KMS relation holds on the code and its modular flow descends to the corner via the corner-inverse $V\rho^{-1}V^\dagger$ (\texttt{finite\_KMS}, \texttt{encoded\_modConj\_corner}); and \textbf{dynamically faithful} — a \emph{supplied} field flow is preserved through the encoding, $\beta(\iota_V(A))=\iota_V(\alpha(A))$, with two-time correlators intact (\texttt{encoded\_heisenberg\_intertwining}, \texttt{dynamical\_twopoint\_preserved}). Finally the equiprobable typicality measure (P5) is \emph{isolated}: a permutation-invariant outcome law is forced uniform (\texttt{uniform\_of\_permInvariant}), so envariance fixes the measure uniquely. \textbf{Scope, enforced in the statements:} these are statements about the field's \emph{records and states}, not a construction of the field or its dynamics from microstates; capacity remains a \emph{constraint, not a generator} (it never derives the electron or photon); and they remove neither (P5) nor the \texttt{HolographicCapacityBound} postulate. The continuum real-time modular flow ($\rho^{it}$, Type III$_1$) and the Gaussian/Williamson entropy–area derivation are the cited research-grade frontiers, honestly checkpointed.

\textbf{Free Standard-Model field content and finite proto-spacetime, in the same corner (2026).} Two further axiom-free modules extend the corner construction along the two directions a reader naturally asks about — \emph{other quantum fields} and \emph{emergent spacetime} — while holding the same honest line: QIQT-H sits \textbf{on top of} quantum field theory (it \emph{accepts} a field algebra as input and bounds its records), it does \textbf{not} \emph{construct} the field, and capacity is a constraint, not a generator. (i) \texttt{QIQTH/FreeFieldCorner.lean} unifies the electron's CAR and the photon's truncated-CCR transport under one \textbf{graded bracket} $[x,y]_\varepsilon := xy + \varepsilon\,yx$ ($\varepsilon=+1$ fermionic, $\varepsilon=-1$ bosonic), proves it transports into the corner, $[\iota_V(x),\iota_V(y)]_\varepsilon = \iota_V([x,y]_\varepsilon)$ (\texttt{encode\_gradedBracket}), and instantiates it for the \emph{free} Standard-Model field content: quarks/leptons (multi-flavor CAR, $\{\iota_V(a(\alpha)),\iota_V(a^\dagger(\beta))\}=c(\alpha,\beta)\,P$), the $W/Z$/gluon vector bosons and the Higgs scalar (truncated bosonic, each carrying the \emph{explicit, surviving} truncation defect $P-N\,\iota_V(|N{-}1\rangle\langle N{-}1|)$), with the corresponding mode counts holographically bounded ($n\log 2 \le A/4\ell_P^2$ for $n$ fermionic modes, $G\log\binom{d{+}N}{N}\le A/4\ell_P^2$ for $G$ truncated-boson components). The capstone \texttt{sm\_free\_field\_in\_corner} packages it: any free SM sector fitting the microstate space has its algebra faithfully transported into the corner \emph{and} its records area-bounded ($S_{\mathrm{vN}}\le A/4\ell_P^2$). (ii) \texttt{QIQTH/EmergentSpacetime.lean} builds the \textbf{finite proto-spacetime} cores: a no-go guard that a finite-dimensional unitary conjugation cannot \emph{rescale} a nonzero operator (\texttt{finiteDim\_scaling\_forces\_zero} — so no exact finite Borchers dilation/boost, forcing the emergent geometry to be approximate); the machine-checked correction that \textbf{min-cut entanglement area is not a metric} (it violates the triangle inequality, \texttt{minCut\_area\_not\_metric}) together with a provably-metric replacement (\texttt{embedDist\_isPseudometric}); a finite graph Ryu–Takayanagi entropy skeleton with \textbf{purity $S(A)=S(A^c)$} and \textbf{subadditivity} (\texttt{cut\_compl}, \texttt{cut\_union\_le}); the conditional finite RT \emph{inequality} $S_{\mathrm{vN}}(\rho)\le \mathrm{cut}\,(S)$ when the bond dimension fits the cut area (\texttt{entropy\_le\_cut}); and an operational \textbf{causal preorder} from a supplied signalling relation (\texttt{reach\_refl}/\texttt{reach\_trans}, causal cones, mutual-reachability poset). \textbf{Scope, enforced throughout:} the field side is \emph{free-field content only} — interactions, non-abelian gauge dynamics, the Yang–Mills mass gap, confinement, chirality, and spontaneous symmetry breaking are cited frontiers (open mathematics), not in scope; the spacetime side delivers \emph{finite proto-geometry with explicit error bounds} — a background-independent $3{+}1$ Lorentzian manifold, and the continuum Borchers/half-sided-modular route, are the cited (open-physics) frontiers. The Jacobson/Bisognano–Wichmann/Sakharov material above \emph{assumes} a smooth spacetime and is therefore emergent \emph{gravity} (dynamics), not emergent \emph{spacetime} (the manifold) — a distinction these modules keep sharp.

\textbf{Emergent geometry from entanglement: the finite, background-general core (2026).} A further axiom-free extension imports — finitely, and \emph{off} anti-de Sitter — the genuinely transferable core of the AdS/CFT holographic dictionary, toward the emergent-\emph{spacetime} frontier (\texttt{QIQTH/EmergentSpacetime.lean}, "Track C"). Its honest centrepiece is the \textbf{entanglement $\to$ distance} reconstruction, which \emph{repairs} the earlier roadmap error that min-cut area is a distance (it violates the triangle inequality, \texttt{minCut\_area\_not\_metric}): the \textbf{weighted L¹/Crofton distance} $d(x,y)=\sum_i \omega_i\,|\chi_i(x)-\chi_i(y)|$ is a provable pseudometric, and a genuine finite metric exactly when the (supplied) probe family separates points (\texttt{weightedCutDist\_isFiniteMetric\_iff}), with a second, complementary \textbf{graph-geodesic} reconstruction (\texttt{graphDist\_isFiniteMetric}, via Mathlib's \texttt{SimpleGraph.dist}). On the area/entropy side it adds the finite \textbf{Ryu–Takayanagi inequality} (\texttt{flow\_weak\_duality} — the easy half of max-flow/min-cut) and \textbf{conditional-mutual-information screening} (\texttt{cmi\_nonneg\_of\_SSA}, \texttt{markov\_screening\_of\_locality}: a graph separator is an approximate Markov blanket). On the dynamics side it adds the finite \textbf{entanglement first law} $\Delta S\le\Delta\langle K\rangle$, and $\delta S\le\eta\,\delta A$ as its area instance (\texttt{deltaEntropy\_le\_eta\_deltaArea}). On the Lorentzian side it adds a \textbf{causal-order skeleton} — strict causal precedence, Alexandrov intervals $J^+(x)\cap J^-(y)$, and a capacity-as-volume proxy (\texttt{intervalCapacityVolume}, the Malament/Hawking–King–McCarthy route) — and a quantitative \textbf{approximate} no-go strengthening the exact one (\texttt{approx\_scaling\_gap\_mul\_norm\_le}: a finite isometry cannot even \emph{approximately} rescale a nonzero operator). All reconstructions are bundled in a provenance- and noncircularity-tagged \textbf{certificate} (\texttt{ApproxMetricCert}, \texttt{reconstruction\_certs\_are\_finite\_proto}). \textbf{Scope, enforced throughout:} these are \emph{finite proto-geometry objects with explicit error tags} — a metric/causal skeleton reconstructed from \emph{supplied} entanglement/connectivity data — NOT a background-independent $3{+}1$ Lorentzian manifold (the open frontier, which AdS/CFT itself realises only in its anti-de Sitter laboratory); a \emph{directed} causal order requires a supplied time orientation (a reversible unitary runtime provides none); min-cut stays the \emph{area} primitive, never a distance; and QIQT-H derives the holographic area law its \emph{own}, background-general way (the P4-MICRO/Sakharov route, valid off AdS) — from AdS/CFT it borrows only the emergent-geometry \emph{techniques}, not the boundary dictionary — with one further machine-checked \textbf{correspondence} (\texttt{QIQTH/HolographicBridge.lean}): under the granularity reframing's induced $G = 1/(N\Lambda_s^2)$, the AdS/CFT boundary Cardy microstate count coincides with QIQT-H's bulk capacity exponent $(A/4)\,N\Lambda_s^2$ (\texttt{btz\_cardy\_eq\_qiqth\_capacity}), so the two holographic bookkeepings \emph{agree} — still a correspondence, not the boundary dictionary or an independent cross-check.

\textbf{Macroscopic Definiteness (H2) — retired as a category error (2026).} Earlier formulations made a further, \emph{load-bearing} conjectural claim (the "Macroscopic Definiteness Conjecture", §7.6): that a regional content carrying two or more macroscopically distinct records has summed cost $I_\Sigma \approx K\cdot I_0$ exceeding $Q_R$, so that finite capacity \emph{itself} forbids the multi-record state and thereby forces a single outcome. \textbf{We retire this claim.} It conflates two different entropies. The capacity $Q_R$ bounds the \emph{von Neumann} entropy $S(\rho_R) = H(\{p_i\}) + \sum_i p_i\,S(\rho_{R,i})$ of the pre-selection state $\Phi$ (the constraint above), in which a second macroscopic branch raises only the \emph{branch-count} term $H(\{p_i\})$ — by $\sim\!\log K$, a handful of bits — while the enormous \emph{within-branch} term $\sum_i p_i\,S(\rho_{R,i})$ is essentially the same whether $K=1$ or $K=2$. A two-record regional state therefore does \textbf{not} overflow $Q_R$. The "$I_\Sigma \approx K\cdot I_0$" branch-summed cost is precisely the additive fiction that drove the false conclusion; the Lean development itself records that the branch-summed cost is \emph{not} bounded by — and is not — the entropy (\texttt{BranchLedger.branchSummed\_not\_bounded\_by\_Shannon}). More fundamentally, in a unitarily-evolving theory no \emph{kinematic} entropy inequality can \emph{select} an outcome: the branch weights $|c_k|^2$ are conserved (\texttt{NoConcentration}), and a capacity (cardinality) bound \emph{underdetermines} the record basis (\texttt{RealmSelection.capacity\_underdetermines\_realm} — "the area budget is not the metaselector"). \textbf{Single outcomes are the work of $\lambda$ (P1) — decoherence fixes the record basis, $\lambda$ selects one — not of the capacity bound.} What survives, and is load-bearing, is the \emph{kinematic} role of $Q_R$: it bounds $S(\rho_R)$ (the now-derived area floor, above) and feeds the field equations. Sections §4, §7.3, and §7.6 are to be read accordingly — their single-record conclusions stand via (P1) + decoherence, but every clause there asserting that capacity \emph{forbids} a multi-record state, or that a multi-record content is \emph{not instantiable for capacity reasons}, is \textbf{superseded} by this note. The genuinely open problems are therefore the \textbf{canonical IC measure} (P5) and $\lambda$'s \textbf{dynamical, Lorentz-covariant realization} — \emph{not} Macroscopic Definiteness, and \emph{not} the typicality$\to$Born step itself. That step \textbf{is} machine-checked and axiom-free: the \textbf{Zurek amplitude$\to$count bridge} $\|\psi_k\|^2/\|\psi\|^2 = \mathrm{count}/|I|$ holds by orthonormality alone (Pythagoras), so the uniform (equiprobable) measure over an equal-amplitude orthonormal fine-graining has outcome-marginal \emph{exactly} $|c_k|^2$ (\texttt{BornEquiprobable.born\_from\_equiprobability}), discharging the previously-posited multiplicity $\mathrm{count}=M\,w_k$. What survives is \textbf{not} the Born derivation but the \textbf{canonicity of that equiprobable measure} — envariance (equal-amplitude $\Rightarrow$ equiprobable) and refinement-additivity (which excludes the $\alpha$-family $w_k^{\alpha}$) — the single named Born-strength premise (P5), stated in its most honest, amplitude-free form in §11.4.

\textbf{Lorentz naturalness of finite capacity — an adversarially-tested open frontier (2026).} A dedicated stress-test program (\texttt{scripts/qg/}, \texttt{QIQTH/QG/}) confronts the sharpest objection to a finite-information foundation: does the finite-capacity postulate (FQ/P4) radiatively generate \emph{Lorentz violation}? The diagnostic is the Collins–Perez–Sudarsky–Urrutia–Vucetich (CPSUV) one-loop speed splitting $\Delta c^2 = Z_s/Z_t - 1$. The findings — machine-verified where Lean-formalizable, reproducible-numerical otherwise: \textbf{(i)} a \emph{naive} "finite capacity = local Lorentz-violating cutoff" is excluded — a sharp three-momentum cutoff radiatively generates an \emph{unsuppressed} dimension-4 violation $\Delta c^2 \to \tfrac43\,g^2/16\pi^2 \neq 0$, while a Lorentz-invariant regulator gives zero, and the escape condition is machine-checked to be the single scalar $\Delta c^2 = 0 \Leftrightarrow$ the matter kernel is Lorentz-scalar (\texttt{QIQTH/QG/WardSpeedSplitting.lean}); \textbf{(ii)} QIQT-H's \emph{actual} capacity is \textbf{not} such a cutoff — it bounds the distinguishable-record/entropy content per causal diamond ($\mathrm{card}\le e^{Q_D}$, $Q_D=A(\partial D)/4\ell_P^2$), with no field-momentum or modular-energy truncation — so the naive failure does \textbf{not} directly apply. \textbf{(iii) The honest dilemma, sharpened by an adversarial review.} Whether this constitutes a genuine \emph{escape} is \textbf{not established}, and the strong claim — \emph{literal finite per-region matter capacity together with exact Lorentz invariance} — is assessed at low confidence ($\approx$10–20\%). The crossed-product (Type II) reconciliation faces a fork: \textbf{either} matter remains ordinary covariant (Type III$_1$) field theory and "finite capacity" is only a finite \emph{renormalized entropy} (consistent, but then the finiteness does no matter-UV work, and the "finite information" content is finite \emph{entropy}, not a finite matter Hilbert space); \textbf{or} the finiteness is made literal for matter, colliding with structural facts — Type III$_1$ has no atoms or finite trace, and the non-compact Lorentz group has no non-trivial finite-dimensional unitary representations. Crucially, **this pressure is confined to the literal finite-\emph{matter} reading; it does not threaten the $(\Phi,\lambda)$ record-selection ontology or the holographic \emph{entropy} bound, both of which are entropy-level and Lorentz-safe.** Resolving which fork QIQT-H occupies — and restating "finite capacity" as finite \emph{entropy} rather than finite \emph{matter} if, as expected, that is the consistent horn — is active work (\texttt{FINITE\_MATTER\_OR\_ENTROPY\_PLAN.md}) and is the most consequential current open problem after $\lambda$'s law. The honest discipline is recorded in the repository: the earlier intermediate "escapes CPSUV" verdict was explicitly \textbf{retracted} under the same adversarial review.

\subsection{Realisations: minimal toy models, SymPy-verified}

The boundary-theory reading (§1.1a) is made concrete by two small, fully executable toy \emph{realisations} — generative models whose regional data, encoded into a finite microstate space, are checked symbolically (in SymPy) to satisfy QIQT-H's boundary conditions. They draw the exact line between what a \emph{finite} realisation gives and what the Einstein equations require.

\textbf{The minimal finite realisation} (\texttt{scripts/qiqth\_minimal\_realisation.py}) is a single fermionic CAR mode — a qubit code $C=\mathbb{C}^2$ — isometrically embedded, $V:\mathbb{C}^2\hookrightarrow\mathbb{C}^N$, into a finite microstate memory $\mathcal{H}_R=\mathbb{C}^N$. The script constructs the objects and verifies, exactly, every boundary condition (each tagged with its Lean theorem): the CAR relation $\{a,a^\dagger\}=1$ fits a finite sector exactly; the corner $P=VV^\dagger$ is an orthogonal projector and the encoding is unital \emph{onto it}, $\iota_V(1)=P\neq 1_{\mathcal{H}_R}$ (the audit tripwire); record statistics read back faithfully, $\operatorname{Tr}\big((V\rho V^\dagger)(VOV^\dagger)\big)=\operatorname{Tr}(\rho O)$; the CAR algebra transports to the corner, $\{\iota_V(a),\iota_V(a^\dagger)\}=P$; the equiprobable measure over an equal-amplitude orthonormal fine-graining reproduces the Born weights $|c_k|^2$; the area floor $S_{\mathrm{vN}}(\rho)\le\log\dim C\le\log\dim\mathcal{H}_R$ holds; and the photon no-go is respected — a truncated bosonic mode necessarily carries the explicit defect $[a,a^\dagger]=1-N\,|N{-}1\rangle\langle N{-}1|$. Three composition checks confirm consistency under combination: \textbf{(a) locality} — disjoint regions' observables commute and the encoding preserves it, $[\iota_V(O_A),\iota_V(O_B)]=\iota_V([O_A,O_B])=0$; \textbf{(b) a tensor-network RT toy} — a $3$-node graph whose cut is the \emph{area} primitive, with purity $S(A)=S(A^c)$, subadditivity, the RT inequality $S_{\mathrm{vN}}(\rho_A)\le\mathrm{cut}(A)$, and min-cut violating the triangle inequality ($\lambda_{02}=5>4=\lambda_{01}+\lambda_{12}$, so it is an area, never a distance); and \textbf{(c) the thermal half} — a finite modular/Gibbs state with a genuine temperature (finite KMS), the exact entanglement first law $\delta S=\delta\langle K\rangle$, and, with the modular Hamiltonian $\propto$ area, the entropy–area variation $\delta S=\eta\,\delta A$. This is as far as a \emph{finite} realisation reaches: it carries the kinematic \emph{and} thermal boundary conditions, but yields no Einstein equations.

**The minimal \emph{gravitational} realisation** (\texttt{scripts/qiqth\_gr\_realisation.py}) is the continuum object that \emph{does} reach GR: the free Klein–Gordon field on a pp-wave background, $ds^2=H(u,x,y)\,du^2+2\,du\,dv+dx^2+dy^2$ — matching the Lean capstone \texttt{qiqt\_gr\_ppwave\_showcase}. The script computes the geometry symbolically (Christoffel $\to$ Ricci $\to$ Einstein tensor: $R=0$, $G_{uu}=-\tfrac12\nabla^2_\perp H$, all other $G_{\mu\nu}=0$); verifies the contracted Bianchi identity $\nabla^\mu G_{\mu\nu}=0$ (conservation, which fixes $\Lambda$ and lets the local Clausius relation become a field equation); proves Jacobson's null-cone lemma (a symmetric $S$ with $S_{\mu\nu}k^\mu k^\nu=0$ for all null $k$ is $S=f\,g$); and assembles Jacobson's equation of state $\delta Q=T\,\delta S$ — with the QIQT-H inputs that \emph{are} theorems ($\delta S=\eta\,\delta A$ from \texttt{DifferentialAreaLaw}, the $1/4$ from \texttt{SakharovRatio}) and the Unruh temperature \emph{discharged} for the free field (\texttt{Fock.OneParticleBW}) — into the explicit Einstein equation $-\tfrac12\nabla^2_\perp H = 8\pi G\,(\partial_u\phi)^2$. So GR is a property of the \emph{continuum} realisation, which supplies the smooth metric and the Raychaudhuri area derivative $\delta A\leftrightarrow R_{kk}$ that no finite toy contains — conditional on the labelled residuals (matter equation of motion, (P4), the localization map, the reference-state identification, and the value of $G$). The two scripts together exhibit, executably, the central claim of §1.1a: **a realisation produces GR only if it also realises the modular/thermal \emph{and} smooth-background structure; the finite toys realise the boundary conditions and deliberately stop short of it.**

\subsection{What this paper does and does not claim}

\textbf{Does.} Introduces (FQ) in its rigorous algebraic form using Type II crossed-product algebras: the finite-entropy bound $S_{\rm ren}(\omega_\Psi) \le Q_R$ on the regional reduced state, with finite renormalized entropy realized via the Type II trace. Distinguishes formal vs per-run wave functions. Discusses how a single actual record per run arises — decoherence fixes the record basis and the non-dynamical selector $\lambda$ makes one record actual (single outcomes are the work of $\lambda$, \textbf{not} of the capacity bound; the earlier "capacity forbids a multi-record state" reading is retired as a category error, §1.1a). States four formal results. Identifies four explicit open problems (§11.4).

\textbf{Does not.} Provide a fully rigorous typicality theorem for Born weights. Modify the Schrödinger equation, posit Bohmian particles, MWI branches, modal value-rules, or any actuality primitive beyond the non-dynamical selector $\lambda$ of the $(\Phi,\lambda)$ ontology.

\subsection{Relation to existing programs}

The framework is positioned in relation to:

\begin{itemize}
  \item \textbf{Algebraic QFT} (Haag, Kastler 1964; Haag 1992): local algebras for bounded regions; we adopt this as the mathematical foundation.
  \item \textbf{Chandrasekaran-Penington-Witten Type II algebras} (CPW 2022; Witten 2022; Jensen-Sorce-Speranza 2023): the crossed-product construction giving Type II regional algebras with finite renormalized entropy; we use this as the rigorous home for (FQ).
  \item \textbf{Holographic bounds} (Bekenstein 1981; 't Hooft 1993; Susskind 1995; Bousso 2002; Banks 2025): the bound that (FQ) imposes; the entropy reading we adopt is supported by these works' information-theoretic motivations.
  \item \textbf{Decoherence} (Zurek 2003; Joos et al. 2003): einselection of pointer-basis records; we add the Type II resolution structure that makes below-threshold branches physically equivalent to nonexistent ones.
  \item \textbf{Ensemble interpretation} (Ballentine 1970): we agree the standard wave function is ensemble-descriptive; we supply the mathematical mechanism via the Type II regional structure.
  \item \textbf{Typicality programs} (statistical mechanics; Wallace 2012; Carroll-Sebens 2014; Dürr-Goldstein-Zanghì 1992 for Bohmian equivariance): we deploy a typicality argument for Born weights.
\end{itemize}

\subsection{Roadmap}

§2 reviews standard QM preliminaries and introduces the formal/per-run distinction. §3 introduces the algebraic QFT framework and the CPW Type II crossed-product construction. §4 states (FQ) in the Type II framework and discusses its motivation. §5 derives the finite regional entropy bound $S(\rho_R) \le Q_R$ as a consequence of the Type II structure. §6 develops the mechanism by which decoherence + the selector $\lambda$ + microscopic initial conditions produce single-record per-run wave functions. §7 states the four formal results. §8 discusses the holographic motivation. §9 sketches possible phenomenology. §10 compares with existing approaches. §11 concludes.

\medskip\hrule\medskip

\section{Formal Preliminaries}

\subsection{Standard QM axioms}

\textbf{(A1) State space.} Pure states are unit rays in a separable complex Hilbert space $\mathcal{H}$; mixed states are density operators in $\mathcal{B}(\mathcal{H})$.

\textbf{(A2) Closed-system evolution.} Unitary $U(t) = e^{-iHt/\hbar}$ with self-adjoint Hamiltonian $H$.

\textbf{(A3) Composite systems.} $\mathcal{H}_{AB} = \mathcal{H}_A \otimes \mathcal{H}_B$.

\textbf{(A4) Observables.} Self-adjoint operators or POVMs.

\textbf{(A5) Born rule.} $\Pr(\omega) = \mathrm{Tr}(\rho E_\omega)$.

\subsection{One wave function per run; subsystem vs universal description}

\textbf{The central commitment: there is exactly ONE wave function per run.} No formal/per-run distinction in the problematic sense of "the same preparation gives different per-run wave functions." Each actual run of an actual experiment corresponds to \emph{one} universal wave function evolving from \emph{one} set of actual initial conditions of the universe.

What varies across runs is \emph{the actual initial conditions themselves}: different particles emitted by the source, different microscopic states of the apparatus, different environmental configurations, different photon backgrounds, different thermal fluctuations. These differences are not hidden variables added to the framework, they are simply the actual physical differences between actual physical universes in different runs. Standard QM already accommodates this: the universal wave function is a function of all degrees of freedom in the universe, which trivially differ across runs.

What is sometimes loosely called the "formal wave function" in textbook accounts is the \emph{subsystem wave function}, the wave function of a subsystem (typically the particle of interest) treated as if it were isolated from the apparatus and environment. This subsystem wave function is the textbook abstraction $\alpha|0\rangle + \beta|1\rangle$. It is not the per-run physical state of the universe; it is a subsystem description that abstracts away the apparatus + environment.

\textbf{Definition (Subsystem wave function).} \emph{The subsystem wave function $|\psi\rangle_{\rm sub}$ for a subsystem $S$ of the universe is the wave function obtained by tracing the universal wave function over all degrees of freedom outside $S$ and (when $S$ is approximately uncorrelated with its complement) recovering an effectively pure state on $S$. For the prepared particle in the double-slit, $|\psi\rangle_{\rm sub} = \alpha|0\rangle + \beta|1\rangle$ is the standard textbook expression.}

\textbf{Definition (Per-run universal wave function).} \emph{The per-run universal wave function $|\Psi\rangle_{\rm run}$ is the actual physical wave function of the universe in a specific individual run, evolving unitarily under the universal Hamiltonian from the specific actual initial conditions of the universe in that run. Its physical content in any bounded region $R$ is given by its values on the regional algebra $\hat{\mathcal{A}}(R)$ (defined in §3), constrained by (FQ).}

The relationship across runs: there is \emph{one} per-run universal wave function per run. Different runs have \emph{different actual} per-run universal wave functions because the actual initial conditions of the universe differ. Born statistics across runs (§7.4) arise from the distribution of these actual initial conditions, not from any "alternative" wave functions for the same preparation. The "ensemble" is over actual different physical universes, not over hypothetical alternatives.

**The framework adds no second \emph{dynamical} substance beyond the wave function** (no Bohmian particles, no extra fields, no guidance law). The framework's contributions are:
\begin{enumerate}
  \item The recognition that the textbook subsystem wave function $|\psi\rangle_{\rm sub}$ is \emph{not} the relevant complete dynamical quantum state of a run; the per-run universal wave function $\Phi = |\Psi\rangle_{\rm run}$ is. (The complete per-run \emph{ontology}, however, is not $\Phi$ alone but the pair $(\Phi, \lambda)$; see the note below.) This is standard QM at the universal level.
  \item The (FQ) axiom limiting the physical information content of the universal wave function per region.
  \item A \emph{non-dynamical run-index} $\lambda$: a global, atemporal actuality fact about \emph{which} of $\Phi$'s macroscopic realizations the whole run is (fixed for the entire 4-dimensional history; \emph{not} past-localized "initial data," to avoid any setting-correlation / Conway–Kochen reading; see §6.9), that selects which single record is realized (§7.6).
\end{enumerate}

\textbf{A note on "ψ-monism" as used throughout.} The framework is \textbf{ψ-monist in the weak (dynamical) sense}: the wave function $\Phi$ is the only thing carrying a law of motion (exactly unitary), and no second \emph{dynamical} ontology is added. It is \textbf{not} ψ-monist in the strong sense that $\Phi$ alone is the complete per-run state, the single realized record per run is fixed by $\lambda$, a \emph{non-dynamical actuality fact} (a broad-sense beable, structurally analogous to the configuration in Bohmian mechanics or the actual history in modal interpretations, but with no guidance law and no back-reaction on $\Phi$; see §7.6). In the Spekkens–Harrigan taxonomy the framework is therefore \textbf{ψ-ontic, weak-ψ-monist, and formal-ψ-incomplete}: $\Phi$ is real and is the sole dynamical ontology, but the complete per-run description is the pair $(\Phi, \lambda)$. Every later use of "ψ-monist" in this paper is shorthand for this weak/dynamical sense.

\subsection{Empirical content}

Empirical content is exhausted by Born statistics across runs. Standard QM is empirically complete in this sense; interpretations differ on the per-run physical content.

\medskip\hrule\medskip

\section{Algebraic Framework: Type II Regional Algebras}

\subsection{Local algebras in QFT}

In algebraic QFT (Haag-Kastler), one associates to each spacetime region $R$ a von Neumann algebra $\mathcal{A}(R) \subset \mathcal{B}(\mathcal{H})$ of observables localized in $R$. The assignment $R \mapsto \mathcal{A}(R)$ satisfies isotony, locality, covariance, and (for the vacuum) the Reeh-Schlieder property.

A foundational result (Buchholz, Wichmann, Borchers, Longo): for any double cone or causal-diamond region $R$ in a generic QFT, $\mathcal{A}(R)$ is a \textbf{Type III$_1$ factor} in the Murray-von Neumann classification. This has consequences directly relevant to foundations:

\begin{itemize}
  \item No trace exists on $\mathcal{A}(R)$.
  \item No tensor factorization $\mathcal{H} = \mathcal{H}_R \otimes \mathcal{H}_{\bar R}$, the algebra of $R$ and its complement do not split the Hilbert space.
  \item No notion of "number of degrees of freedom" in $R$.
  \item Entanglement entropy $S(\rho_R)$ is UV-divergent.
  \item No pure states with respect to $\mathcal{A}(R)$.
\end{itemize}

The Type III$_1$ structure is the formal obstruction that has historically prevented the holographic bound from being applied directly to regional entanglement entropy: there's nothing well-defined to bound.

\subsection{The crossed-product construction (Chandrasekaran-Penington-Witten)}

Recent work (CPW 2022; Witten 2022; Jensen-Sorce-Speranza 2023; Chandrasekaran-Longo-Penington-Witten 2022) has shown that \textbf{gravitational dressing}, implemented via the crossed product with the modular flow, converts Type III$_1$ local algebras into \textbf{Type II algebras} with a semifinite trace structure and a renormalized entropy whose differences reproduce the generalized entropy.

The construction, in outline:

\begin{enumerate}
  \item Start with the Type III$_1$ local algebra $\mathcal{A}(R)$ and a cyclic-separating reference state $\Omega$ (e.g., the vacuum, or a state corresponding to a chosen observer's perspective in the semiclassical setting). The choice of $\Omega$ enters the construction and we discuss reference-state dependence in §3.4.
  \item Let $\Delta_\Omega$ be the modular operator and $\sigma_t^\Omega$ the corresponding modular automorphism of $\mathcal{A}(R)$.
  \item In the presence of gravity (perturbatively in $G_N$ or $1/N$), impose the constraint that physical observables commute with the boost generator of the modular flow. Equivalently: dress observables by the modular flow.
  \item The resulting algebra is the \textbf{crossed product} $\hat{\mathcal{A}}(R) := \mathcal{A}(R) \rtimes_{\sigma^\Omega} \mathbb{R}$, naturally represented on the enlarged Hilbert space $\mathcal{H} \otimes L^2(\mathbb{R})$.
\end{enumerate}

\textbf{Theorem (CPW / Witten).} \emph{The crossed-product algebra $\hat{\mathcal{A}}(R)$ is of Type II in the Murray-von Neumann classification. For general bounded regions, $\hat{\mathcal{A}}(R)$ is Type II$_\infty$ with semifinite trace (the identity having infinite trace). In special gravitational sectors with a finite phase-space volume (most prominently the de Sitter static patch, treated in CLPW 2022), the algebra is Type II$_1$ with finite trace $\tau(\mathbf{1}) < \infty$.}

The Type II structure provides what was missing in Type III$_1$:

\begin{itemize}
  \item A \textbf{trace or weight} $\tau$ exists on $\hat{\mathcal{A}}(R)$ (semifinite in the general Type II$_\infty$ case, finite in the special Type II$_1$ case).
  \item For the class of normal states $\omega$ on $\hat{\mathcal{A}}(R)$ with a well-defined trace-class density operator $\rho_\omega$, the \textbf{renormalized entropy} $S_{\rm ren}(\omega) := -\tau(\rho_\omega \log \rho_\omega)$ is defined.
  \item \emph{Entropy differences} $S_{\rm ren}(\omega) - S_{\rm ren}(\omega')$ between such states reproduce the generalized entropy differences $S_{\rm gen}(\omega) - S_{\rm gen}(\omega')$ with $S_{\rm gen} = A/(4\ell_P^2) + S_{\rm matter}$. The absolute value of $S_{\rm ren}$ depends on trace normalization and additive constants.
  \item "Number of effective degrees of freedom" is well-defined as a renormalized quantity, with the modular spectrum providing concrete structure.
\end{itemize}

\subsection{The regional algebra and the holographic bound, what CPW gives and what we postulate}

What the CPW/Witten framework rigorously provides:

\begin{itemize}
  \item A semifinite Type II algebra $\hat{\mathcal{A}}(R)$ for bounded regions in semiclassical gravity.
  \item A trace structure with respect to which renormalized entropy differences are well-defined.
  \item \emph{Entropy differences} matching the generalized entropy $A/(4\ell_P^2) + S_{\rm matter}$, with the area term arising naturally from the algebraic structure.
  \item In particular semiclassical gravitational sectors, finite area-controlled generalized entropy.
\end{itemize}

What CPW/Witten do \emph{not} establish in full generality:

\begin{itemize}
  \item A theorem of the form "for every bounded region $R$, $\max_{\omega} S_{\rm ren}(\omega) = A(\partial R)/(4\ell_P^2)$."
  \item A finite Hilbert dimension or finite number of distinguishable normal states.
  \item An \emph{absolute} bound $S(\rho_R) \le Q_R$ on the regional entropy (CPW match entropy \emph{differences}, not an absolute ceiling).
\end{itemize}

The holographic bound applied to the regional Type II algebra,

\[
S_{\rm ren}(\omega) \le Q_R := \frac{A(\partial R)}{4 \ell_P^2},
\]

is therefore not a theorem of CPW. \textbf{We postulate it as a foundational axiom} in this algebraic setting (it is the content of part (ii) of (FQ) below). The Type II framework provides the rigorous mathematical language in which the bound can be stated without UV divergences and in a basis-independent way; the bound itself is our axiom, motivated by the holographic principle of quantum gravity but not derived from CPW.

This sharp separation, CPW gives the algebraic language, we postulate the bound, is intentional. The Type II framework is borrowed; the foundational application to QM measurement is our contribution.

\subsection{State extension and reference-state dependence}

Two technical issues require comment.

\textbf{State extension.} The crossed-product algebra $\hat{\mathcal{A}}(R) = \mathcal{A}(R) \rtimes_{\sigma^\Omega} \mathbb{R}$ is naturally represented on $\mathcal{H} \otimes L^2(\mathbb{R})$, not on the original QFT Hilbert space $\mathcal{H}$. A per-run wave function $|\Psi\rangle \in \mathcal{H}$ does not automatically define expectation values on operators in $\hat{\mathcal{A}}(R)$ that act on the $L^2(\mathbb{R})$ factor (the modular-flow / "clock" degrees of freedom). Following CPW, we assume an explicit extension prescription: per-run universal wave functions are taken to live in the enlarged Hilbert space, $|\Psi\rangle_{\rm run} \in \mathcal{H} \otimes L^2(\mathbb{R})$, with the $L^2(\mathbb{R})$ factor encoding the dressing degrees of freedom required to make the gravitational constraint well-defined. The state $\omega_\Psi$ on $\hat{\mathcal{A}}(R)$ is then the standard vector state of $|\Psi\rangle_{\rm run}$ restricted to $\hat{\mathcal{A}}(R)$.

\textbf{Reference-state dependence.} The crossed product uses the modular flow $\sigma^\Omega$ of a reference state $\Omega$. Different reference states give different crossed-product algebras; however, the continuous-core construction is known to be state-independent up to isomorphism (Connes-Takesaki), and physical predictions (entropy differences, generalized-entropy comparisons between states) are reference-state-independent in the semiclassical regime where CPW operates. We adopt the same convention as CPW: a natural reference state is chosen (e.g., the vacuum or a global Bunch-Davies-like state), and physical content is read off via entropy differences and algebra-state equivalences that are reference-independent in the relevant sense. A fully reference-state-independent formulation is a technical refinement we defer.

\medskip\hrule\medskip

\section{The (FQ) Axiom in the Type II Framework}

\subsection{Statement}

\textbf{Axiom (FQ), finite regional entropy budget.} \emph{For every per-run wave function $|\Psi\rangle_{\rm run}$ of the universe (regarded as an element of the enlarged Hilbert space $\mathcal{H} \otimes L^2(\mathbb{R})$, cf. §3.4) and every bounded region $R$ of space:}

\textbf{Local von Neumann net and normality.} Work in the GNS representation $(\mathcal{H}_\sigma, \pi_\sigma, \Omega_\sigma)$ of the reference state $\sigma$ (§7.6). For each bounded region $R$, let

\[
\mathcal{M}(R) := \pi_\sigma(\mathcal{A}_{\rm loc}(R))''
\]

be the local von Neumann algebra. The hatted algebra $\hat{\mathcal{A}}(R)$ denotes the corresponding semifinite Type II crossed-product / core algebra, represented in the standard crossed-product representation induced by $\mathcal{M}(R)$ and $\sigma_R$. (If no crossed-product refinement is being used, set $\hat{\mathcal{A}}(R) = \mathcal{M}(R)$.) The net is assumed isotonic and local: $R' \subset R \Rightarrow \hat{\mathcal{A}}(R') \subset \hat{\mathcal{A}}(R)$, and $R_1 \perp R_2 \Rightarrow [\hat{\mathcal{A}}(R_1), \hat{\mathcal{A}}(R_2)] = 0$.

\emph{(i) \textbf{Local normality.} The physical content of $|\Psi\rangle_{\rm run}$ in $R$ is given by the regional restriction $\omega_\Psi|_R$. This restriction must be a normal state on $\hat{\mathcal{A}}(R)$, equivalently, a positive unital $\sigma$-weakly continuous linear functional, $\omega_\Psi|_R \in \hat{\mathcal{A}}(R)_*^+$ with $\|\omega_\Psi|_R\| = 1$.}

\emph{(ii) \textbf{Finite holographic modular budget.} Define the regional information functional as the Araki relative entropy with respect to the canonical sector reference state $\sigma_R$ (§7.6):}

\[
\chi_R(\omega_\Psi|_R) \;:=\; S^{\hat{\mathcal{A}}(R)}_{\rm Araki}(\omega_\Psi|_R \,\|\, \sigma_R).
\]

\emph{The (FQ) postulate (ii) is the modular-local bound:}

\[
\chi_R(\omega_\Psi|_R) \;\le\; Q_R \;\equiv\; C(R) \;:=\; \frac{A(\partial R)}{4\ell_P^2}.
\]

\emph{(The symbols $Q_R$ and $C(R)$ are used interchangeably throughout for the regional holographic capacity. We retain both because $Q_R$ emphasizes "information capacity" while $C(R)$ emphasizes "modular-cost ceiling"; the underlying quantity is identical.)}

\emph{Notational caveat: we will also write this functional as $S_{\rm ren}(\omega_\Psi)$ for emphasis on "modular-local information content", and (FQ)(ii) then reads $S_{\rm ren}(\omega_\Psi) \le Q_R$. This $S_{\rm ren}$ is \textbf{not} the same functional as the CPW/Witten "renormalized entropy" of the state on the Type II crossed-product algebra, they are distinct functionals related by the modular identity}

\[
\chi_R(\omega) \;=\; \Delta_\omega \langle K_R^\sigma \rangle - \Delta_\omega S_R^{\rm CPW}
\]

\emph{where $K_R^\sigma$ is the modular Hamiltonian of $\sigma_R$ and $S_R^{\rm CPW}$ is the CPW renormalized entropy. \textbf{The (FQ) bound is on the relative-entropy functional $\chi_R$, not on the CPW renormalized entropy.} The two functionals can disagree substantially (a state may satisfy one bound and violate the other; see the formal verification module \texttt{lean/mathlib/QIQTH/EntropyBridge.lean} for a classical counterexample making this concrete).}

\textbf{Machine-checked modular substrate (companion formalization).} A companion Lean 4/Mathlib formalization (§11.4b; project repository \texttt{lean/mathlib/QIQTH/}) independently verifies the bounded modular-theoretic and coherent-state relative-entropy \emph{calculus} on which $\chi_R$ rests, for the free-field coherent-state case: the finite Araki/Umegaki convention lock $S_{\rm Araki}(\rho\,\|\,\sigma)=\operatorname{tr}\rho(\log\rho-\log\sigma)$; the bounded Rieffel-Van Daele standard-subspace Tomita-Takesaki construction (modular conjugation $J$, the relation $JRJ=2-R$, the continuum modular group $\Delta^{it}$ and its strong continuity); the Casini-Grillo-Pontello one-particle relative entropy as a bounded scalar spectral integral together with its \emph{positivity} $S(\xi)\ge 0$ for localized $\xi$; the free-field (Fock) modular flow $\Gamma(\Delta^{it})$ with $\sigma_t(W(u))=W(\Delta^{it}u)$ and the vacuum as the modular state; the coherent-state relative modular operator $\Delta^{it}_{W(f)\Omega\,|\,\Omega}=W(f)\,\Gamma(\Delta^{it})\,W(f)^*$ (valid for $W(f)$ in the local algebra) with its Connes cocycle $[D\omega_{W(f)\Omega}:D\omega_\Omega]_t=W(f)W(-\Delta^{it}f)$ and cocycle chain rule; and the \emph{entropy reduction} $S_{\rm Araki}(\omega_{W(f)\Omega}\,\|\,\omega_\Omega)=S_{\rm CGP}(f)$ identifying the coherent-state Araki relative entropy with the one-particle CGP value. The development carries no \texttt{sorry} and, as reported by \texttt{\#print axioms}, depends only on the standard classical foundations of Lean/Mathlib (\texttt{propext}, \texttt{Classical.choice}, \texttt{Quot.sound}); it introduces no further axioms.

\emph{Scope.} What is verified is the modular/relative-entropy \emph{calculus} motivating $\chi_R$, in the finite and free-field coherent sectors, \textbf{not} the dressed Type II regional content map itself, the CPW/Witten crossed product, the holographic capacity axiom (FQ), the Donald/Fano argument, or the \textbf{canonical IC measure that grounds Born} (P5), all of which remain explicit physical assumptions or open problems of this paper (the typicality$\to$Born \emph{step} itself, by contrast, is machine-checked and axiom-free — the Zurek amplitude$\to$count bridge, §1.1a / §11.4 — so the residue is the \emph{canonicity} of the equiprobable measure, not the Born derivation). (The \emph{Macroscopic Definiteness Conjecture} — that capacity forbids a $\ge 2$-record state — is \textbf{retired} as a category error, §1.1a.) The single-record results (Theorem 4) rest on $\lambda$-selection (P1) + decoherence, not on capacity. Their \emph{dynamical, Lorentz-covariant} realization (that unitary measurement evolution produces the decoherent record structure $\lambda$ selects from) and the \textbf{canonical IC measure (P5)} that grounds Born — \emph{not} the typicality$\to$Born derivation, which the formalization does close — are the load-bearing gaps the formalization does not close.

\emph{(iii) \textbf{Superseded reading (retained for the record).} An earlier version of (FQ) added a third clause reading the bound as a finite physical specification precision on the amplitudes of $|\Psi\rangle_{\rm run}$ in $R$ (a finite-resolution-of-amplitudes postulate). That amplitude-precision reading is \textbf{not load-bearing} and is superseded (§1.1a). The load-bearing content of (FQ) is part (ii): the von Neumann / Araki entropy bound $S_{\rm ren}(\omega_\Psi) \le Q_R$, i.e. $Q_R$ bounds the entropy $S(\rho_R)$ of the pre-selection state $\Phi$ in $R$ (equivalently, its effective modular rank $e^{S(\rho_R)}$). It is not a bound on a finite matter Hilbert-space dimension and not a finite-precision reading of amplitudes.}

Part (i) defines the mathematical content of "wave function in $R$": it is the state on the Type II regional algebra introduced by CPW/Witten. Part (ii) is the holographic information bound on regional renormalized entropy, postulated (not derived from CPW): it is the entropy bound $S_{\rm ren}(\omega_\Psi) \le Q_R$ on $S(\rho_R)$ of $\Phi$. Part (iii), the finite-precision-of-amplitudes clause, is the superseded reading noted above and is retained only for the intellectual record.

The entropy reading is a precise statement of what Witten's algebraic framework supports: a bound on the regional renormalized entropy. It, together with the $(\Phi,\lambda)$ ontology (P1), is our foundational content. The relationship to Witten's framework is made explicit in §4.3.

\textbf{Important: (FQ) is a modular-local kinematic constraint, not a single global budget.} Each bounded region $R$ has its own capacity $Q_R$ set by \emph{its} boundary area; the bound is imposed independently for each $R$ as a constraint on the algebra-state pair $(\hat{\mathcal{A}}(R), \omega_\Psi|_R)$. There is no single shared information budget over all spacetime, and no joint-region cutoff is applied to the spacelike-combined algebra $\hat{\mathcal{A}}(D_A) \vee \hat{\mathcal{A}}(D_B)$.

\textbf{Spacelike combination by meet of local predicates.} For spacelike-separated regions $D_A, D_B$:

\[
\mathrm{Adm}(D_A \cup D_B) \;=\; \mathrm{Adm}(D_A) \wedge \mathrm{Adm}(D_B).
\]

That is: a state is admissible on the spacelike pair iff its restriction to \emph{each} local algebra is admissible. The joint algebra $\hat{\mathcal{A}}(D_A) \vee \hat{\mathcal{A}}(D_B)$ may remain vacuum-entangled and Type III non-factorizing, the admissibility predicate imposes nothing on the joint structure beyond the two local restrictions. The framework's admissibility is therefore a \textbf{local subfunctor} of the AQFT state functor.

This modular-local formulation makes no-signaling a \emph{theorem}, not an extra axiom: it follows automatically from AQFT microcausality (§7.7 below). The naive alternative, imposing the bound on the joint spacelike algebra as well, would produce operational signaling, because the vacuum is not a product state across spacelike-separated boundaries.

\subsection{Motivation: the entropy reading of the bound}

We adopt (FQ) as a foundational axiom. The motivation rests on:

\begin{itemize}
  \item \textbf{Bekenstein's original information-theoretic argument.} The original bound was derived as a bound on the information needed to specify any physical system, not specifically as an entanglement-entropy inequality on a reduced state.
  \item \textbf{'t Hooft's dimensional reduction.} All degrees of freedom in a region are encoded in boundary data; this is a statement about \emph{total} physical content.
  \item \textbf{Susskind's holographic principle.} Physics in a region is holographically described by boundary data of bounded information content. The boundary holds at most $A/(4\ell_P^2)$ bits, bounds \emph{all} the information to specify the physics, not just one measure of it.
  \item \textbf{Banks 2025.} Finite entropy in quantum gravity implies finite Hilbert-space dimension for bounded subsystems. (We invoke this only as motivation for a finite \emph{entropy} budget: QIQT-H's (P4) is the entropy bound $S(\rho_R) \le Q_R$ on the pre-selection $\Phi$, \textbf{not} an endorsement of a finite matter Hilbert-space dimension.)
  \item \textbf{CPW 2022 and successors.} The Type II crossed-product structure provides the rigorous algebraic realization compatible with the entropy reading.
\end{itemize}

The entropy reading takes the bound seriously as a physical limit on the regional entropy $S(\rho_R)$ of the pre-selection state $\Phi$: the bound caps the von Neumann / renormalized entropy (equivalently the effective modular rank) of the regional content, region by region. It is not a bound on amplitude precision and not a finite matter Hilbert-space dimension; the wave function retains continuous amplitudes.

\subsection{What Witten gives us and what we postulate on top}

The relationship between our framework and the Witten/CPW algebraic infrastructure is precisely the following.

\textbf{What Witten / CPW provide (mathematical infrastructure):}
\begin{itemize}
  \item The Type III$_1$ classification of local QFT algebras.
  \item The crossed-product construction $\hat{\mathcal{A}}(R) = \mathcal{A}(R) \rtimes_{\sigma^\Omega} \mathbb{R}$ giving Type II algebras.
  \item A semifinite (or finite, in Type II$_1$) trace $\tau$ on $\hat{\mathcal{A}}(R)$.
  \item A well-defined renormalized entropy $S_{\rm ren}(\rho) = -\tau(\rho \log \rho)$ for suitable states.
  \item Matching of entropy \emph{differences} with generalized entropy differences $\Delta S_{\rm gen}$.
  \item Observer-frame degrees of freedom built in via the $L^2(\mathbb{R})$ clock factor.
\end{itemize}

\textbf{What we postulate as a foundational interpretation on top of this infrastructure:}
\begin{enumerate}
  \item \emph{(FQ) part (i)}: The mathematical content of "regional wave function content" is the state on the CPW Type II algebra.
  \item \emph{(FQ) part (ii)}: The holographic bound $S_{\rm ren} \le Q_R$ holds as an absolute bound, not just as a matching of differences.
  \item \emph{The $(\Phi,\lambda)$ ontology (P1)}: the regional content is the pre-selection state $\Phi$ (whose entropy $S(\rho_R)$ the bound of part (ii) constrains) together with the non-dynamical selector $\lambda$ that marks one decohered record actual per run. (An earlier version added an amplitude-precision clause as part (iii); that finite-resolution reading is superseded, §1.1a.)
\end{enumerate}

Our distinctive contribution is the pairing of the entropy bound (part (ii)) with the $(\Phi,\lambda)$ ontology (P1), \emph{not} a precision floor on amplitudes. Witten's theorem supplies the continuous normal-state space and the exact algebraic equivalence; on top of it we assert (a) that $Q_R$ is an absolute bound on the regional entropy $S(\rho_R)$ of $\Phi$, and (b) that a non-dynamical actuality selector $\lambda$ fixes the single realized record. The Type II algebra is the mathematical scaffold; the entropy bound and the selector $\lambda$ are the physics laid on top.

We are not contradicting Witten. We are extending the foundational reading of the bound beyond what Witten directly proves. Witten's theorem is about the rigorous mathematical home for finite renormalized entropy on bulk regions; our postulate is about how this entropy bound translates into physical precision on the wave function as a physical state of spacetime.

The two readings can be compared:

\begin{center}
\small
\begin{tabularx}{\linewidth}{|Y|Y|Y|}
\hline
\textbf{} & \textbf{Narrow algebraic reading (Witten)} & \textbf{Entropy reading (ours)} \\
\hline
What the bound limits & Renormalized entropy on Type II algebra & The entropy $S(\rho_R)$ of the pre-selection $\Phi$ in $R$ (its effective modular rank $e^{S(\rho_R)}$) \\
\hline
Equivalence relation on states & Exact equality on algebra & Exact algebraic-state equality (no separate precision floor) \\
\hline
Near-zero amplitudes & Distinct from zero (continuous state space) & Distinct from zero; the bound limits $S(\rho_R)$, not amplitude precision \\
\hline
Single-record per run & Requires separate Concentration Conjecture & Decoherence + the selector $\lambda$ (not the capacity bound) \\
\hline
Status w.r.t. Witten's theorem & Direct mathematical reading & Stronger interpretive postulate; compatible but not derived \\
\hline
\end{tabularx}
\end{center}

Whether the absolute entropy reading is the correct foundational reading of the bound is a substantive philosophical / physical question, not a mathematical theorem. Our position is that reading $Q_R$ as an absolute bound on the regional entropy $S(\rho_R)$ of $\Phi$ (not merely a matching of entropy differences) is the natural reading of "Bekenstein-Bousso as a Lorentz-invariant information limit on bounded regions of spacetime".

\medskip\hrule\medskip

\section{Finite Regional Entropy: the $Q_R$ Bound on $S(\rho_R)$}

\subsection{The entropy consequence of the bound}

Under (FQ), the load-bearing consequence in $R$ is an \emph{entropy} bound: the pre-selection state $\Phi$ has regional von Neumann / renormalized entropy bounded by the holographic capacity,

\[
S_{\rm vN}(\rho_R) \;\le\; Q_R \;=\; \frac{A(\partial R)}{4\ell_P^2},
\]

equivalently, the effective modular rank of $\rho_R$ is at most $e^{S(\rho_R)} \le e^{Q_R}$. This is a bound on the entropy of $\Phi$ in $R$, not a finite-resolution reading of amplitudes and not a finite matter Hilbert-space dimension.

\emph{Superseded reading (retained for the record).} An earlier version drew instead a finite \emph{amplitude-resolution} consequence, that amplitudes are physically specified only up to a floor $\epsilon(R)$. This amplitude-precision reading is \textbf{not load-bearing} and is superseded (§1.1a); it is recorded here only for the intellectual history. In that reading one posited:

\textbf{Definition (Physical resolution floor) [superseded].} \emph{Let $\epsilon(R)$ denote the smallest amplitude difference physically distinguishable in region $R$ under (FQ). Two physical wave functions in $R$ whose physical instantiations differ by less than $\epsilon(R)$ in their content on $\hat{\mathcal{A}}(R)$ are the same physical state.}

\textbf{Lemma 1 (Resolution implies indistinguishability of near-extremes) [superseded].} \emph{Under the superseded amplitude-resolution reading, for any per-run wave function $|\Psi\rangle_{\rm run}$ in $R$:}
\begin{itemize}
  \item \emph{An amplitude $c$ with $|c| < \epsilon(R)$ in the algebraic decomposition of $\omega_\Psi$ on $\hat{\mathcal{A}}(R)$ corresponds to a physical state indistinguishable from amplitude exactly $0$ on that component.}
  \item \emph{An amplitude $c$ with $|c|^2 > 1 - \epsilon(R)$ for a single component is physically indistinguishable from that component having amplitude exactly $1$.}
\end{itemize}

The load-bearing statement of this section is the entropy bound $S_{\rm vN}(\rho_R) \le Q_R$ above; the $\epsilon(R)$ floor is retained only as superseded history. Where later sections invoke "the (FQ) resolution-equivalence of §5.1", this refers to the subsidiary role of erasing the operationally-negligible residual coherence, not to any load-bearing amplitude-precision claim.

\textbf{Stability requirement (recoherence).} For the resolution-$\epsilon$ relation to be a well-defined identity of physical states, it must be \emph{stable} under admissible subsequent evolution: two regional contents identified as the same physical state at resolution $\epsilon(R)$ must not be separable by later dynamics into distinct macroscopic records. In the measurement/decoherence setting this holds because the relevant sub-$\epsilon$ content is the residual record-coherence of an already-decohered, effectively irreversible interaction with $\sim 10^{20}$ environmental degrees of freedom; recoherence of distinct macroscopic records is dynamically excluded on any physical timescale. The general characterization of which contexts guarantee this stability is recorded as a subsidiary part of Open Problem 3.

\subsection{The Mandelbrot rendering analogy}

\emph{Superseded illustration.} The following analogy illustrates the retired finite-amplitude-resolution reading (§1.1a); it is retained for the intellectual record and is not part of the load-bearing entropy framing (§5.1).

The Mandelbrot set has unbounded fractal detail as a mathematical object. Any physical rendering, pixels, etched plates, printed images, has finite resolution determined by the physical substrate. Below pixel scale, mathematically distinct points are the \emph{same physical pixel}. The pixel does not partially exist or fractionally exist; it is one specific color.

The wave function, on the present view, is similar:
\begin{itemize}
  \item As an abstract Hilbert-space vector it can have arbitrary continuous amplitudes.
  \item As a \emph{physical} state of spacetime, it is rendered with finite resolution determined by the holographic bound on the region.
  \item Below resolution, fractional amplitudes are not physically realized; they are physically equivalent to exact $0$ or exact $1$.
\end{itemize}

This is the \emph{physical-instantiation reading} of the wave function, distinct from the abstract Hilbert-space reading. The latter treats amplitudes as freely specifiable continuous parameters; the former takes seriously the requirement that those amplitudes be physically encoded in the regional substrate, which has finite information capacity.

\subsection{Defense against the qubit objection}

\emph{Superseded subsection.} This defense addresses an objection to the retired finite-precision-on-amplitudes reading (§1.1a). Under the load-bearing entropy reading (§5.1) there is \textbf{no} amplitude-precision claim to defend: $Q_R$ bounds the regional entropy $S(\rho_R)$ of $\Phi$, and a qubit's continuous Bloch-sphere amplitudes are untouched. The material below is retained for the intellectual record only.

A standard objection to a finite-precision-on-amplitudes reading of the holographic bound runs as follows: "Consider a qubit. It has Hilbert dimension 2 and entropy bound $\log 2$. But its pure states form a continuum, the Bloch sphere, with arbitrary $\alpha, \beta$. So finite entropy does not imply finite amplitude precision."

This objection treats the qubit as an abstract Hilbert-space entity. The framework's response distinguishes the mathematical idealization from the physical instantiation. \textbf{The Bloch sphere stands in the same relation to a physical qubit as Euclidean geometry stands to a physical measuring rod}: an indispensable, extraordinarily accurate mathematical idealization, but not a literal description of what the physical object instantiates with infinite precision. No physical measuring rod realizes exact Euclidean points; no physical qubit realizes exact Bloch-sphere amplitudes. The continuum is the mathematical scaffolding; the physical substrate has finite information capacity.

A physical qubit is always implemented in a substrate, an ion, an atom, a photon's polarization, a superconducting circuit, a nuclear spin, a quantum dot, occupying some bounded spatial region. The amplitudes $\alpha, \beta$ must be physically encoded in that substrate's quantum state. The holographic bound on the region containing the qubit gives $Q_R$ bits of physical information capacity. The amplitudes can be specified to a precision corresponding to this capacity.

For a qubit in a 1m region with $Q_R \sim 10^{70}$ bits, amplitude precision is astronomically fine, the Bloch sphere is effectively continuous for all practical purposes. \textbf{This is precisely why standard QM works at lab scales: the framework reproduces ordinary continuous-qubit phenomenology because the physical precision floor is far below experimental resolution.} For a qubit in a Planck-scale region, $Q_R$ is order unity and the Bloch sphere is genuinely discretized.

Quantum error correction does not evade this. A logical qubit encoded in many physical qubits has a larger substrate and therefore larger information capacity, but still finite. The continuum is never physically recovered; it is only postponed into a bigger idealization.

For \emph{macroscopic} systems, where each macroscopic record component, considered as the full quantum state of the apparatus+environment region, uses a substantial fraction of the regional capacity, and superpositions of multiple records require encoding all of them simultaneously, the operationally relevant precision for amplitudes between records is finite. This is the regime where the precision floor has structural foundational consequences (Theorem 4).

The qubit objection therefore identifies the \emph{limit where the framework reproduces standard QM} (laboratory qubit regime); it does not show that the framework is incoherent in the regime where it bites (macroscopic record superpositions). The Bloch sphere is recovered to astronomical precision in lab regimes; finite precision matters foundationally only at the macroscopic-record scale, where the measurement problem actually lives.

\subsection{Relationship to the algebraic indistinguishability}

The Witten/CPW algebraic framework gives a basis-independent \emph{exact} equivalence relation on regional states: $\omega_\Psi = \omega_\Phi$ as states on $\hat{\mathcal{A}}(R)$. This exact algebraic-state equality is the operative equivalence relation on physical regional states. The (FQ) content is a \emph{separate} statement, an entropy bound $S(\rho_R) \le Q_R$ on that state, not a coarsening of the equivalence relation. (An earlier version added an $\epsilon$-coarsened "physical-resolution" equivalence $\omega_\Psi \approx_\epsilon \omega_\Phi$ as a postulate; that amplitude-precision reading is superseded, §1.1a, and no longer figures as an equivalence relation.)

Algebraic equivalence (Witten) is what is established by the Type II algebra structure; it is basis-independent and intrinsic to the regional geometry, and it is the relation relevant to the measurement problem. The (FQ) entropy bound is imposed on the state, not on the equivalence.

The macroscopic record structure enters naturally: macroscopically distinct records $\{|s_i\rangle|A_i\rangle|E_i\rangle\}$ are those distinguished by the environmental algebra $\hat{\mathcal{A}}(E)$ (the algebraic formulation of einselection). After decoherence these records are non-interfering, and $\lambda$ (P1) selects the one that is actual.

\medskip\hrule\medskip

\section{The Mechanism: Decoherence and $\lambda$-Selection (the $Q_R$ Entropy Bound as Arena)}

\subsection{Decoherence in algebraic terms}

In the algebraic framework, decoherence corresponds to the dynamical suppression of off-diagonal expectation values $\omega_\Psi(O_{ij})$ for operators $O_{ij}$ connecting macroscopically distinct record components. Standard decoherence theory (Zurek, Joos, Hartle) gives the dynamical mechanism for the suppression. The algebraic formulation makes this precise: after decoherence, $\omega_\Psi$ on the apparatus + environment algebra approaches a classical-record state, a convex combination $\sum_i p_i \omega_i$ of mutually decohered records.

This is the standard decoherence result, restated algebraically. It does not, by itself, select a single outcome.

\subsection{What decoherence actually does: exponentially-classical mixtures}

A common but misleading reading treats the post-measurement state as a "naive superposition" of macroscopic records, say, "$70\%$ amplitude at spot A, $30\%$ at spot B sitting there as accessible amplitudes." This picture treats the system as if its amplitudes could be directly read off; the puzzle then becomes "how do $70/30$ amplitudes become $100/0$ in a specific run?"

This picture is wrong. We never have direct access to "the amplitudes of the wave function." We have access only to \emph{macroscopic records} on detectors, screens, and observers, which are themselves complex quantum systems with $\sim 10^{20}$–$10^{25}$ degrees of freedom, entangled with the system being measured.

What decoherence actually does is the following.

After interaction with the apparatus and environment, the joint state of system + apparatus + environment is

\[
|\Psi\rangle = \sum_k c_k \,|s_k\rangle_S \otimes |A_k\rangle_A \otimes |E_k\rangle_E,
\]

where $\{|s_k\rangle\}$ are macroscopically distinguishable system states, $\{|A_k\rangle\}$ are macroscopically distinct apparatus configurations correlated with each $|s_k\rangle$, and $\{|E_k\rangle\}$ are environmental record states.

Decoherence theory (Zurek 2003; Joos et al. 2003) establishes that for any pair of macroscopically distinct records:

\[
\langle E_j | E_k \rangle \sim \exp(-\Gamma t) \to \exp(-N)
\]

where $\Gamma$ is the decoherence rate and $N \sim 10^{20}$ for typical macroscopic systems on any reasonable timescale. The cross-overlap is \textbf{exponentially small}, by factors like $e^{-10^{20}}$.

Translating to the regional algebra-state on $\hat{\mathcal{A}}(R)$ for $R$ containing the apparatus + environment: the induced state is exponentially close to a \emph{classical-looking mixture}:

\[
\omega_\Psi^{R} = \sum_k |c_k|^2 \omega_k^R + O(\exp(-N)),
\]

where $\omega_k^R$ are the macroscopic record states on $\hat{\mathcal{A}}(R)$ (defined more precisely in §5) and the off-diagonal coherence terms are exponentially suppressed.

This is what decoherence delivers. The diagonal weights $|c_k|^2$ remain. \textbf{The off-diagonal coherence terms do not vanish exactly under standard unitary evolution; they become exponentially small but mathematically nonzero.}

\subsection{The role of decoherence: rendering the records non-interfering}

The off-diagonal coherence terms after decoherence are exponentially small but mathematically nonzero. In standard QM, this is the source of the measurement problem: even if the coherence terms are $\sim e^{-10^{20}}$, they are still nonzero, and one is forced to either treat all branches as real (MWI) or invoke a separate selection mechanism (collapse, modal value-rule).

\textbf{Decoherence drives the records to be non-interfering.} On any physical timescale the off-diagonal coherence between macroscopically distinct records is suppressed like $e^{-10^{20}}$; the records become, for all practical purposes, mutually exclusive alternatives. The capacity bound (FQ) is \textbf{not} an amplitude eraser: it bounds the von Neumann entropy $S(\rho_R)$ of the pre-selection $\Phi$ (§1.1a), it does not set exponentially-small amplitudes to exact zero. What converts "one of these records is observed" into "exactly one record is actual" is the non-dynamical selector $\lambda$ (P1), applied to the decohered (non-interfering) record set — not the entropy bound.

For a typical macroscopic detection scenario the off-diagonal coherence is driven to $\sim e^{-10^{20}}$ dynamically, and recoherence of the distinct macroscopic records is excluded on any physical timescale; the residual sub-threshold coherence has no operational consequence.

Therefore: after decoherence, the regional state on $\hat{\mathcal{A}}(R)$ is, for all operational purposes, the classical-looking mixture

\[
\omega_\Psi^{R} = \sum_k |c_k|^2 \omega_k^R,
\]

with the off-diagonal coherence terms dynamically suppressed to negligibility. The macroscopic records are effectively exclusive in $R$.

This is the cooperative structure of the mechanism:
\begin{itemize}
  \item \textbf{Decoherence} drives off-diagonal coherence to exponentially small values dynamically, so the records no longer interfere (standard QM)
  \item \textbf{The selector $\lambda$ (P1)} marks exactly one of the non-interfering records actual; the capacity bound (FQ) is the kinematic \emph{arena} (it bounds $S(\rho_R)$), not the agent that removes coherence
\end{itemize}

The result: the operational regional state on $\hat{\mathcal{A}}(R)$ is a classical mixture over non-interfering macroscopic records, of which $\lambda$ selects one as actual.

\subsection{Per-run branch selection via microscopic initial conditions}

The classical mixture $\sum_k |c_k|^2 \omega_k^R$ is \emph{not} a single record. It is a weighted list of records carrying the Born weights $|c_k|^2$, weights that the framework reads as the \emph{empirical frequency pattern across runs} (§7.4), not as a fundamental per-run probability (QIQT-H has no fundamental probabilities).

Two distinct questions must be separated. \emph{Why is the resident content a single record at all (rather than a multi-record mixture)?} and \emph{which} record is it in this run? The first is answered by the $(\Phi,\lambda)$ ontology (P1): once decoherence renders the records non-interfering, the non-dynamical selector $\lambda$ marks exactly one of them actual. (An earlier version answered it by a \emph{finite-information restriction} — that a $\ge 2$-record content exceeds $Q_R$ and is "not instantiable"; that Macroscopic-Definiteness reading is \textbf{retired} as a category error, §1.1a: a second branch raises $S(\rho_R)$ by only $\sim\!\log K$, far below $Q_R$, so capacity does not forbid it.) The second is answered by the run's actual microscopic initial conditions:

**Per-run, the universe carries one specific record of the decohered regional content. \emph{Which} record is the per-run microscopic initial conditions of the apparatus + environment — this \emph{is} the selector $\lambda$ (P1), propagated through the unitary dynamics from these specific microstates. Decoherence makes the records non-interfering so that a definite record can be actual; $\lambda$ makes exactly one of them the lived one. (Earlier drafts attributed the single-ness instead to a "finite-information restriction" of capacity; that reading is \textbf{retired} as a category error, §1.1a — capacity is kinematic and does not select outcomes, which is $\lambda$'s role.)**

This is analogous to how, in classical statistical mechanics, a system at temperature $T$ has specific molecular configurations in any individual realization. The thermodynamic ensemble $e^{-\beta H}$ is a probability distribution over microstates; any specific realization sits at one microstate. We do not say "the system is in a quantum superposition of all microstates with thermal weights"; we say "the system is in some specific microstate, the distribution over runs follows the thermal ensemble."

In the quantum case, the same structure obtains with one additional ingredient: the per-run microstate must be consistent with the (FQ)-decohered structure on the regional algebra. After decoherence + (FQ), the regional algebra-state is strictly classical (no remaining off-diagonal coherence); the per-run universe occupies one branch of this strictly-classical mixture; which branch is determined by the specific microscopic initial conditions.

\textbf{This is not a magic transition from "70\% A + 30\% B as a superposition" to "100\% A as a single outcome".} It is the recognition that:
\begin{enumerate}
  \item We never have direct access to a "70\%/30\% superposition" of macroscopic alternatives, we have access only to macroscopic records, which after decoherence are exponentially exclusive
  \item (FQ) makes "exponentially exclusive" into "strictly exclusive": the residual sub-$\epsilon$ record-coherence has no physical referent
  \item After decoherence renders the records non-interfering, the non-dynamical selector $\lambda$ (P1) makes exactly one record actual — this is what makes the outcome single (an earlier version attributed single-ness to capacity forbidding a $\ge 2$-record content; that reading is \textbf{retired} as a category error, §1.1a — capacity does not forbid multiplicity)
  \item \emph{Which} single record the per-run universe carries is indexed by its microscopic IC (the IC does not produce the single-ness: step 3 does)
\end{enumerate}

The per-run "microscopic initial conditions" are \emph{not} hidden variables in any ontological sense. They are simply \emph{the actual physical state of the apparatus + environment in the specific run}. Different runs of an experiment correspond to different actual physical universes, different particles emitted, different microscopic apparatus states, different environmental configurations. These differences are not added structure; they are the standard fact that \emph{actual physical universes differ across actual physical experiments}. Standard QM, at the universal-wave-function level, already accommodates this trivially: the universal wave function is a function of all degrees of freedom in the universe; those degrees of freedom take different actual values in different runs.

The framework is \textbf{ψ-monist in the weak (dynamical) sense} (§2.2): the wave function is the only ontology carrying a law of motion; the only additional fact is the non-dynamical run-index $\lambda$, not a second dynamical substance. What the framework does, and what standard textbook QM elides, is to take seriously that:
\begin{enumerate}
  \item The universal wave function (not the textbook subsystem wave function) is the actual physical state in each run
  \item Different runs are different actual physical universes, not different versions of "the same" universe
  \item The Born statistics across runs come from the distribution of \emph{actual} initial conditions across \emph{actual} runs, not from any "alternative" wave functions for a single preparation
\end{enumerate}

In Bohmian language: there are no Bohmian particles and no second \emph{dynamical} ontology; the wave function $\Phi$ is the only thing carrying a law of motion. What plays a role structurally analogous to the Bohmian configuration is the run-index $\lambda$, the run's actual microscopic initial conditions, a \emph{non-dynamical} actuality fact (no guidance law, no back-reaction on $\Phi$) that selects which single record is realized. Different runs carry different $\lambda$. This is weak-ψ-monist standard QM with (FQ) added as a finite-information constraint, plus the non-dynamical actuality selector $\lambda$.

\subsection{Theorem: Single-record per-run wave functions}

\textbf{Theorem 1 (Single-record per-run wave functions).} \emph{Under (FQ) and standard decoherence dynamics of unitary QFT, the per-run physical state on the regional algebra $\hat{\mathcal{A}}(R)$ for $R$ containing apparatus + environment after measurement is, with regional physical content given by the algebraic-state-modulo-(FQ)-precision equivalence class, a single-record state $\omega_{k_{\rm run}}^R$.}

\emph{Proof.} By standard decoherence theory, after measurement interaction, the joint state's induced regional state on $\hat{\mathcal{A}}(R)$ has off-diagonal coherence terms suppressed exponentially: $\omega_\Psi^R = \sum_k |c_k|^2 \omega_k^R + O(\exp(-N))$. By Lemma 1 (§5.1), components below the (FQ) precision floor $\epsilon(R)$ are physically equivalent to zero. Since $e^{-N} \ll \epsilon(R)$ for any reasonable macroscopic decoherence regime, the off-diagonal coherence terms are physically exactly zero. The resulting regional state is the strict classical mixture $\sum_k |c_k|^2 \omega_k^R$, with macroscopic records physically exclusive (no remaining coherence between distinct $\omega_k^R$).

This strict classical mixture still carries every record; decoherence has removed interference, not multiplicity. The single-record conclusion is supplied by the selector $\lambda$ (P1): on the decohered (non-interfering) record set, $\lambda$ marks exactly one record actual, $\omega_{k_{\rm run}}^R$ for some specific $k_{\rm run}$, with the per-run microscopic initial conditions indexing which $k_{\rm run}$ (the ICs index the run; $\lambda$, not capacity, is what makes the content single). (An earlier version derived single-ness from a \emph{finite-information restriction} — summed cost $I_\Sigma \approx K\cdot I_0 > Q_R$; that is \textbf{retired} as a category error, §1.1a: the branch-summed cost is not the von Neumann entropy $Q_R$ bounds — \texttt{branchSummed\_not\_bounded\_by\_Shannon} — and a second branch raises $S(\rho_R)$ by only $\sim\!\log K$.) The global evolution remains exactly unitary throughout; no amplitude is trimmed. $\blacksquare$

\emph{(The single-record step is now carried by $\lambda$-selection (P1) + decoherence; the former §7.6 Macroscopic-Definiteness / $I_\Sigma > Q_R$ step is \textbf{retired} as a category error, §1.1a — not an open problem. The (FQ) resolution-equivalence of §5.1 still removes the residual sub-$\epsilon$ off-diagonal coherence.)}

\textbf{Remark on the role of each ingredient:}
\begin{itemize}
  \item \textbf{Decoherence} drives the regional algebra-state to be exponentially close to a classical mixture, so the records no longer interfere (standard QM, no postulate)
  \item \textbf{The selector $\lambda$ (P1)} marks exactly one of the non-interfering records actual; the capacity bound (FQ) plays only the kinematic role of bounding $S(\rho_R)$, it does not render coherence exactly zero and does not select a branch
  \item \textbf{The per-run wave function itself} carries one $\lambda$-selected record of the decohered classical mixture, this is the $(\Phi,\lambda)$ ontology: no second dynamical ontology, the actual per-run state is simply different from the textbook formal ensemble descriptor
  \item \textbf{Born statistics} across runs follow from the typicality of microscopic IC (open: Concentration / Born-typicality program)
\end{itemize}

\textbf{What standard linear QM does not provide:} a single actual record per run. Decoherence renders the records non-interfering (off-diagonal coherence exponentially small), and the non-dynamical selector $\lambda$ (P1) marks exactly one actual; the per-run microscopic IC index which one. Neither ingredient requires modifying the Schrödinger evolution; the capacity bound (FQ) is the kinematic arena, not a selection agent.

\subsection{What is preserved, and what algebraic restriction is (and is not)}

\textbf{Preserved exactly:} the unitary Schrödinger / Heisenberg evolution of the underlying field algebra. No modification of dynamics, no stochastic term, no projection operator, no restriction on the admissible Hamiltonians. The evolution law is exactly standard QM; what the framework adds operates not on the dynamics but on which \emph{regional contents} count as physically instantiable (next paragraph).

\textbf{Constrained at the kinematic level:} what (FQ) constrains is the regional \emph{entropy}, $S(\rho_R) \le Q_R$ for the pre-selection state $\Phi$. This is a kinematic ceiling on the entropy (equivalently the effective modular rank) of the regional content, \textbf{not} an exclusion of $\ge 2$-record contents: a second macroscopic branch raises $S(\rho_R)$ by only $\sim\!\log K$, far below $Q_R$, so capacity does not forbid multi-record contents (the "cost exceeding $Q_R$ / no per-run referent" reading is retired as a category error, §1.1a). \textbf{This is not a restriction on the Hamiltonians and not a modification of the dynamics:} the global Schrödinger / Heisenberg evolution law is exactly unitary and unconstrained, including ordinary measurement interactions. Single-record definiteness is supplied separately by decoherence + $\lambda$ (P1), not by the entropy bound.

\textbf{The relationship between global evolution and regional physical content.} On $\mathcal{H}_{\rm phys}$ the universal state evolves unitarily under a physical Hamiltonian: $|\Psi_t\rangle = U(t)|\Psi_0\rangle$. The induced regional algebra-state evolves linearly: $\omega_t(O) = \omega_0(U(t)^\dagger O U(t))$ for $O \in \hat{\mathcal{A}}(R)$. The map $|\Psi\rangle \mapsto \omega_\Psi$ from Hilbert vectors to regional algebra-states is a structural restriction (analogous to forming a reduced state). It is mathematically linear on density operators and quadratic in vectors.

The operative equivalence relation on physical regional states is exact algebraic-state equality $\omega_\Psi = \omega_\Phi$ on $\hat{\mathcal{A}}(R)$ (§5.4). (An earlier version introduced instead a coarser "physical-precision quotient" $\omega_\Psi$ modulo resolution-$\epsilon$ equivalence and made it do the foundational work; that amplitude-precision reading is superseded, §1.1a.)

The foundational work is done not by any precision quotient but by decoherence and the selector $\lambda$: decoherence drives the off-diagonal coherence between macroscopic records to exponentially small values, so the records no longer interfere, and $\lambda$ (P1) marks exactly one record actual, with the per-run microscopic IC indexing which one. The capacity bound (FQ) plays only the kinematic role of bounding $S(\rho_R)$; it does not render coherence exactly zero and does not select a branch.

\subsection{Physics is the macroscopic observable content; branches are not part of it}

A potential objection to the framework: even after decoherence renders the off-diagonal coherence negligible, the resulting classical mixture $\sum_k p_k \omega_k^R$ on $\hat{\mathcal{A}}(R)$ is still a mixture over multiple macroscopic records. To get a single-record per run, the framework does need a selector: this is exactly the role of $\lambda$ (P1), and it is neither MWI (all branches real) nor a dynamical hidden-variable mechanism, but a non-dynamical actuality fact.

The framework's response cuts deeper than introducing a selection mechanism. It refuses the question of "which Everett branch is realized" by reframing what counts as physical content:

\begin{qiqtquote}
\textbf{Physics is what is encoded in the macroscopic observable algebra. The full universal wave function (with all its mathematical branch structure) is a calculational/formal apparatus; the physical content of any region is the state on the macroscopic record subalgebra $\mathcal{C}(R) \subset \hat{\mathcal{A}}(R)$. We do not have direct physical access to "the universal wave function's branch structure", only to the macroscopic record content. (FQ) bounds the entropy $S(\rho_R)$ of that content but does not by itself single out one record. Single-record per run arises because decoherence renders the records non-interfering and the non-dynamical selector $\lambda$ (P1) marks exactly one of them actual. The microscopic branch question is not part of the physics because it does not appear in the macroscopic observable content.}
\end{qiqtquote}

This is the operational-realist + algebraic + holographic move. Physics is the content of macroscopic observable algebras (algebraic QFT lineage). These algebras are constrained by (FQ). The "Everett branches" of the universal wave function are mathematical artifacts of the underlying Hilbert-space formalism; they don't enter the macroscopic observable content; they're not part of physics.

This reframing originally motivated the \emph{Macroscopic Definiteness Conjecture} (§7.6); that conjecture is now \textbf{retired} as a category error (§1.1a) — capacity is kinematic and does not select outcomes, which is $\lambda$'s role.

\subsection{The macroscopic record subalgebra and the spectrum of records}

\begin{qiqtquote}
\textbf{Read with the §1.1a retraction.} The macroscopic record subalgebra and spectrum below are current and stand. But the \emph{cost-counting} development in this subsection (the multi-record "information cost" conjecture and its number-bound) was the engine of the now-\textbf{retired} Macroscopic Definiteness Conjecture: it does \textbf{not} establish that capacity forbids a multi-record state (capacity bounds the von Neumann entropy $S(\rho_R)$, which a second branch raises only by $\sim\!\log K$), and single-record definiteness comes from the selector $\lambda$ + decoherence, not from a capacity exclusion (§1.1a). The finite \emph{number-bound} Lean theorem (\texttt{SBSBridge}, \texttt{CoreNoCollapse}) remains valid as the storage cost of the $\lambda$-selected record — it bounds an \emph{additive} cost that is provably \textbf{not} the holographic entropy (\texttt{branchSummed\_not\_bounded\_by\_Shannon}).
\end{qiqtquote}

\textbf{Definition (Macroscopic record subalgebra).} \emph{Let $\hat{\mathcal{A}}(R)$ be the regional Type II crossed-product algebra for region $R$ containing apparatus + environment. The \textbf{macroscopic record subalgebra} $\mathcal{C}(R) \subset \hat{\mathcal{A}}(R)$ is the maximal commutative subalgebra of observables that are decoherence-stable under the dynamics of the apparatus + environment, that is, the subalgebra selected by einselection (Zurek) and that supports definite macroscopic records over the relevant timescales.}

For the double-slit screen, $\mathcal{C}(R)$ is generated by the position-localized record projectors $\{P_k\}$ corresponding to spots at different positions. For a generic measurement, $\mathcal{C}(R)$ is the algebra of pointer-basis observables. The subalgebra is approximately commutative (off-diagonal terms between distinct records are exponentially suppressed by decoherence, so the records are non-interfering; cf. §6.2–6.3).

\textbf{Spectrum of records.} $\mathrm{Spec}(\mathcal{C}(R))$ is the set of macroscopic record configurations, for the screen, the set of distinct spot positions; for a Stern-Gerlach detector, the set of distinct deflection records; etc. By the Gelfand representation theorem, $\mathcal{C}(R) \cong C(\mathrm{Spec}(\mathcal{C}(R)))$ (continuous functions on the spectrum), and states on $\mathcal{C}(R)$ are probability measures on $\mathrm{Spec}(\mathcal{C}(R))$.

For a finite or countable spectrum (which is the relevant case for macroscopic records on a finite-resolution screen), states on $\mathcal{C}(R)$ are probability distributions $\{p_k\}$ over records.

\textbf{Thickened-state construction.} For each record $r \in \mathrm{Spec}(\mathcal{C}(R))$, the \textbf{thickened state} $\tilde{\delta}_r$ on the full algebra $\hat{\mathcal{A}}(R)$ is the state corresponding to "record $r$ realized", including all the microscopic structure of the apparatus + environment configuration that produces macroscopic record $r$. This thickened state uses approximately $Q_R$ bits of physical information (it is a specific microscopic configuration of the full apparatus + environment region; macroscopic records consume approximately the full regional holographic capacity).

For a multi-record state, a probability measure $\mu = \sum_k p_k \delta_{r_k}$ on $\mathrm{Spec}(\mathcal{C}(R))$ with multiple records, the \emph{thickened state} is $\tilde{\mu} = \sum_k p_k \tilde{\delta}_{r_k}$ on $\hat{\mathcal{A}}(R)$. This thickened state must encode each constituent macroscopic record's full microscopic configuration plus the probability weights.

\textbf{Conjecture (Information cost of multi-record thickened states):} \emph{The renormalized entropy $S_{\rm ren}(\tilde{\mu})$ of a thickened multi-record state $\tilde{\mu} = \sum_k p_k \tilde{\delta}_{r_k}$ scales as $\sum_k p_k S_{\rm ren}(\tilde{\delta}_{r_k}) + H(\{p_k\})$ where $H$ is the Shannon entropy, when the constituent records $\{r_k\}$ are physically distinct macroscopic configurations with mutually exclusive microscopic specifications. When each $S_{\rm ren}(\tilde{\delta}_{r_k}) \approx Q_R$ (records saturate the regional capacity), the thickened multi-record state has $S_{\rm ren}(\tilde{\mu}) \approx Q_R + H(\{p_k\})$, which exceeds $Q_R$ when $H(\{p_k\}) > 0$, i.e., for genuinely multi-record states.}

This is the cost-counting argument that \emph{originally} motivated the Macroscopic Definiteness Theorem of §7.6 — now \textbf{retired} as a category error (§1.1a; see the banner opening this subsection). The framework no longer rests on it for single outcomes.

\textbf{Formalization note, the number-bound is a machine-checked conditional theorem (the strong form is not).} The cost-counting argument above has a \emph{continuum} part (the renormalized-entropy scaling on the Type II regional algebra, which remains the open Macroscopic Definiteness Conjecture, Open Problem 3) and a \emph{finite-dimensional core} that we have now mechanized. In a tractable finite model the argument is no longer a conjecture but an axiom-free Lean theorem (modules \texttt{CoreNoCollapse.lean}, \texttt{CapacityModel.lean}, \texttt{SBSBridge.lean}; standard Lean axioms only). Concretely: realizing the constituent records as an objective Quantum-Darwinist record, a Spectrum Broadcast Structure (Korbicz, arXiv:2007.04276) in which the $n$-outcome pointer is redundantly broadcast to $R$ environment fragments, forces the broadcast support to have Hilbert dimension $\ge n^{R}$ (the fragment Hilbert spaces tensor, so dimensions multiply), hence an \emph{information} storage cost $\ge R\log n$ (\texttt{SBSBridge.redundancy\_le\_logStorage}). Under a finite additive storage capacity $Q_{\max}$, a "macroscopic" record (large redundancy $R$) then has cost exceeding $Q_{\max}/2$, so a region cannot carry two of them: at most one such record is active (\texttt{CoreNoCollapse.coactual\_subsingleton}, \texttt{SBSBridge.sbs\_single\_outcome}); the selector $\lambda$ supplies which one. The capacity threshold is \emph{derived} from orthonormality of the pointer states ($\sum_j\dim_j\le D$, \texttt{CapacityModel.capacity\_total}), not stipulated, and the chain is robust to \emph{approximate} decoherence, a near-orthonormal fragment family (pairwise overlap $<1/(n-1)$) is still linearly independent, so the dimension bound survives imperfect distinguishability (\texttt{fragment\_finrank\_ge\_approx}), with an amplification lemma showing that weak per-collision distinguishability $\gamma<1$ compounds to exponentially reliable records (block overlap $\le\gamma^{L}\to 0$). This mechanizes the \emph{number-bound} (at most one macroscopic record) form of the conjecture; it does \textbf{not} establish the strong superposition-exclusion form (the redundant cat state $\alpha|0\rangle^{\otimes N}+\beta|1\rangle^{\otimes N}$ remains the obstruction, §7.6, Open Problem 3). It is moreover a \emph{conditional} theorem whose one physical input is the factorized, branch-dependent scattering premise ($\gamma<1$); \textbf{for the standard collisional/QND monitoring model that premise is itself now a machine-checked theorem.} Taking the textbook interaction $H_{\rm int} = g\,\sigma_z^S \otimes \sigma_x^E$ (Zurek), the branch-conditioned environment records overlap by $\langle E_+|E_-\rangle = \cos^2\theta - \sin^2\theta = \cos 2\theta$, so the per-collision distinguishability $\gamma = |\cos 2\theta| < 1$ for generic coupling $\theta = g\tau$ is \emph{derived}, not assumed (module \texttt{CollisionalGamma.lean}: \texttt{branch\_overlap}, \texttt{gamma\_lt\_one}). The propagator $U_s(\theta) = \exp(-i\theta\,sA)$ is a verified one-parameter unitary group ($A = \sigma_x$ a self-adjoint involution; \texttt{collisionU\_group}, \texttt{hamiltonian\_isSymmetric}), the pointer basis is einselected because $[H_{\rm int},\sigma_z^S]=0$, and over $L$ independent collisions the overlap factorizes to $\gamma^L \to 0$ (\texttt{collisional\_overlap\_tendsto\_zero}), supplying the previously-assumed amplification input. The residual open physics is the \emph{field-theoretic} version, deriving uniform $\gamma<1$, and the QND/no-recoherence structure itself, from a realistic system–environment Hamiltonian rather than positing the collisional form, a sharply isolated form of the Strasberg–Winter "branch selection problem" (arXiv:2601.19703). Everything downstream of the premise is machine-checked.

\subsection{Bell: the framework violates Parameter Independence, not measurement independence}

A natural objection: any deterministic single-world theory reproducing standard quantum predictions must confront Bell's theorem. The framework is \textbf{ψ-monist in the weak (dynamical) sense} (no extra dynamical particles or fields; §2.2), the complete per-run description is the pair $(\Phi, \lambda)$, with $\lambda$ a non-dynamical actuality selector, it reproduces standard quantum predictions, and it has single outcomes per run. \emph{No theory} reproducing the quantum correlations can keep all of Bell's premises; the only question is which one it gives up. We answer that question precisely.

\textbf{The framework's commitment: it keeps measurement independence (free settings) and violates Bell local causality, specifically Parameter Independence, at the ontic level, while preserving operational no-signaling.}

Bell's \emph{local causality} factorizes (Jarrett; Shimony) into two sub-conditions, conditional on the complete ontic state:

\begin{itemize}
  \item \textbf{Outcome Independence (OI):} Alice's outcome distribution is independent of Bob's \emph{outcome}, given the ontic state.
  \item \textbf{Parameter Independence (PI):} Alice's outcome distribution is independent of Bob's \emph{setting}, given the ontic state.
\end{itemize}

The framework's commitments fix the choice with no remaining freedom:

\begin{enumerate}
  \item It keeps \textbf{outcome definiteness} (a single record per run): so the Everett escape (denying definite outcomes) is unavailable.
  \item It keeps \textbf{measurement independence / free settings}: the settings $a,b$ are statistically independent of the run's actuality fact $\lambda$, $\;\rho(\lambda\mid a,b)=\rho(\lambda)$. The framework therefore \textbf{rejects superdeterminism}: this escape is deliberately closed.
  \item Because the per-run selection is \textbf{deterministic} in $(\Phi,\lambda)$, \textbf{Outcome Independence holds trivially} (a deterministic outcome is automatically independent of the distant outcome given $\lambda$).
\end{enumerate}

With (1), (2), and (3) fixed, the only Bell premise left to deny is \textbf{Parameter Independence}. There is no fourth door. The framework is therefore \textbf{Bell-nonlocal by necessity}, exactly as every single-world, non-superdeterministic, definite-outcome theory must be, and it is honest about \emph{which} horn that is.

\textbf{This is not a special property of holography; it is the generic price.} Orthodox quantum mechanics \emph{also} violates Bell local causality. What the present framework adds on top of orthodox QM is \emph{definite $\lambda$-selected outcomes}, and that addition is exactly what re-incurs a burden orthodox (statistical, no-signaling) QM avoids: the ontic-level PI violation must wash out, after averaging over the actuality measure $\mu$, to \textbf{operational no-signaling},

\[
P(A\mid a,b) = P(A\mid a),
\]

even though at the ontic level $P(A\mid a,b,\lambda)\neq P(A\mid a,\lambda)$. Operational no-signaling is established for the framework's instruments (Theorem 7, §7.7; Lean \texttt{Theorem7.Setup.no\_signaling}). That the \emph{required} cancellation is not fine-tuned, i.e. that the same $\mu$ delivering Born weights also enforces no-signaling without an experiment-by-experiment conspiracy, is a genuine open burden (the Wood–Spekkens fine-tuning problem), recorded as part of the open Born-typicality program (§7.4, §11.4) and in \texttt{PROGRAM\_STATUS.md}.

\textbf{What holographic nonseparability does and does not contribute.} It is tempting to say the Bell correlations are "explained by holographic nonseparability." We resist this: Bell local causality can be formulated for \emph{commuting spacelike algebras} and does \textbf{not} require a Hilbert-space tensor factorization, so the Type III$_1$ non-factorization of local algebras (Reeh–Schlieder, vacuum entanglement) is shared with \emph{orthodox} AQFT and does not by itself explain definite-outcome correlations or supply a local completion. Algebraic nonseparability explains why the \emph{state} is nonseparable; it does not select outcomes. The genuine, non-redundant content is the ontic PI violation above; "nonseparability" is the \emph{register} in which that violation lives, not an independent escape. (We also avoid leaning on "no tensor factorization" as the escape while simultaneously invoking the split property elsewhere to model local instruments; the split property supplies an \emph{approximate} Type I factorization adequate for ordinary Bell experiments, and the two registers must not be conflated.)

\textbf{No-signaling ≠ locality ≠ Lorentz invariance.} These are distinct, and we claim only what is established: PR-box correlations are no-signaling yet \emph{more} nonlocal than quantum mechanics, so operational no-signaling (which we have) does not by itself secure Bell-correct correlations \emph{or} Lorentz invariance of the beable $C_R$. Lorentz covariance of the selection map $A_R[\Phi,\lambda]$, that it carries no hidden foliation, is \textbf{not} established and is recorded as an open problem of equal rank to the Born-typicality measure (§11; \texttt{PROGRAM\_STATUS.md}).

\textbf{Position with respect to Kochen–Specker.} If per-run structure determines outcomes, value assignments must be contextual, which the framework accepts. The per-run wave function's content on a regional algebra is naturally contextual: different measurement contexts probe different algebra-projections, with no requirement of joint-value consistency across noncommuting contexts.

\textbf{Position with respect to PBR.} The framework is \textbf{ψ-ontic and ψ-monist in the weak (dynamical) sense}: $\Phi$ is physically real and is the sole \emph{dynamical} ontology. It adds no hidden particle positions, no hidden fields, no second \emph{dynamical} layer. It does include one non-dynamical actuality fact, the selector $\lambda$ fixing which single record is realized (§2.2, §7.6), so it is \textbf{formal-ψ-incomplete}: the actual per-run ontic state is the pair $(\Phi, \lambda)$, finer-grained than the textbook formal wave function (which describes the across-run distribution). PBR rules out certain ψ-epistemic positions; the framework is ψ-ontic and not one of those, $\Phi$ is real, not a state of knowledge.

\textbf{Contrast (not alliance) with Palmer's RaQM / Invariant Set Theory.} Palmer (2025 and prior work) takes the \emph{other} horn: a principled rejection of \textbf{measurement independence} via the geometry of a fractal invariant set, formally a (non-conspiratorial) measurement-dependence / superdeterminism-family escape. The present framework deliberately takes the opposite horn: it \emph{keeps} measurement independence and pays at Parameter Independence instead. The two programs share only single-world realism and a finite-information motivation; on the Bell question they are contrasting, not allied. Detailed comparison is in §10.8.

\subsection{Tunneling: standard dynamics, record-level reading}

Quantum tunneling is a useful test of what QIQT-H does and does not modify. The under-barrier amplitude is real physical content of the formal field state, and \textbf{all tunneling amplitudes and rates are exactly those of standard quantum theory}, coherent tunneling (ammonia inversion, Josephson oscillation, double-well coherent oscillation), WKB barrier penetration, $\alpha$-decay and fusion rates, instanton / vacuum-decay amplitudes, and Quantum Zeno / anti-Zeno effects are all unchanged, because (FQ) does not modify the Schrödinger evolution of the underlying field algebra.

**The information limit constrains how the wave function \emph{resides}; it does not trim it.\textbf{ This point must be stated carefully, because QIQT-H is ψ-monist and deterministic: there are }no fundamental probabilities** in the theory, only the wave function and its unitary evolution. The holographic / (FQ) bound is a constraint on the \emph{physical specification capacity} of a region, on how much amplitude/phase structure the wave function can be instantiated with in $R$, not a dynamical rule that sets small amplitudes to zero. Concretely:

\begin{itemize}
  \item The under-barrier evanescent tail is \textbf{genuine wave-function content} and stays so. (FQ) never deletes it, never truncates it to zero, and never alters its unitary evolution. The Schrödinger equation runs untouched.
  \item (FQ) does \emph{not} act pointwise on $\psi$ in the forbidden region. The pointwise value $\psi(x)$ is not even a well-defined physical target for a floor: it is basis-, normalization-, and coarse-graining-dependent, and in the continuum it is dimensionful. The bound is a statement about the regional algebra-state's information content, not about individual amplitudes.
  \item There is therefore \textbf{no kinematic cutoff on tunneling amplitudes}. Any reading on which "a tunneling amplitude below $\varepsilon(R)$ is excluded" is mistaken. The decisive objection is dynamical and basis-theoretic: a fixed pointwise floor would be ill-defined (a transition amplitude over a time slice scales with the arbitrary slice duration; the pointwise value is basis-, normalization-, and coarse-graining-dependent) and, applied to the dynamics, would break unitarity and wrongly freeze ordinary decay. Independently, such a reading misreads (FQ): (FQ) is a capacity constraint on regional residence, not an amplitude-trimming operation: and, since the theory has no fundamental probabilities, "excluding a low-amplitude outcome" is not even a move the framework can make.
\end{itemize}

What (FQ) \emph{does} constrain is the regional \textbf{record structure}: $Q_R$ caps \emph{how many} distinguishable decohered records a region can hold (its genuine kinematic role; it does \textbf{not} forbid a superposition of records — §1.1a). Which single tunneling \emph{outcome} resides as the definite regional record is then fixed by decoherence + the selector $\lambda$, not by a capacity exclusion. This is where the framework's content enters — not in the coherent tunneling amplitude, but in how a tunneling outcome comes to reside as a definite regional record.

The correct statement is therefore: QIQT-H leaves the tunneling wave function and its evolution exactly as in standard quantum theory, and differs only in what it says about the \textbf{regional record once decoherence has occurred} (the particle is found, or not, on the far side). Per run, the deterministic evolution of the full wave function drives the records to be non-interfering, and the non-dynamical selector $\lambda$ (P1) leaves a single record actual in $R$ (the microscopic IC index which one; the capacity bound only bounds $S(\rho_R)$). The textbook Born numbers $|\psi|^2$ are recovered as the \emph{empirical frequency pattern} of records across many runs (the typicality reading of §7.4, §11.4), never as a fundamental per-run probability. This is the same mechanism as for any other measurement (§6); tunneling is not a special case. \textbf{Standard dynamics (barrier penetration, unchanged); QIQT-H reading of the regional record only, and even there, by $\lambda$-selection, not by amplitude trimming.}

\textbf{Horizon tunneling (Hawking radiation).} The Parikh-Wilczek (1999) tunneling derivation of Hawking emission expresses the emission rate as $\Gamma \sim \Gamma_{\rm grey}\, e^{\Delta S_{\rm BH}}$, where $\Delta S_{\rm BH} = \Delta A/(4\ell_P^2)$ is the change in Bekenstein-Hawking entropy (in nats). For an emitted quantum the horizon area decreases, $\Delta A < 0$, so the exponential factor \emph{suppresses} emission; writing the positive entropy loss $\Delta S_{\rm loss} = S_i - S_f > 0$, $\Gamma \sim \Gamma_{\rm grey}\, e^{-\Delta S_{\rm loss}}$ (greybody prefactor $\Gamma_{\rm grey}$ aside). The same entropy is what QIQT-H reads as the finite information content of the horizon region (in bits, $Q_{\rm hor} = S_{\rm BH}/\ln 2 = A/(4\ell_P^2 \ln 2)$). This is a clean conceptual alignment, the same quantity governs lab-scale per-run wave functions and the horizon's information budget, but it is a \emph{reinterpretation}, not a new prediction: QIQT-H does not by itself modify or re-derive the Parikh-Wilczek emission rate.

\medskip\hrule\medskip

\section{Formal Consequences}

\subsection{Theorem (FQ restricts physically realized wave functions)}

\textbf{Theorem 2.} \emph{The class of per-run wave functions consistent with (FQ) is a proper subclass of formally describable wave functions in standard $T_{QM}$. Specifically, formal wave functions whose induced state on some $\hat{\mathcal{A}}(R)$ violates the renormalized entropy bound $S_{\rm ren}(\omega_\Psi) > Q_R$ are excluded as physically realized per-run states.}

\emph{Proof.} (FQ) part (ii) imposes $S_{\rm ren}(\omega_\Psi) \le Q_R$ on per-run wave functions. Formal wave functions in standard $T_{QM}$ are not subject to this constraint; we can construct formal states whose induced regional states violate the bound. Such formal states are excluded from the physical model class. $\blacksquare$

\subsection{Theorem (Finite regional entropy)}

\textbf{Theorem 3 (Finite regional entropy).} \emph{Under (FQ), the per-run pre-selection state $\Phi$ has, in any bounded region $R$, regional von Neumann / renormalized entropy bounded by the holographic capacity, $S_{\rm vN}(\rho_R) \le Q_R = A(\partial R)/(4\ell_P^2)$; equivalently, the effective modular rank of $\rho_R$ is at most $e^{Q_R}$.}

\emph{Proof sketch.} This is (FQ) part (ii) applied to the pre-selection state: the regional information functional $\chi_R = S_{\rm ren}(\omega_\Psi)$ is bounded by $Q_R$ by postulate. $\blacksquare$

\emph{Remark (superseded subsidiary reading).} An earlier version stated Theorem 3 as a finite amplitude-\emph{resolution} floor $\epsilon(R) > 0$ (amplitudes within $\epsilon(R)$ of $0$ or $1$ physically equivalent to exact $0$ or $1$). That amplitude-precision reading is subsidiary and superseded (§1.1a); the load-bearing statement is the entropy bound $S(\rho_R) \le Q_R$ above. The quantitative form of any $\epsilon(R)$ is not needed for the framework's conclusions.

\subsection{Theorem (Single-record per-run wave functions)}

\begin{qiqtquote}
\textbf{Superseded (2026) — see the Macroscopic Definiteness retraction in §1.1a.} The single-record conclusion stands, but it follows from $\lambda$-selection (P1) + decoherence — \textbf{not} from the finite capacity forbidding a $\ge 2$-record state. The "summed information cost $I_\Sigma \approx K\cdot I_0 > Q_R$ / not instantiable" step (and its labelling as "the load-bearing step") is the \textbf{retired category error}: the branch-summed cost is not the entropy (\texttt{BranchLedger.branchSummed\_not\_bounded\_by\_Shannon}), and a second branch adds only $\sim\!\log K$ to $S(\rho_R)$, far below $Q_R$. Throughout this theorem and its remarks, read "the load-bearing single-ness step" as \textbf{$\lambda$-selection}, and the capacity bound as playing only its kinematic role.
\end{qiqtquote}

\textbf{Theorem 4 (Single-record per-run).} \emph{Under (FQ) (the finite regional entropy bound), together with (i) standard decoherence (which renders the macroscopic records non-interfering) and (ii) the non-dynamical selector $\lambda$ of the $(\Phi,\lambda)$ ontology (P1), which marks exactly one decohered record actual per run, the per-run wave function in any region $R$ containing apparatus + environment after measurement is physically a single-record state. Which record is realized is indexed by the run's microscopic initial conditions; the unrealized records have no physical referent in $R$ (they are not instantiable, not "real elsewhere"). No amplitude is trimmed and the global evolution is exactly unitary throughout.}

\emph{Proof.} By standard (unitary) decoherence, the regional state on $\hat{\mathcal{A}}(R)$ is exponentially close to a strict classical mixture $\sum_k |c_k|^2 \omega_k^R$ of macroscopic records, interference is gone, but the mixture still carries every record (branch weights $|c_k|^2$ are conserved; cf. \texttt{NoConcentration.lean}). By Theorem 3 the (FQ) resolution floor $\epsilon(R) > 0$ makes the residual sub-$\epsilon$ off-diagonal coherence have no physical referent (a resolution-equivalence on regional states, §5.1), so the resident content is exactly that strict mixture. By $\lambda$-selection (P1), on the decohered (non-interfering) record set exactly one record is marked actual; hence the resident regional content is single-record. (The earlier "summed cost $I_\Sigma \approx K\cdot I_0 > Q_R$ / not instantiable" route is \textbf{retired} as a category error, §1.1a.) Which record is realized is indexed by the run's microscopic initial conditions. The global evolution remains exactly unitary; no amplitude is trimmed. $\blacksquare$

\emph{(The single-record step is supplied by $\lambda$-selection (P1) + decoherence; the former §7.6 Macroscopic-Definiteness / $I_\Sigma > Q_R$ step is \textbf{retired} as a category error, §1.1a — not an open problem. The (FQ) resolution-equivalence of §5.1 removes the residual sub-$\epsilon$ off-diagonal coherence. The machine-checked \texttt{CoreNoCollapse}/\texttt{CapacityModel}/\texttt{SBSBridge} development (§6.8) proves a conditional single-record statement for a finite \textbf{additive}-cost model; but since the additive (branch-summed) cost is provably not the holographic von Neumann entropy \texttt{branchSummed\_not\_bounded\_by\_Shannon}, this models the storage of the $\lambda$-selected record, and does \textbf{not} revive the retired "capacity forbids two records" reading.)}

\emph{Remark on what does the work (not "amplitude concentration").} It would be a mistake to read Theorem 4 as "decoherence + microscopic IC drive the per-run amplitudes dynamically to $0$ or $1$." They do not: by unitarity the branch weights $|c_k|^2$ are conserved (the \texttt{NoConcentration.lean} audit makes this precise), so decoherence removes \emph{interference} but never \emph{selects} a record, the decohered diagonal content still carries every record. The single-record conclusion comes instead from the selector $\lambda$ (P1): on the decohered record set, $\lambda$ marks exactly one record actual. (The earlier "finite-information restriction / summed cost $> Q_R$" route is \textbf{retired}, §1.1a — capacity is kinematic, not exclusionary.) The (FQ) resolution floor plays the subsidiary role of removing the residual sub-$\epsilon$ off-diagonal coherence (as a resolution-equivalence on regional states, not by trimming any amplitude). Microscopic IC then merely index which single record the run carries. This is the cooperative structure: decoherence (interference) + $\lambda$-selection (single-ness, the load-bearing step) + IC (which one).

\emph{Remark on how a single record arises (decoherence is necessary but not sufficient).} A machine-verified audit (\texttt{lean/mathlib/QIQTH/NoConcentration.lean}) confirms that linear unitary measurement-decoherence \textbf{alone} cannot produce a single record: branch weights $|c_k|^2$ are conserved by unitarity, not concentrated toward $0$ or $1$. The equal-superposition input $\psi = (|0\rangle + |1\rangle)/\sqrt{2}$ leaves post-measurement weights exactly $(1/2, 1/2)$. Decoherence solves the \emph{interference} problem (it makes distinct records exponentially distinguishable), but it does \textbf{not} solve the \emph{selection} problem: the decohered regional state $\sum_k |c_k|^2 \omega_k^R$ is diagonal yet still carries $K \ge 2$ macroscopic records. The single-record-per-run content comes from the \textbf{selector $\lambda$}, in two cooperating steps with the second doing the load-bearing work:

\begin{enumerate}
  \item \textbf{Decoherence (unitary, standard, dynamical).} Entangling a record with $N$ apparatus + environment particles drives the off-diagonal coherence between macroscopically distinct records down like $e^{-N}$. This is ordinary Schrödinger evolution: no modification, no probability, no amplitude trimming. The records become non-interfering.
  \item \textbf{$\lambda$-selection (P1, the single-ness source).} On the decohered (non-interfering) record set, the non-dynamical selector $\lambda$ marks exactly one record as the actual regional content. The capacity bound (FQ) enters here only kinematically, it bounds the \emph{entropy} $S(\rho_R) \le Q_R$ (equivalently the effective modular rank, a cardinality/entropy ceiling), and does \textbf{not} forbid a $\ge 2$-record content: a second macroscopic branch raises $S(\rho_R)$ by only $\sim\!\log K$, far below $Q_R$. (An earlier version derived single-ness from a branch-summed cost $I_\Sigma \approx K\cdot I_0 > Q_R$ making a multi-record content "not a point of the region's physical state space at all", the Branch-Summed Holographic Bound / Macroscopic Definiteness Conjecture; that reading is retired as a category error, §1.1a.) The resident regional physical content is single-record because $\lambda$ selects one.
\end{enumerate}

This is decisively \textbf{not} many-worlds with regional coarse-graining: the unrealized records are not "real elsewhere", they have \emph{no physical referent}, because $\lambda$ (P1) marks exactly one decohered record as the actual regional content (the unselected records remain in $\Phi$ but are not regionally instantiated). It is also \textbf{not} a modification of the dynamics: the Schrödinger / Heisenberg evolution \emph{law} is exactly unitary throughout; what selects one record is $\lambda$ (P1), on the decohered set. A formal multi-record Hilbert-space vector remains a perfectly good \emph{calculational} descriptor; it simply has no per-run \emph{actual} referent in $R$ because $\lambda$ selects only one, this is the formal/per-run distinction plus $\lambda$-selection, not a ban on Hamiltonians and not a non-unitary projection. \textbf{There are no fundamental probabilities anywhere:} the weights $|c_k|^2$ are the empirical frequency pattern of records \emph{across runs} (§7.4, §11.4), and the per-run microscopic initial conditions merely \emph{index which} single-record universe a given run is (distinct runs are distinct actual universes), IC does not \emph{produce} single-ness; the selector $\lambda$ does.

\subsection{Theorem (Conditional Born typicality)}

\textbf{Reframing.} QIQT-H has no fundamental probabilities. The framework is deterministic per run: microscopic initial conditions of that run determine which macroscopic record is realized. "Born" in QIQT-H refers to the \emph{empirical frequency pattern} across many runs, a statistical regularity that must emerge from the IC distribution, not a primitive probability assignment. Accordingly Theorem 5 is best understood not in the probability-axiomatic mode of Gleason's theorem, but in the \textbf{typicality} mode of Bohmian Dürr-Goldstein-Zanghì equivariance (adapted to QIQT-H's no-particle, no-branching ontology).

\textbf{Theorem 5 (Conditional Born typicality).} \emph{Let $\rho$ be a preparation state, $\Omega_\rho$ the QIQT-H microscopic-IC space, and $O_M : \Omega_\rho \to \mathrm{Outcomes}(M)$ the deterministic outcome map for a measurement protocol $M$. Suppose:}

\emph{(i) \textbf{Canonicality.} There exists a measure $\mu_\rho$ on $\Omega_\rho$ defined from QIQT-H primitives (not fitted per measurement).}

\emph{(ii) \textbf{Born pushforward.} $(O_M)_* \mu_\rho(i) = \mathrm{tr}(\rho E_i)$ for every allowed outcome effect $E_i$ in the measurement protocol $M$.}

\emph{(iii) \textbf{Equivariance / stationarity.} $\mu_\rho$ is preserved (up to allowed parameter-dependence on $\rho$) by the (FQ)-restricted physical Hamiltonian.}

\emph{(iv) \textbf{Repeated-trial typicality.} Repeated trials of the prepare-measure procedure are represented by a $\mu_\rho$-typical iid (or stationary-ergodic) process on $\Omega_\rho$.}

\emph{Then for $\mu_\rho$-typical microscopic-IC sequences, the empirical relative frequency of macroscopic outcome $k$ across many runs converges to the Born weight $\mathrm{tr}(\rho E_k)$ (= $|c_k|^2$ for pure states with $\rho = |\psi\rangle\langle\psi|$ and projective measurements).} $\blacksquare$

\textbf{Status of the four hypotheses.} Hypotheses (ii) and (iv) are standard: (ii) is what defines $\mu_\rho$ as "the Born measure" for $\rho$; (iv) is a routine LLN/ergodicity input. Hypotheses (i) and (iii), \emph{canonicality} and \emph{equivariance}, are the \textbf{framework's load-bearing open commitments}. They are not delivered by FQ + AQFT + holography alone; see machine-verified audits below.

The conditional shape of Theorem 5 is now formalized in \texttt{lean/mathlib/QIQTH/BornTypicality.lean} (theorem \texttt{qiqth\_born\_typicality\_conditional}). The deterministic core (mean per-run frequency equals the outcome marginal) is rigorously proved; the LLN step is taken as a standard probability black box at the interface layer.

\textbf{Important, QIQT-H does not derive the Born measure.} It identifies the exact canonical-measure / equivariance principle required for Born frequencies, and reduces the Born problem to that principle. This is the typicality-paradigm analog of Bohmian $|\psi|^2$-equivariance, adapted for QIQT-H's no-particle ontology.

\textbf{Audit remark, universal realizability (the negative companion).} A machine-verified audit (\texttt{lean/mathlib/QIQTH/NoBornFromNothing.lean}) establishes that for \textbf{any} target outcome distribution $p$ (not just $|c_k|^2$) and any surjective outcome map from the IC space, there exists a measure $\mu_p$ whose outcome-marginal equals $p$. Construction: choose a section of the outcome map and place mass $p_k$ on $s(k)$. Hence the framework's structural axioms (FQ, microcausality, Donald, holographic bound) do \emph{not} select Born over any other distribution. Hypothesis (i) of Theorem 5 is therefore \emph{genuinely independent} of the other QIQT-H postulates, it must be added as a separate principle, not derived.

\textbf{Audit remark, support preservation $\neq$ Born equivariance.} A second machine-verified audit (\texttt{lean/mathlib/QIQTH/EquivarianceGap.lean}) establishes that the framework's claim that the exactly-unitary dynamics carries instantiable regional contents to instantiable regional contents (the closure property of §7.6, preservation of the \emph{support} of the admissible set) is strictly weaker than the measure-preservation hypothesis (iii) of Theorem 5. Concrete counterexample: a bijection on a 2-point space preserves the support trivially yet shuffles a non-uniform measure. Hypothesis (iii) is therefore a \emph{genuine additional commitment} beyond the framework's current postulates.

\textbf{Candidate principles for hypothesis (i).} Four candidate physical mechanisms could supply the Canonical IC Measure Principle. A 2026-06 circularity audit (GPT-5.5-pro consultation, cross-checked against the Lean Born modules) sharpened their relative status. The discriminating ingredient that every route must supply is one and the same: \textbf{additivity / non-contextuality of the selector measure} (equivalently, no-signaling of the outcome marginal under record refinement). This is not a neutral regularity condition: on the quantum effect algebra, positivity + normalization + that additivity is \emph{equivalent} to the Born trace form $\mu(E)=\mathrm{tr}(\rho E)$ (the machine-checked \texttt{EffectGleason.finite\_effect\_gleason}, all finite dimensions), and the $\alpha$-deformation $q_k \propto w_k^{\alpha}$ that satisfies every $\alpha$-blind structural axiom fails \emph{exactly} additivity (\texttt{RefinementBorn.sq\_not\_additive}). The candidates therefore differ not in \emph{whether} they need a Born-strength premise, but in whether that premise is stated transparently or is smuggled through structure that already encodes $|c_k|^2$.

*(α) \textbf{Canonical tracial typicality from CPW Type II structure.} The crossed-product Type II algebra $\hat{\mathcal{A}}(R)$ carries a canonical normal semifinite trace $\tau_R$ (unique up to scale), and it is tempting to read $\tau_R$ as the canonical IC measure. The audit flags this as \textbf{high-circularity}, and we accordingly retract the earlier "leading candidate" status. The trace alone yields the relative-dimension / maximally-mixed measure $\tau_R(P)/\tau_R(\mathbf 1)$, \textbf{not} the Born measure for a given $\Phi$; recovering Born requires the $\Phi$-normal state $\omega_\Phi(P)=\tau_R(h_\Phi P)$ with Haagerup/Radon–Nikodym density $h_\Phi=d\omega_\Phi/d\tau_R$, and the identification $\mu_\rho := \omega_\Phi$ \textbf{is} the Born postulate in algebraic dress (precisely the empirical-selection step isolated in §11.4.2). Moreover the existence and uniqueness of $\tau_R$ are $\alpha$-blind: they do not exclude $q_k \propto \omega_\Phi(P_k)^{\alpha}$. So (α) supplies the most QIQT-H-native \emph{language} but does not by itself force additivity.*

*(β) \textbf{Symmetric equiprobability} on a natural decomposition of $\Omega_\rho$ into record-fibers, with Born weights emerging from fiber-volume ratios (formalized for an equal-amplitude orthonormal fine-graining in \texttt{BornEquiprobable.lean}). The Zurek amplitude→count bridge is machine-checked: the \emph{uniform} measure over an equal-\textbf{norm} orthonormal fine-graining has outcome-marginal exactly $|c_k|^2$, by orthonormality alone (Pythagoras), discharging the previously-posited multiplicity $\mathrm{count}=M\,w_k$. But the load-bearing premises survive untouched: (i) \textbf{envariance} — why the measure is uniform over the equal-amplitude atoms — and (ii) \textbf{canonicity of the equal-norm fine-graining} (refinement-additivity, which is exactly what excludes the $\alpha$-family); the unequal-amplitude step itself \emph{is} refinement-additivity. So (β) relocates, but does not eliminate, the additivity premise.*

\emph{(γ) \textbf{Holographic / modular construction} from the canonical sector reference state $\sigma_R$ (vacuum on Minkowski, KMS on stationary thermal, Bunch-Davies on de Sitter, etc.) via a modular-theoretic construction yielding $\mu_\rho$. Physically attractive but the least developed; the audit notes it inherits (α)'s circularity, since the modular flow / $\Phi$-normal state already carries the $|c_k|^2$ content.}

*(δ) \textbf{Actuality Projective Consistency} (a selector-level principle, and the one \emph{syntactically} Born-free option). Upgrade $\lambda$ from a flat index over histories to a \textbf{projective/sheaf-like assignment over record refinements}, and impose: a local nondemolition (e.g. Spectrum-Broadcast) refinement $C'$ of a recorded coarse observable $C$ only \emph{refines} the selected history — it does not change its coarse restriction — pointwise in the IC, $s_C(\omega)=\pi\!\big(s_{C'}(\omega)\big)$ for the forgetful map $\pi:H_{C'}\to H_C$. Ordinary within-context additivity of any measure $\mu$ on the IC then \emph{pushes forward} to refinement-additivity of the outcome marginal, hence, with positivity, to Born (via \texttt{RefinementBorn.refinementNatural\_additive} composed with \texttt{EffectGleason}). The distinctive feature is that this premise mentions \textbf{no amplitudes, no $|c_k|^2$, no trace}: it is \emph{syntactically} Born-free, an ontic statement purely about the consistency of $\lambda$ under coarse/fine-graining of records. It is \textbf{not} thereby Born-\emph{weak} — it is the ontic form of no-refinement-signaling, and on the effect algebra it is Born-strength — but it isolates the missing premise at the selector level without smuggling the Born numbers, and is therefore the \textbf{most honest} statement of what must be assumed. Two caveats keep it honest: (a) it must be restricted to \emph{pure refinements of the same recorded coarse observable}; extended to arbitrary settings it would collide with the Parameter-Independence horn of §6.9 (deterministic parameter independence for all settings plus measurement independence is Bell-local and cannot reproduce the quantum correlations); and (b) its own \emph{weak} grounding fails — microcausality of $\Phi$ together with an inert, non-back-reacting $\lambda$ does \textbf{not} force no-signaling of $\mu$. The $\alpha=2$ rule is the explicit counterexample: it leaves $\Phi$ exactly unitary and $\lambda$ inert, yet a local binary split of a remote-correlated branch changes the remote marginal (a $\tfrac12\!\to\!\tfrac13$ shift). So (δ) names the irreducible premise sharply; it does not derive it.*

None of (α)–(δ) is derivable from FQ + AQFT + holography alone; each supplies — transparently in (δ), implicitly in (α)–(γ) — the single additivity / non-contextuality bridge that the \texttt{NoBornFromNothing} and \texttt{RefinementBorn.sq\_not\_additive} countermodels show to be \textbf{necessary and independent}. The position the framework adopts (§11.4) is accordingly that \emph{one} such bridge must be posited, the relativistic analogue of Bohmian quantum equilibrium and of the Everettian probability postulate; the contribution of this audit is to identify that bridge as additivity/non-contextuality, to retract the over-optimistic reading of the tracial route (α), and — via (δ) — to state the bridge in its most honest, amplitude-free, selector-level form. Open Problem 1 (§11.4) sharpens the canonical-measure problem with explicit sub-conditions.

\emph{Remark on Gleason.} Gleason's theorem derives the Born rule from probability-axiomatic constraints (noncontextuality + orthogonal additivity) on a projection lattice. It is \emph{not} the appropriate primary derivation for QIQT-H, which is deterministic per run and has no primitive probability assignment. Gleason functions as a consistency check on the \emph{target} empirical frequencies (Born is the unique probability rule consistent with the noncontextuality axioms), but cannot select the microscopic IC measure $\mu_\rho$, many distinct $\mu_\rho$ push forward to the same Born outcome distribution. Closing the gap requires the typicality argument above, not Gleason.

\subsection{What §7.1–7.4 establishes}

(FQ) restricts the model class properly (Theorem 2); it bounds the regional entropy $S(\rho_R) \le Q_R$ (Theorem 3); a single actual record per run then follows from decoherence (which fixes the record basis and removes interference) together with the non-dynamical selector $\lambda$ (P1), which makes exactly one record actual (Theorem 4); Born statistics emerge as an across-run typicality / relative-frequency result (Theorem 5, schematic). \textbf{(Revised 2026.} Earlier versions derived single-recordness instead from a "finite-information restriction" that allegedly made a $\ge 2$-record regional content \emph{non-instantiable} (§7.6) — that reading is \textbf{retired} as a category error, §1.1a: $Q_R$ bounds the von Neumann entropy $S(\rho_R)$, which a second branch raises by only $\sim\!\log K$, so capacity does \textbf{not} forbid a multi-record state, and no kinematic bound can \emph{select} an outcome in a unitary theory. The single-record conclusion of Theorem 4 stands, but via $\lambda$ + decoherence, \textbf{not} via capacity.)

The framework's structural skeleton: (FQ) is the entropy reading of the Bekenstein-Bousso bound, $S(\rho_R) \le Q_R$, stated rigorously in the CPW Type II algebraic framework; decoherence fixes the regional record basis and the selector $\lambda$ makes one record actual, with the run's microscopic initial conditions indexing \emph{which} record and the global evolution left exactly unitary. (The capacity bound $S(\rho_R)\le Q_R$ remains load-bearing kinematically — it yields the area floor and feeds the field equations — but it is \textbf{not} the mechanism of single outcomes.)

\subsection{The Modular-Local Holographic Superselection Rule (RETIRED — historical record)}

\begin{qiqtquote}
\textbf{Retired (2026) — see the Macroscopic Definiteness retraction in §1.1a.} This section's central postulate — that a $\ge 2$-record regional content has summed cost exceeding $Q_R$ and is therefore \emph{not an instantiable physical state} (the Branch-Summed Holographic Bound / Macroscopic Definiteness Conjecture) — is \textbf{withdrawn as a category error}: $Q_R$ bounds the von Neumann entropy $S(\rho_R)$, which a second macroscopic branch raises by only $\sim\!\log K$ (not by a further $\approx Q_R$), so a multi-record state does \textbf{not} overflow capacity; and no kinematic entropy bound can \emph{select} an outcome in a unitarily-evolving theory. Single-record definiteness comes from the selector $\lambda$ (P1) together with decoherence, \textbf{not} from this rule. The material below is retained for the historical record and for the parts that \emph{do} stand — the kinematic bound $S(\rho_R)\le Q_R$ and the no-signaling corollary (§7.7) — but its single-ness / superselection / "not instantiable" claims are superseded.
\end{qiqtquote}

\textbf{Reference state sector.} Fix once and for all a faithful locally normal reference state $\sigma$ on the quasilocal algebra, and choose the hatted local algebras coherently with this reference sector. In the Minkowski-vacuum sector $\sigma$ is the vacuum state; on stationary thermal backgrounds it is the geometrically preferred KMS state; and on de Sitter / cosmological backgrounds it is the Bunch–Davies, Hartle–Hawking, or other specified invariant reference state. For every bounded diamond $R$, write

\[
\sigma_R := \sigma|_{\hat{\mathcal{A}}(R)}.
\]

The family $\{\sigma_R\}$ is required to be \textbf{restriction-compatible}: whenever $R' \subset R$, isotony gives $\hat{\mathcal{A}}(R') \subset \hat{\mathcal{A}}(R)$, and

\[
\sigma_{R'} = \sigma_R |_{\hat{\mathcal{A}}(R')}.
\]

All relative entropies and modular Hamiltonians below are taken with respect to this $\sigma_R$. The notation $\Omega_R$, when it appears in vacuum-sector examples, denotes only the GNS vector representative of $\sigma_R$, not a second reference state.

The entropy quantity denoted $S_{\rm ren}$ in §4.1(ii) is the modular-local information functional used here:

\[
S_{\rm ren}(R; \omega) \;\equiv\; \chi_R(\omega_R) \;:=\; S^{\hat{\mathcal{A}}(R)}_{\rm Araki}(\omega_R \,\|\, \sigma_R).
\]

When the Type II crossed-product representation supplies a canonical trace and a modular Hamiltonian $K_R^\sigma$, this is equivalently written as

\[
\chi_R(\omega_R) \;=\; \Delta_\omega \langle K_R^\sigma \rangle - \Delta_\omega S_R.
\]

The holographic condition (FQ)(ii) is precisely the modular-local capacity bound $\chi_R(\omega_R) \le C(R)$.

The above theorems establish the framework's structure within standard renormalized entropy. They leave a critical question: even after decoherence + (FQ) deliver a strict classical mixture $\sum_k p_k \omega_k^R$ on the regional algebra, this is \emph{still a mixture over multiple macroscopic records}. The modular-local capacity bound controls the record-entropy / effective support of this mixture: for $k$ equiprobable records of per-record cost $I_0$, one typically has $\chi_R \simeq I_0 + \log k$, so raw cardinality is not excluded without additional branch-summed assumptions. To establish single-record per-run as a structural consequence at the level of the macroscopic observable content, the framework commits to a stronger principle than standard entropy bounds: a \textbf{modular-local holographic superselection rule}.

\emph{(As flagged in the retirement box above, the single-ness / superselection reading of this rule is withdrawn, §1.1a. What survives is the kinematic entropy bound $\chi_R \le C(R)$, one of the framework's five postulates, (P4); the irreducible new physics is (P4) + (P5) on the (P1) $(\Phi,\lambda)$ ontology, not this rule as an outcome-selector.)} The bound is formulated \emph{intrinsically} on the algebra-state pair via Araki / Type II core relative entropy, with no joint-region cutoff applied to spacelike-combined algebras.

\textbf{Definition (Regional information functional, restated).} \emph{The regional information functional $\chi_R$ is the Araki relative entropy on $\hat{\mathcal{A}}(R)$ with respect to the fixed reference $\sigma_R$ defined above. Equivalently, in the Type II crossed-product core with canonical semifinite trace $\tau_R$ and Haagerup densities $h_{\omega, R}$, $h_{\sigma, R}$:}

\[
\chi_R(\omega) \;=\; \tau_R\!\left[ h_{\omega, R} \left( \log h_{\omega, R} - \log h_{\sigma, R} \right) \right] + \text{counterterms},
\]

\emph{finite by the CPW/Witten construction (no UV divergence; no density-matrix assumption).}

\textbf{Postulate (Modular-Local Holographic Bound, the finite regional entropy bound (P4)).} \emph{The regional information functional is bounded by the holographic capacity, region by region:}

\[
\boxed{\chi_R(\omega) \;\le\; C(R) \;=\; \frac{A(\partial R)}{4\ell_P^2} \quad \text{for every bounded region } R \text{ and every physical state } \omega.}
\]

\emph{For spacelike-separated regions $D_A, D_B$, admissibility combines as the meet of local predicates:}

\[
\mathrm{Adm}(D_A \cup D_B) \;=\; \mathrm{Adm}(D_A) \wedge \mathrm{Adm}(D_B).
\]

\emph{No bound is imposed on the joint spacelike algebra $\hat{\mathcal{A}}(D_A) \vee \hat{\mathcal{A}}(D_B)$ beyond what the local restrictions already require.}

This is a \textbf{strengthening of the standard Bekenstein-Bousso bound}, not a derivation from it. Standard holographic entropy bounds limit the renormalized entropy of regional reduced states; the modular-local bound limits the \emph{relative} entropy with respect to a reference state, intrinsically defined via Araki/Connes machinery in Type III local QFT.

**Scope: what the bound does \emph{not} forbid, and what therefore does the work (important).** It is essential to be precise about the reach of any entropy-type bound, because a naive reading overclaims. The naive argument, "a two-record content must encode the joint information of both records, which exceeds $Q_R$", \emph{does not follow from the Bekenstein–Bousso bound, and we explicitly disavow it.} An entropy bound limits the number of mutually \textbf{distinguishable} records a region can carry ($\# \le e^{Q_R}$); it does \textbf{not} forbid a \textbf{superposition} of records. The decisive counterexample is a redundant cat state,

\[
|\Psi\rangle = \alpha\,|0\rangle^{\otimes N} + \beta\,|1\rangle^{\otimes N},
\]

which encodes two macroscopically distinct, redundantly copied records, fits in $N$ qubits, and is a \textbf{pure state of zero entropy}, violating no entropy bound whatsoever. More generally, by linearity of the exactly-unitary dynamics, if $|R_0\rangle$ and $|R_1\rangle$ are each individually instantiable in $R$, then so is $\alpha|R_0\rangle + \beta|R_1\rangle$, \emph{unless an additional superselection-type principle is imposed}, and any such principle stands in tension with the exact unitary linearity this framework otherwise preserves. (This is precisely why the single-record claim is a \emph{postulated superselection rule} above, a strengthening of standard holography, and not a corollary of it.)

Consequently the honest \textbf{division of labor} is: decoherence + einselection supply \emph{stable, robust, effectively Boolean} candidate records (the classical stage; this gives robustness, not uniqueness); the modular-local capacity bound makes the set of \emph{distinguishable} records \emph{finite} (cardinality, not superposition-exclusion); and the \textbf{uniqueness of the actual record is supplied by the non-dynamical actuality fact $\lambda$} (§1.0a, §11.3), not derived from $Q_R$. The framework's single-record-per-run content rests on $\lambda$ regardless of whether the superselection rule below can be established. A positive resolution of that rule would upgrade $\lambda$-selection from a postulate toward a structural consequence; a negative resolution leaves all empirical claims intact, with the selection mechanism explicitly postulated. In its strong "finite capacity \emph{forbids} macroscopic superposition" form, the rule is not merely unproven but, in light of the cat-state obstruction above, currently \emph{implausible} without additional nonstandard structure; this is recorded honestly as Open Problem 3.

\textbf{Branch-summed cost as derived approximation.} When $\omega_R$ is a strict classical mixture $\bar\omega_R = \sum_k p_k \omega_{k,R}$ over decoherent macroscopic records (as after decoherence + (FQ) at lab scale, §6.2–6.3), Donald's identity for Araki relative entropy gives

\[
\chi_R(\bar\omega_R) \;=\; \sum_k p_k \, c_R(r_k) - I_{\rm Hol}^R,
\]

where $c_R(r_k) = \chi_R(\omega_{k,R})$ is the per-record cost (relative entropy of branch $k$ to the reference $\sigma_R$) and $I_{\rm Hol}^R = \sum_k p_k S_{\rm Araki}(\omega_{k,R} \| \bar\omega_R)$ is the Holevo information of the mixture. In the strict distinguishable-record limit, $I_{\rm Hol}^R \approx H(\{p_k\})$. Effective branch multiplicity is thus controlled by the Holevo term that reappears in Theorem 6, not by an independent spacelike joint cutoff. \textbf{The earlier "branch-summed cost" $I_\Sigma^\epsilon$ is a derived classical-mixture approximation in this sense, useful for operational calibration (the $I_0$ parameter, Schrödinger-cat scale, exclusion curves) but not the fundamental statement of the bound.}

This formulation matters for two reasons:

\begin{enumerate}
  \item \textbf{Type III compatibility.} $\chi_R$ is defined directly in Type III via Araki; it does not require selecting a record subalgebra or smoothing parameter to be fundamentally defined. Branch counting is a coarse-graining of $\chi_R$ in the classical-mixture limit, not a new entropy concept.
  \item \textbf{No-signaling becomes a theorem.} Because $\chi_R$ depends only on the algebra-state pair $(\hat{\mathcal{A}}(R), \omega_R)$ and not on any joint-region structure, and because the admissibility predicate for spacelike-separated regions is the meet of local predicates, no-signaling follows automatically from AQFT microcausality. See §7.7.
\end{enumerate}

\textbf{The framework's commitment: the modular-local capacity bound is a superselection/admissibility rule, not a dynamical collapse rule and not a universal branch-summed cardinality postulate.}

Under this commitment, the framework's structure is the following.

\textbf{Definition (Stagewise causal admissibility).} \emph{Fix a global time function with Cauchy slices $\{\Sigma_t\}$. Let $\{\mathcal{F}_t\}$ be the classical filtration generated by outcome records whose localization diamonds lie in $J^-(\Sigma_t)$. For a branch/history $h \in \mathcal{F}_t$, let $\mathcal{O}_t(h)$ be the set of record diamonds realized on that branch up to time $t$.}

\emph{A bounded diamond $R$ is \textbf{causally instantiated} at stage $(t, h)$ if $R \Subset J^-(\Sigma_t)$ and either:}
\emph{(a) \textbf{Conditioned case} ($Y \ne \varnothing$): there exists a finite nonempty set $Y \subset \mathcal{O}_t(h)$ of records whose alternatives are being jointly constrained in $R$, with $R \Subset \bigcap_{O \in Y} J^+(O)$. This is the joint-record-comparison case where $R$ sits causally downstream of all records in $Y$.}
\emph{(b) \textbf{Unconditioned local case} ($Y = \varnothing$): $R$ has no past records constraining it. This case gives the ordinary local capacity constraint on $R$ alone, in force whenever $R$ has been physically instantiated (i.e., it is a finite causal diamond contained in $J^-(\Sigma_t)$ at the current stage).}
\emph{Denote the resulting collection of instantiated regions by $\mathfrak{R}_t(h)$.}

\emph{The \textbf{stagewise physical state space} is}

\[
\mathcal{S}_{\rm phys}(t, h) \;:=\; \left\{ \omega_t^h : (\omega_t^h)_R \in \hat{\mathcal{A}}(R)_*^+,\; \|(\omega_t^h)_R\| = 1,\; \chi_R((\omega_t^h)_R) \le C(R) \text{ for all } R \in \mathfrak{R}_t(h) \right\}.
\]

\emph{A \textbf{physical process} is an adapted family $h \mapsto \omega_t^h$ such that $\omega_t^h \in \mathcal{S}_{\rm phys}(t, h)$ for every stage $t$ and $\mathcal{F}_t$-almost every branch $h$.}

\emph{The unqualified notation $\mathcal{S}_{\rm phys}$ denotes this stagewise family unless a terminal completed-history limit is explicitly specified. Future joint-diamond constraints are therefore imposed only when the relevant region has become causally instantiated; they do not retroactively delete branches that were admissible at earlier stages.}

\emph{If a universal Hilbert-space representation is used, let $\pi(|\Psi\rangle) = \omega_\Psi$ be the induced algebraic state. Then $\mathcal{H}_{\rm phys}$ denotes the set of $|\Psi\rangle$ whose induced stagewise family $(\omega_\Psi)_t^h$ satisfies the stagewise condition above, modulo the equivalence $|\Psi\rangle \sim |\Phi\rangle \iff \omega_\Psi|_R = \omega_\Phi|_R$ for every $R$. This is a set of physical equivalence classes, not a linear subspace of $\mathcal{H}$.} States outside $\mathcal{S}_{\rm phys}$ are \emph{kinematically forbidden}, they are mathematically writeable in the unrestricted formalism but are not physically realizable.

\textbf{Definition (Physical dynamics and instruments).} \emph{A deterministic evolution is represented in the Heisenberg picture by a normal unital completely positive map $\alpha_t: \hat{\mathcal{A}} \to \hat{\mathcal{A}}$; for closed reversible dynamics, $\alpha_t$ is an automorphism. It is \textbf{physical} iff}

\[
\alpha_{t*}(\mathcal{S}_{\rm phys}) \subseteq \mathcal{S}_{\rm phys}, \qquad (\alpha_{t*}\omega)(A) := \omega(\alpha_t(A)).
\]

\emph{Here $\mathcal{S}_{\rm phys}$ is the set of physically instantiable regional contents (those satisfying the finite-information admissibility bound). The condition says only that physical evolution carries instantiable contents to instantiable contents; it is \textbf{not} a restriction on the admissible Hamiltonians. The global Schrödinger / Heisenberg law is exactly unitary and unconstrained, a formal global state whose regional content leaves $\mathcal{S}_{\rm phys}$ remains a valid calculational descriptor, it simply has no per-run physical referent in that region.}

\emph{A measurement instrument $\mathcal{I} = \{\mathcal{I}_a\}_a$ is \textbf{physical} iff $\sum_a \mathcal{I}_a$ is unital and, for every $\omega \in \mathcal{S}_{\rm phys}$:}
\emph{(i) the non-selective state $\omega \circ (\sum_a \mathcal{I}_a)$ lies in $\mathcal{S}_{\rm phys}$; and}
\emph{(ii) \textbf{branchwise}: for each outcome $a$, the corresponding durable classical record is localized in a bounded diamond $O_a$. Whenever $p_a = \omega(\mathcal{I}_a(\mathbf{1})) > 0$, the normalized conditional state}

\[
\omega_a(B) := \frac{\omega(\mathcal{I}_a(B))}{p_a}
\]

\emph{must satisfy the admissibility bound on every bounded diamond $R \Subset J^+(O_a)$:}

\[
(\omega_a)_R \in \hat{\mathcal{A}}(R)_*^+, \qquad \chi_R((\omega_a)_R) \le C(R).
\]

\emph{At finite process stage $t$, this condition is imposed only for those $R \Subset J^+(O_a) \cap J^-(\Sigma_t)$ that lie in the instantiated collection $\mathfrak{R}_t(a)$. Conditional feed-forward operations must satisfy the same condition on each branch.}

Equivalently: physical evolution preserves the set of instantiable regional contents. This is a closure property of $\mathcal{S}_{\rm phys}$ under the (exactly unitary, unconstrained) dynamics, not a selection rule on Hamiltonians, every standard Hamiltonian is admissible; what the framework restricts is which \emph{regional contents} count as physically instantiated, not which evolutions are allowed.

\textbf{Axiom (Causal application of admissibility).} \emph{Modular-local admissibility is imposed on the state of a region when that region is physically instantiated in the causal order. If earlier alternatives are separately admissible in their respective past algebras, a later joint diamond $D^+$ constrains only the state restricted to $\hat{\mathcal{A}}(D^+)$ after the relevant systems enter $D^+$. Failure of $\chi_{D^+} \le C(D^+)$ makes the proposed future joint state or transition inadmissible; it does \textbf{not} retroactively delete previously admissible past branches.}

\textbf{Theorem 6 (Effective Macroscopic Definiteness via Donald's identity, conditional form).} \emph{Throughout this theorem, $\chi_R(\varphi) := S_{\rm Araki}^{\hat{\mathcal{A}}(R)}(\varphi \| \sigma_R)$ denotes Araki relative entropy to the \textbf{fixed} regional reference state $\sigma_R$ of §7.6, and all states are normal on $\hat{\mathcal{A}}(R)$.}

\emph{Let $\omega \in \mathcal{S}_{\rm phys}$ and let $R$ contain a decohered macroscopic record algebra with mutually exclusive record predicates $\{P_k\}_{k \in \mathcal{A}_\epsilon}$ forming the $\epsilon$-smoothed active set, with $|\mathcal{A}_\epsilon| < \infty$ (finite effective support). Let $p_k = \omega_R(P_k)$, $q := \sum_{k \in \mathcal{A}_\epsilon} p_k > 0$, $\tilde p_k := p_k/q$. Let $\omega_{k,R}$ be the conditional regional state of branch $k$, $\bar\omega_R := \sum_k \tilde p_k \, \omega_{k,R}$, and}

\[
I_{\rm Hol}^R \;:=\; \sum_{k \in \mathcal{A}_\epsilon} \tilde p_k \, S^{\hat{\mathcal{A}}(R)}_{\rm Araki}(\omega_{k,R} \,\|\, \bar\omega_R), \qquad H_\epsilon \;:=\; -\sum_k \tilde p_k \log \tilde p_k.
\]

\emph{The exact \textbf{Donald's identity} holds for Araki relative entropy:}

\[
\sum_{k \in \mathcal{A}_\epsilon} \tilde p_k \, \chi_R(\omega_{k,R}) \;=\; \chi_R(\bar\omega_R) + I_{\rm Hol}^R. \tag{$\star$}
\]

\emph{Assume:}
\emph{(H1) \textbf{Branchwise admissibility.} Each active branch satisfies $\chi_R(\omega_{k,R}) \le C(R)$.}
\emph{(H2) \textbf{Record-instantiation cost} (independent framework postulate). The mean regional state has cost $\chi_R(\bar\omega_R) \ge I_0 - \eta_0$, where $I_0$ is the per-record cost parameter and $\eta_0 \ge 0$ a tolerance.}
\emph{(H3) \textbf{Operational distinguishability.} There exists a normal measurement instrument on $\hat{\mathcal{A}}(R)$ decoding the active record index $K \in \mathcal{A}_\epsilon$ with average error probability $P_e$. Standard information theory gives $H_\epsilon = I(K;Y) + H(K \mid Y)$ where $Y$ is the decoded outcome. The Holevo bound yields $I(K;Y) \le I_{\rm Hol}^R$, and Fano's inequality gives $H(K \mid Y) \le h_2(P_e) + P_e \log(|\mathcal{A}_\epsilon| - 1) =: \eta_{\rm def}$. Hence $H_\epsilon \le I_{\rm Hol}^R + \eta_{\rm def}$.}

\emph{Then Donald's identity $(\star)$ together with (H1) and (H2) gives}

\[
I_{\rm Hol}^R \;=\; \sum_k \tilde p_k \chi_R(\omega_{k,R}) - \chi_R(\bar\omega_R) \;\le\; C(R) - I_0 + \eta_0.
\]

\emph{(Since $I_{\rm Hol}^R \ge 0$, this is informative only when $C(R) - I_0 + \eta_0 \ge 0$. If $C(R) < I_0 - \eta_0$, the hypotheses (H1) and (H2) are jointly unsatisfiable, the region cannot support even a single record at its specified cost, so no admissible record-bearing state on $R$ exists.)}

\emph{and with (H3),}

\[
\boxed{H_\epsilon \;\le\; C(R) - I_0 + \eta_0 + \eta_{\rm def}, \qquad N^{(\epsilon)}_{\rm eff} := \exp H_\epsilon \;\le\; \exp(C(R) - I_0 + \eta_0 + \eta_{\rm def}).}
\]

\emph{Here $\eta_{\rm def} = h_2(P_e) + P_e \log(|\mathcal{A}_\epsilon| - 1)$ is the explicit Fano residual.}

\emph{Moreover, for every $0 < \delta < 1$, there exists $T_\delta \subseteq \mathcal{A}_\epsilon$ with normalized active probability $\ge 1 - \delta$ and}

\[
|T_\delta| \;\le\; \exp\!\left(\frac{C(R) - I_0 + \eta_0 + \eta_{\rm def}}{\delta}\right).
\]

\emph{If the leakage sector $\mathcal{A}_\epsilon^c$ is retained, the theorem bounds only the conditional active entropy $H_\epsilon$. The full unnormalized entropy decomposes as $H(p) = h_2(q) + q H_\epsilon + (1-q) H(p \mid \mathcal{A}_\epsilon^c)$, so an additional finite-support or entropy assumption on the leakage sector is required to bound $H(p)$ itself.}

\emph{At \textbf{exact} saturation $I_0 = C(R)$ and $\eta_0 = \eta_{\rm def} = 0$, one obtains $H_\epsilon = 0$: a single record has normalized active probability one. For \textbf{finite} tolerances, the conclusion weakens to $H_\epsilon \le \eta_0 + \eta_{\rm def}$, one record dominates with probability $\ge 1 - O(\eta_0 + \eta_{\rm def})$. For approximate saturation, the conclusion is \textbf{effective / probability-weighted} rather than a raw cardinality bound.}

\emph{Proof.} Donald's identity $(\star)$ is a standard property of Araki relative entropy on a von Neumann algebra (Donald 1986; cf. Ohya-Petz 1993, Thm 5.21). Under (H1), the LHS of $(\star)$ is bounded by $C(R)$. Under (H2), the term $\chi_R(\bar\omega_R)$ in the RHS of $(\star)$ is at least $I_0 - \eta_0$. Rearranging gives the bound on $I_{\rm Hol}^R$. Applying (H3) gives the bound on $H_\epsilon$. The smoothed support bound follows by Markov's inequality on $-\log \tilde p_k$, and $\tilde p_{\max} \ge e^{-H_\epsilon}$ gives the saturation conclusion. $\blacksquare$

\textbf{Status of the assumptions.} (H1) is a direct consequence of $\omega \in \mathcal{S}_{\rm phys}$ provided each $\omega_{k,R}$ inherits admissibility, natural under the branchwise instrument condition of §7.6. (H3), operational distinguishability, is the formal expression of macroscopic record decodability, established under decoherence + Quantum Darwinism for the macroscopic-record subalgebra (§6.4–6.7). (H2) is the \emph{independent framework postulate}: it asserts that physical instantiation of a macroscopic record on $R$ requires modular cost $\ge I_0$. This is \emph{not} a theorem of standard AQFT or holography; it is the framework's central commitment, calibrated empirically against the observed quantum-to-classical transition scale.

\textbf{Sharpening of (H2), reference-weight form.} A machine-verified audit (see \texttt{lean/mathlib/QIQTH/H1H2Audit.lean} in the project repository) confirms that (H2) is genuinely independent of (H1), Donald's identity, Klein positivity, and DPI: a classical KL countermodel with $\sigma=(1/2,1/2)$ and $\rho=\delta_0$ satisfies all the structural hypotheses but violates (H2) with $I_0=1$. The same audit also yields the following sharp structural reformulation of (H2). For a perfect-record state on event $E_{\rm record}$,

\[
\chi_R(\bar\omega_R) \ge I_0 - \eta_0 \quad \Longleftrightarrow \quad \sigma_R(E_{\rm record}) \le e^{-(I_0 - \eta_0)}.
\]

Equivalently: (H2) is the assertion that \textbf{macroscopic record sectors occupy exponentially small reference-weight} under the canonical sector reference state $\sigma_R$. This identifies the load-bearing physical content of (H2): not the positivity of relative entropy (which Klein already provides), but a reference-weight bound on pointer sectors that a future first-principles derivation must obtain from modular/holographic structure on macroscopic records.

\textbf{Conditional modular estimate of Born-weight deviations.} Let $\mu$ be the Born measure over conditional record states $r \mapsto \omega_r$, and define the inadmissible set

\[
F_R := \{r : \chi_R(\omega_r) > C(R)\}, \qquad \delta_R := \mu(F_R).
\]

Since $\chi_R \ge 0$, Markov's inequality gives

\[
\delta_R \;\le\; \frac{\mathbb{E}_\mu[\chi_R(\omega_r)]}{C(R)}.
\]

The conversion of $\chi_R$ into a stress-energy estimate is \emph{not} assumed for a generic QFT ball. It is used only when the reference modular Hamiltonian is known geometrically, or when $R$ is compared to such a region by monotonicity of relative entropy.

\emph{Wedge case (Bisognano-Wichmann).} In the Minkowski-vacuum sector, i.e., with $\sigma_R$ taken to be the vacuum / geometric-modular reference state corresponding to a Rindler wedge $W \supset R$, the Bisognano-Wichmann theorem gives, for states normal to the vacuum representation and in the domain of the boost / modular Hamiltonian, the \emph{exact first-moment expression}

\[
\Delta\langle K_W^\sigma\rangle \;=\; \frac{2\pi}{\hbar c} \int_W x^1 \, \Delta\langle T_{00}\rangle(x) \, d^{d-1}x.
\]

\emph{(This formula is specific to the vacuum / geometric-wedge modular setting; for arbitrary $\sigma_R$ no such universal local form is available.)}
The simplified estimate $\Delta\langle K_W^\sigma\rangle \le 2\pi L E_W/(\hbar c)$ requires: (i) finite renormalized energy; (ii) finite first moment; (iii) the excitation has support effectively in $0 \le x^1 \le L$ (split / smeared localization to avoid UV pathologies of sharp localization); (iv) nonnegative renormalized energy density / positive energy measure in that region. If these conditions fail, use the exact first-moment expression rather than the $L E_W$ bound.

By monotonicity of Araki relative entropy under restriction,

\[
\chi_R(\omega_R) \;\le\; \chi_W(\omega_W) \;=\; \Delta\langle K_W^\sigma\rangle - \Delta S_W.
\]

Hence $\chi_R(\omega_R) \le 2\pi L E_W/(\hbar c) + s_W$ provided the entropy shift satisfies $\Delta S_W \ge -s_W$. The commonly used estimate $\chi_R(\omega_R) \le 2\pi L E_W/(\hbar c)$ requires the additional hypothesis $\Delta S_W \ge 0$ along with (i)–(iv) above.

\emph{CFT-vacuum ball case.} For a ball $B_L$ in a CFT vacuum, the conformal modular Hamiltonian gives

\[
K_{B_L}^\sigma \;=\; \frac{2\pi}{\hbar c} \int_{|x| < L} \frac{L^2 - |x|^2}{2L} T_{00}(x) \, d^{d-1}x,
\]

and the same conclusion holds, with the corresponding entropy-shift qualification.

\emph{Generic QFT ball.} For a non-conformal QFT ball, no universal local formula for $K_R^\sigma$ is assumed; the bound is obtained only by enclosing-wedge monotonicity as above.

Under these hypotheses, the Born-deviation estimate reads

\[
\delta_R \;\lesssim\; \frac{2\pi L \, \mathbb{E}_\mu[E_R]/(\hbar c) + s_R}{A(\partial R)/(4\ell_P^2)}.
\]

If $s_R$ is negligible, $A(\partial R) \simeq 4\pi L^2$, $L \simeq 1\,\mathrm{m}$, and $E_R \simeq 1\,\mathrm{kg}\,c^2$, this gives the \textbf{conditional} order-of-magnitude estimate

\[
\delta_R \;\sim\; 10^{-27}.
\]

This number is therefore \textbf{not a generic QFT consequence}; it is a wedge / CFT-ball modular-energy estimate with the stated entropy-shift hypothesis. The deviation theorem (§7.4) then implies that Born probabilities under modular-local admissibility deviate from standard Born probabilities by at most $\delta_R$ on bounded observables, operationally invisible at lab scale, conditional on the hypotheses above.

The single-record-per-region structure of physical states is therefore a \emph{structural consequence} of the modular-local capacity bound. \textbf{Effective form (Theorem 6, the precise statement).} Multi-record states are not absent because of dynamical collapse; under approximate saturation $I_0 \approx C(R)$ they are \emph{exponentially suppressed} on the normalized active distribution (effective multiplicity $N^{(\epsilon)}_{\rm eff} = \exp H_\epsilon \le \exp(C(R) - I_0 + \eta_0 + \eta_{\rm def})$), and at exact saturation $I_0 = C(R), \eta_0 = \eta_{\rm def} = 0$ they are \emph{kinematically forbidden} (single record with active probability one). The strong-language passages below ("not in the physical state space", "$N_{\max} \approx Q_R/I_0$") are to be read in the exact-saturation limit; the precise effective form for finite tolerances is Theorem 6.

\textbf{The adopted reading: a restriction on the physical state space, not on the dynamics.} Standard Hilbert-space QM is the unconstrained formalism. The constrained set $\mathcal{H}_{\rm phys}$, regional contents with $I_\Sigma \le Q_R$, is a \emph{nonlinear set of admissible regional contents} (the possible values of $C_R$), \textbf{not} an invariant Hilbert subspace in which the global $\Phi(t)$ evolves; the constraint $I_\Sigma \le Q_R$ is nonlinear in the state. The restriction is on \emph{which regional contents are instantiable physical states}, \textbf{not} a restriction on the allowed Hamiltonians and \textbf{not} a claim that "certain unitaries do not occur." The Schrödinger / Heisenberg evolution law is exactly unitary and entirely standard, including ordinary measurement interactions. What the formalism's linearity produces, the global multi-record vector $\sum_k c_k |A_k\rangle|E_k\rangle$, remains a perfectly good \emph{calculational} descriptor; the framework's claim is that its regional content in $R$, carrying $\ge 2$ macroscopic records at summed cost $I_\Sigma \approx K \cdot I_0 > Q_R$, \textbf{cannot be instantiated in its entirety as the actual content of $R$}.

It is essential to be precise about what this does and does not deliver. The capacity bound does \emph{negative} work only: it forbids a $\ge 2$-record regional content. It does \textbf{not}, by itself, pick out which single record is the actual one, and decoherence cannot do so either (it conserves the branch weights $|c_k|^2$; \texttt{NoConcentration.lean}). The actual single-record content of $R$ in a given run is the derived beable

\[
C_R = A_R[\Phi, \lambda],
\]

the single record sector selected by $\lambda$, the run's actual microscopic initial conditions of apparatus + environment. Here $A_R$ is a \emph{fixed, non-dynamical reading functional}, part of the interpretive postulate (not a per-run free object): given $\Phi(t)$, the run-index $\lambda$, and the decoherence-defined record decomposition of region $R$, it returns one admissible record sector $C_R(t)$. It carries no run-dependent degrees of freedom beyond $\lambda$, does not alter $\Phi$, and is required to be consistent across overlapping/nested regions and across time. Thus the per-run ontology is the pair $(\Phi, \lambda)$ (the exactly-unitary global $\Phi$ plus the non-dynamical actuality selector $\lambda$), and the fixed theoretical structure comprises the unitary law, the capacity bound, the typicality measure, and the reading functionals $\{A_R\}$, lawlike structure, not a further per-run substance. $C_R$ is not a third substance and is not $\Phi$ restricted to $R$ (that restriction is the full $K$-record mixture); it is the $\lambda$-selected sector, eliminable in favour of $(\Phi, \lambda, A_R)$ but \textbf{not} in favour of $\Phi$ alone. Two honest consequences follow. First, $\lambda$ is a \emph{broad-sense} beable, an actuality fact not contained in bare $\Phi$, so the framework is ψ-monist only in the \emph{weak (dynamical)} sense ($\Phi$ is the sole dynamical ontology), not in the strong sense that $\Phi$ alone is the complete per-run state. Second, the selection role of $\lambda$ and the typicality measure it carries are genuine additional commitments (consistent with the \texttt{EquivarianceGap.lean} audit), not consequences of (FQ) alone. This is (FQ) part (iii) plus the formal/per-run distinction, a statement about physical instantiability of regional contents and their $\lambda$-selection, with the global evolution law left untouched.

\textbf{Response to the linear measurement obstruction.} A standard objection: "By linearity, $U(\alpha|0\rangle + \beta|1\rangle)|A_{\rm ready}\rangle = \alpha|A_0\rangle + \beta|A_1\rangle$, exactly the forbidden multi-record state."

This objection assumes that the unrestricted-Hilbert-space wave function $\alpha|A_0\rangle + \beta|A_1\rangle$ is the physically primary object, the actual physical state that we have direct access to. The framework denies this. \textbf{We do not have direct physical access to unrestricted-Hilbert-space states; we have access only to the physical observable content of bounded regions, the state on the macroscopic record subalgebra $\mathcal{C}(R)$, with the (FQ) precision floor.}

The screen is a complex quantum system with $\sim 10^{20}$–$10^{25}$ degrees of freedom that becomes entangled with the particle during detection. Decoherence between distinct macroscopic record states $|A_0\rangle$ and $|A_1\rangle$ produces environmental cross-overlaps $\langle E_0 | E_1 \rangle \sim e^{-10^{20}}$ on physically realistic timescales (§6.2), so the records become non-interfering (§6.3); the selector $\lambda$ (P1) then marks one actual.

So at the level of physical observable content, which is what physics is about, the formal $\alpha|A_0\rangle + \beta|A_1\rangle$ structure decoheres to a classical-looking expression, and the residual sub-$\epsilon$ coherence has no physical referent under (FQ) (a resolution-equivalence on regional states). The physical content on the macroscopic observable algebra is a strict classical mixture over records. The finite-information restriction then says: a regional content carrying $\ge 2$ macroscopically distinct records (summed cost $I_\Sigma \approx K\cdot I_0$ approaching/exceeding $Q_R$) is \emph{not an instantiable physical state of the region}, it has no per-run physical referent.

It is essential that this is a restriction on which \emph{regional contents} are physically instantiable, \textbf{not} a modification of the dynamics. The formal global vector $\alpha|A_0\rangle|E_0\rangle + \beta|A_1\rangle|E_1\rangle$ produced by ordinary linear unitary evolution is a perfectly good \emph{calculational descriptor}, it is the cross-run ensemble descriptor, not the per-run physical state. What is physically real per run is the regional content, and by the finite-information restriction that content is single-record. \emph{Which} record a given run carries is fixed by that run's actual microscopic initial conditions (distinct runs are distinct actual universes; the IC index the run, they do not produce the single-ness). The global Schrödinger / Heisenberg evolution law is left exactly unitary and entirely standard throughout; no Hamiltonian is restricted and no amplitude is trimmed.

The "linear measurement obstruction" thus disappears once one recognizes that (a) physical content is the regional algebra-state content of bounded regions, not the unrestricted global Hilbert-space vector; (b) decoherence removes interference and (FQ) removes the residual sub-$\epsilon$ coherence as a resolution-equivalence; (c) the finite-information restriction makes a $\ge 2$-record \emph{regional content} non-instantiable, the multi-record global vector remains a valid ensemble descriptor but has no single-region per-run physical referent.

\textbf{The per-record cost $I_0$ as an experimental parameter.} The framework's central physical postulate, the Branch-Summed Holographic Bound $I_\Sigma^\epsilon[\omega_R] \le Q_R$, combines two quantities:
\begin{itemize}
  \item $Q_R = A(\partial R)/(4\ell_P^2)$: the \textbf{regional holographic capacity}, a geometric quantity determined by the boundary area in Planck units. Set by quantum gravity; not a free parameter.
  \item $c_R(r) \approx I_0$: the \textbf{per-record physical cost}, the information cost of specifying a macroscopic record's full microscopic configuration. This is an \textbf{experimental parameter} of the theory.
\end{itemize}

The number of coexisting macroscopic records allowed per region is $N_{\max} \approx Q_R/I_0$. The empirically observed boundary at which macroscopic superpositions cease to occur (the quantum-to-classical transition scale) corresponds to the scale at which $N \cdot I_0$ approaches $Q_R$ for the relevant macroscopic record cost $I_0$.

\textbf{This is directly analogous to GRW's parameter structure.} GRW has a collapse rate $\lambda$ that is a free parameter calibrated empirically: small enough that microscopic quantum coherence is preserved, large enough that macroscopic superpositions decohere on observed timescales. QIQT-H has $I_0$ playing an analogous role: small enough that microscopic records (atoms, qubits, small molecules) have $N \cdot I_0 \ll Q_R$ and standard QM behavior is recovered, large enough that macroscopic records have $N \cdot I_0$ approaching $Q_R$ so the superselection enforces single-record per run.

Unlike GRW, where $\lambda$ is ad hoc, in QIQT-H $I_0$ has a \textbf{physical interpretation}: the information cost of specifying a macroscopic record's full microscopic configuration in the regional Type II algebra. This is in principle calculable from microphysics (Zurek physical entropy, decoherent-history weights, Quantum Darwinist redundancy), though current best estimates give $I_0 \sim 10^{25}$ bits for typical macroscopic records, far below $Q_R \sim 10^{70}$ for a 1m region.

The framework's empirical content includes the prediction: \textbf{the per-record cost $I_0$ must be calibrated such that $N \cdot I_0$ approaches $Q_R$ at the observed quantum-to-classical transition scale.} If standard Zurek entropy gives $I_0 \sim 10^{25}$ bits and laboratory macroscopic records are still well below the regional capacity, then the framework predicts the observable quantum-to-classical boundary involves macroscopic records whose cost is much larger than naive Zurek estimates, including all entangled environmental degrees of freedom out to the relevant decoherence horizon. The framework's specific prediction is that this calibration converges to a definite value of $I_0$ for macroscopic record types, and this value is testable against the observed scale at which Schrödinger-cat-like experiments fail.

The two-parameter structure of the theory:
\begin{center}
\small
\begin{tabularx}{\linewidth}{|Y|Y|Y|}
\hline
\textbf{Parameter} & \textbf{Origin} & \textbf{Status} \\
\hline
$Q_R = A(\partial R)/(4\ell_P^2)$ & Holographic principle of quantum gravity & Geometric; not adjustable \\
\hline
$I_0 = c_R(r)$ for typical macroscopic records & Per-record physical cost in regional Type II algebra & Experimental; calibrated empirically \\
\hline
\end{tabularx}
\end{center}

The single-outcome enforcement threshold occurs at $N \cdot I_0 \approx Q_R$. Below this threshold (microscopic regime, $N \cdot I_0 \ll Q_R$): constraint vacuous, standard QM behavior. Above this threshold (macroscopic regime): constraint enforces single-record per run.

A clarifying contrast with gauge theory: in gauge theory the physical state space is a constrained submanifold (gauge-invariant states) and one \emph{also} restricts the admissible Hamiltonians to gauge-preserving ones. QIQT-H is \textbf{not} of this form. The restriction here is on which \emph{regional contents} count as instantiable physical states (those with $I_\Sigma \le Q_R$); the global evolution law is left exactly unitary and unconstrained, and the formal global wave function, including formal multi-record vectors, evolves by ordinary standard QM. The Branch-Summed Bound constrains what can physically \emph{reside} in a region, not which Hamiltonians may act.

\textbf{The Schrödinger law is exactly unitary and unmodified.} There is no restriction on Hamiltonians and no non-unitary or non-linear term anywhere. What the framework adds is a kinematic statement about regional physical instantiability: a regional content carrying $\ge 2$ macroscopically distinct records has no per-run physical referent. This is (FQ) part (iii) plus the formal/per-run distinction, not a dynamical modification.

\textbf{The price the framework pays.} The commitment is a new \emph{kinematic postulate}, the Branch-Summed Holographic Bound on instantiable regional content, not a change to the dynamics. It is substantive at the foundational level, but only in regimes where the branch-summed cost approaches the holographic capacity. At lab scales, where macroscopic records use only $\sim 10^{25}$ bits of an available $\sim 10^{70}$-bit holographic capacity per cubic meter, the restriction is operationally vacuous and standard QM is recovered exactly. The restriction becomes operationally relevant only when branch-summed cost approaches $Q_R$, i.e., for macroscopic measurement records, where its content is to make $\ge 2$-record regional content non-instantiable, single-record per run at the kinematic level, with the dynamics still exactly unitary.

\textbf{Mathematical work needed:}

\begin{enumerate}
  \item \textbf{Precise specification of $\mathcal{C}(R)$} as the einselected/Darwinistic record subalgebra (drawing on decoherent histories, Quantum Darwinism).
  \item \textbf{Per-record cost $c_R(r)$} rigorously defined via Zurek physical entropy.
  \item \textbf{Branch-summed bound $I_\Sigma^\epsilon \le Q_R$} as new physical principle, not derivable from existing holography, but conjecturally connectable to deeper quantum-gravity arguments about distinguishable record content per region.
  \item \textbf{Characterization of instantiable regional content}: a precise statement of which regional algebra-states satisfy $I_\Sigma \le Q_R$, and the proof that $\ge 2$-record content violates it (the Macroscopic Definiteness Conjecture, Open Problem 3). This is a constraint on regional content, not on Hamiltonians; the dynamics is left exactly unitary.
  \item \textbf{Born statistics from typicality}: the typicality theorem for which single-record content is realized per run, with the across-run realization frequency reproducing Born weights $|c_k|^2$ (probability emergent, not a per-run primitive).
\end{enumerate}

Each of these is a concrete open mathematical problem; together they constitute the framework's explicit research program beyond the borrowed CPW/Witten scaffolding.

\textbf{Why this dissolves the MWI tension.} Under the Branch-Summed Superselection Postulate (Theorem 6), the framework escapes the MWI-without-many-worlds problem cleanly:

\begin{itemize}
  \item The unconstrained Hilbert space contains states with multi-record macroscopic content; these are mathematically writeable but not in $\mathcal{H}_{\rm phys}$
  \item Physical states (those in $\mathcal{H}_{\rm phys}$) have $I_\Sigma \le Q_R$ for every bounded region; with $I_0 \approx Q_R$ at macroscopic scales, this means single-record per region
  \item The "Everett branches" are mathematical artifacts of considering unrestricted Hilbert-space states; the actual physical state space excludes them by the superselection rule
  \item We don't need to "select" one branch from many physically real branches, there are no multi-branch physical states to select from
\end{itemize}

The framework's single-world per run is therefore a \emph{kinematic structural feature} of the physical state space, not a dynamical selection event. The superselection rule is what makes this true.

\textbf{Why this is genuinely new physics beyond Witten/CPW.} Witten/CPW provide the Type II algebraic infrastructure for regional generalized entropy. They do \emph{not} establish:
\begin{itemize}
  \item The branch-summed bound as a strengthening of standard holographic entropy
  \item The kinematic non-instantiability of $\ge 2$-record regional content (a restriction on which regional contents are physical, not on the dynamics)
  \item The exclusion of multi-record regional content as kinematically forbidden
\end{itemize}

These are the framework's specific new physical principles, building on the Witten/CPW scaffolding. They constitute a concrete research program: define the branch-summed cost rigorously; postulate the Branch-Summed Bound on instantiable regional content as new physics; prove that $\ge 2$-record content is excluded; derive Born statistics as an across-run typicality / relative-frequency result. The global evolution law remains exactly unitary throughout.

\subsection{Theorem (No-signaling from modular-local admissibility)}

\textbf{Non-selective-instrument convention.} All operations used in the no-signaling argument below are deterministic \emph{non-selective} channels, i.e., instruments enter only through their sum $\mathcal{I} = \sum_a \mathcal{I}_a$. Selective conditional states $\omega_a$ are admissibility-checked branchwise (cf. §7.6 Definition (Physical instruments)), but \textbf{postselection on the outcome label $a$ is not an operation available for spacelike signaling}; it becomes operational only once the outcome record is classically available in the common future. Without this restriction, even ordinary Born probabilities can be made to "signal" by postselecting on rare outcomes; with it, no-signaling reduces cleanly to algebraic commutation.

\textbf{Theorem 7 (No-signaling).} \emph{Let $D_A, D_B$ be spacelike-separated causally complete regions with $[\hat{\mathcal{A}}(D_A), \hat{\mathcal{A}}(D_B)] = 0$. Let $\{\Phi_a^x\}$ be a normal CP instrument localized in $\hat{\mathcal{A}}(D_A)$ with setting $x$ and outcomes $a$, and let $\{\Psi_b^y\}$ be similarly localized in $\hat{\mathcal{A}}(D_B)$. Under the modular-local finite-information admissibility restriction (§7.6) and the non-selective-instrument convention above, Alice's marginal probability is independent of Bob's setting:}

\[
P_{\rm QIQT}(a \mid x, y) \;=\; P_{\rm QIQT}(a \mid x).
\]

\emph{Proof.} Bob's nonselective channel is $\Psi^y = \sum_b \Psi_b^y$, localized in $\hat{\mathcal{A}}(D_B)$ by construction. By locality of this deterministic channel (i.e., $\Psi^{y*}$ is implemented by operators in $\hat{\mathcal{A}}(D_B)$, equivalently it acts trivially on $\hat{\mathcal{A}}(D_B)'$) together with microcausality $[\hat{\mathcal{A}}(D_A), \hat{\mathcal{A}}(D_B)] = 0$, the Heisenberg dual $\Psi^{y*}$ acts as the identity on Alice's algebra:

\[
\Psi^{y*}(X) \;=\; X \qquad \forall X \in \hat{\mathcal{A}}(D_A).
\]

Alice's outcome effect is $E_a^x = \Phi_a^{x*}(\mathbf{1}) \in \hat{\mathcal{A}}(D_A)$. Therefore:

\[
P_{\rm QIQT}(a \mid x, y) \;=\; \omega\!\left(\Psi^{y*}(E_a^x)\right) \;=\; \omega(E_a^x),
\]

which is manifestly independent of $y$.

The admissibility predicate does not alter this conclusion, because (§7.6) admissibility on spacelike-separated regions is the \emph{meet of local predicates}: $\mathrm{Adm}(D_A \cup D_B) = \mathrm{Adm}(D_A) \wedge \mathrm{Adm}(D_B)$. Alice's local admissibility predicate $\mathrm{Adm}(D_A)$ depends only on $\omega|_{\hat{\mathcal{A}}(D_A)}$, which Bob's nonselective operation does not affect. Hence Alice's branch admissibility cannot depend on $y$. $\blacksquare$

\textbf{Remarks.}

\begin{enumerate}
  \item The proof uses \emph{only} microcausality $[\hat{\mathcal{A}}(D_A), \hat{\mathcal{A}}(D_B)] = 0$, which is the bedrock of relativistic AQFT. It does \emph{not} assume tensor factorization $\mathcal{A}(D_A \cup D_B) \cong \mathcal{A}(D_A) \,\bar\otimes\, \mathcal{A}(D_B)$, which Type III local QFT generally lacks because of vacuum boundary entanglement.
  \item The key structural feature is that the admissibility predicate is a \textbf{local subfunctor} of the AQFT state functor: $\mathrm{Adm}(D) \subseteq \mathrm{States}_{\rm normal}(\hat{\mathcal{A}}(D))$ is defined entirely from data on $\hat{\mathcal{A}}(D)$, with no joint-region cutoff. Imposing a hard cutoff on the joint relative entropy $S_{\hat{\mathcal{A}}(D_A) \vee \hat{\mathcal{A}}(D_B)}(\omega \| \Omega)$ would \emph{not} satisfy the modular-local form and would re-introduce signaling, because the vacuum is not a product state across spacelike boundaries (the difference between joint and additive relative entropies is precisely the mutual-information / boundary-entanglement contribution).
  \item \textbf{Joint causal diamonds containing record-comparison events} ($D_{AB}$ containing both Alice's and Bob's records, encountered when records meet in a common future) carry their own capacity $C(D_{AB})$ and their own modular-local admissibility predicate $\mathrm{Adm}(D_{AB})$, which constrains the \emph{single} algebra-state pair $(\hat{\mathcal{A}}(D_{AB}), \omega|_{D_{AB}})$. No additional cross-coupling is imposed. The admissibility of the joint record is the relative entropy of the joint state on the joint algebra, not a sum of separate Alice-cost plus Bob-cost.
  \item \textbf{Operational rule on instruments}, not on branches: a physical instrument $\{\Phi_a^x\}$ must be branchwise admissibility-preserving, if $\omega \in \mathrm{Adm}(D)$, then for every nonzero branch $a$ the post-outcome state $\omega \circ \Phi_a^* / \omega(\Phi_a^*(\mathbf{1}))$ must also be in $\mathrm{Adm}(D)$. Inadmissible outcomes are not produced because the dynamics that would produce them is not a physical operation, \emph{not} because branches are postselected away after the fact. This preserves the superselection (kinematic-exclusion) flavor of the framework.
\end{enumerate}

This theorem retires the earlier no-signaling concern arising from joint-comparison-diamond admissibility. An explicit Bell-style counter-example with hard joint-diamond cutoff and asymmetric record costs motivated the modular-local reformulation; the meet-of-local-predicates structure resolves it.

\medskip\hrule\medskip

\section{The Bekenstein-Bousso Bound and the Entropy Reading}

\subsection{The bound's pedigree}

Bekenstein (1981), 't Hooft (1993), Susskind (1995), Bousso (2002), and most recently Banks (2025) and the CPW (2022) Type II construction. The bound has multiple independent motivations across the quantum-gravity literature.

The numerical scale: $Q_R \sim 10^{70}$ nats for a one-meter region, $\sim 10^{122}$ for the observable universe. Astronomically large, but finite, and area-law.

\subsection{The entropy reading as the natural reading}

The historical conflation of "the holographic bound" with "an entanglement-entropy inequality on a fictitious tensor factorization" is an artifact of working in finite-dimensional toy models or in spin-chain QFTs where Hilbert-space factorization is artificially imposed.

In genuine continuum QFT, local algebras are Type III$_1$. There is no entanglement entropy of a region's reduced state; there is no reduced state. The narrow interpretation is mathematically ill-defined.

The entropy interpretation, that the bound limits the renormalized von Neumann entropy $S(\rho_R)$ of regional states in the algebraic sense, is the mathematically natural reading once the Type II crossed-product construction is recognized as the rigorous home of regional generalized entropy.

(FQ) is then a precise mathematical statement: the renormalized entropy on $\hat{\mathcal{A}}(R)$ is bounded by $Q_R$.

\subsection{Wald and higher-derivative corrections}

For modified gravitational theories, the bound generalizes via the Wald entropy. The Type II framework extends naturally: the trace on the regional algebra is determined by the gravitational theory, and $Q_R$ takes the Wald form $A f'(R)/(4\ell_P^2)$ for $f(R)$ gravity.

\medskip\hrule\medskip

\section{Possible Phenomenological Consequences}

\subsection{Laboratory regime: no deviations}

At ordinary laboratory scales, $Q_R$ is astronomically large, and the resolution structure of the Type II algebra is fine relative to the macroscopic record structure of any laboratory apparatus. The framework predicts no deviations from standard QM in standard laboratory experiments. The role of (FQ) at lab scales is structural, it provides the resolution that makes per-run states physically single-record.

\subsection{Long-baseline regimes}

Long-baseline neutrino oscillation and other extended-coherence systems may approach regimes where the Type II resolution structure becomes operationally relevant. The current IceCube and Super-Kamiokande bounds on quantum-gravity-induced decoherence ($\Gamma_Q/E \lesssim 10^{-32}\,\mathrm{GeV}^{-1}$) are consistent with the framework. Rigorous derivation of phenomenological predictions from the Type II structure is deferred to subsequent work.

\subsection{Cosmological and horizon regimes}

The de Sitter static patch and black-hole horizon regions are the natural arenas where Type II$_1$ algebras (with finite trace $\tau(\mathbf{1}) < \infty$) directly apply. CPW developed the Type II$_1$ structure precisely for de Sitter. The cosmological extensions of (FQ) are the natural next step and are deferred to companion papers.

\medskip\hrule\medskip

\section{Comparison with Existing Approaches}

\subsection{Copenhagen}

Invokes collapse without dynamics. We invoke no collapse. The single-record per-run regional content is a consequence of decoherence (removing interference) + the non-dynamical selector $\lambda$ (P1) marking exactly one decohered record actual, not of dynamical reduction. (Single-ness is $\lambda$'s work, not a capacity exclusion — §1.1a.)

\subsection{Many-Worlds}

Posits all formal components as equally real branches. We are single-world per run: the per-run regional content is single-record because the non-dynamical selector $\lambda$ (P1) marks exactly one decohered record actual; the unrealized records have no physical referent (they are not "real elsewhere"), in contrast to Many-Worlds.

\subsection{Bohmian mechanics}

Bohm adds particle positions in $3N$-dim configuration space as a \emph{second dynamical layer} beyond the wave function, with a guidance equation connecting the two, the configuration evolves under a law driven by $\Phi$. The framework adds \emph{no second dynamical layer and no guidance law}: $\Phi$ is the sole dynamical ontology (weak ψ-monism; §2.2). The structural counterpart of the Bohmian configuration is the run-index $\lambda$, the run's actual microscopic initial conditions, which select the realized record, but $\lambda$ is \emph{non-dynamical}: it carries no equation of motion and exerts no back-reaction on $\Phi$. So the difference from Bohm is sharp: same single-world realism and a run-indexing actuality fact, but no extra guided substance and no modification of the dynamics. The complete per-run state is the pair $(\Phi, \lambda)$; the formal wave function is the textbook ensemble descriptor.

\subsection{Objective collapse (GRW, CSL, DP, OR)}

Modify Schrödinger with collapse terms. We modify nothing dynamical. The Schrödinger evolution of the underlying field algebra is exactly unitary; the regional physical content is single-record because the non-dynamical selector $\lambda$ (P1) marks exactly one decohered record actual — a fact about actuality, not a modification of the dynamics.

\subsection{QBism, relational quantum mechanics}

Treat the wave function as epistemic / perspectival. We are objectively realist: the per-run wave function is real, the (FQ) constraint is real, the Type II algebra structure is real, the regional states are real physical states.

\subsection{Modal interpretations, decoherent histories}

Introduce a modal value-rule or history-selection structure. We require neither. Decoherence fixes the regional record basis and the non-dynamical selector $\lambda$ (P1) marks exactly one record actual, yielding single records without external value-rules (capacity is kinematic, not exclusionary — §1.1a).

\subsection{Ensemble interpretations (Ballentine)}

The closest neighbor. Ballentine treats the formal wave function as ensemble-descriptive with per-run states underdetermined. We supply the mechanism: the Type II regional algebra structure makes per-run states physically single-record despite the formal wave function's superposed appearance.

\subsection{Palmer's Rational Quantum Mechanics / Invariant Set Theory (nearest neighbour on motivation, opposite horn on Bell)}

Palmer's Rational Quantum Mechanics (RaQM 2025, PNAS) and the broader Invariant Set Theory (IST) program is the closest neighbor in the space of \emph{single-world, structural-principle, finite-information} frameworks. Both treat the wave function as not exhausting the per-run physical content, and both seek to ground the measurement problem's resolution in a deep structural feature of physics rather than in collapse, branching, or particle ontology. \textbf{On the Bell question, however, the two take opposite horns:} Palmer rejects Bell's measurement-independence assumption (a measurement-dependence / superdeterminism-family escape, principled rather than conspiratorial), whereas the present framework \emph{keeps} measurement independence and rejects superdeterminism, paying instead at Parameter Independence at the ontic level (§6.9). The comparison below is therefore a contrast on the central Bell move, not an alliance.

\textbf{Shared structural commitments:}
\begin{itemize}
  \item Single-world realism per run
  \item Recognition that the textbook formal wave function does not exhaust the per-run physical content
  \item Motivation by deep structural principles rather than ad hoc modification of QM
\end{itemize}

\textbf{The decisive difference (Bell horn):} Palmer denies measurement independence (a measurement-dependence / superdeterminism-family escape, with a non-conspiratorial rejection of metaphysical counterfactual recombination); the present framework \emph{keeps} measurement independence and rejects superdeterminism, denying Parameter Independence at the ontic level instead (§6.9). This is the opposite Bell move, not a shared one.

(The frameworks differ in ontology: Palmer treats the wave function as derived/emergent from invariant-set statistics, with the underlying ontology being points on the fractal invariant set $I_U$; our framework is weak-ψ-monist, with $\Phi$ the sole dynamical ontology plus a non-dynamical run-index $\lambda$ selecting the realized record, the per-run state $(\Phi,\lambda)$ being finer-grained than the textbook formal ensemble descriptor. Interestingly Palmer's $I_U$-point and our $\lambda$ play structurally similar run-indexing roles.)

\textbf{Differences in mathematical apparatus and motivating principle:}

\begin{center}
\small
\begin{tabularx}{\linewidth}{|Y|Y|Y|}
\hline
\textbf{} & \textbf{Palmer's RaQM / IST} & \textbf{Present framework} \\
\hline
Mathematical apparatus & p-adic measures; Niven's theorem; fractal invariant set $I_U$ in cosmological state space & CPW Type II crossed-product algebra; finite-precision physical-instantiation postulate \\
\hline
Bell horn & Denies measurement independence (counterfactual settings off the invariant set not realizable; Impossible Triangle Corollary), superdeterminism-family & Keeps measurement independence; denies Parameter Independence at the ontic level; rejects superdeterminism (§6.9) \\
\hline
Status of wave function & Derived / emergent from invariant-set statistics; not fundamentally physically real & Physically real but finite-precision (ψ-ontic) \\
\hline
Status of Schrödinger evolution & Emergent from chaotic dynamics on invariant set & Preserved exactly at the underlying field-algebra level \\
\hline
Connection to quantum gravity & Indirect: chaos-based, cosmological motivation & Direct: Bekenstein-Hawking entropy, Witten/CPW Type II structure \\
\hline
Empirical predictions & Cosmological (CMB anomalies, fine-tuning, dark matter); explicit Bell-correlation models & Empirically conservative at lab scales; possible long-baseline signatures; cosmological extensions deferred \\
\hline
\end{tabularx}
\end{center}

\textbf{Relationship: same motivation, opposite Bell horn.} The two frameworks share single-world realism and a finite-information motivation, but they are \emph{not} the same Bell escape: Palmer pays at measurement independence, the present framework pays at Parameter Independence (§6.9). They are therefore contrasting positions on the central Bell move, not interchangeable. (They are not in direct conflict as research programs, one could imagine a hybrid in which an invariant-set structure characterizes which per-run configurations are realized while a holographic bound constrains regional record capacity, but on the Bell question itself they diverge.)

\textbf{Where the frameworks complement each other's weaknesses:} Palmer's program is more developed (PNAS-published, decade of prior literature, explicit cosmological predictions, concrete Bell-model with Niven's theorem) but uses a more idiosyncratic mathematical machinery (p-adic measures, rational/irrational distinction). Our framework uses mainstream algebraic QFT + holographic bound (lower entry cost for the foundations community) but is at the proposal stage with the Concentration Conjecture and Born-typicality program as open work.

\textbf{For readers familiar with Palmer:} do \emph{not} read the present framework as "Palmer with holography", that conflates the two on Bell. Palmer takes the measurement-dependence horn; the present framework takes the Parameter-Independence horn while keeping free settings (§6.9). The shared content is single-world realism plus a finite-information motivation; the Bell move is opposite.

\subsection{The algebraic-foundational program}

Our framework can be read as an algebraic foundational program: the macroscopic world is the regional Type II algebra structure, the dynamics is the algebraic evolution, and the (FQ) bound is the holographic constraint on regional information capacity. This connects naturally to:

\begin{itemize}
  \item Algebraic QFT (Haag-Kastler)
  \item Generalized entropy program (Wall, Engelhardt-Wall, CPW)
  \item Quantum-gravity ideas about emergence (Verlinde, Jacobson, Padmanabhan)
\end{itemize}

\subsection{Summary comparison}

\begin{center}
\small
\begin{tabularx}{\linewidth}{|Y|Y|Y|Y|Y|Y|}
\hline
\textbf{Approach} & \textbf{Modifies dynamics} & \textbf{Hidden particle state} & \textbf{Branches} & \textbf{Resolution structure} & \textbf{Per-run state} \\
\hline
Copenhagen & Yes (collapse) & No & No & None & Formal + collapse \\
\hline
Bohm & No & Yes (positions) & No & None & WF + particles \\
\hline
MWI & No & No & Yes (all real) & None & All formal branches \\
\hline
GRW/CSL/DP/OR & Yes (stochastic) & No & No & None & Formal until stoch collapse \\
\hline
QBism / RQM & No & No & No & None & Epistemic / perspectival \\
\hline
Modal & No & No & No & None & Formal + value rule \\
\hline
Decoherent histories & No & No & No & None & Formal + history selection \\
\hline
Ballentine ensemble & No & No & No & None & An ensemble member, unspecified \\
\hline
Palmer RaQM / IST & No (emergent dynamics) & No particles; per-run microstate on invariant set & No & Invariant-set discreteness (p-adic) & Specific microstate on $I_U$; measurement-dependent \\
\hline
\textbf{(FQ) entropy-bound framework} & \textbf{No} & No second \emph{dynamical} ontology (weak ψ-monist); one non-dynamical run-index $\lambda$, formal-ψ-incomplete & \textbf{No} & \textbf{Finite regional entropy $S(\rho_R) \le Q_R$ + decoherence + $\lambda$-selection} & \textbf{One single-record regional content $C_R = A_R[\Phi,\lambda]$: $\lambda$ selects the single actual record (capacity is kinematic, not exclusionary — §1.1a); dynamics exactly unitary} \\
\hline
\end{tabularx}
\end{center}

The (FQ) entropy-bound framework is characterized by a \textbf{finite regional entropy budget from the holographic bound}: the pre-selection state $\Phi$ has regional von Neumann entropy $S(\rho_R) \le Q_R$; combined with decoherence (which renders the records non-interfering), the non-dynamical selector $\lambda$ (P1), and microscopic initial conditions (which index the run), this yields a single-record per-run outcome, without collapse, without modified dynamics, without hidden particles, without branches. The CPW Type II algebra framework provides the rigorous mathematical infrastructure; the entropy bound plus the $(\Phi,\lambda)$ ontology is the additional foundational content that does the work.

\medskip\hrule\medskip

\section{Conclusion}

\subsection{Summary}

We have developed a foundational framework for quantum mechanics that combines two ingredients: (a) the algebraic-QFT-plus-gravitational-dressing framework of Chandrasekaran-Penington-Witten and Witten, providing the rigorous mathematical infrastructure for finite renormalized entropy on bounded regions; and (b) the \emph{finite regional entropy reading} of the Bekenstein-Bousso bound (postulated as axiom (FQ)), that the pre-selection state $\Phi$ has regional von Neumann entropy $S(\rho_R) \le Q_R$ (its effective modular rank $e^{S(\rho_R)}$), together with the $(\Phi,\lambda)$ ontology in which the non-dynamical selector $\lambda$ fixes the single actual record.

The structure rests on:

\begin{enumerate}
  \item \textbf{The CPW/Witten algebraic infrastructure (borrowed).} Local algebras in QFT with gravitational dressing are Type II with semifinite trace and well-defined renormalized entropy differences matching generalized entropy. This provides the rigorous algebraic home for the bound.
  \item \textbf{The (FQ) axiom as a finite regional entropy budget (our postulate).} The pre-selection state $\Phi$ in any bounded region $R$ has regional von Neumann / renormalized entropy bounded by the holographic capacity, $S(\rho_R) \le Q_R = A(\partial R)/(4\ell_P^2)$ (equivalently, its effective modular rank is at most $e^{Q_R}$). This is an entropy bound on $\Phi$, not a finite matter Hilbert-space dimension and not a finite-precision reading of amplitudes. Witten supplies the mathematical scaffold; the absolute entropy bound is the additional postulate.
  \item \textbf{Finite regional entropy (Theorem 3).} $S(\rho_R) \le Q_R$: the load-bearing statement is this von Neumann / renormalized entropy bound. (An earlier version stated a subsidiary finite amplitude-resolution floor $\epsilon(R) > 0$; that reading is superseded, §1.1a.)
  \item \textbf{The decoherence-plus-$\lambda$-selection mechanism.} Under standard unitary dynamics, decoherence removes interference between macroscopic record components (their weights $|c_k|^2$ are conserved, not driven to $0$ or $1$), and it is robust; the capacity bound (FQ) provides the kinematic arena, bounding the entropy/cardinality of records $S(\rho_R) \le Q_R$; and the non-dynamical selector $\lambda$ (P1) makes the resident regional content single-record, with the run's microscopic initial conditions indexing which record (Theorem 4). Capacity does not forbid a $\ge 2$-record content (§1.1a).
  \item \textbf{The Born typicality program (open).} A conjectural program for deriving Born statistics from typicality of microscopic initial conditions; analogous to but distinct from Bohmian $|\psi|^2$-equivariance, identified as the central open quantitative problem.
\end{enumerate}

We have proved four theorems: model-class restriction (Theorem 2), finite regional entropy (Theorem 3), single-record per-run wave functions (Theorem 4), and Born from typicality (Theorem 5, schematic).

\subsection{What the framework does and does not deliver}

\textbf{Delivers:}
\begin{itemize}
  \item A foundational position in which the macroscopic world emerges from the holographic entropy bound + decoherence + the non-dynamical selector $\lambda$ + actual run-specific initial conditions, without collapse, hidden particles, branching, or modal value-rule; the ontology is the pair $(\Phi,\lambda)$
  \item A precise foundational axiom (FQ) as a finite regional entropy bound $S(\rho_R) \le Q_R$ on the pre-selection $\Phi$
  \item A rigorous mathematical home for the bound via the CPW/Witten Type II algebra framework
  \item A finite regional entropy bound $S(\rho_R) \le Q_R$ (Theorem 3)
  \item A single-record per-run consequence (Theorem 4) following from decoherence + the selector $\lambda$ (P1), with the microscopic initial conditions indexing which record
  \item A clearly identified open program for Born statistics from typicality
\end{itemize}

\textbf{Does not deliver:}
\begin{itemize}
  \item A quantitative form of the (superseded, §1.1a) subsidiary resolution floor $\epsilon(R)$ as an explicit function of regional geometry; this is not needed for the framework's conclusions, which rest on the entropy bound $S(\rho_R) \le Q_R$
  \item A rigorous typicality theorem reproducing Born weights from explicit measure
  \item A proof that the absolute entropy reading is the only defensible reading of the bound (it is, we argue, the natural reading; alternative narrow readings are possible and have different structural consequences)
  \item Lab-scale predictions deviating from standard QM (the framework is empirically conservative)
\end{itemize}

\subsection{The mathematical status of the framework: what we have and what we don't}

A central question for any foundations paper: do we have the mathematics, or are we hand-waving? We answer honestly.

\textbf{Rigorous mathematical infrastructure (borrowed from the QFT / quantum-gravity literature):}

\begin{center}
\small
\begin{tabularx}{\linewidth}{|Y|Y|Y|}
\hline
\textbf{Component} & \textbf{Status} & \textbf{Source} \\
\hline
Local algebras in QFT as Type III$_1$ factors & Rigorous theorem & Buchholz, Wichmann, Borchers, Longo \\
\hline
Algebraic QFT framework (Haag-Kastler nets) & Rigorous & Haag 1992; Haag-Kastler 1964 \\
\hline
Crossed-product Type II construction with gravitational dressing & Rigorous (in CPW's setting) & Witten 2022; CPW 2022; CLPW 2022 \\
\hline
Semifinite trace on Type II crossed-product algebras & Rigorous & Murray-von Neumann; standard operator algebras \\
\hline
Renormalized entropy on Type II algebras matching generalized entropy differences & Rigorous (in CPW's setting) & CPW 2022; Jensen-Sorce-Speranza 2023 \\
\hline
\end{tabularx}
\end{center}

This is the \emph{mathematical scaffolding} on which the framework is built. It is real mathematical physics, currently under active development in the quantum-gravity literature.

\textbf{The framework's foundational axiom (our postulate, not derived from the above):}

The (FQ) axiom is stated precisely in §4.1. It postulates a finite regional entropy budget: the pre-selection state $\Phi$ has regional von Neumann / renormalized entropy bounded by the holographic capacity, $S(\rho_R) \le Q_R = A(\partial R)/(4\ell_P^2)$. The binding ontological reading is that $Q_R$ is an entanglement entropy of the pre-selection $\Phi$, and geometry responds to $\Phi$, never to the actualized branch selected by $\lambda$. This is a foundational postulate, not a theorem; it is not a finite-precision reading of amplitudes and not a finite matter Hilbert-space dimension. (Per the P4-MICRO refinement, §1.1a, the irreducibly postulated \emph{content} is just that the regional capacity is \emph{finite} — the area \emph{floor} $S_{\rm vN}(\rho_R)\le Q_R$ is itself a derived theorem, the area \emph{form} $\propto A$ is from the Sakharov bridge, and only the \emph{value} of $G$ is carried.)

\textbf{Qualitative consequences worked out:}

\begin{itemize}
  \item The entropy bound $S_{\rm vN}(\rho_R) \le Q_R$ (§5.1; the subsidiary amplitude-resolution floor $\epsilon(R) > 0$ is superseded, §1.1a)
  \item Decoherence → the records become non-interfering (argued qualitatively with numerical scaling: decoherence gives $\sim e^{-10^{20}}$ off-diagonal coherence)
  \item Per-run single-record outcome from decoherence + the selector $\lambda$ (P1), with the microscopic initial conditions indexing which record (Theorem 4)
\end{itemize}

These are stated with reasoning but not proved as explicit theorems with formal proofs.

\textbf{What we do not have (open mathematical work):}

\begin{center}
\small
\begin{tabularx}{\linewidth}{|Y|Y|}
\hline
\textbf{Open problem} & \textbf{What's needed} \\
\hline
Explicit form of $\epsilon(R)$ & A function $\epsilon(R) = f(Q_R, \text{record subalgebra dimension}, \text{geometry})$ \\
\hline
Decoherence + (FQ) → strict mixture & A theorem in algebraic QFT showing that, under (FQ), decoherent suppression below $\epsilon$ gives a strict classical mixture on the regional algebra \\
\hline
Born-typicality theorem & A measure $\mu$ on the space of universal initial conditions, plus an equivariance-type theorem $\mu\{\lambda : k(\lambda, \Psi) = k\} = \vert c_k\vert ^2$ \\
\hline
FQ as operational equivalence (not literal thresholding) & A clean mathematical formulation: states differing on the regional algebra by less than $\epsilon$ (in some operational metric like trace distance) are physically equivalent \\
\hline
Connection to Bell correlations & An explicit $(\Phi,\lambda)$ model reproducing CHSH $=2\sqrt2$ via ontic Parameter-Independence violation, keeping free settings and yielding operational no-signaling without fine-tuning (§6.9) \\
\hline
\end{tabularx}
\end{center}

\textbf{Honest characterization of math status:} the framework provides the \emph{scaffolding} (Type II algebras from CPW/Witten) and the \emph{axiom} ((FQ), the finite regional entropy bound). It establishes \emph{qualitative} consequences. It does not yet have \emph{explicit theorems} for the central claims. This is the status of a research program, not a completed mathematical theory.

\textbf{Updated formal-verification status.} The characterization above remains correct \emph{about the physics}: the central \emph{physical} claims (the (FQ) entropy bound, the canonical IC/Born-typicality measure, and $\lambda$'s dynamical Lorentz-covariant law) are not theorems and remain the research agenda. (Macroscopic Definiteness is \emph{not} among them — it is retired as a category error, §1.1a, not an open target.) What has changed since earlier drafts is the \emph{formal} status of the \textbf{deductive core}, which is now machine-checked and \textbf{axiom-free} (0 project-specific axioms, no \texttt{sorry}; CI budget 0, verified by \texttt{scripts/axiom\_budget\_check.sh}). Three layers must be held apart:

\begin{center}
\small
\begin{tabularx}{\linewidth}{|Y|Y|}
\hline
\textbf{Layer} & \textbf{Formal status} \\
\hline
Conditional/structural deductive core (Theorems 3, 6, 7; Lemma 1; Donald's identity; the no-signaling chain; the finite no-collapse Born representation; the negative audits) & Machine-checked, \textbf{axiom-free} (standard Lean axioms only) \\
\hline
Physical postulates ((FQ)/P4; Canonical IC / Born Principle (P5); Lorentz covariance of $A_R[\Phi,\lambda]$) & \textbf{Open} — not theorems; axiom-free Lean does not bear on their truth (the Macroscopic Definiteness Conjecture is \emph{retired}, §1.1a — not an open postulate) \\
\hline
Continuum realization (Type II crossed-product entropy scaling; Fock / Weyl / CCR modular flow) & \textbf{Formalization frontier}, in active development (\texttt{QIQTH/Fock}, \texttt{Spectral}, \texttt{Entropy}) \\
\hline
\end{tabularx}
\end{center}

A subtlety the axiom \emph{count} alone cannot see: a vacuous hypothesis (an interface axiom whose antecedent is trivially true) can trivialize a theorem without any axiom being declared. The project guards against this with a vacuity lint (\texttt{scripts/vacuity\_lint.sh}) alongside the axiom-budget check; one historical instance (a \texttt{True}-antecedent locality axiom) was caught and converted to an explicit hypothesis, and the current lint reports a single benign site (an indiscrete-preorder definition).

The open problems are \emph{concrete and well-defined}: each can be attacked in principle by an analyst willing to engage with both algebraic QFT and foundations of QM. We identify them not as defects of the framework but as the explicit research agenda the framework opens.

\textbf{Comparison with other foundations programs:} Bohmian mechanics has a complete mathematical theory (configuration space, guidance equation, $|\psi|^2$-equivariance) but adds a second ontology and faces relativistic incompatibility. Many-Worlds keeps standard QM but multiplies branches and lacks a derivation of Born from amplitudes. GRW modifies dynamics with specific parameters and faces empirical bounds on those parameters. QIQT-H, by comparison, has a partial mathematical theory (scaffolding + axiom + qualitative consequences) with explicit open theorems; in exchange, it preserves Schrödinger evolution exactly, adds no extra ontology, and is grounded in holography and quantum gravity. The mathematical work needed to complete it is well-defined, not vague.

\subsection{The central thesis, in one paragraph}

\begin{qiqtquote}
The Bekenstein-Bousso holographic bound, taken literally as a Lorentz-invariant physical information limit on bounded regions of spacetime, bounds the number of mutually \emph{distinguishable} macroscopic records a region can carry ($\# \le e^{Q_R}$). We are explicit that this is a \textbf{capacity (cardinality) bound, not a superposition-exclusion rule}: by linearity of the exactly-unitary dynamics it does \emph{not}, on its own, forbid a macroscopic superposition of records, the redundant zero-entropy cat state $\alpha|0\rangle^{\otimes N}+\beta|1\rangle^{\otimes N}$ is the obstruction (§7.6). Whether finite capacity plus einselection can be strengthened into a genuine macroscopic \emph{superselection} rule consistent with unitarity is the open Macroscopic Definiteness Conjecture (Open Problem 3); its strong "forbids superposition" form is currently \emph{implausible}, and the framework's single-outcome content does \textbf{not} depend on its resolution. The Schrödinger / Heisenberg evolution law of the underlying field algebra is left \textbf{exactly unitary and entirely standard} (including ordinary measurement interactions); no amplitude is trimmed and no Hamiltonian is restricted. \textbf{The per-run ontology is the pair $(\Phi, \lambda)$:} an exactly-unitarily-evolving global wavefunction $\Phi(t)$, together with a fixed \emph{non-dynamical, atemporal actuality selector} $\lambda$, a single global fact about which macroscopic realization the run is, fixed for the whole 4-dimensional history (not past-localized "initial data," to avoid any Conway–Kochen / setting-correlation reading; see §6.9). For each region $R$ the actual regional content is the derived beable $C_R(t) = A_R[\Phi(t), \lambda]$, the single $\lambda$-selected decoherent record sector of $\Phi(t)$; the unselected sectors remain in $\Phi(t)$ but are not regionally instantiated. The mechanism has a clean \textbf{division of labour}: decoherence removes interference and einselects stable, robust, effectively Boolean records (conserving the weights $|c_k|^2$; proved \emph{not} to select an outcome, \texttt{NoConcentration.lean}), this gives robustness, not uniqueness; the capacity bound makes the set of distinguishable records \emph{finite}, cardinality, not selection; and \textbf{$\lambda$ supplies which single record is actual}, this is where outcome-uniqueness genuinely comes from. No collapse dynamics, no Bohmian particle-position ontology or guidance law, no Everettian parallel \emph{actual} branches (the unrealized records are not a parallel reality, they remain in $\Phi$ but are not regionally instantiated), and \textbf{no modification of the dynamics} are added; the only actuality rule is the fixed $\lambda / A_R$ reading rule, not a stochastic or modal dynamics. The framework is \textbf{ψ-monist in the weak (dynamical) sense}: $\Phi$ is the only thing that dynamically evolves and the only ontology subject to a law of motion, but the complete description of a run is the pair $(\Phi, \lambda)$, with $\lambda$ a non-dynamical actuality fact (a broad-sense beable, structurally analogous to the configuration in Bohmian mechanics or the actual history in modal interpretations, but with no guidance law and no back-reaction on $\Phi$). There are \textbf{no fundamental single-run probabilities}, no per-run stochastic or collapse chance is postulated; the Born weights $|c_k|^2$ are to be recovered, \emph{if at all}, as an emergent relative frequency across the $\lambda$-indexed ensemble of runs, conditional on a typicality measure $\mu$ over actuality selectors that is \textbf{currently unconstructed / interface-level} (the Born-typicality program, §7.4, §11.4, the framework's central open problem, together with the requirement that the same $\mu$ enforce operational no-signaling without fine-tuning, §6.9). The CPW Type II crossed-product algebra framework provides the rigorous mathematical infrastructure; the holographic capacity bound is the framework's new physical input on top. The single-record-per-region claim is therefore an \emph{architecture} in which decoherence + holography furnish a finite stage of robust classical records and $\lambda$ makes one actual, \textbf{a decomposition of the measurement problem into a robustness part, a finiteness part, and two clearly-stated primitives ($\lambda$ actuality, $\mu$ typicality), not a completed derivation of single outcomes from the bound alone.}
\end{qiqtquote}

\subsection{Open problems}

\begin{qiqtquote}
\textbf{Audit status (current).} The discharge narrative in this section is historical. As of the latest audit the project is \textbf{axiom-free}: 0 project-specific axioms, no \texttt{sorry}, CI budget 0 (\texttt{scripts/axiom\_budget\_check.sh}), 2,263 \texttt{\#print axioms} directives over 259 modules, within a full library of 299 files and roughly 3,300 theorems, full build green. Every interface axiom mentioned below as "remaining" has since been discharged in its finite-dimensional realization — the Goldstein-Struyve Schur classification (\texttt{schur\_classification\_real}) and tensor multiplicativity, Klein's inequality, Donald's identity, the data-processing inequality, the Araki/\texttt{EntropyBridge} relative-entropy layer, the Tsirelson-attainability bound, and the Fano / \texttt{RelEntPositivity} finite steps — and the \texttt{TypicalityMackeyGleason} packaging was retired in favour of the axiom-free \texttt{EffectGleason}. The figures "57 → 40 → 37 → 35" below trace that effort; it has since reached \textbf{0}. What remains genuinely open is the \emph{continuum} (von Neumann-algebraic) version of these results and the \emph{physical} postulates (Open Problems 1–3b); neither is touched by the finite-dimensional formal discharge.
\end{qiqtquote}

The Lean audit work (see \texttt{lean/mathlib/QIQTH/} and the formal-verification footnote in the abstract) has sharpened the boundary between what the framework derives from FQ + AQFT + holography and what remains as load-bearing additional commitments. The open problems are reorganized here to reflect that sharpened boundary, with the Born / canonical-measure problem promoted to first position as the most central remaining commitment.

\textbf{Open Problem 1, Canonical IC Measure Principle (the central Born/typicality problem).}

This problem decomposes naturally into six layers, each requiring different kinds of work. The decomposition is based on a closing audit (\texttt{lean/mathlib/QIQTH/AxiomAudit.lean} + GPT-5.5-pro consultation) that distinguishes what is mathematically derivable from what is genuinely open physics from what is irreducible empirical calibration.

\textbf{Goal.} Specify, for each preparation state $\rho$, a measure $\mu_\rho$ on the QIQT-H microscopic-IC space $\Omega_\rho$ (alternatively: a typicality structure on the FQ-restricted physical Hilbert space) satisfying:
   (a) \textbf{Canonicality}, defined from QIQT-H primitives, not fitted per measurement;
   (b) \textbf{Born pushforward}, $(O_M)_* \mu_\rho(i) = \mathrm{tr}(\rho E_i)$ for every allowed measurement protocol $M$ with effects $\{E_i\}$;
   (c) \textbf{Equivariance / stationarity}, preserved by the (FQ)-restricted physical Hamiltonian;
   (d) \textbf{Repeated-trial typicality}, supports iid / exchangeable / stationary-ergodic frequency theorems;
   (e) \textbf{Measurement-setting independence}, $\mu_\rho$ does not depend on the measurement context (unless the theory explicitly accepts contextual / measurement-dependent typicality);
   (f) \textbf{Uniqueness / stability}, robust under coarse-graining and equivalent IC descriptions.

Machine-verified audits flag the difficulty: \texttt{NoBornFromNothing.lean} (any distribution realizable; the structural axioms do not pick out Born), \texttt{EquivarianceGap.lean} (support preservation $\neq$ Born equivariance), \texttt{OperationalNoGo.lean} (operational frequency data alone insufficient). The independence package \texttt{BornMinimalityTable.lean} re-exports these three countermodels and shows that the locality premise is \emph{not} an independent Born-selection sub-axiom, its marginal-level content is a theorem holding for any measure (\texttt{MarginalLocality.pushforward\_marginal\_local}), modulo a named Hilbert→set-level bridge axiom (§11.4.5). Relative to the current finite formal decomposition, the Canonical IC Measure Principle's irreducible Born-selection content is three named sub-axioms (canonical measure, equivariance, operational sufficiency), each provably necessary by concrete finite countermodel, alongside the two acknowledged operator-algebra interface axioms (Schur classification, tensor multiplicativity) that the Born representation itself rests on.

\emph{\textbf{The recommended route: a Gleason spine (GPT-5.5-pro μ-options consultation).}} Of the six candidate constructions for $\mu$, (i) POVM-Gleason / Mackey-Gleason representation, (ii) decoherent-histories quantum measure + a typicality postulate, (iii) envariance / self-locating symmetry (Zurek; Sebens–Carroll), (iv) Bohmian quantum-equilibrium equivariance (Dürr–Goldstein–Zanghì), (v) frequency-operator (Hartle; Farhi–Goldstone–Gutmann), (vi) Deutsch–Wallace decision theory, only \textbf{(i)} plausibly yields an \emph{objective, Lorentz-equivariant} $\mu$ meeting all of (a)–(f), with operational no-signaling following structurally from \emph{one} canonical $\mu$ rather than per-experiment tuning. The recommended architecture is therefore: \textbf{Gleason as the spine} (the representation theorem that \emph{forces} the weights), the \textbf{decoherence functional as the cylinder construction} ($\mu_\Phi(\mathrm{Cyl}\,\alpha) = D_\Phi(\alpha,\alpha) = \langle\Phi|C_\alpha^\dagger C_\alpha|\Phi\rangle$, extended by Kolmogorov–Carathéodory), and a \textbf{frequency theorem as the LLN corollary}; envariance enters only as optional symmetry motivation.

The decisive conceptual point: the Born content is \emph{not} contained in any single record/Boolean algebra, a single measurement context admits \textbf{any} probability vector on its outcomes (this is exactly the \texttt{NoBornFromNothing} countermodel). What forces Born is \textbf{non-contextuality across overlapping effect contexts} together with state-certainty, the Gleason hypotheses. Concretely, the target is a \emph{record-effect Gleason theorem}: a normalized, additive, non-contextual effect-weight $w$ that is certain on the state's ray ($w(P_\Phi)=1$) must equal the Born functional $w(E)=\langle\Phi|E|\Phi\rangle$; hence on any complete record family $\{C_k\}$ ($\sum_k C_k^\dagger C_k = I$) the selector measure is \emph{forced} to $\mu(k) = \langle\Phi|C_k^\dagger C_k|\Phi\rangle$. This is \textbf{uniqueness/canonicity of $\mu$ from structural hypotheses}, not "probability from nothing".

\emph{Lean status (\texttt{GleasonSelector.lean}), Born forced from positivity.} The architecture is formalized in the finite-dimensional record model. A first version named the Busch / POVM-Gleason step as an interface axiom; a soundness review (GPT-5.5-pro) showed that axiom was \textbf{false as stated}, it omitted \emph{positivity}, and the explicit weight $w(E) = E_{00} + E_{01}$ on $\mathbb{C}^2$ satisfies normalization, additivity, homogeneity, and ray-certainty yet is not the Born functional. The axiom was \textbf{retired} and replaced by \emph{proved} content (every theorem standard-axioms-only, verified by \texttt{AxiomAudit.lean}, except where the two standard linear-algebra interface axioms below are explicitly invoked):
\begin{itemize}
  \item \texttt{naive\_gleason\_premises\_insufficient}: the $\mathbb{C}^2$ counterexample, proved: the positivity-free premises do \emph{not} force Born (a soundness red-team result that documents why positivity is indispensable);
  \item \texttt{rankOne\_sandwich}: proved (no normalization needed): $|\psi\rangle\langle\psi|\, E\, |\psi\rangle\langle\psi| = \langle\psi|E|\psi\rangle \cdot |\psi\rangle\langle\psi|$ (the algebraic heart of "ray-certainty $\Rightarrow$ Born");
  \item \texttt{support\_of\_positive\_certain}: proved: a \textbf{positive}, additive, homogeneous, ray-certain functional is \textbf{ray-supported} ($w(E) = w(P_\psi E P_\psi)$), via the Cauchy–Schwarz null-radical argument ($Q = I - P_\psi$ has $w(Q) = 0$, so all off-ray terms vanish). This \emph{derives} the support property from positivity, the genuine Gleason/Busch bridge, rather than assuming it;
  \item \texttt{positive\_ray\_certain\_forces\_born}: \textbf{the capstone, proved}: a positive, normalized, ray-certain weight $w$ \textbf{is} the Born functional. Born now follows from \emph{positivity + normalization + ray-certainty}, with positivity (not a conclusion-equivalent support premise) as the substantive hypothesis;
  \item \texttt{history\_measure\_is\_born} / \texttt{history\_measure\_total}: the record measure is the Born / decoherence-functional measure, normalized;
  \item \texttt{no\_signaling\_marginal}: requirement (e)/no-signaling, honest form: one bilinear correlation gives spacelike-marginal independence for \emph{all} of Bob's settings.
\end{itemize}

The capstone originally rested on two \emph{standard} linear-algebra interface axioms; \textbf{both have since been discharged}, so the entire finite-dimensional Gleason module is now \textbf{axiom-free} (every theorem, including the capstone, depends only on Lean's standard axioms \texttt{propext}, \texttt{Classical.choice}, \texttt{Quot.sound}, verified by \texttt{AxiomAudit.lean}). Specifically: \texttt{positive\_functional\_hermitian} (a positive functional is a $*$-functional) is proved by the standard polarization argument (realness of $w((A+X)^\dagger(A+X))$ and $w((A+iX)^\dagger(A+iX))$ forces conjugate symmetry); and \texttt{psd\_null\_radical} (a null vector of a positive-semidefinite sesquilinear form lies in its radical, the Cauchy–Schwarz core) is proved by deriving first-slot conjugate-linearity from conjugate symmetry, then applying the real-quadratic discriminant lemma \texttt{quadratic\_nonneg\_forall\_linear\_zero} (also proved) at $c = t$ and $c = it$ to force $\mathrm{Re}\,(B\,Q\,X) = \mathrm{Im}\,(B\,Q\,X) = 0$. The honest residual is now only the \emph{continuum} Type II / history-net version (the finite-dimensional case is complete and unconditional).

\emph{Honest verdict, and the position adopted.} We do \textbf{not} claim to derive $\mu$ from unitarity plus the bare existence of $\lambda$; on structural grounds \textbf{one} bridge principle is required (here: \emph{positivity} of the selector functional + state-certainty), the same status that quantum equilibrium has in Bohmian mechanics, and the analogue of the Everettian probability problem. (We do not assert a formal no-go that such a derivation is \emph{impossible}, only that this framework, like every single-world programme, supplies one bridge principle rather than deriving probability from nothing.) The adopted position is: \emph{QIQT-H posits a selector functional constrained by positivity, finite-record non-contextuality, covariance, and state-certainty; a Gleason-type theorem (\texttt{positive\_ray\_certain\_forces\_born}, proved) then shows the only canonical $\mu$ is the decoherence-functional measure $\mu_\Phi(\mathrm{Cyl}\,\alpha) = D_\Phi(\alpha,\alpha)$; the one remaining postulate is that the actual $\lambda$ is $\mu$-typical}, the relativistic analogue of quantum equilibrium. Strict Wood–Spekkens "no fine-tuning" in the causal-faithfulness sense is unachievable for any ontic-Parameter-Independence-violating model (a known no-go); what \emph{is} achieved is the defensible weaker form, one canonical $\mu$ enforces operational no-signaling for every spacelike instrument pair, with no per-experiment tuning (\texttt{no\_signaling\_marginal}).

\subsubsection{Formal Born representation theorem (Mackey-Gleason + noncommutative Radon-Nikodym)}

\textbf{The mathematical core.} For a normal probability measure on the projection lattice of a von Neumann algebra without Type $I_2$ summand, Mackey-Gleason (Bunce-Wright 1990s, extending classical Gleason 1957) gives a unique normal state extension. Combined with the noncommutative Radon-Nikodym theorem (Sakai 1971; Takesaki vol. II), this yields the trace-density form

\[
\mu_\rho(P) = \tau(D_\rho \cdot P)
\]

on semifinite vN algebras with faithful normal trace $\tau$.

\textbf{Caveats to be explicit about:}
\begin{itemize}
  \item \textbf{No bare qubit Gleason}: classical Gleason has a $d=2$ exception. The fix is Busch's POVM-Gleason theorem (2003) extending the result to all effects (Busch, Caves-Fuchs-Manne-Renes 2004).
  \item Normality / σ-additivity vs finite additivity assumptions.
  \item Complex Hilbert spaces (real/quaternionic variants require separate treatment).
  \item Factor vs non-factor center (superselection sectors).
  \item PVM-only vs POVM/effect-level Born rule.
  \item Exclusion of singular states.
\end{itemize}

\textbf{Lean status:} the finite-dimensional content is now the axiom-free \texttt{EffectGleason.finite\_effect\_gleason} (Born forced from positivity); the earlier \texttt{TypicalityMackeyGleason} packaging was retired and deleted. The full von Neumann-algebraic Mackey-Gleason + noncommutative Radon-Nikodym version remains a Mathlib-formalization task (multi-week).

\subsubsection{Existence and canonicality in finite factors (Type $II_1$)}

\textbf{Standard mathematical existence (NOT open).} For a finite Type II factor $M$ with faithful normal tracial state $\tau$, the construction

\[
\mu_\rho(p) = \tau(\rho \cdot p), \quad \rho \in L^1(M, \tau)_+, \quad \tau(\rho) = 1
\]

yields a normal probability measure on the projection lattice $\mathrm{Proj}(M)$. Positivity, normalization, orthogonal additivity, normality, and unitary covariance are all standard theorems. \textbf{In a $II_1$ factor, the trace $\tau$ is the unique normal state invariant under all inner unitaries}, so canonicality under unitary invariance is \emph{not} an open problem; it is a consequence of factor structure.

\textbf{What remains open in $II_1$:} none of the \emph{mathematical} existence; only the (genuinely open) question of canonical \emph{physical selection}, what makes a particular $\rho$ correspond to a particular preparation. That is the empirical calibration step (§11.4.4 below).

\textbf{Lean status:} This is a Lean formalization gap, \emph{not} an open math problem. The construction can be formalised directly given the requisite vN-algebra infrastructure in Mathlib.

\subsubsection{Infinite trace and Type $II_\infty$, a sharp obstruction}

\textbf{The obstruction.} For a Type $II_\infty$ factor with semifinite trace $\tau$, the trace satisfies $\tau(1) = \infty$. Therefore the trace itself is a \emph{weight}, not a probability state. There is \textbf{no normalized normal state on a $II_\infty$ factor invariant under all unitaries}.

Concretely: one cannot define a "canonical uniform" probability measure $\mu(p) = \tau(p)$ in $II_\infty$; the candidate is not normalizable.

\textbf{Options for the QIQT-H Type II setting:}
\begin{enumerate}
  \item Choose a density $\rho$: $\mu_\rho(p) = \tau(\rho \cdot p)$ with $\rho \in L^1(M, \tau)_+$ and $\tau(\rho) = 1$. The measure exists; canonicality requires an additional principle to select $\rho$ for a given physical preparation.
  \item Restrict to a finite corner $eMe$ with $\tau(e) < \infty$: recovers the $II_1$ analysis above on the corner.
  \item Use semifinite \emph{typicality weight} rather than probability measure: appropriate if the framework adopts a non-probability typicality reading.
  \item Provide additional physical selection principles (canonical sector reference state, modular structure, holographic origin).
\end{enumerate}

\textbf{Lean status:} Sub-theorem A's interface (Mackey-Gleason + RN) applies to $II_\infty$; the obstruction is \emph{physical canonicality}, not formal existence of $\mu_\rho$.

\subsubsection{Type III caveat (AQFT local algebras)}

\textbf{A non-trivial caveat.} AQFT local algebras of bounded regions are typically Type $III_1$, not Type II (Buchholz, Borchers, Longo; see foundations paper §3). The QIQT-H framework uses the CPW Type II crossed-product construction (CPW 2022, Witten 2022) precisely to \emph{escape} Type III to a well-defined entropy structure. But:

\begin{itemize}
  \item \textbf{There is no trace-density picture in Type III.} Normal states still give projection probabilities $\phi(p)$, but the formula $\phi(p) = \tau(\rho \cdot p)$ has no analog in Type III (no trace).
  \item The relationship between the underlying Type III local algebras and the crossed-product Type II algebras is mediated by the modular flow of the canonical sector reference state $\sigma_R$.
  \item The Born representation problem in Type III is structurally different and requires either passing through the Type II crossed product or using Connes' spatial derivative / Haagerup's $L^p$-space construction.
\end{itemize}

\textbf{Open:} make explicit how the Type II crossed-product trace-density picture relates to the underlying Type III local-algebra structure, particularly whether the canonical IC measure $\mu_\rho$ on the Type II core has a natural Type III pullback.

\subsubsection{Operational calibration, split into derivable and irreducible parts}

\textbf{Splitting the calibration step.} The "measurement-calibration" step (often informally called "the Born content bridge") is actually two parts:

\begin{itemize}
  \item \textbf{(4a) Mathematical restriction (derivable).} If $A \cong M_d(\mathbb{C})$ is a finite-dimensional measurement subalgebra of a Type II factor and $\phi$ is a normal state, then $\phi|_A(a) = \mathrm{tr}_d(\rho_A \cdot a)$ for a unique ordinary $d \times d$ density matrix $\rho_A$. If the embedding is trace-preserving, the Type II trace restricts to the normalized matrix trace on $A$. \emph{This is standard mathematics.}
  \item \textbf{(4b) Empirical calibration (irreducible).} The identifications:
  \item this detector $\leftrightarrow$ this projection / effect;
  \item this preparation $\leftrightarrow$ this density;
  \item this pointer event $\leftrightarrow$ this subalgebra;
  \item this finite-dimensional operational model $\leftrightarrow$ the right model of the actual lab.
\end{itemize}
  
  Operational reconstructions (Hardy 2001; Chiribella-D'Ariano-Perinotti 2011; Masanes-Müller 2011; Brukner-Zeilinger) derive the Born pairing from operational axioms, but the axioms themselves are physical/empirical inputs. \emph{Calibration is not eliminated; it is relocated.}

\textbf{The irreducible part is small but real.} Any deterministic theory needs (4b). It is the analog of the bridge from Liouville measure to thermodynamic ensembles in classical statistical mechanics.

\textbf{POVMs and instruments.} Modern operational treatments use POVMs (effects) rather than PVMs (projections). For sequential measurements and state update, Davies-Lewis / Ozawa instruments are the appropriate framework. The QIQT-H paper should include POVM formulations in §7.6 and §7.7 to avoid the bare-qubit Gleason exception.

\subsubsection{Dynamics, equivariance, and locality}

The conditional Born typicality theorem (\texttt{BornTypicality.lean}) requires three additional physical inputs:

\begin{itemize}
  \item \textbf{Equivariance theorem.} If $\rho_t = U_t \rho_0 U_t^*$ under FQ-restricted unitary dynamics, then $\mu_{\rho_t}(p) = \mu_{\rho_0}(U_t^* p U_t)$. For QFT-style Hamiltonians, this requires careful treatment of domain / self-adjointness issues. The \texttt{EquivarianceGap.lean} audit establishes that support preservation (the framework's current claim) is strictly weaker than measure preservation.
  \item \textbf{Projection locality.} Need precise assumptions on the local algebraic structure: tensor product factorisation (or its absence in Type III), commuting local algebras (microcausality), split property, finite corners. The \texttt{CompressionLocality.lean} audit isolates the compression-locality leakage identity and identifies when the FQ projection preserves microcausality at the restricted level. \emph{Post A1 strengthening} (\texttt{MarginalLocality.pushforward\_marginal\_local}): given set-level locality of Bob's dynamics ($r\circ T = r$), the Alice-marginal is invariant for \textbf{any} IC measure: with no equivariance assumption, so locality of the marginal is a theorem, not a separate Born-selection postulate. This reduction is, however, \emph{conditional} on a named Hilbert→set-level locality bridge axiom (\texttt{set\_level\_locality\_from\_unitary\_dilation}, a standard Heisenberg↔Schrödinger correspondence): locality has been relocated into that explicit bridge, not eliminated. Subject to that caveat, locality is no longer counted among the independent Born-selection sub-axioms of the Canonical IC Measure Principle.
  \item \textbf{Empirical frequency bridge.} If μ is interpreted as typicality (rather than primitive probability), one needs a DGZ-style step from typical initial configurations to observed frequencies. The \texttt{BornTypicality.lean} module formalizes this conditionally via standard finite LLN.
\end{itemize}

\subsubsection{Typicality vs probability}

\textbf{The conceptual framing.} QIQT-H is fully deterministic; the "measure" $\mu_\rho$ is a \emph{typicality structure} on the microscopic-IC space (à la Dürr-Goldstein-Zanghì 1992 in Bohmian mechanics), not a primitive probability assignment. The candidate routes for canonicality (§7.4 α/β/γ) are typicality principles, not probability axioms.

\textbf{Caveats and references:}
\begin{itemize}
  \item \textbf{Frequency concentration (PROVED, single-trial + independence step).} The Lean module \texttt{BornConcentration.lean} proves Chebyshev's tail bound on finite probability spaces ($\Pr(|X-\mu| \ge \varepsilon) \le \mathrm{Var}/\varepsilon^2$), computes the Bernoulli variance ($p(1-p)$), proves the variance-addition (independence) lemma \texttt{variance\_add\_of\_product} (variance of a sum of independent variables equals the sum of variances), and combines these to give the single-trial concentration inequality and the two-trial Bernoulli variance $2p(1-p)$. \emph{What is proved:} the single-trial Chebyshev bound and the reusable variance-addition lemma the $N$-fold induction would iterate. \emph{What remains axiomatic:} the full $N$-trial product-measure scaling $\Pr(|\mathrm{freq}_N - p| \ge \varepsilon) \le p(1-p)/(N\varepsilon^2)$ rests on the LLN interface axiom in \texttt{BornTypicality.lean}. Accordingly the framework's frequency statement ("Born \emph{frequencies}", not merely "Born \emph{means}") is established up to that single named scaling axiom: for which the per-trial variance and the independence step are now both in hand.
  \item \textbf{Goldstein-Struyve uniqueness} (J. Stat. Phys. 128, 1197, 2007): in Bohmian mechanics, $|\psi|^2$ is the unique local equivariant typicality density. The adaptation to QIQT-H's Type II + (FQ) setting is precisely the multi-week mathematical task of Sub-theorem C.
  \item \textbf{Valentini nonequilibrium} (1991, 2002; Valentini-Westman 2005): there exist non-Born typicality measures whose dynamics relaxes (or fails to relax) to equilibrium. This is the principled objection to "canonical = equivariant": equivariance is necessary but may not be sufficient without an additional relaxation argument.
  \item \textbf{Noncontextuality} in this deterministic-with-typicality setting refers to the equilibrium probability assignment to quantum effects (Spekkens 2005 generalised noncontextuality), \emph{not} to noncontextual pre-existing values (which Kochen-Specker rules out in $d \ge 3$).
\end{itemize}

\subsubsection{Independence / minimality of the remaining sub-axioms}

Per GPT-5.5-pro's sixth audit, the honest Born claim is not "Born is derived from holography alone", \texttt{NoBornFromNothing} refutes that, but rather:

\begin{qiqtquote}
\emph{Born is the unique admissible measure given a NAMED minimal set of additional postulates, and each postulate is INDEPENDENT, dropping any one yields a finite countermodel where Born fails.}
\end{qiqtquote}

The independence package \texttt{BornMinimalityTable.lean} makes this explicit. Relative to the current finite formal decomposition, the Canonical IC Measure Principle has three irreducible Born-selection sub-axioms, each with a concrete finite countermodel witnessing its necessity (with the locality premise reducible, modulo the named bridge axiom of §11.4.5):

\begin{center}
\small
\begin{tabularx}{\linewidth}{|Y|Y|Y|Y|}
\hline
\textbf{Tag} & \textbf{Premise} & \textbf{Minimality witness (Lean module)} & \textbf{Countermodel size} \\
\hline
P1 & Canonical IC measure (some structurally-distinguished $\mu$ is selected) & \texttt{NoBornFromNothing.any\_anti\_born\_realizable} & $\Gamma$ arbitrary finite \\
\hline
P2 & Measure equivariance under FQ-restricted dynamics & \texttt{EquivarianceGap.support\_preservation\_does\_not\_imply\_…} & $\mathrm{Fin}\,2$ swap \\
\hline
P3 & Operational sufficiency (click-statistics determine IC marginal) & \texttt{OperationalNoGo.operational\_data\_insufficient} & $\mathrm{Fin}\,3 \to \mathrm{Fin}\,2$ \\
\hline
P4 & Locality of Alice's marginal under Bob's local dynamics & \textbf{REDUCIBLE}, proved as a theorem holding for \emph{any} measure (\texttt{MarginalLocality.pushforward\_marginal\_local}), conditional on the named Hilbert→set-level bridge axiom (§11.4.5) & (not an independent Born-selection axiom) \\
\hline
\end{tabularx}
\end{center}

This converts the standing concern "you still assume X" into a sharper audit conclusion: X is \emph{necessary}. Relative to the current finite formal decomposition, the irreducible Born-selection premise set is $\{P1, P2, P3\}$ (down from four, once the marginal-locality step is theoremized modulo its bridge axiom); each premise has a finite-dimensional countermodel; and Born follows uniquely when all three hold, together with the operator-algebraic Mackey-Gleason / Radon-Nikodym input (the finite Schur classification and tensor multiplicativity, once acknowledged interface axioms, are now proved axiom-free). This is the strongest honest statement available without resolving the \emph{continuum} Sub-theorems A and C as theorems-in-Mathlib; it is a \emph{minimality relative to this decomposition}, not a proof that no coarser axiom set could suffice.

\subsubsection{Combined: what is and is not open}

\begin{center}
\small
\begin{tabularx}{\linewidth}{|Y|Y|Y|}
\hline
\textbf{Layer} & \textbf{Item} & \textbf{Status} \\
\hline
11.4.1 & Mackey-Gleason + RN representation theorem & Standard math; Mathlib formalisation gap (multi-week) \\
\hline
11.4.2 & Existence in finite Type II ($II_1$) & Standard math; Lean formalisation gap (constructive) \\
\hline
11.4.3 & Canonical normalised probability state on $II_\infty$ & \textbf{Sharply obstructed}; requires density or finite corner or additional principle \\
\hline
11.4.3a & Type III ↔ Type II crossed-product correspondence for IC measure & Open; mathematically non-trivial \\
\hline
11.4.4a & Type II → finite-dim restriction (math) & Standard math \\
\hline
11.4.4b & Empirical lab-to-formalism calibration & Irreducible empirical bridge (no derivation in any framework eliminates it) \\
\hline
11.4.5 & Equivariance theorem under FQ-restricted Hamiltonian & Open; requires concrete FQ-Hamiltonian + self-adjointness analysis \\
\hline
11.4.6 & Justification of typicality framing (Goldstein-Struyve adaptation; Valentini relaxation) & Open; multi-year research direction \\
\hline
11.4.6 & Born \textbf{frequency} concentration (single-trial Chebyshev) & \textbf{PROVED} in \texttt{BornConcentration.born\_chebyshev\_single\_trial}; $N$-trial scaling at LLN axiom interface \\
\hline
11.4.7 & Independence/minimality of P1, P2, P3 by finite countermodel & \textbf{PROVED} in \texttt{BornMinimalityTable.lean} \\
\hline
11.4.7 & Reducibility of locality (P4): marginal-locality step proved for \emph{any} measure & \textbf{PROVED} in \texttt{MarginalLocality.pushforward\_marginal\_local}; reduction conditional on the named Hilbert→set bridge axiom \\
\hline
\end{tabularx}
\end{center}

\textbf{Bottom line.} The Canonical IC Measure Principle is \textbf{not} a monolithic mystery. The hard parts are: the $II_\infty$ canonicality obstruction (11.4.3); the Type III correspondence (11.4.3a); the equivariance theorem for FQ-restricted dynamics (11.4.5); and the typicality justification (11.4.6). The mathematical core (11.4.1, 11.4.2, 11.4.4a) is standard given operator-algebra infrastructure; the empirical calibration (11.4.4b) is irreducible. \textbf{After the A1+A2+A4+A6 strengthening pass:} the marginal-locality step is theoremized (for \emph{any} measure, no equivariance assumed) so locality is no longer an independent Born-selection sub-axiom, modulo a named Hilbert→set-level bridge axiom into which the locality content is relocated, not eliminated; relative to the current decomposition the irreducible Born-selection set has three named members (P1/P2/P3), each provably necessary by concrete finite countermodel; and single-trial Born frequency concentration plus the variance-addition (independence) lemma are proved, with the $N$-trial scaling to full frequency convergence still resting on the named product-measure LLN axiom. These are minimality and reduction results \emph{relative to this formal decomposition}, not claims that the underlying operator-algebra interface axioms (Schur classification, tensor multiplicativity, Mackey-Gleason/Radon-Nikodym) have been discharged.

\textbf{Open Problem 2, Reference-weight bound (sharpened H2).}
The record-instantiation-cost postulate (H2 in Theorem 6) is \emph{equivalent} to the reference-weight bound $\sigma_R(E_{\rm record}) \le \exp(-(I_0 - \eta_0))$ on macroscopic pointer sectors (see §7.6 Sharpening of (H2); formalized in \texttt{lean/mathlib/QIQTH/H1H2Audit.lean}). Derive this reference-weight bound from modular / holographic first principles, as a strengthening of standard Bekenstein-Bousso. Donald's identity + Klein positivity + DPI alone are insufficient. This would be genuinely new physics. \textbf{(Revised 2026.} This problem was framed as feeding the former Open Problem 3 (Macroscopic Definiteness), which is now \textbf{withdrawn} as a category error (§1.1a). It is therefore \emph{no longer load-bearing} for single-outcome definiteness; it survives only as an optional sharpening of the pointer-sector Bekenstein–Bousso weight bound, not as a route to a capacity exclusion of multi-record states.)

\textbf{Open Problem 3 — WITHDRAWN (2026): the Macroscopic Definiteness Conjecture is retired as a category error.}
This problem formerly asked to \emph{prove the Branch-Summed bound} — that a regional content carrying $K \ge 2$ macroscopically distinct records has summed information cost $I_\Sigma \approx K\cdot I_0 > Q_R$ and is therefore not an instantiable physical state, so that finite capacity itself \emph{forbids} the multi-record state. \textbf{We withdraw it} (§1.1a). $Q_R$ bounds the \emph{von Neumann} entropy $S(\rho_R) = H(\{p_i\}) + \sum_i p_i S(\rho_{R,i})$ of the pre-selection state $\Phi$, and a second macroscopic branch raises only the branch-count term — by $\sim\!\log K$, far below $Q_R$ — so a $\ge 2$-record state does \textbf{not} overflow capacity; the branch-summed cost $I_\Sigma$ is provably not the entropy $Q_R$ bounds (\texttt{BranchLedger.branchSummed\_not\_bounded\_by\_Shannon}); and in a unitarily-evolving theory a kinematic bound cannot in any case \emph{select} an outcome (\texttt{NoConcentration}, \texttt{RealmSelection.capacity\_underdetermines\_realm} — "the area budget is not the metaselector"). \textbf{Single-record definiteness is the work of the selector $\lambda$ (P1) together with decoherence, not of a capacity exclusion.} What remains genuinely open is \emph{not} this conjecture but the \textbf{dynamical, Lorentz-covariant realization} of $\lambda$-selection (Open Problem 3b below) and $\lambda$'s law (Born from typicality); the subsidiary quantitative questions — recoherence stability, $\epsilon(R)$, $I_0$, and no-signaling consistency (Open Problem 10) — are retained under their own headings.

\emph{Status (2026-06), the finite additive-cost number-bound is mechanized — but it does \textbf{not} revive the withdrawn capacity-exclusion.} In a tractable finite-dimensional \textbf{additive-cost} model the number-bound is an axiom-free theorem; however the additive (branch-summed) cost it bounds is provably \emph{not} the holographic von Neumann entropy $Q_R$ bounds (\texttt{branchSummed\_not\_bounded\_by\_Shannon}), so it models the storage of the $\lambda$-selected record, not a capacity exclusion of multi-record states. Concretely: realizing records as a redundant Spectrum Broadcast Structure makes the summed cost a genuine \emph{information} quantity $\ge R\log n$ (broadcasting tensors the fragments, so dimensions multiply), and a finite additive capacity then admits at most one macroscopic record, an axiom-free Lean theorem (\texttt{SBSBridge.sbs\_single\_outcome}, \texttt{CoreNoCollapse.coactual\_subsingleton}), with the threshold \emph{derived} from pointer-state orthonormality (\texttt{CapacityModel.capacity\_total}) and robust to approximate decoherence (\texttt{fragment\_finrank\_ge\_approx}). This is the \emph{number-bound} (one-record) content; what remains genuinely open is (i) the continuum realization, reproducing this scaling as a renormalized-entropy bound on the Type II regional algebra rather than a finite register, and (ii) the \emph{field-theoretic} origin of the scattering premise. Input (ii) has been \textbf{discharged for the standard collisional/QND toy Hamiltonian} $H_{\rm int}=g\,\sigma_z^S\otimes\sigma_x^E$ (\texttt{CollisionalGamma.lean}): the branch overlap is $\cos 2\theta$, so $\gamma=|\cos 2\theta|<1$ is derived and amplifies to $\gamma^L\to0$ (§6.8); what stays open is deriving the \emph{uniform-$\gamma$, factorized, no-recoherence} structure itself from a realistic field-theoretic Hamiltonian (the constructive form of the Strasberg–Winter branch-selection problem, arXiv:2601.19703). The strong superposition-exclusion form (forbidding the cat state) is \emph{not} claimed and is not needed for the single-record content.

\textbf{Open Problem 3b, Lorentz covariance of the selection structure (of equal rank to the Born/typicality problem).}
Establish that the actuality selector $\lambda$ and the regional reading map $A_R[\Phi,\lambda]$ introduce \textbf{no hidden preferred foliation}, i.e. that the single-world structure is genuinely Lorentz-covariant, not merely operationally no-signaling. This is the relativistic counterpart of the Born/typicality problem (Open Problem 1) and is its equal in rank; it is \emph{not} secured by the operational no-signaling theorem (§7.7), since no-signaling $\neq$ Lorentz invariance of the beable (§6.9, §11). The framework is structurally \emph{better placed} than Bohmian mechanics here, because $\Phi$ evolves by exactly-unitary, already-covariant dynamics and $\lambda$ is \emph{non-dynamical} (there is no collapse or guidance event that must be time-ordered, hence nothing that requires a "now"), but "better placed" is not "proved." Five explicit requirements, in increasing difficulty:

\begin{enumerate}
  \item \textbf{$\lambda$ as a genuinely 4-dimensional object.} $\lambda$ must be defined geometrically over the whole spacetime history: e.g. a choice of one globally consistent decoherent history of $\Phi$, or a section of an appropriate bundle over spacetime, with the Poincaré group acting on it \emph{geometrically} and with no reference to a Cauchy slice. Any definition that reduces $\lambda$ to "initial data on $\Sigma_0$" reintroduces a preferred frame and fails. (This is the substance of the relabel from "microscopic initial conditions" to "atemporal actuality selector," §2.2, §6.9, here it is made a covariance \emph{requirement}, not just terminology.)
  \item \textbf{Covariant region structure.} $Q_R = A(\partial R)/4\ell_P^2$ (and the subsidiary resolution floor $\epsilon(R)$, superseded — §1.1a) must be defined on \textbf{causal diamonds} via the \textbf{covariant (Bousso light-sheet) entropy bound}, \emph{not} on spatial slices: the area of $\partial R$ on a spacelike slice is frame-dependent. The CPW/Witten Type II scaffolding already lives on causal diamonds with microcausality built in (and machine-checked), so the scaffolding is covariant; the burden is to make $Q_R$ and the reading map respect it. \emph{This requirement is concrete and fixable now} and should be discharged first.
  \item \textbf{The covariance identity (the core theorem).} Prove
\end{enumerate}

\[
A_{gR}[\,U_g\Phi,\; g\!\cdot\!\lambda\,] \;=\; g\cdot A_R[\Phi,\lambda] \qquad \text{for every Poincaré } g,
\]

i.e. selecting the record in the boosted diamond $gR$ from the boosted state $U_g\Phi$ and boosted selector $g\lambda$ yields the boost of the originally-selected record. This is the precise statement that "every observer in every frame agrees on the physical facts." Open.

\begin{enumerate}
  \item \textbf{Foliation-free global consistency.} Show that the family $\{C_R = A_R[\Phi,\lambda]\}$ over all diamonds is globally consistent: restriction-compatible on nested diamonds (extending the $\sigma_R$ restriction-compatibility of §7.6 to the $\lambda$-selected content), and jointly consistent on spacelike-separated diamonds, \emph{with the consistency condition statable purely in the causal partial order}, requiring no global time function. The danger is that a "global history" selection covertly needs a foliation to define "consistent across time"; operational no-signaling (§7.7) is the observable shadow of this consistency, but the ontic statement must be proved. Open.
  \item **Poincaré-\emph{equivariant} typicality.** The measure $\mu$ over $\lambda$ (Open Problem 1) must transform \emph{covariantly with the state}, $g_*\mu_\Phi = \mu_{U_g\Phi}$: \textbf{not} strict invariance $\mu(g\cdot S)=\mu(S)$, which is too strong since the Born weights depend on $\Phi$ (strict invariance is the special case where $\Phi$ is itself Poincaré-invariant, e.g. the vacuum). This couples the relativistic problem to the Born problem: an equivariant $\mu_\Phi$ is strictly harder than a single-frame $\mu$, and a non-equivariant family would re-import a preferred frame through the statistics even if $\lambda$ itself is clean. This is the relativistic analogue of the difficulty Bohmian mechanics has with frame-dependence of $|\psi|^2$-equilibrium. Open, and entangled with Open Problem 1.
\end{enumerate}

We note candidly that requirements 3–5 are genuine open theorems, requirement 5 the hardest; that the program's Lorentz-friendliness is \emph{bought with} an atemporal / all-at-once (block-universe, mildly retrocausal-flavored) reading of $\lambda$, which is owned openly rather than hidden; and that this places QIQT-H in the Lorentz-friendly family of single-world proposals (Kent final-boundary; Wharton/Sutherland all-at-once) rather than the foliation-bound Bohm/GRW family, a structural advantage that is, as yet, a promissory note.

\emph{\textbf{Proposed construction (a concrete attack, not yet a theorem).}} Two design choices make the five requirements simultaneously well-posed and, we argue, jointly natural; they are recorded here as the intended line of attack.

\emph{(i) Relativistic dynamics: $\Phi$ is a quantum field, not a Schrödinger wavefunction.} Throughout this paper "Schrödinger / Heisenberg evolution" is shorthand for the manifestly covariant dynamics of the underlying field theory; for matter the relevant equation is the \textbf{Dirac equation} $(i\gamma^\mu\partial_\mu - m)\psi = 0$ (second-quantized: the Dirac quantum field), not its non-relativistic limit. This is not cosmetic but \emph{required}: requirement 3's covariance identity $A_{gR}[U_g\Phi, g\!\cdot\!\lambda] = g\cdot A_R[\Phi,\lambda]$ is only statable once $\Phi$ carries a genuine Poincaré action $U_g$, which the Dirac field provides through the spinor representation $S(\Lambda)$, and a Galilean Schrödinger equation does not. The Dirac field is also the natural matter content of the CPW/Witten AQFT scaffolding (the local algebras are field algebras), so this choice merely names the QFT the algebras already presuppose. (The single-particle Dirac equation's negative-energy / Klein pathologies are why the \emph{field}, not the one-particle wave equation, is the carrier; $\Phi$ is the global QFT state.)

\emph{(ii) $\lambda$ as holographic boundary data on causal-diamond screens.} We propose that $\lambda$ be \textbf{encoded on boundaries, not in the bulk}: for each causal diamond $D$, the $\lambda$-selected record sector $C_D = A_D[\Phi,\lambda]$ is reconstructed from $\lambda$'s restriction to the diamond's edge / holographic screen $\partial D$. The motivation is internal and, we think, forcing: the \emph{capacity} that constrains records is itself a boundary quantity, $Q_R = A(\partial R)/4\ell_P^2$, so the actuality fact that \emph{spends} that capacity should live where the budget does. This single choice addresses three requirements at once. (1) The boundary of a causal diamond is a Lorentz-covariant geometric object, it transforms under boosts with no preferred slicing, so boundary data is the slice-free 4-D object requirement 1 demands (boundary data $\neq$ Cauchy data on $\Sigma_0$). (2) $Q_R$ and $\lambda$ then live on the same screen, automatically aligning the covariant capacity bound with the selector. (4) Boundary data on nested / overlapping diamonds carries a natural restriction-compatibility (agreement on shared edges), turning "foliation-free global consistency" into an order-theoretic gluing condition on the causal-diamond poset rather than a statement requiring a global time function.

\emph{Honest status of the proposal.} "Boundary data determines the bulk record sector" is the holographic principle / bulk reconstruction, and in controlled settings (AdS/CFT: HKLL, entanglement-wedge reconstruction) that reconstruction is \emph{subtle and not literal}; for general causal diamonds in non-AdS spacetimes it is an open question in its own right. So this proposal does not \emph{solve} requirements 3–4, it \emph{relocates} them to the sharper, more tractable question of whether holographic boundary data on a causal-diamond screen uniquely and covariantly fixes the $\lambda$-selected bulk record. That is the theorem to attempt. The merit of the proposal is that it gives $\lambda$ a concrete geometric home (a covariant holographic screen) and a manifestly covariant dynamics (Dirac/QFT $U_g$), so both halves of the pair $(\Phi,\lambda)$ transform geometrically, with nothing referencing a "now."

\emph{\textbf{Reduction, corrections, and the linchpin theorem (GPT-5.5-pro adversarial sharpening).}} Subjecting the proposal to an adversarial review converts the slogan into a precise target and corrects three mistakes in its naive form. We record both the corrections and the resulting linchpin theorem, because a sharply-posed target is itself a contribution.

\emph{Correction A, the gluing is over bulk overlaps, not screen intersections.} For nested diamonds $K \subset D$ the boundary $\partial K$ is in general \textbf{not} a subset of $\partial D$, so literal "agreement on shared edges" is too weak. The correct, frame-free consistency condition is \textbf{bulk-overlap agreement}: writing $\rho_{D,L}$ for the restriction of a $D$-record to a subdiamond $L$,

\[
\rho_{D,L}(C_D) = \rho_{E,L}(C_E) \qquad \text{for every } L \subset D \cap E .
\]

Equivalently, organize the admissible record sectors as a contravariant functor $X_\Phi(D) = \mathrm{Stone}(\mathsf{B}_\Phi(D))$ on the causal-diamond poset $\mathsf{Diam}$ (where $\mathsf{B}_\Phi(D) \subset \mathrm{Proj}(\widehat{\mathcal A}(D))$ is the finite Boolean algebra of decoherent macroscopic record propositions), with restriction maps $\rho_{D,K} : X_\Phi(D) \to X_\Phi(K)$ for $K \subset D$. Then $\lambda$ is a \textbf{global section},

\[
\lambda \in \Gamma(X_\Phi) = \varprojlim_{D \in \mathsf{Diam}} X_\Phi(D), \qquad \rho_{D,K}(\lambda_D) = \lambda_K ,
\]

and requirement 4 ("foliation-free global consistency") \emph{is} the existence of such a global section. This is the precise, order-theoretic, time-function-free statement we wanted.

\emph{Correction B, the gluing obstruction is AQFT net cohomology, and the split property does not trivialize it.} When local record-sector labels are defined only up to automorphism, gluing requires transition functions $g_{ij} \in \mathrm{Aut}(X_\Phi)$ on overlaps satisfying the cocycle condition $g_{ij}g_{jk}g_{ki} = 1$ on triple overlaps; the obstruction is a class $[g_{ij}] \in \check{H}^1(\mathcal U, \mathrm{Aut}(X_\Phi))$ (rising to an $H^2$ / gerbe class if agreement is only projective). In AQFT this is precisely \textbf{Roberts net cohomology}, the same machinery that classifies DHR superselection sectors under Haag duality, so the obstruction can carry charge sectors, gauge flux, and edge modes; it is not mere topological Čech cohomology. The Type III$_1$ \textbf{split property} (Type-I intermediate factors for separated inclusions) and \textbf{Haag duality} \emph{help}, they let one define finite local pointer algebras and identify causal complements, but they do \textbf{not} by themselves trivialize the obstruction: the split inclusion is \emph{non-canonical}, and a non-canonical choice breaks covariance unless one proves a coherent, Poincaré-natural split. (Gauge theories further complicate naive Haag duality via edge/flux degrees of freedom, and the Kochen–Specker theorem forbids $\lambda$ from being a global valuation on the \emph{full} projection lattice, $\lambda$ must select one compatible recorded-history \emph{realm}, not assign values to all contexts.)

\emph{Correction C, requirement 5 is equivariance, not invariance} (already folded into requirement 5 above): $g_*\mu_\Phi = \mu_{U_g\Phi}$, not $\mu(g\cdot S) = \mu(S)$.

\emph{The covariance identity becomes automatic from equivariant naturality.} Requirement 3 is cleanest categorically, but not as "$A_R$ is a natural transformation of nets" ($A_R$ is not a $*$-homomorphism; it selects a classical sector). Rather: build the record-sector functor $X_\Phi$ \emph{naturally} from the covariant Haag–Kastler net $(\mathcal A, U)$ over the action category $\mathsf{Diam} \rtimes G$ (with $G = \mathcal{P}^\uparrow_+$, or its spin cover $\mathrm{ISpin}(1,3)$ for the Dirac field), with natural isomorphisms $\gamma_g : X_\Phi(D) \to X_{U_g\Phi}(gD)$, and take $\lambda$ to be an \textbf{equivariant global section}. Then requirement 3 is one line:

\[
A_{gD}[U_g\Phi, g\lambda] = (g\lambda)_{gD} = \gamma_g(\lambda_D) = g \cdot A_D[\Phi,\lambda].
\]

So requirement 3 reduces to proving that $X_\Phi$ is (i) functorial under inclusions, (ii) Poincaré-equivariant, (iii) compatible with the boundary-reconstruction maps. Bisognano–Wichmann / Borchers supply modular covariance ($\Delta_W^{it} = U(\Lambda_W(-2\pi t))$ for wedges); what \emph{breaks} naturality is exactly the list of frame-dependent choices to avoid: a non-canonical pointer basis, a non-canonical split inclusion, state-dependent modular crossed products not respecting inclusions, a frame-dependent decoherence threshold $\epsilon(D)$, gauge-fixing / missing edge modes, and tie-breaking rules in any "most decoherent branch" prescription.

\emph{The measure: the decoherence functional is the Born-reproducing, covariant $\mu$.} Of the candidates, the canonical Type II trace $\tau_D$ is a \emph{carrier} (it writes local Born weights as $p_i^D = \omega_{\Phi,D}(P_i^D) = \tau_D(h_{\Phi,D}\,P_i^D)$ via the Haagerup density $h_{\Phi,D}$) but its uniqueness-up-to-scale does \textbf{not} by itself pin $\mu$ (the bare trace gives uniform counting, not Born for arbitrary $\Phi$); Bisognano–Wichmann / KMS is geometric only for wedges and CFT diamonds, not general interacting diamonds. The winning candidate is the \textbf{Gell-Mann–Hartle decoherence functional}: for a medium-decoherent recorded-history family $\{\alpha\}$ with class operators $C_\alpha$,

\[
\mu_\Phi(\mathrm{Cyl}(\alpha)) = D_\Phi(\alpha,\alpha) = \langle \Phi | C_\alpha^\dagger C_\alpha | \Phi \rangle = \tau_D(h_{\Phi,D}\,P_\alpha),
\]

extended from finite compatible cylinder events to global histories by a Kolmogorov–Carathéodory extension theorem. Its covariance is immediate from covariant class operators, $C_{g\alpha} = U_g C_\alpha U_g^{-1} \Rightarrow D_{U_g\Phi}(g\alpha,g\alpha) = D_\Phi(\alpha,\alpha)$, giving exactly the equivariance $g_*\mu_\Phi = \mu_{U_g\Phi}$ of requirement 5, and it reproduces Born by construction (the Kolmogorov sum rules hold on decoherent families). This unifies requirements 1 and 5: $\lambda$ = a maximal consistent recorded decoherent history, $\mu$ = the decoherence-functional measure on history space, implemented locally through Type II trace densities.

\emph{No-signaling survives, conditional on one lemma.} Boundary-data $\lambda$ does not threaten the proved operational no-signaling \textbf{provided} one proves a \textbf{screen-local marginal lemma}: for every local AQFT instrument, the $\mu$-pushforward of $\lambda$ to its outcome records equals the AQFT Born state $\omega_\Phi$ and is independent of spacelike-separated instrument choices, i.e. $\sum_y \omega_\Phi(E_x^a F_y^b) = \omega_\Phi(E_x^a)$ using microcausality $[E_x^a, F_y^b]=0$ and completeness $\sum_y F_y^b = 1$. The named dangers to exclude are: screen-dependence (two screens reconstructing different records for the same diamond), remote-setting dependence of \emph{local} marginals, and any access to sub-$\lambda$ boundary data beyond ordinary QFT observables (which would enable postselection-style signaling). \textbf{One reformulation is forced here:} if $\lambda$ is a \emph{whole-history} object that already contains the setting records, then Bell measurement independence $\rho(\lambda \mid a,b) = \rho(\lambda)$ is ill-posed as stated; MI must be re-expressed for the \textbf{past / common-cause component} of $\lambda$ (the conditional history measure), not for the global history. This is the relativistic refinement of the §6.9 Bell stance, and is itself an open sub-task.

\emph{Borrowed machinery and what breaks.} The reduction "$\lambda$ = maximal consistent recorded decoherent history (Gell-Mann–Hartle; Isham History Projection Operator formalism, already 4-D and time-neutral) + $\mu$ = decoherence-functional measure" lets QIQT-H import the spacetime-QM / decoherent-histories apparatus wholesale, but \textbf{seven repairs are required}, and naming them is part of the honest agenda: (1) the realm-selection / Dowker–Kent "many incompatible quasiclassical sets" problem, there is no automatic unique maximal realm; (2) contrary-inference pathologies, controlled only by restricting to \emph{recorded} quasiclassical histories; (3) relativistic frame-dependence of time-ordered class operators $C_\alpha = P_{\alpha_n}(t_n)\cdots P_{\alpha_1}(t_1)$, must be replaced by genuinely spacetime / Schwinger–Keldysh / Tomonaga–Schwinger coarse-grainings; (4) Type III$_1$ obstruction to exact projection-valued local histories (needs split / Type-II-core / coarse-grained pointer algebras); (5) exact decoherence is rare, approximate decoherence needs explicit $\epsilon(D)$ bounds \emph{stable under restriction and gluing}; (6) \textbf{the holographic capacity bound $\log \#\mathrm{Atoms}(\mathsf{B}_\Phi(D)) \le Q_D$ is NOT automatic in Gell-Mann–Hartle}, it is QIQT-H's genuinely new ingredient and must be proved compatible with refinement; (7) the measurement-independence reformulation of the previous paragraph.

\emph{\textbf{The linchpin: a Poincaré-equivariant holographic recorded-history sheaf theorem.}} The single theorem that, if proved, discharges the most requirements at once. For a Poincaré-covariant Haag–Kastler net $(\mathcal A, U)$ satisfying locality, isotony, additivity, time-slice, Haag duality, split/nuclearity, and Bisognano–Wichmann modular covariance, and for every admissible global state $\Phi$, there exist: (1) finite Boolean record algebras $\mathsf{B}_\Phi(D) \subset \mathrm{Proj}(\widehat{\mathcal A}(D))$ for every causal diamond $D$; (2) boundary record algebras $\mathsf{B}_{\Phi,\partial}(D)$ with natural reconstruction isomorphisms $\mathsf{B}_{\Phi,\partial}(D) \cong \mathsf{B}_\Phi(D)$; (3) capacity and decoherence bounds $\log\#\mathrm{Atoms}(\mathsf{B}_\Phi(D)) \le Q_D$ and $\mathrm{DecErr}(D) \le \epsilon(D)$; (4) functorial restriction maps making $X_\Phi(D) = \mathrm{Stone}(\mathsf{B}_\Phi(D))$ a sheaf/stack over $\mathsf{Diam}$; (5) a nonempty global-section space $\Lambda_\Phi = \Gamma(X_\Phi)$; (6) a probability measure $\mu_\Phi$ on $\Lambda_\Phi$ given on cylinder events by the decoherence-functional / Born weights $\mu_\Phi(\lambda_D = i) = \omega_{\Phi,D}(P_i^D) = \tau_D(h_{\Phi,D}P_i^D)$; (7) Kolmogorov consistency $\sum_{j : \rho_{D,K}(j)=i} \mu_\Phi(\lambda_D = j) = \mu_\Phi(\lambda_K = i)$; (8) Poincaré equivariance $X_{U_g\Phi}(gD) \cong X_\Phi(D)$ and $g_*\mu_\Phi = \mu_{U_g\Phi}$; (9) operational no-signaling marginals for all local AQFT instruments. The selector is then mere evaluation, $A_D[\Phi,\lambda] = \lambda_D$, and the covariance identity (requirement 3) follows immediately: $A_{gD}[U_g\Phi, g\lambda] = g\cdot A_D[\Phi,\lambda]$. \textbf{This theorem is the real target of Open Problem 3b.} Its nine hypotheses are exactly the framework's five requirements made precise, plus the AQFT regularity conditions; proving it (even for free fields, as a first case) would convert the Lorentz-covariance promissory note into a structural result.

\emph{Formalization note (Lean), now partially carried out.} The linchpin's discrete skeleton is more amenable to machine verification than its analytic core, and the project's existing strategy, prove the order-theoretic / combinatorial scaffold, name the functional-analytic inputs as interface axioms (cf. \texttt{AxiomAudit.lean}), applies directly. Four Lean modules now realize this, all compiling and axiom-audited (full build $\approx$ 2994 jobs; standard axioms only on every headline theorem):

\begin{itemize}
  \item \textbf{\texttt{LorentzSelection.lean}}: the poset/sheaf layer: the causal-diamond poset $\mathsf{Diam}$, the record-sector presheaf $X_\Phi$ with functorial restriction, the global section $\lambda \in \Gamma(X_\Phi)$, \texttt{bulk\_overlap\_agreement} (Correction A; \emph{no axioms}), and the \textbf{evaluation-gives-covariance theorem} $A_{gD}[U_g\Phi,g\lambda] = g\cdot A_D[\Phi,\lambda]$, proved from equivariant naturality, depending on standard axioms only and on \emph{none} of the (now-retired) AQFT axioms. The Poincaré action is an \texttt{OrderIso} of diamonds; this makes the pushed-forward section's consistency a \emph{proved theorem} (\texttt{actSection\_consistent}) rather than an axiom, the earlier transport-bookkeeping axiom was thereby eliminated. The four AQFT analytic inputs are no longer opaque axioms but an explicit \texttt{RecordedHistoryNet} structure (holographic finiteness bound, boundary-reconstruction natural iso, decoherence-functional probability weights, projective/no-signaling marginal), over which the covariance result becomes the conditional theorems \texttt{covariant\_selection\_of\_net} / \texttt{covariant\_selection\_exists} and no-signaling becomes the theorem \texttt{net\_no\_signaling}, leaving the module \textbf{axiom-free}.
  \item \textbf{\texttt{FiniteModularTheory.lean}}: the finite-matrix modular engine: the inner conjugation automorphism, the KMS \emph{boundary} identity $\omega(xy)=\omega(y\,\sigma(x))$ from trace cyclicity, and, the genuine modular \emph{data}, the \textbf{real-time modular flow} $\sigma_t(x)=\rho^{it}x\rho^{-it}$ in the diagonal case, proved to be a real one-parameter group ($\sigma_s\circ\sigma_t=\sigma_{s+t}$). This is honestly the \emph{finite Type-I} modular skeleton, not the continuum Tomita–Takesaki data (GNS $\Delta$, $J$, $S=J\Delta^{1/2}$, the implementation theorem, KMS strip-analyticity remain future work).
  \item \textbf{\texttt{FreeFieldRecord.lean}}: the free-field finite-mode instance: the holographic record count $\log_2 \#\mathrm{Atoms} = N \le Q_R$; the \textbf{decoherence factorization} $\langle E_\alpha|E_\beta\rangle = \prod_i \langle E_{\alpha_i}|E_{\beta_i}\rangle$ with the \emph{derived} bound $|\langle E_\alpha|E_\beta\rangle| \le q^{\,d_H(\alpha,\beta)} \to 0$ (Gaussian/Wick suppression, derived from the mode product rather than postulated); and the finite-mode Lorentz action on record sectors (a one-parameter group of bijections, capacity-covariant).
  \item \textbf{\texttt{LorentzSelectionStrong.lean}}: the \emph{strengthening} pass that answers the adversarial-review caveats (i)–(ii) above with genuine content (the review's items A–G), keeping the interface \emph{conditional} but making it \textbf{rigid} and \textbf{connected} rather than vacuous, and still axiom-free: \textbf{(A)} an \emph{externalized} \texttt{GeometrySpec} fixes the area-law budget $N$ and the boundary presheaf \emph{before} a net is chosen, so the holographic bound and reconstruction can no longer be satisfied by the cheap $N:=\#X_\Phi$, $P_b:=P$ choices (consuming theorems: \texttt{card\_le\_of\_le}, using holographic monotonicity, and \texttt{reconSection}, the screen-encoded selected history as a consistent boundary section); \textbf{(B)} a genuine Poincaré \emph{group} action \texttt{GroupAction G P} with identity/composition laws, giving covariance for \emph{every} $g$ (\texttt{group\_evaluation\_covariance}) and law-consuming identities (\texttt{act\_one\_diam}, \texttt{act\_mul\_diam}), replacing the single order-automorphism; \textbf{(C)} \texttt{measure\_pushforward\_total}, a theorem that \emph{consumes} the measure-covariance hypothesis (a $g$-covariant weight has $g$-invariant total mass, by reindexing along the sector equivalence), the measure-level content the bare covariance theorem ignored; and \textbf{(G)} the \textbf{Born link} (the anti-vacuity lock): a \texttt{BornData} pins the weights to Born values $\omega_\Phi(P_i^D)=\langle\psi_D|E_x^D|\psi_D\rangle$ of an actual unit state on a resolution of identity, whence \emph{normalization is a theorem} (\texttt{bornω\_sum\_one} / \texttt{bornωRe\_sum\_one}: $\sum_x\omega = \langle\psi|\,{\textstyle\sum_x}E_x\,|\psi\rangle = \langle\psi|\psi\rangle = 1$) derived from the already-discharged Born functional (\texttt{GleasonSelector.born\_add}/\texttt{born\_one}), tying the Lorentz strand to the axiom-free Gleason strand. A second pass (after a follow-up adversarial review) closes the three remaining gaps that review named: \textbf{(probabilities, not just affine normalization)} a \texttt{PVMData} adds the projection hypotheses $E_x^\dagger=E_x$, $E_x^2=E_x$ that \texttt{BornData} lacked, whence \texttt{born\_posSemidef\_nonneg} ($\langle\psi|E|\psi\rangle\ge 0$ for positive-semidefinite $E$, via \texttt{Matrix.PosSemidef.dotProduct\_mulVec\_nonneg}) gives \texttt{pvm\_bornωRe\_nonneg} and \texttt{pvm\_bornω\_im\_zero}, so \texttt{pvm\_isProbability} certifies the weights are a genuine finite probability distribution (nonnegative \emph{and} summing to $1$), not the signed/complex "weights" the bare $\sum_x E_x=1$ would have permitted; \textbf{(a real representation)} the γ-cocycle laws \texttt{IsRepOne}/\texttt{IsRepMul} (with the fibre transport \texttt{fibCast}) upgrade the family $\{\gamma_g\}$ from a per-element natural equivalence to a genuine action on the record fibres, with \texttt{γ\_cocycle\_apply} ($\gamma_{g_1 g_2}=\gamma_{g_2}\circ\gamma_{g_1}$ up to the diamond cast) as the consuming theorem; and \textbf{(per-cell covariance)} \texttt{measure\_pushforward\_cell} strengthens the total-mass statement to per-sector equality $\omega_{gD}(y)=\omega_D(\gamma_g^{-1}y)$. A third pass then \textbf{discharges the measure-covariance hypothesis itself} (\texttt{hcov}, previously the marquee \emph{assumption}): the unconditional theorem \texttt{born\_unitary\_invariant} ($\langle U\psi|UEU^\dagger|U\psi\rangle=\langle\psi|E|\psi\rangle$ for $U^\dagger U=1$, proved by the adjoint identities \texttt{star\_mulVec}/\texttt{dotProduct\_mulVec}/\texttt{vecMul\_vecMul}) is the engine; a \texttt{UnitaryCovariance} records the \emph{lower-level} equivariance, the state transforms $\psi_{gD}=U_{g,D}\psi_D$ and the effects conjugate $E_{gD}(\gamma_g x)=U_{g,D}E_x^D U_{g,D}^\dagger$, from which \texttt{ubornω\_covariant} \emph{derives} $\omega_{gD}(\gamma_g x)=\omega_D(x)$ as a \textbf{theorem}, and \texttt{ubornω\_pushforward\_cell} / \texttt{ubornω\_total\_invariant} then make the per-cell and total-mass covariances \textbf{unconditional on \texttt{hcov}}. This converts caveats (i)–(ii) from "the analytic fields are ignored / non-rigid" into "the fields are externally rigid, consumed by theorems, the weights are certified probabilities, and their Poincaré covariance is \emph{derived from unitarity} rather than assumed." A fourth pass then threads the γ-cocycle onto the \emph{selection itself}, \texttt{selection\_cocycle} proves $\gamma_{g_1 g_2}(\lambda_D)=\gamma_{g_2}\big(\text{selector}(g_1\!\cdot\!\lambda)(g_1 D)\big)$ (up to the diamond cast), i.e. multi-step transformed selections agree with the one-step transform, consuming both \texttt{IsRepMul} and \texttt{group\_evaluation\_covariance}, and fuses the two strands into a single object: \texttt{UniformPVMData} carries the projection hypotheses \emph{and} the uniform-dimension Born data, and \texttt{upvm\_covariant\_probability} certifies the weights on each diamond are simultaneously (i) nonnegative, (ii) summing to $1$, and (iii) Poincaré-covariant $\omega_{gD}(\gamma_g x)=\omega_D(x)$, a genuine \textbf{covariant probability distribution} whose covariance is \emph{derived}. A coherence check is also discharged as a theorem: \texttt{proj\_conj\_unitary} (unitary conjugation preserves a Hermitian projection) gives \texttt{E\_cov\_preserves\_proj}, the boosted effects $E_{gD}(\gamma_g x)=U E_x^D U^\dagger$ remain a PVM, so \texttt{UnitaryCovariance} over a \texttt{UniformPVMData} is internally consistent (not over-determined); and \texttt{covariantProbability\_of\_unitaryPVM} packages the three facts as a \texttt{CovariantProbability} bundle. A final pass discharges the two remaining reachable in-interface statements. \textbf{(Section-object group action)} \texttt{LorentzSelection} now exposes a public spec \texttt{actSection\_val} (the \texttt{actVal} field is private) plus \texttt{GlobalSection.ext} and \texttt{GlobalSection.val\_cast}; on top of these \texttt{actSection\_one} ($1\cdot\lambda=\lambda$) and \texttt{actSection\_mul} ($(g_1 g_2)\cdot\lambda = g_2\cdot(g_1\cdot\lambda)$) are \emph{proved}, so \texttt{actSection} is a genuine \textbf{group action on $\Gamma(X)$}, of which the earlier selector-level \texttt{selection\_cocycle} was only the evaluation shadow. The transport casts are handled cleanly by \texttt{actSection\_val\_act'} (the cast-free evaluation $(g\cdot\lambda)(gD)=\gamma_g(\lambda_D)$), \texttt{val\_cast}, and \texttt{fibCast\_symm\_castSector}, no raw \texttt{HEq}, the route the adversarial review outlined. \textbf{(Full PVM preservation)} \texttt{unitary\_preserves\_resolution} proves $U(\sum_x E_x)U^\dagger = U\,U^\dagger = 1$ (using $UU^\dagger=1$ derived from the finite-square $U^\dagger U=1$), so the unitary transport preserves the \emph{resolution of identity}, not just individual projections, the boosted effects are a \emph{full} PVM. What remains genuinely open is now only: the boundary-action/reconstruction compatibility, and, the principal gap, the continuum \emph{realization} (existence of such a net, with its unitary Poincaré transport, from a genuine relativistic QFT: the Type III$_1$ / Tomita–Takesaki problem). Every reachable assumption-to-theorem conversion in the finite conditional interface is complete: the interface is rigid, connected, a \emph{covariant probability distribution} with covariance derived from unitarity, coherent under transport (a full PVM preserved by the action), and carries a genuine group action on its global sections; its \emph{instantiation} from QFT is the open problem.
  \item \textbf{\texttt{LorentzWitness.lean}}: a \textbf{concrete non-trivial model} of the strengthened interface, settling the \emph{vacuity} question the review raised (the structure also admits the trivial one-point net, so "a model exists" alone is empty). Over a single causal diamond with a two-element record fibre \texttt{Fin 2}, a two-dimensional Hilbert space, the rational unit state $\psi=(3/5,4/5)$, and the diagonal projective measurement, it builds a full \texttt{UniformPVMData} + \texttt{UnitaryCovariance} and hence (via \texttt{covariantProbability\_of\_unitaryPVM}) a \texttt{CovariantProbability}, with \texttt{witness\_nondegenerate} proving the Born weights are $\omega(0)=9/25,\ \omega(1)=16/25 \in (0,1)$, a genuinely \emph{spread} two-outcome distribution rather than a point mass, and \texttt{witness\_fibre\_card} confirming the fibre has two elements (not one). So the conditional interface has real content: it is satisfiable by non-degenerate models, not only the trivial net (standard Lean axioms only). A second witness exercises the \emph{covariance} machinery on a genuine orbit: two causal diamonds with the indiscrete preorder (so every permutation is an order-isomorphism), the group $G=\mathrm{Perm}$ permuting them, and homogeneous PVM data, \texttt{witness2\_covariantProbability} gives a \texttt{CovariantProbability} over this non-trivial group, and \texttt{witness2\_action\_nontrivial} proves the swap moves diamond $d_0$ to $d_1$ (the action is genuinely non-trivial, unlike Witness A's trivial group). \emph{Honest scope:} these are finite, rational (no $\sqrt 2$), homogeneous models that prove the interface is non-vacuous and its covariance applies to a non-trivial orbit; the continuum QFT realization (Type III$_1$ / Tomita–Takesaki) remains the genuine open content, untouched.
\end{itemize}

What was previously deferred to \emph{four opaque interface axioms} in \texttt{LorentzSelection.lean}, corresponding to the existence of the finite record algebras with the holographic bound (linchpin hypotheses 1–3), the boundary-reconstruction isomorphism (2), the Born/decoherence-functional measure with its Kolmogorov consistency and $\sigma$-additive extension (6), and the no-signaling marginal lemma (9), has been \textbf{removed from the axiom budget}: those four \texttt{axiom \_ : Prop} placeholders (which asserted nothing and were used by no theorem) are \textbf{retired} and replaced by an explicit \texttt{RecordedHistoryNet} structure that \emph{writes the content down} as type-checked data and propositions (a holographic finiteness/cardinality bound, a boundary-reconstruction natural isomorphism commuting with restriction, decoherence-functional probability weights, and the projective/no-signaling marginal). The covariance result is then a \emph{conditional theorem over this structure} (\texttt{covariant\_selection\_of\_net}, \texttt{covariant\_selection\_exists}), and the projective marginal yields the convenience theorem \texttt{net\_no\_signaling} (the marginal onto a sub-diamond $K$ is the same regardless of which larger diamond $D\supseteq K$ it is computed from). The whole \texttt{LorentzSelection} module now adds \textbf{zero project axioms} (every headline theorem depends on the standard Lean axioms only, per \texttt{AxiomAudit.lean}; the project-wide raw axiom count drops $44\to 40$).

\emph{What this does and does not establish (adversarial-review caveat, GPT-5.5-pro).} This is the corpus's \emph{interface-as-hypothesis, not axiom} discipline applied to the AQFT inputs, and as such it is a genuine improvement in formal hygiene, but it must not be overstated, and three honest limitations are recorded here. \textbf{(i) The covariance theorem is evaluation-equivariance, not covariance of the physics.} \texttt{covariant\_selection\_of\_net} consumes only the presheaf field \texttt{net.P}; its proof is the structural identity \texttt{evaluation\_covariance}, and the analytic fields (\texttt{N}, \texttt{recon}, \texttt{ω}, \texttt{ω\_marg}) play no role in it, so what is machine-checked is "a transformed global section, evaluated at the transformed diamond, equals the transform of the evaluation," not covariance of the decoherence measure or of a \emph{fixed} selection rule. \textbf{(ii) Several fields are non-rigid as written, so the structure admits degenerate models.} Because \texttt{N}, \texttt{Pb}/\texttt{recon}, and \texttt{ω} are chosen \emph{within} the same structure, the holographic bound is satisfiable by taking $N(D) := \#X_\Phi(D)$, boundary reconstruction by $P_b := P,\ \mathrm{recon} := \mathrm{id}$, and the weights by any finite PMF; in particular the one-point net ($X_\Phi(D) := \mathbf 1$) satisfies every field. **Consequently the bare existence statement "a \texttt{RecordedHistoryNet} exists" is \emph{not} the open problem, it is trivially true.** The genuine Open Problem 3b is the \emph{realization} problem: existence of a net whose data are \emph{extracted from a fixed relativistic QFT and geometry}, the area-law $N(D)=\lfloor e^{Q_D}\rfloor$ and boundary algebra fixed externally, and the weights pinned to Born values $\omega_\Phi(P_i^D)$ of an actual global state $\Phi$, together with a genuine Poincaré \emph{group} action (the current \texttt{PoincareAction} is a single order-automorphism, not a representation) under which the measure is equivariant. \textbf{(iii) \texttt{net\_no\_signaling} is a two-line rewrite of the assumed marginal}, not a derivation of operational no-signaling from microcausality; the real content remains the \emph{assumption} \texttt{ω\_marg}. The honest status, then: the four opaque axioms are gone and the finite combinatorial skeleton (functorial restriction, global sections, evaluation-equivariance, the finite Type-I modular flow, and the free-field finite-mode record structure) is machine-checked and axiom-free; but the substantive Lorentz-covariance content, Born-pinned covariant measure, externally-rigid holographic/boundary data, a true group action, and the Type III$_1$ / Tomita–Takesaki / Haagerup-$L^p$ \emph{realization}, remains open, and is the agenda for the next formalization pass (§11.4 strengthening targets: externalize the geometry spec, replace the single automorphism by a group action with equivariant weights, and tie $\omega$ to the already-formalized Born functional).

\textbf{Open Problem 4, Quantitative form of $\epsilon(R)$.}
\emph{(Subsidiary / superseded, §1.1a.)} The subsidiary resolution-floor quantity satisfies $\epsilon(R) > 0$ (formalized in \texttt{Resolution.lean}). One could derive an explicit functional form $\epsilon(R) = f(\text{geometry}, Q_R, \text{macroscopic record dimension})$, but this amplitude-resolution reading is no longer load-bearing; the framework's conclusions rest on the entropy bound $S(\rho_R) \le Q_R$ (Theorem 3). Downstream of Open Problems 2 and 3.

\textbf{Open Problem 5, Empirical calibration of $I_0$.}
The per-record physical cost $I_0$ is a phenomenological parameter, analogous to GRW's collapse rate $\lambda$. Current best estimates (Zurek physical entropy for typical macroscopic records) give $I_0 \sim 10^{25}$ bits. The framework's prediction: a definite calibrated value of $I_0$ such that the saturation condition $I_0 \approx C(R)$ matches the observed quantum-to-classical transition. Downstream of Open Problem 4.

\medskip\hrule\medskip

\textbf{Algebraic / mathematical-physics infrastructure.} \emph{Finite-classical discharge note (Lean).} The pieces of items 6 and 9 that are \emph{finite-classical}, and were previously carried as named axioms with an in-file proof sketch, have now been \textbf{proved} and removed from the axiom budget (project raw axiom count $40 \to 37$, machine-verified, standard Lean axioms only): the \textbf{Fano-step bound} behind item 6 (\texttt{ShannonFano.H\_bound\_imp\_max\_lb}: the Rényi-$\infty \le$ Shannon relation $\max_i p_i \ge e^{-H(p)}$, with \texttt{H\_zero\_imp\_dirac} now a corollary), and the \textbf{finite relative-entropy positivity} behind item 9 (\texttt{RelEntPositivity.KL\_classical\_nonneg}: Gibbs' inequality $\sum_i p_i\log(p_i/q_i)\ge 0$ from $\log x \le x-1$). What remains axiomatic in this stack is exactly the genuinely \emph{vN-algebraic / continuum} content, Klein's inequality at the von Neumann level (operator convexity of $-\log$), Donald's identity, the quantum data-processing inequality, the CPW entropy bridge, and Araki relative entropy, none of which Mathlib yet reaches; these are the true open analytic inputs, not finite bookkeeping.

\begin{enumerate}
  \item \textbf{Operational distinguishability axiomatization (H3).} Make precise, under decoherence + Quantum Darwinism, when the macroscopic-record subalgebra admits a normal measurement instrument decoding the record index with small error, so that the Fano-type bound $H_\epsilon \le I_{\rm Hol}^R + \eta_{\rm def}$ holds. \emph{(The finite-classical Fano step: \texttt{H\_bound\_imp\_max\_lb}, is now proved, §above; the open part is the operational existence of the small-error decoding instrument at the AQFT level.)} Existing ingredients: spectrum broadcast structure (Zurek, Brandão-Piani-Horodecki), redundancy plateaus, decoherent-histories medium-decoherence conditions.
  \item \textbf{Stagewise adapted process and refinement compatibility.} Construct rigorously the stagewise process $h \mapsto \omega_t^h$ of §7.6, including: existence under Cauchy-slice refinement; compatibility of $\mathfrak{R}_t(h)$ with the AQFT net's isotony; conditions under which the adapted family extends to a completed-history limit.
  \item \textbf{Functorial / isotonic crossed-product local nets.} Establish the functorial structure $D \mapsto \hat{\mathcal{A}}(D)$ as a covariant net of Type II crossed-product algebras with normal embeddings under inclusion. Needed for the modular-local bound to behave consistently under region refinement / coarse-graining.
  \item \textbf{Modular-Hamiltonian estimates beyond wedges and CFT balls.} Generalize the wedge / CFT-ball modular-Hamiltonian estimate to broader classes of regions, either via tighter enclosing-wedge monotonicity arguments or new geometric forms of $K_R^\sigma$. Without this the conditional $\delta_R \sim 10^{-27}$ estimate cannot be extended to generic laboratory geometries. \emph{(The finite-classical relative-entropy positivity that anchors the modular bound: \texttt{KL\_classical\_nonneg}, Gibbs, is now proved, §above; the open part is the continuum modular geometry and vN-level Klein inequality.)}
  \item \textbf{State-extension and reference-state issues} in the crossed-product algebra formulation: extension of states across causal-diamond boundaries, choice of $\sigma_R$ in cosmological / non-stationary backgrounds.
  \item \textbf{Locality of the (FQ)-restricted dynamics.} The (FQ) projection onto $\mathcal{H}_{\rm phys}$ must commute with the local algebras' observables for Theorem 7's no-signaling argument to apply to the \emph{restricted} (physical) theory rather than only the ambient (unrestricted) theory. The compression-locality leakage identity (formalized in \texttt{CompressionLocality.lean}) isolates this implicit constraint. Construct concrete (FQ) projections satisfying it; characterize the class of admissible projections.
\end{enumerate}

\textbf{Empirical / phenomenological:}

\begin{enumerate}
  \item \textbf{Concrete measurement-apparatus models realizing finite $\chi_R$ budgets.} Build explicit toy models (Stern-Gerlach, interferometer, optomechanical) where $\chi_R$ for the apparatus-record state is computable and the saturation regime can be approached. Without these the framework's content cannot be confronted with experiment.
  \item \textbf{Phenomenological predictions.} With $I_0$ calibrated, the framework predicts: maximum macroscopic-superposition scale (testable against Schrödinger-cat experiments with progressively larger systems); long-baseline coherence limits; specific signatures distinguishing the framework from GRW-style stochastic collapse (the framework predicts kinematic exclusion, not stochastic events).
  \item \textbf{Cosmological / horizon applications}, extend to the de Sitter static patch and black-hole horizon regions, where $C(R)$ becomes finite and saturation may be physically relevant.
\end{enumerate}

\medskip\hrule\medskip

\textbf{Strategic note.} Of the central open problems, the audit work suggests two bottleneck chains:

  • \textbf{Foundations bottleneck:} Open Problem 1 (canonical $\mu$-selection / equivariance). Now decomposed into three sub-theorems A, B, C (see above), with the A1+A2+A4+A6 strengthening pass further narrowing the load-bearing content. Two of the three sub-theorems are direct applications of existing operator-algebra machinery (Mackey-Gleason + noncommutative Radon-Nikodym); the third requires adapting the Goldstein-Struyve 2007 uniqueness theorem from Bohmian dynamics to QIQT-H's Type II / (FQ) setting. Concrete progress: the permutation-symmetry collapse component of Goldstein-Struyve Step 1 sub-lemma 1c is \textbf{PROVED} under its necessary hypothesis (\texttt{step1c\_collapse\_of\_perm\_symmetric}) along with the foundational matrix-conjugation identities (\texttt{permutation\_conj\_matrixUnit}, \texttt{diagonalU\_conj\_matrixUnit}), while the full finite Schur classification has since been \textbf{PROVED} outright (\texttt{schur\_classification\_real}); the marginal-locality step is \textbf{PROVED} as a theorem holding for any measure (\texttt{MarginalLocality.pushforward\_marginal\_local}), reducing locality to a named Hilbert→set-level bridge axiom rather than eliminating it; single-trial Chebyshev frequency concentration and the variance-addition (independence) lemma are \textbf{PROVED} (\texttt{BornConcentration}), with the $N$-trial scaling still on the LLN axiom; and the duplicate axiomatization of Goldstein-Struyve sub-steps in \texttt{FQEquivarianceUniqueness} is \textbf{CONSOLIDATED} to route through the concrete \texttt{GoldsteinStruyveFinDim} proof (\texttt{canonical\_ic\_measure\_principle} was consolidated to two project-specific axioms, down from nine — both since discharged, so it is now axiom-free), with the five false/unused Step-1 placeholder sub-axioms deleted (project axiom total 57 → 40, since driven to 0). Progress on the remaining load-bearing axioms transforms the Born section from conditional to substantive, and is the most pressing problem for foundations-of-QM defensibility.

  • \textbf{Quantitative-emergence bottleneck:} Open Problem 2 (reference-weight bound) $\Rightarrow$ Open Problem 4 ($\epsilon(R)$) $\Rightarrow$ Open Problem 5 ($I_0$ calibration). Progress on Open Problem 2 unlocks the quantitative QIQT-H pipeline. (The former Open Problem 3, Macroscopic Definiteness / Branch-Summed bound, is \textbf{withdrawn} — see §11.4; single-record definiteness is $\lambda$'s role, not a capacity exclusion.)

These are the two genuinely open research directions left after the Lean audit work has clarified the deductive boundary. The remaining items (6–14) are infrastructure and phenomenology that follow once the bottlenecks are addressed.

\subsection{Claim-to-Lean theorem matrix}

The following table maps each major paper claim to its Lean theorem, with explicit status. Status labels:

  • \textbf{U}: unconditional (no project-specific axioms; depends only on standard Lean axioms \texttt{propext}, \texttt{Classical.choice}, \texttt{Quot.sound}).
  • \textbf{C}: conditional on a standard external theorem (Mathlib-citable; axiomatized at clean interface in the project).
  • \textbf{P}: programmatic interface axiom (specific framework or AQFT-level axiom; concrete Lean discharge is a multi-day to multi-week task).
  • \textbf{N}: negative audit / counterexample (proves a \emph{non}-derivation).

\textbf{Current status (audit): every former C and P row has been discharged to U.} The project is now axiom-free (0 project-specific axioms; CI budget 0, \texttt{scripts/axiom\_budget\_check.sh}). The C/P labels and the discharge notes below are retained as the historical record of how each interface axiom was either proved in a concrete finite model or recast as an explicit hypothesis. The continuum versions of these results remain the cited/open frontier (Open Problems 3/3b), and the \emph{physical} postulates are unaffected by the formal discharge.

\begin{center}
\small
\begin{tabularx}{\linewidth}{|Y|Y|Y|Y|}
\hline
\textbf{Paper claim} & \textbf{Lean theorem} & \textbf{Status} & \textbf{Notes} \\
\hline
Theorem 3 (subsidiary resolution floor; superseded, §1.1a) & \texttt{Resolution.eps\_pos} & U & $\varepsilon(R) > 0$ from finite $Q_R$ \\
\hline
Lemma 1 (near-extreme indistinguishability) & \texttt{Resolution.lemma1\_zero} & U & Discrete-bin formalization \\
\hline
Theorem 6 (effective definiteness) & \texttt{Theorem6.effective\_definiteness} & U & Donald-bound chain \\
\hline
Theorem 6 inner step & \texttt{Theorem6.BranchData.holevo\_le\_capacity} & U & $I_{\rm Hol}^R \le C(R)$ from $\chi_R \le C(R)$ and Klein \\
\hline
Theorem 7 (no-signaling) & \texttt{Theorem7.Setup.no\_signaling} & U & From microcausality + locality \\
\hline
Donald's identity & \texttt{Donald.donald\_identity} & C & Conditional on 3 cross-entropy axioms (A1, A2, A3) \\
\hline
Microcausality ⇒ locality (unitary) & \texttt{UnitarityLocality.locality\_of\_conjugation} & U & StarRing argument \\
\hline
Microcausality ⇒ locality (Kraus) & \texttt{KrausLocality.kraus\_channel\_aliceFixing} & U & Generalization to CPTP channels \\
\hline
Bell/CHSH inequality (LHV bound) & \texttt{Bell.LHVModel.chsh\_le\_two} & U & Finite probability + ±1 algebra \\
\hline
Tsirelson saturation by singlet & \texttt{Tsirelson.singlet\_chsh\_abs\_gt\_two} & U & Explicit 4D real construction; attains $2\sqrt{2}$ exactly (does \emph{not} prove the universal upper bound) \\
\hline
(H1) ⇏ (H2) & \texttt{H1H2Audit.H1\_does\_not\_imply\_H2} & N & Classical KL countermodel \\
\hline
(H2) ⇔ reference-weight bound & \texttt{H1H2Audit.H2\_iff\_reference\_weight} & U & Sharp structural reformulation \\
\hline
Decoherence ⇏ concentration & \texttt{NoConcentration.audit\_conclusion} & N & Linear unitary preserves branch weights \\
\hline
$S_{\rm ren}$ (paper) $\neq$ $S_{\rm ren}^{\rm CPW}$ & \texttt{EntropyBridge.fq\_ambiguity\_counterexample} & N & Classical KL counterexample distinguishing the two \\
\hline
Branch-summed cost > Shannon possible & \texttt{BranchLedger.branchSummed\_not\_bounded\_by\_Shannon} & N & Binary uniform witness \\
\hline
FQ-literal finite ⇒ trivial dynamics & \texttt{FQDynamicsNoGo.finite\_admissible\_flow\_fixed} & U & Continuous flow on finite T2 space (only, not arbitrary finite-info dynamics) \\
\hline
Compression-locality leakage identity & \texttt{CompressionLocality.compressed\_commutator\_with\_commute} & U & Pure ring algebra \\
\hline
Any $p$ realizable (no Born from nothing) & \texttt{NoBornFromNothing.exists\_probability\_realizing} & N & Concrete section-based construction \\
\hline
Support preservation $\neq$ Born equivariance & \texttt{EquivarianceGap.support\_preservation\_does\_not\_imply\_measure\_preservation} & N & Concrete \texttt{Fin 2} counterexample \\
\hline
Born typicality (mean form) & \texttt{BornTypicality.born\_mean\_conditional} & U & Conditional on canonical $\mu_\rho$ with Born marginal \\
\hline
Operational data ⇏ unique IC measure & \texttt{OperationalNoGo.operational\_data\_insufficient} & N & Concrete \texttt{Fin 3} witness \\
\hline
Sub-thm A (Born from non-contextuality) & \texttt{EffectGleason.finite\_effect\_gleason} & U & \textbf{The \texttt{TypicalityMackeyGleason} module was retired and deleted.} Its Mackey-Gleason packaging is superseded by the axiom-free finite effect-Gleason theorem (Born forced from positivity + normalization + ray-certainty) \\
\hline
Sub-thm C (FQ-equivariance uniqueness) & \texttt{FQEquivarianceUniqueness.canonical\_ic\_measure\_principle} & U & \textbf{Now axiom-free:} the former dependencies (Schur classification, tensor multiplicativity) are both proved, so \texttt{canonical\_ic\_measure\_principle} depends only on the standard Lean axioms \\
\hline
GS Step 2 (normalization) & \texttt{GoldsteinStruyveFinDim.step2\_normalization} & U & Concrete \texttt{Matrix.trace} computation \\
\hline
GS Step 4 (non-degeneracy) & \texttt{GoldsteinStruyveFinDim.step4\_nondegeneracy} & U & Direct case analysis \\
\hline
GS Step 3 algebraic core & \texttt{GoldsteinStruyveStep3.step3\_algebraic\_core} & U & Polynomial identity via \texttt{ring} \\
\hline
GS Step 3 Kronecker bridge & \texttt{GoldsteinStruyveKronecker.step3\_kronecker\_bridge} & U & Concrete \texttt{d = 2} witness with E₁₁ ⊗ E₁₁ \\
\hline
GS Step 1 (Schur classification) & \texttt{GoldsteinStruyveStep1.step1\_via\_sub\_lemmas} & U & \textbf{The \texttt{step1\_schur\_classification} interface axiom is RETIRED} (Step 1 is now proved outright, see \texttt{schur\_classification\_real} below). The five former placeholder sub-axioms (1a–1e), three literally false, all unused, were \textbf{deleted}; \texttt{permutation\_conj\_matrixUnit}, \texttt{diagonalU\_conj\_matrixUnit}, and \texttt{step1c\_collapse\_of\_perm\_symmetric} are the proved building blocks \\
\hline
GS Step 1c (coefficient unification, partial) & \texttt{GoldsteinStruyveStep1.step1c\_collapse\_of\_perm\_symmetric} & U & \textbf{NEW (A2):} the \emph{permutation-symmetry collapse component} of Step 1c, proved via \texttt{Equiv.swap} transitivity under its necessary permutation-symmetry hypothesis. The prior bare-coefficient axiom (no hypothesis) was literally false and has been \textbf{deleted}. The \emph{full} Schur classification (Step 1) remains the single named interface axiom \\
\hline
Permutation conjugation of matrix units & \texttt{GoldsteinStruyveStep1.permutation\_conj\_matrixUnit} & U & \textbf{NEW (A2):} $P_\sigma \cdot E_{ij} \cdot P_\sigma^* = E_{\sigma(i),\sigma(j)}$ via direct matrix-entry computation \\
\hline
Diagonal-unitary conjugation of matrix units & \texttt{GoldsteinStruyveStep1.diagonalU\_conj\_matrixUnit} & U & \textbf{NEW (A2):} $D(z)\cdot E_{ij}\cdot D(z)^* = (z_i\,\overline{z_j})\,E_{ij}$, self-contained diagonal parameterization (no \texttt{Complex.exp}) \\
\hline
Combined GS finite-dim theorem & \texttt{GoldsteinStruyveFinDim.goldstein\_struyve\_findim} & U & Proved by composition; Steps 1 and 3 are now proved axiom-free, so it inherits no axioms \\
\hline
Marginal locality (pure pushforward) & \texttt{MarginalLocality.pushforward\_marginal\_local} & U & \textbf{NEW (A1):} $r\circ T = r \Rightarrow r_*(T_*\mu) = r_*\mu$ for \textbf{any} $\mu$, no equivariance assumption; reindexing argument \\
\hline
Marginal locality (equivariant corollary) & \texttt{MarginalLocality.alice\_marginal\_unchanged\_by\_bob\_dynamics} & U & \textbf{NEW (A1):} thin corollary of the pure version; discharges the \emph{marginal-locality step} (not the Hilbert→set bridge, which stays a named axiom) \\
\hline
Chebyshev tail bound on finite probability space & \texttt{BornConcentration.chebyshev\_tail\_bound} & U & \textbf{NEW (A4):} abstract Chebyshev via variance-dominates-restriction \\
\hline
Bernoulli variance $p(1-p)$ & \texttt{BornConcentration.bernoulli\_variance} & U & \textbf{NEW (A4):} direct two-term sum computation \\
\hline
Variance addition (independence) & \texttt{BornConcentration.variance\_add\_of\_product} & U & \textbf{NEW (A4):} $\mathrm{Var}(X_1+X_2)=\mathrm{Var}\,X_1+\mathrm{Var}\,X_2$ on a product measure, the reusable lemma the $N$-fold LLN induction would iterate \\
\hline
Two-trial Bernoulli variance $2p(1-p)$ & \texttt{BornConcentration.two\_trial\_bernoulli\_variance} & U & \textbf{NEW (A4):} concrete instance of variance-addition \\
\hline
Born single-trial frequency concentration & \texttt{BornConcentration.born\_chebyshev\_single\_trial} & U & \textbf{NEW (A4):} single-trial bound; the $N$-trial scaling to full "Born frequencies" still rests on the LLN interface axiom \\
\hline
Minimality witnesses for P1, P2, P3 & \texttt{BornMinimalityTable.\{P1, P2, P3\}\_…\_necessary} & N + U & \textbf{NEW (A6):} unified independence package showing each Born sub-axiom is necessary by concrete finite countermodel (minimality \emph{relative to the current decomposition}) \\
\hline
Locality (P4) is reducible & \texttt{BornMinimalityTable.P4\_locality\_reducible\_to\_equivariance} & U & \textbf{NEW (A6):} marginal-locality theoremized (modulo the named bridge axiom); not an independent Born-selection sub-axiom \\
\hline
\textbf{GS Step 1 (Schur classification)} & \texttt{GoldsteinStruyveStep1.schur\_classification\_real} & U & \textbf{2026-06 UPDATE:} the former interface axiom \texttt{step1\_schur\_classification} is \textbf{RETIRED}, Step 1 is now fully proved axiom-free (Hadamard relation → linear assembly → reality via \texttt{IsHermitianPreserving}). Step 3 likewise proved. The finite Goldstein–Struyve chain is axiom-free; the project axiom budget has since reached \textbf{0} \\
\hline
Finite effect (Busch) Gleason & \texttt{EffectGleason.finite\_effect\_gleason} & U & normalized + positive + coexistent-additive functional on effects $=\mathrm{tr}(\rho\,\cdot)$, Born from positivity, finite-dim \\
\hline
Finite weak LLN (Chebyshev) & \texttt{BornTypicalityFinite.chebyshev\_freq} / \texttt{chebyshev\_freq\_union\_le} & U & $P(\lvert\mathrm{freq}-p_k\rvert\ge\varepsilon)\le p_k(1-p_k)/(N\varepsilon^2)$; union bound $\le 1/(N\varepsilon^2)$ \\
\hline
Quantum bridge (tensor trace) & \texttt{BornTypicalityQuantum.quantumWeight\_eq\_w} & U & $N$-copy product-measurement weight $=$ classical product weight of the Born vector \\
\hline
Product-measure uniqueness & \texttt{BornMeasureUniqueness.product\_born\_measure\_unique} & U & the product Born measure is the unique additive history measure with the Born marginals (independence explicit) \\
\hline
\textbf{C1} one record → one \textbf{value} & \texttt{PointerValue.existsUnique\_actualValue} / \texttt{existsUnique\_actualHistory} & U & in the finite additive-cost model, ≤1 distinct-valued coactual record (λ selects it); many same-value records may coexist \\
\hline
\textbf{C2} one-site Born (vector state) & \texttt{OneSiteBorn.vectorState\_eq\_weight} & U & $\mathrm{Re}\langle\psi,E_a\psi\rangle=\lVert E_a\psi\rVert^2$ on a PVM projection \\
\hline
\textbf{Single-trial Born from non-contextuality} & \texttt{OneSiteGleason.oneSite\_forced} & U & non-contextual \texttt{EffectMeasure} $\Rightarrow$ $\mu(P_a)=\mathrm{Re}\,\mathrm{tr}(\rho P_a)$, \textbf{forced} by effect-Gleason (strong premise; not Born by hand) \\
\hline
Non-contextual $=$ Born (converse) & \texttt{OneSiteGleason.traceEffectMeasure} & U & every density matrix is a non-contextual \texttt{EffectMeasure} \\
\hline
Subadditive capacity exclusion & \texttt{CoreNoCollapse.joint\_coactual\_subsingleton}; \texttt{OrthogonalCapacity.pair\_exceeds} & U & monotone joint cost + pairwise overflow (overflow \textbf{derived} from orthogonality; no additivity assumed) \\
\hline
\textbf{Finite no-collapse Born representation} & \texttt{BornJoin.finite\_noCollapseBornRepresentation}; \texttt{BornJoinGleason.finite\_noCollapseBorn\_fromNoncontextuality} & U & \textbf{the join:} unique actual history + Born product law + Chebyshev typicality; single-trial law \textbf{forced Born} from non-contextuality. \emph{Conditional representation theorem, assumes non-contextuality + product preparation, NOT Born from $Q_{\max}$.} See \texttt{FINITE\_BORN\_REPRESENTATION.md} \\
\hline
Independence $=$ product preparation & \texttt{BornTypicalityFinite.w\_history\_factorizes} & U & product world-measure $\Rightarrow$ trial independence; independence is the irreducible product-preparation input \\
\hline
World-measure observationally free & \texttt{BornJoin.ActualEnsemble.history\_law\_unique} & U & same single-trial law $\Rightarrow$ identical actual-history statistics, independent of the posited world-measure on $\Omega$ \\
\hline
\end{tabularx}
\end{center}

\textbf{Finite no-collapse Born representation (2026-06).} The rows above marked C1/C2/non-contextuality/the
join formalize, \textbf{axiom-free}, a finite no-collapse Born representation theorem: capacity forces a
unique actual pointer value; finite effect-Gleason forces a non-contextual outcome assignment to be the
Born weight $\mathrm{tr}(\rho P_a)$; product preparation gives independent trials; the actual-value
histories carry the Born product law and are Chebyshev-typical. Two GPT-5.5-pro verification passes
confirm it is **sound, non-vacuous, but a *conditional representation theorem***, Born weights and
factorization are \emph{derived} from non-contextuality + product preparation, \textbf{not} from the capacity
bound alone. Full claim→theorem map and honest scope caveats: \texttt{FINITE\_BORN\_REPRESENTATION.md}.

\textbf{No-signaling clarification.} All "no-signaling" results in this paper are about \emph{nonselective} CPTP instruments, the unconditional marginal probability after Bob's measurement, marginalized over outcomes (equivalently, the trace-preserving sum over Kraus operators). Selective post-measurement conditional states can change remotely once Bob's outcome is known; this is not "signaling" because the conditional update requires classical communication.

\textbf{Tsirelson clarification.} The Lean module \texttt{QIQTH.Tsirelson} proves \emph{attainability} of the value $2\sqrt{2}$ by an explicit singlet-state construction in 4-dimensional real Euclidean space (which is the achievability content of Tsirelson's theorem). The \emph{upper bound} statement that CHSH $\le 2\sqrt{2}$ for \emph{every} quantum state and \emph{every} measurement choice, the full Tsirelson theorem, requires operator-norm machinery on C\textbackslash\{\}*-algebras and is \emph{not} proved in this formalization. The framework relies on the standard Tsirelson upper bound as a cited result.

\textbf{Markov suppression clarification.} The \texttt{CapacityPacking} module formalizes only the \emph{Markov-style} multi-record suppression (polynomial tail in the modular slack), explicitly \emph{not} exponential suppression. Where the paper says "exponentially suppressed", the reference is to \emph{decoherence-driven} exponential suppression of off-diagonal coherence (which is genuinely exponential, e.g., $e^{-10^{20}}$ for macroscopic environment dimensions), not to Markov-style multi-record suppression from the entropy bound alone.

\subsection{Machine-checked modular substrate (companion formalization)}

A companion Lean 4/Mathlib development (pinned commit \texttt{adf50bf}, toolchain \texttt{leanprover/lean4:v4.30.0}) machine-checks the standard modular / relative-entropy \emph{calculus} underlying the regional cost functional $\chi_R$, for the free-field coherent-state sector. Every result below has status \textbf{U} (no project-specific axioms; depends only on \texttt{propext}, \texttt{Classical.choice}, \texttt{Quot.sound}, verified in \texttt{QIQTH/AxiomAudit.lean}); none formalizes the holographic axiom (FQ), the Macroscopic Definiteness Conjecture, or Born-from-typicality. A full statement-level index with \texttt{file:line} references and the companion paper appears in \texttt{paper\_strategy/45\_Theorem\_Paper\_Index.md}.

\emph{Finite Araki relative entropy.}

  • \texttt{arakiEntropy\_eq\_relEntropy}: $S_{\rm Araki}(\rho\,\|\,\sigma)=\operatorname{tr}\rho(\log\rho-\log\sigma)$ (the Umegaki form).

\emph{Bounded Tomita-Takesaki (standard subspace).}

  • \texttt{modConj\_rvdRC\_modConj}: $JRJ=2-R$.
  • \texttt{modConj\_rvdT\_of\_mem\_K}: $J(T\xi)=(2-R)\xi$ for $\xi\in\mathcal{K}$.
  • \texttt{modUnitary} (\texttt{\_add}, \texttt{\_unitary}, \texttt{\_stronglyContinuous}): $\Delta^{it}$, with group law, unitarity, and strong continuity.

\emph{One-particle CGP relative entropy.}

  • \texttt{cgpEntropy}: $S(\xi)=-\int\log((2-r)/r)\,d\mu^R_\xi$.
  • \texttt{rvdSpec\_balance}: the CGP spectral balance.
  • \texttt{cgpEntropy\_nonneg}: $S(\xi)\ge 0$ for localized $\xi\in\mathcal{K}$.

\emph{Free-field (Fock) modular flow.}

  • \texttt{secondQuantModFlowH}: $\Gamma(\Delta^{it})$ (a one-parameter group, vacuum-fixing, strongly continuous on coherent vectors).
  • \texttt{secondQuantModFlowH\_weylH}: $\sigma_t(W(u))=W(\Delta^{it}u)$.

\emph{Coherent-state relative modular operator and reduction.}

  • \texttt{relModFlowH}: $\Delta^{it}_{W(f)\Omega\,|\,\Omega}=W(f)\,\Gamma(\Delta^{it})\,W(f)^*$.
  • \texttt{connesCocycleH}, \texttt{connesCocycleH\_chain}: $[D\omega_{W(f)\Omega}:D\omega_\Omega]_t=W(f)W(-\Delta^{it}f)$, together with the cocycle chain rule.
  • \texttt{hasDerivAt\_relModFlow\_vacuum}: $S_{\rm Araki}(\omega_{W(f)\Omega}\,\|\,\omega_\Omega)=S_{\rm CGP}(f)$ (the Casini-Grillo-Pontello identity).

\emph{The free graviton (kinematics, classical field, canonical quantization).} All linearized, free-field, flat-background — standard free-field content, machine-checked; \textbf{not} a claim of quantum gravity.

  • \texttt{tt\_decomposition} + \texttt{polarizations\_not\_gauge}: the physical polarization space (transverse-traceless modulo gauge) is \textbf{exactly 2-dimensional} (the $D(D{-}3)/2=2$ count, via the explicit gauge quotient).
  • \texttt{eR\_helicity} / \texttt{eL\_helicity}: the circular polarizations $e_\pm=e_+\!\pm i e_\times$ are eigenvectors of the rotation conjugation with eigenvalue $e^{\mp 2i\theta}$ — helicity $\pm 2$ as explicit eigenvalues.
  • \texttt{kUp\_null}, \texttt{physProj\_*}, \texttt{graviton\_null\_wave}: masslessness ($k^2=0$), the physical-state projector (the harmonic-gauge propagator numerator; idempotent, kills gauge and trace), and the classical wave equation $\partial_t^2 h=\partial_z^2 h$ for null profiles (genuine calculus).
  • \texttt{ccr}, \texttt{numberOp\_pow}, \texttt{hamiltonian\_vacuum}, \texttt{helicityOp\_plus/minus}, \texttt{annih\_coherent}, \texttt{twoPoint}: canonical quantization of the two helicity modes on the Bargmann-Fock space $\mathbb{C}[X_0,X_1]$ — the CCR $[a_i,a_j^\dagger]=\delta_{ij}$, bosonic occupation spectrum, the Hamiltonian $\omega(N_0{+}N_1{+}1)$ with zero-point energy, one-graviton helicity $\pm 2$, coherent states $a|\alpha\rangle=\alpha|\alpha\rangle$, and the vacuum two-point function $\langle 0|a_ia_j^\dagger|0\rangle=\delta_{ij}$.

\emph{The linearized bridge (conditional assembly; \texttt{BRIDGE\_PLAN.md}, 9/9 increments).} The entanglement$\,\to\,$linearized-Einstein template (Van Raamsdonk; Faulkner et al. 2013; Jacobson 2015) assembled from the pieces above: every derived step a theorem, every physical input an \textbf{explicit hypothesis} (never an axiom) — the Clausius/area law $\delta S=\delta A/4G$, the Iyer-Wald identity, the Bisognano-Wichmann/CHM identifications, scattering genericity, and the value of $G$. \textbf{Conditional; not a derivation of gravity} (background independence, the nonlinear completion, and the area law from microstate counting remain open).

  • \texttt{graviton\_solves\_linearized\_einstein} + \texttt{einstein\_iff\_dispersion}: the quantized graviton's polarization content solves linearized vacuum Einstein, and conversely $\delta G=0\Leftrightarrow k^2=0$ (Einstein forces light-cone propagation); \texttt{bianchi\_einsteinSymbol}: the linearized Bianchi identity $k^\mu(\delta G)_{\mu\nu}=0$, identically.
  • \texttt{couple\_gauge\_invariant\_iff\_conserved}: gauge invariance of the matter coupling $\int h_{\mu\nu}T^{\mu\nu}$ $\Leftrightarrow$ $\partial_\mu T^{\mu\nu}=0$; \texttt{einstein\_source\_conserved}: an Einstein-sourced $T$ is automatically conserved (the Bianchi payoff).
  • \texttt{soft\_gauge\_invariant\_iff\_ward} + \texttt{equivalence\_principle}: longitudinal decoupling of the soft graviton $\Leftrightarrow$ the Weinberg sum rule $\sum_i\eta_ig_ip_i^\mu=0$, and (for generic momenta) \textbf{all couplings equal} — Weinberg's equivalence-principle theorem at the algebraic level.
  • \texttt{boost\_flux\_unique} / \texttt{ball\_flux\_unique}: the wedge and per-ball Clausius data $\delta\langle K\rangle=-\delta S$ are \textbf{forced} (derivative uniqueness) given the carried BW/CHM identifications, riding the derived modular flow; \texttt{chmWeight\_edge\_slope}: the CHM ball kernel meets the entangling surface with unit slope (the wedge-ball $2\pi$ consistency); \texttt{cke\_*}: the diamond conformal Killing equation by real calculus.
  • \texttt{area\_probes\_separate}: geometric area probes \textbf{separate} symmetric perturbations (the FGHMVR separating-family hypothesis becomes a theorem); \texttt{screenArea\_eq\_bg\_add\_areaVar}: the holographic-screen-code area charge varies by exactly the geometric $\delta A_\Sigma(h)$ under the supplied identification.
  • \texttt{bridge\_firstLaw\_iff\_einstein} / \texttt{bridge\_conditional}: the assembled capstone — given the carried inputs, the entanglement first law $\delta S=\delta K$ at every probe $\Leftrightarrow$ the emergent perturbation satisfies linearized vacuum Einstein.

\emph{The microtheory earns its gravity (the E1--E5 campaign).} Joins between held theorems, upgrading the bridge from "assembled" to "earned in-model" — every derived step a theorem, every physical input a named hypothesis; \textbf{not} quantum gravity (the calibration, the continuum trace, and background independence remain the open walls).

  • \texttt{freeFieldWedgePackage} / \texttt{freeField\_clausius\_unconditional}: the Bisognano-Wichmann identification is \textbf{discharged} for the free field (wired from the unconditional one-particle BW theorem) — the wedge Clausius datum $\textbackslash\{\}delta\textbackslash\{\}langle K\_\{
m boost\}
angle=-\textbackslash\{\}delta S$ is forced with \emph{no external BW premise}.
  • \texttt{reconstruct} / \texttt{reconstruct\_areaVar}: the explicit decoder $h_{ii}=2A(e_i)$, $h_{ij}=A(e_i{+}e_j)-A(e_i)-A(e_j)$ — \textbf{the metric perturbation is a function of the code's own area data} (the emergence map inverted; pointwise, basis-level, symmetric sector).
  • \texttt{calibrated\_entanglement\_cut\_area\_law} / \texttt{uniform\_realizes\_area\_law}: with the area \emph{induced} from the calibrated entanglement cut (\texttt{inducedScreenArea} $:=4G\textbackslash\{\}cdot\textbackslash\{\}mathrm\{cut\}(w\_\{
m Ent\},S)$; no separate area label), the microstate count obeys $\textbackslash\{\}log\textbackslash\{\}\#\{
m microstates\}=\{
m screenArea\}/(4G)$ under the single named calibration $\textbackslash\{\}log D\_e=w\_\{
m Ent\}(e)$, realized exactly by the maximum-entropy record. The calibration carries the physical content — an explicit hypothesis, not a derivation of area from entanglement.
  • \texttt{rayFamily\_firstLaw} / \texttt{code\_equilibrium\_einstein}: a code whose per-ray state paths sit at relative-entropy equilibrium satisfies the entanglement first law at \textbf{every} probe (stationarity, per ray), hence — through the assembled skeleton — \textbf{linearized vacuum Einstein}: Jacobson's equation of state with the state the code's equilibrium (per-ray BW/analytic data and Iyer-Wald carried; explicit $K\mapsto-K$ sign adapter).
  • \texttt{gravStress\_conserved} / \texttt{deser\_selfcoupling\_consistent}: the graviton's own radiation-form stress $T\textasciicircum{}\{\textbackslash\{\}mu
u\}\_\{
m GW\}=k\textasciicircum{}\textbackslash\{\}mu k\textasciicircum{}
u\textbackslash\{\}langle e,e
angle\_\textbackslash\{\}eta$ is conserved **on-shell** (masslessness $\textbackslash\{\}Rightarrow$ conservation), so its coupling to \emph{its own} stress is gauge-invariant — the first order of the Deser bootstrap toward nonlinear GR (the full iteration and its quantum completion are not built).

\emph{The Type II dual-weight trace on the crossed-product core (the W1--W4 wall ladder).} The Chandrasekaran-Penington-Witten trace, constructed on an honest algebraic core with its three defining laws EXACT — the deepest cut into the Type II frontier; the von Neumann closure stays a named carried hypothesis. \textbf{The wall is not crossed}; the trace exists on the core.

  • \texttt{dualAction\_matter} / \texttt{dualAction\_clock}: the Takesaki dual action $\theta_s = \mathrm{Ad}(V_s)^{-1}$ on the crossed product — fixes the matter $\pi(a)$, phases the clock $\theta_s(\lambda_t)=e^{ist}\lambda_t$ (the vector-valued Weyl relation), with group law and multiplicativity.
  • \texttt{Iexp\_dualShift}: the log-clock weight density $I_{\exp}f=\int e^x f(x)\,dx$ scales EXACTLY, $I_{\exp}(f(\cdot{+}s))=e^{-s}I_{\exp}f$ — the $\tau\circ\theta_s=e^{-s}\tau$ mechanism (the density on the log-clock spectral variable, per the verifier's binding correction).
  • \texttt{zWeight\_shift\_quasiInvariant} / \texttt{zWeight\_dualCircle\_invariant}: the $\mathbb{Z}$-clock regression — the $e^{-1}$ scaling belongs to the SHIFT while the true dual (circle) action leaves the weight invariant; the distinction is machine-checked.
  • \texttt{tauMonomial\_dual}: the trace on normal-ordered core monomials $\pi(a)\lambda_t f(L)$ obeys $\tau_0\circ\theta_s=e^{-s}\tau_0$ exactly (the Weyl phase cancels against the density shift).
  • \texttt{eigen\_tau\_trace} / \texttt{eigen\_tau\_star\_mul\_nonneg}: on the modular eigen-core, \textbf{traciality} — the matter KMS factor $\omega(ab)=e^{\kappa}\omega(ba)$ cancels exactly against the $\int e^x$ change of variables (the Type II mechanism: a KMS state becomes a trace under the log-clock dressing) — and \textbf{positivity} $\tau_0(x^*x)\ge 0$ (the $x^*x$ symbol is a pointwise norm-square). The matter-side KMS-eigen/positivity inputs are carried (provable for the finite corner).
  • \texttt{phase5\_from\_core\_trace} / \texttt{traceCapacity\_from\_core}: the capacity interfaces (\texttt{Phase5Master}, \texttt{TraceCapacity}) instantiated \textbf{by the constructed trace} — the nonnegative JLMS remainder realized as $\tau_0(r^*r)$ and the bound $S\_\{
m ren\}\textbackslash\{\}le Q$ proven rather than assumed; the von Neumann extension is the carried \texttt{DualWeightTraceExtension} typeclass (normal weights/affiliated operators — the genuine remaining frontier), never an axiom.

\emph{The joins — hypothesis deletion (J1--J4).} A dedicated campaign shrinking the carried inputs of the landed results by connecting held theorems; every entry states which named hypothesis became a theorem or shrank to a smaller named input.

  • \texttt{finiteCorner\_tau\_trace} / \texttt{finiteCorner\_tau\_pos} (J1): the eigen-core trace laws hold on the concrete finite corner ($\rho=\mathrm{diag}\,p$, matrix units $E_{ij}$, $\kappa_{ij}=\log p_i-\log p_j$) with \textbf{no matter hypotheses} — the KMS-eigen law, frequency conservation (automatic from the matrix-unit index loop), and positivity are all theorems (\texttt{sigmaDiag\_single} pins the eigen convention). The constructed Type II trace's laws are unconditional in a genuine model; hkms/hfreq/hpos \textbf{deleted} (for this model).
  • \texttt{chmRadialMass3\_eq} / \texttt{CHMSymbolProbe3\_eq} / \texttt{bridge\_conditional\_probe} (J2): the CHM ball kernel's radial mass $\int_0^R 4\pi r^2\beta_R(r)dr=4\pi R^4/15$ (one-variable calculus) and the mass-normalized kernel pairing \textbf{equals} the ray area probe; the bridge assembly refactored to consume the derivable probe. The Iyer--Wald input \textbf{factors} — (physical deficit = kernel probe)$\circ$(kernel probe = area probe) — with the second factor now a theorem; carried residual: \texttt{hDeficit}, stated once (the FGHMVR content).
  • \texttt{hCHM\_of\_conformal\_transport} / \texttt{toBallModularFamily} (J3): the per-ball CHM identification is a \textbf{theorem} of named transport inputs — ONE carried massless wedge-BW datum (the $m>0$ BW theorem is never instantiated at $m=0$), per-ball conformal unitaries with geometric conjugacy, and the modular transport \texttt{hmodVac} carried in its weakest (pointwise-on-vacuum) form. C2b's per-ball \texttt{hCHM} \textbf{shrunk} to wedge datum + functoriality + geometry.
  • \texttt{next\_source\_conserved} / \texttt{extend\_of\_solver} / \texttt{einsteinDeserSystem} (J4): the all-order Deser statement in its honest form — formal-Bianchi consistency \textbf{propagation} (order-indexed operators at the shifted harmonics $n\cdot k$; conservation of the order-$(N{+}1)$ source \emph{derived} from a tower solving through $N$; extension given a solver), instantiated with the held linearized symbols and the proven Bianchi identity at every harmonic. No tower positing conservation per order; order 2 remains the concrete Deser theorem; the nonlinear Einstein coefficients are the cited frontier.

\emph{The Lorentz stress-test gates.} Three falsification gates on the finite-capacity postulate itself (numerics validated against closed forms; the decisive constants and limits machine-checked). The arc's verdict: every \emph{regulator} reading of $Q_D$ is dead; the \emph{state-level} (entropy) reading is provably frame-free; the dynamical-realization gap is the sole remaining Lorentz door.

  • The CPSUV gate (\texttt{cpsuv\_gate\_sharp\_fails}, \texttt{covariantSplit\_eq\_zero}; numerics \texttt{lorentz\_stress\_test.py}): in interacting Yukawa theory a preferred-frame spatial cutoff generates the one-loop speed splitting $\Delta c^2\to g^2/12\pi^2$ — an \textbf{unsuppressed constant} (no $E/\Lambda$ decoupling; the closed form's nonzero limit is a Lean theorem; quadrature matched to $2\times10^{-18}$) — while any O(4)-invariant regulator gives $\Delta c^2=0$ by symmetry. A naive mode cutoff reading of finite capacity is \textbf{falsified} (Collins--Perez--Sudarsky--Urrutia--Vucetich); the free-field pass ($|E_a^2-(m^2+p^2)|\le a^2p^4/12$, also a Lean theorem) is certified \emph{not decisive}.
  • The diamond-tip gate (\texttt{tipSplit\_eq\_zero\_iff}, \texttt{boostAvg\_log\_channel}, \texttt{u0sq\_avg\_diverges}; numerics \texttt{diamond\_tip\_test.py}): a causal diamond's tip vector $u^\mu_D$ \textbf{does} reach the effective action — within the anisotropic family $\Delta c^2=2C\,H_{\rm both}(\sqrt{a/b})$ (exact rationals $11/16$, $37/72$, $-17/18$ validated to $\le0.16\%$) the splitting vanishes \textbf{iff} the regulator is isotropic, with first-order sensitivity $-2C\ne0$; and the rapidity-average escape fails \emph{provably} — the boost-averaged cutoff's null channel equals $W/12$ exactly (linear growth: not a regulator; the boost group's noncompactness is fatal). There is no local Lorentz-invariant finite-capacity cutoff.
  • The state-level gate (\texttt{admissible\_smul\_iff}, \texttt{constraintSet\_invariant}, \texttt{stateLevel\_noDeltaC2}, \texttt{operationalLV\_iff\_not\_invariant}, \texttt{conditioned\_invariant\_iff\_orbit\_constant}, \texttt{equivariant\_enforcement\_preserves\_invariance}): the surviving reading — $Q_D$ bounds the \emph{entropy of the diamond algebra in the covariant vacuum} — makes \textbf{no low-energy Lorentz-violation prediction}: the admissible set is frame-free, the constraint touches no propagator ($\Delta c^2=0$ identically), and the only reopening channels are named theorems — a non-invariant \emph{prepared} state (state breaking, not law-level), background-selecting saturation, a \textbf{non-equivariant enforcement mechanism} (the dynamical-realization gap, now doubling as the last Lorentz gate), or a biased selector (closed by the held covariant-probability theorem). Consistent with the covariant-entropy-bound literature (Bousso; Casini; QNEC): entropy bounds are state-region inequalities, not regulators.

\emph{The operator emergence map (Q1--Q5): "the graviton is a quantized area fluctuation of the screen code," theorem-shaped.} The classical area$\to$metric decoder promoted to an operator map on the polynomial Bargmann--Fock carrier ($\mathrm{End}$ of $\mathbb{C}[X_0,X_1]$; the completion is never used as an operator domain). Two structural honesty rules are machine-enforced: equal-time quantized areas COMMUTE (the naive noncommutativity claim was cut at design time), and the finite code can never carry exact CCR (the trace obstruction: $\mathrm{tr}[Q,P]=0$ vs $\mathrm{tr}(i\mathbf{1})\neq 0$) — the code join is expectation-level, permanently.

  • \texttt{reconstruct\_hHat} (Q1): the decoder inverts the QUANTIZED area map at operator level — $\hat h_{\mu\nu}$ is a function of its own area-fluctuation observables, entrywise in $\mathrm{End}(\mathcal{F})$ (the module-level lift of the classical inversion).
  • \texttt{comm\_area\_mom} / \texttt{vacuum\_area\_pair} (Q2): the honest quantum structure of the area observables — the canonical pair $[\hat A(\Sigma),\hat\Pi(\Sigma')]=i\,\mathrm{areaPair}(\Sigma,\Sigma')\cdot\mathbf{1}$ (built on the operator-level CCR and the master c-number commutator of linear observables) and the vacuum fluctuation $\langle 0|\hat A(\Sigma)\hat A(\Sigma')|0\rangle=\mathrm{areaPair}(\Sigma,\Sigma')$; equal-time areas commute (\texttt{comm\_area\_area = 0}).
  • \texttt{coherent\_hHat} / \texttt{coherent\_area} (Q3): the CLASSICAL bridge is the coherent shadow — $\langle\alpha|\hat h|\alpha\rangle = h(\alpha) = \sum_\lambda 2\,\mathrm{Re}(\alpha_\lambda)\,\mathrm{pol}^\lambda$ and $\langle\alpha|\hat A(\Sigma)|\alpha\rangle = \delta A_\Sigma(h(\alpha))$ — exactly the $\delta A$ the assembled first-law$\Leftrightarrow$Einstein skeleton consumes (expectation rules grounded by the held $a|\alpha\rangle=\alpha|\alpha\rangle$ eigenvalue theorem and cited Bargmann adjointness).
  • \texttt{heis\_q} / \texttt{hHatT\_wave} / \texttt{comm\_areaT} (Q4): the Heisenberg flow DERIVED from the explicit monomial scaling ($U_z\,q\,U_z^{-1}=q(t)$, chain rule proven by induction — no Stone theorem used), the operator WAVE EQUATION $\ddot{\hat h}+\omega^2\hat h=0$ (coefficientwise — the honest ODE on a normless carrier), and the time-separated area commutator $[\hat A_\Sigma(t),\hat A_{\Sigma'}(s)]=2i\sin(\omega(s{-}t))\,\mathrm{areaPair}\cdot\mathbf{1}$ — the causal structure of quantized areas, vanishing at equal times.
  • \texttt{code\_count\_eq\_fock\_area\_expect} (Q5, the join): with the held calibration and ONE named carried input \texttt{hJoin} (the code's induced screen area equals the coherent total-area expectation, $\hat A_{\rm tot}=A_0\cdot\mathbf{1}+\hat A$ — total-to-total), $\log\#\text{microstates} = \langle\alpha|\hat A_{\rm tot}(\Sigma)|\alpha\rangle/4G$: \textbf{the screen code's counting and the graviton's area operator agree as two computations of one number.} Carried residue: \texttt{hJoin}, the calibration $\log D_e = w_{\rm Ent}(e)$, the Bargmann-adjointness grounding. Fixed momentum, linearized, free; NOT quantum gravity — the genuine count and the von Neumann closure remain the frontier.

\emph{The grounding campaign (G1--G5): carried hinges in multiple landed results, deleted at once.} Two files; every deletion states its residue.

  • \texttt{bargmann\_adjoint} / \texttt{coeffFamilyPair\_cohCoeff} / \texttt{cohPair\_X\_mul} (G1): the Bargmann pairing $\langle p,q\rangle_B=\sum_n n!\,\overline{p_n}q_n$ on the polynomial Bargmann--Fock space with \textbf{creation adjoint to annihilation} ($\langle p, X_l q\rangle_B = \langle \partial_l p, q\rangle_B$ — the factorial shift $(m{+}e_l)! = m!(m_l{+}1)$), and the coherent reproducing rule $\langle \mathrm{coh}\,\alpha, p\rangle_B = p(\bar\alpha)$ with the creation rule $\langle\mathrm{coh}\,\alpha, X_l p\rangle = \bar\alpha_l\langle\mathrm{coh}\,\alpha,p\rangle$. The operator-emergence coherent v-rule — previously "standard Bargmann calculus, cited" — is now a polynomial-level theorem (the completion-level identification stays cited: the honest boundary of a polynomial carrier).
  • \texttt{rvdRC\_transport} (G2): the Rieffel--Van Daele operator transports under unitary conjugacy, $R_{S'} = U R_S U^{-1}$, from the uniqueness characterization of real orthogonal projections under the $\mathbb{R}$-isometry (with $i\mathcal{K}$ transporting automatically by $\mathbb{C}$-linearity); carrier hypotheses in membership form.
  • \texttt{cfc\_conjU} / \texttt{specMeasure\_conjU} / \texttt{specProj\_conjU} / \texttt{borelFC\_conjU} (G3 — new tower infrastructure): \textbf{the unitary covariance of the spectral theorem and the bounded Borel functional calculus} — the continuous calculus transports by star-hom functoriality; the Riesz--Markov scalar spectral measures transport as pushforwards (compactly-supported tests Tietze-extended to ambient symbols); the spectral projections transport ($E'(s') = U E(e^{-1}s')U^{-1}$); and the capstone $f(UTU^{-1}) = U\,(f\circ e)(T)\,U^{-1}$ for bounded measurable symbols. The scalar-measure route throughout — the generator-uniqueness shortcut was rejected at design time.
  • \texttt{modUnitary\_transport} and the payoffs (G4): the continuum modular flow transports, $\Delta^{it}_{S'} = U \Delta^{it}_S U^{-1}$, under membership-level carrier conjugacy. Consequences, each an honest deletion or instantiation: \textbf{J3's \texttt{hmodVac} carried field is DELETED} (\texttt{CHMTransportDataOfCarrierMap} builds the ball-transport package from GEOMETRIC carrier conjugacy alone; the residue is the geometry itself plus the massless wedge BW); \textbf{Gate 3's covariance hinge is fed by a derived theorem} (\texttt{modUnitary\_inner\_cov}: the modular correlators are carrier-covariant); and \textbf{the ball-Clausius per-ball modular input is replaced by per-ball geometry} (\texttt{ball\_modUnitary\_cov}). NOT quantum gravity; no wall crossed.

\emph{The keystone — THE COUNT (K0--K6): the holographic count as a theorem in the finite branch.} The campaign derives, rather than posits, that a diamond/screen algebra's entropy against the CONSTRUCTED crossed-product trace $\tau_0$ equals induced area over $4G$ — deleting, in the finite branch, the Clausius input, the geometric input, the calibration $\log D_e = w_{\rm Ent}(e)$, and the emergence join at once (\texttt{QIQTH/Keystone.lean}, \texttt{KeystoneOperator.lean}; axiom-free, std-3).

  • \texttt{vonNeumannEntropy\_maxMixed} / \texttt{vonNeumannEntropy\_le\_log\_card} (K0): the entropy substrate — $S(\text{maximally mixed}) = \log N$ against the UNNORMALIZED counting trace, with the Gibbs/Jensen guard $S(\rho)\le\log N$ for \emph{every} density (the count equality is claimed only where it holds — at maximal mixing).
  • \texttt{K2a\_count\_capstone} (K2a): the standalone finite count — link dimensions $D_e$, microstates $N_C=\prod_e D_e$, the diamond matrix algebra with the unnormalized $\tau(\mathbf{1})=N_C$, record projections with $\tau(P_R)=|R|$ (the trace COUNTS records), and $S(\text{maxMixed}) = \sum_e \log D_e = A_\tau(C)/4G$ as a theorem, with $G$ entering only through the normalization. The honest boundary is stated as a theorem too: \texttt{count\_matches\_external\_weights\_iff} — matching EXTERNAL geometric weights \emph{is} the old calibration, never counted as deleted.
  • \texttt{tauMonomial\_uniform\_eq\_tauCount} / \texttt{wEntTau\_eq\_log\_tau0Dim} / \texttt{K2b\_tau0\_capstone} (K2b, the core result): \textbf{the counting trace IS the restriction of the constructed crossed-product trace $\tau_0$} — the held monomial-trace formula at the uniform matter state and a mass-$N_C$ clock window reproduces $\operatorname{Tr}x$ exactly, so the count is not a new postulate; and \textbf{the calibration is a theorem} because the weight is trace-defined — the link weight IS the log of the link fiber's $\tau_0$-dimension. Capstone: $S(\text{record corner}) = \log\dim_{\tau_0}(\mathcal{D}_C) = A_\tau(C)/4G$ in the finite record corner of the constructed core.
  • \texttt{tauCount\_conj} / \texttt{K5\_dual\_covariant\_count} (K5): the covariance checks — trace-preserving code unitaries preserve the count, and the dual action SCALES it exactly, $S(\theta_s\,\cdot) = S(\cdot) - s$ (the honest transported-area covariance law from the held $e^{-s}$ dual scaling — never naive invariance).
  • \texttt{clockMul} / \texttt{clockTransl\_clockMul} / \texttt{repMonomial} (K1): the operator packaging — bounded-symbol clock multiplication operators on the crossed-product space $L^2(\mathbb{R};H)$, the product law, the \textbf{Weyl covariance} $\lambda_t \circ M_g = M_{g(\cdot+t)}\circ\lambda_t$, and the represented core monomial $\pi(a)\lambda_t M_F$ as a genuine continuous operator.
  • \texttt{tauCount\_norm\_le\_sum\_diag} and the vacuity audit fix (K3): finite-corner boundedness of the counting trace; and a soundness find — the carried von Neumann-extension interface (\texttt{DualWeightTraceExtension}) was VACUOUS (an abelian collapse witness $M=\mathbb{C}$ satisfied it for any algebra), now strengthened with a multiplicative-embedding requirement that provably kills the witness. No fake finite instance was shipped; the genuine $\sigma$-weak/normal-weight extension stands as a named wall, carried non-vacuously.
  • The checkpoint (K6), verbatim: HAVE — "every finite code screen realized as a finite record corner of the constructed crossed-product core has $S_{\tau_0} = \log\dim_{\tau_0}(\mathcal{D}_C) = \sum_e \log\dim_{\tau_0}(P_e) = A_\tau(C)/4G$; in the code instance $\dim_{\tau_0}(P_e) = D_e$, so the calibration is a theorem (trace-defined weight) and the count-built area operator gives the join by construction — no hClausius/hGeom/hCalib/hJoin carried in this branch." HAVE NOT (Walls 1--5, named): continuum QFT diamond algebras ARE these corners; external geometric area $=$ count-built area; Type III$_1$/II$_\infty$ continuum structure in Lean; $\sigma$-weak/normal weights; the value of $G$. NOT quantum gravity solved; no wall crossed.

\emph{The join instance (JI1--JI7): \texttt{hJoin} deleted by construction — the two towers merged at the finite level.} The screen-code tower (the keystone count) and the graviton tower (Q1--Q5) joined into ONE object: the dictionary instance — links $=$ screen elements, code weights defined FROM the geometry — makes the Q5 carried hypothesis a theorem (\texttt{QIQTH/JoinInstance.lean}; axiom-free, std-3). Two-level construction per the consult verdict: generic-exact with REAL trace-dimensions; integer codes only under a named realizability datum.

  • \texttt{sum\_localArea} / \texttt{A0Split} (JI1): the per-element area shares $\beta_a + \delta A_a$ reassemble the total $A_0 + \delta A_\Sigma(h)$ — with the background apportionment carried as NAMED data (a global constant has no canonical per-link split; the uniform split is an optional policy, never pretended-derived).
  • \texttt{hJoin\_tau} / \texttt{hcal\_tau} (JI2): with $w_{\rm Ent}(a) := A^{\rm loc}_a/4G$ and $D^\tau_a := e^{w_{\rm Ent}(a)}$ (a positive REAL — the Type II lesson: trace-dimensions need not be integers), \textbf{the carried \texttt{hJoin} equality is a theorem} — \texttt{inducedScreenArea} $= A_0 + \delta A_\Sigma(h(\alpha))$ — and the calibration is \texttt{Real.log\_exp}. Direction geometry $\to$ code: the theorem is about the instance \emph{built from} the geometric data, so nothing is smuggled.
  • \texttt{exists\_tau0\_corner\_of\_posReal} (JI3): the dictionary lives in the held crossed-product core — every positive real is a \emph{realized} $\tau_0$ corner value with the clock-window witness explicit (the free window mass; never subcorners of one fixed fiber$+$window, which are provably rank-quantized).
  • \texttt{Stau\_eq\_area\_over\_4G} (JI4, the generic replacement): the instance's count $S_\tau(J) = \log\prod_a D^\tau_a = \sum_a w_{\rm Ent}(a) = A_J/4G$ for ARBITRARY graviton data — count and geometry as two computations of one number, with nothing carried (the join is the construction).
  • \texttt{code\_count\_eq\_fock\_area\_expect\_noJoin} (JI5): the old Q5 capstone re-proved \textbf{with no \texttt{hJoin} hypothesis} — for a nat-realizable geometry (\texttt{NatRealizable}, the named integer-dimension datum, where $D^\tau_a = D_a$), $\log\#\text{microstates} = \langle\alpha|\hat A_{\rm tot}(\Sigma)|\alpha\rangle/4G$ with the join supplied by \texttt{hJoin\_tau}.
  • \texttt{Stau\_eq\_capacity\_primitives} and the costs (JI6): substituting the DERIVED $G = 1/(N\Lambda_s^2)$, the count is $S_\tau(J) = (A_J/4)\,N\Lambda_s^2$ — the count-built and induced-Newton normalizations are ONE formula in $\{$area, species, granularity$\}$; per-link capacity $(A^{\rm loc}_a/4)N\Lambda_s^2$, the patch bound, and the area costs: one nat of link entropy costs $4/(N\Lambda_s^2)$, one qubit $4\log 2/(N\Lambda_s^2)$.
  • The checkpoint (JI7), verbatim: HAVE — "after JI1--JI6, \texttt{hJoin} is no longer a hypothesis for the constructed τ join or for nat-realizable finite-code joins, and the count normalization rewrites to $(A_J/4)\cdot N\cdot\Lambda_s^2$ with local capacity corollaries." HAVE NOT — "no theorem says arbitrary external real geometry has exact natural link dimensions, no asymptotic approximation is included, and no canonical $A_0$ split is asserted beyond the named apportionment data/policy." The CCR-isometry obstruction is permanent (the join is expectation-level forever); NOT quantum gravity solved; no wall crossed.

\emph{The embedding (EM1--EM7): the matter-side dictionary — the truncated field diamond IS a counted record corner; the finite-level bridge closed end to end.} The key observation is definitional: the keystone's microstate space $\mathrm{Micro} = \prod_{k} \mathrm{Fin}(D_k)$ IS a multi-mode truncated Fock basis (joint occupation numbers $n_k < D_k$), so no new Hilbert space is built — the campaign gives the already-counted diamond algebra its FIELD structure and the mode$\leftrightarrow$link semantics (\texttt{QIQTH/Embedding.lean}; axiom-free, std-3).

  • \texttt{ModeAssignment} / \texttt{toLinkDims} / \texttt{truncated\_field\_diamond\_entropy} (EM1): a LINK IS A FIELD MODE (its dimension the mode's truncation cutoff, carried as named finite data — the continuum localization map stays the wall); the keystone count reads verbatim as the truncated field diamond's entropy, $S = \sum_k \log D_k = A_\tau(C)/4G$.
  • \texttt{modeOp} and the transport package (EM2): the direct-entry coordinate embedding ($A$ on fiber $k$, delta elsewhere — never Kronecker induction), unital/additive/multiplicative/⋆-preserving/injective via a reusable fiber-sum lemma — each single-mode truncated-oscillator algebra genuinely embeds, and every single-mode theorem transports without unfolding the operators.
  • \texttt{mode\_ladder\_commutator} / \texttt{occupationProj\_joint\_eigen} (EM3): the per-mode oscillators $a_k, N_k$, with the HONEST transported defect $[a_k, a_k^\dagger] = 1 - D_k P_{\rm top,k}$ (exact CCR permanently impossible — stated, not hidden), the finite spectrum reading ($N_k$ eigenvalue $n_k$ on every occupation projector), and $[N_k, a_k] = -a_k$, $[N_k, a_k^\dagger] = a_k^\dagger$.
  • \texttt{modeOp\_commute\_of\_ne} (EM4): ONE generic theorem — coordinate operators at different modes commute — with all cross-mode commutators ($[a_k,a_j] = [a_k,a_j^\dagger] = [N_k,a_j] = 0$, $j\neq k$) as corollaries; the honest scope: this IS the bosonic sector (pi-fiber ladders commute; fermionic CAR needs the held graded layer).
  • \texttt{recordProj\_eq\_sum\_occupationProj} / \texttt{encoded\_mode\_ladder\_commutator} (EM5): RECORDS ARE OCCUPATION POINTER-BASIS SUBSETS as a theorem (the occupation projectors are orthogonal, complete, and every keystone record projector is their sum — welding the record ontology to the field ontology); the field record trace runs through the constructed τ₀; and the corner transport carries the honest unit — $[\iota_V(a_k), \iota_V(a_k)^\dagger] = P - D_k\,\iota_V(P_{\rm top,k})$ with $P = VV^H$, never the ambient 1.
  • \texttt{field\_entropy\_le\_area\_of\_capacity} (EM6): CAPACITY IS A CONSTRAINT, NOT A GENERATOR — the FQ bound $\sum_k \log D_k \le A/4G$ selects admissible mode assignments and implies $S \le A/4G$; saturation gives equality for the CHOSEN assignment (integer-saturation existence never claimed); each mode occupies the local area $4G\log D_k$, and a diamond of two-level modes carries at most $A/(4G\log 2)$ of them.
  • \texttt{truncated\_field\_count\_eq\_fock\_area\_expect\_noJoin} (EM7, the capstone): \textbf{the finite-level bridge, end to end} — for a nat-realizable geometry, $\log\#(\text{truncated Fock basis}) = \langle\alpha|\hat A_{\rm tot}(\Sigma)|\alpha\rangle/4G$, composing field → corner → count → area → graviton on ONE object with no join hypothesis; \texttt{LocalizedModeFrame} supplies the certifying (never constructing) localization witness.
  • The checkpoint (EM7), verbatim: HAVE — "the N-mode truncated free-field diamond algebra IS a counted record corner — the occupation basis is Micro, the per-mode truncated oscillators embed with their honest defect ([a,a†] = 1 − D·P\_top), records are occupation projectors, the count S = Σ log D\_k = A/4G reads as the truncated field diamond's entropy, capacity bounds/saturates the cutoffs as a constraint, and the mode dictionary composes with the join instance end-to-end (field → corner → count → area → graviton expectation)." HAVE NOT — "no exact finite CCR (the truncation defect is permanent); no Type III₁ finite corner (the cutoff→continuum limit is THE wall, never claimed); no construction of continuum-localized modes from the standard subspace (mode membership is named finite data, at most CERTIFIED by a supplied localization witness); capacity is a constraint, not a generator." NOT quantum gravity solved; no wall crossed.

\emph{The dynamics (DY1--DY7): the code's time evolution, the independent cross-check, and the correspondence conjecture.} The microscopic side gets what a definition of a theory requires — a time evolution — and the campaign closes with the QIQT-H analogue of the Brown--Henneaux $=$ Cardy check in its honest finite form, plus the sharp continuum conjecture stated (never assumed) in Lean (\texttt{QIQTH/Dynamics.lean}, \texttt{CrossCheck.lean}, \texttt{Conjectures.lean}; axiom-free, std-3).

  • \texttt{alpha} and the entry formula (DY1): the diagonal code Hamiltonian $H = \sum_k \omega_k N_k$ with its Heisenberg flow defined by explicit phase unitaries (no Stone theorem, no matrix exponential; every proof runs through $\alpha_t(A)(n,m) = e^{it(E(n)-E(m))}A(n,m)$) — a one-parameter group of ⋆-automorphisms under which \textbf{records are stationary} (the honesty point: $H$ is a function of the $N_k$, so what $\lambda$ selects among does not move) while the ladders rotate at their mode frequencies, $\alpha_t(a_k) = e^{-i\omega_k t}a_k$.
  • \texttt{gibbsDensity} (DY2): the thermal state as an EXPLICIT product diagonal density (per-mode Boltzmann weights, normalized by the occupation-basis product-sum interchange) — a genuine density for every $\beta$, stationary under the flow.
  • \texttt{sigmaDiag\_gibbs\_eq\_alpha\_rescale} (DY3): \textbf{the Gibbs state's modular flow IS the rescaled physical flow}, $\sigma_s^{\rho_\beta} = \alpha_{-\beta s}$ (the partition function cancels; the flow is never \emph{defined} by modular theory — the bridge runs from the independent dynamics to the KMS certificate); the held finite Tomita--Takesaki then certifies $\rho_\beta$ as the KMS state, and at $\beta = 0$ the Gibbs state IS the keystone's maximally mixed counting state — the thermal and counting towers share their ground floor.
  • \texttt{marginal\_gibbsWeight} / \texttt{entropy\_gibbs\_region} (DY4--DY5): regions are subsets of MODE labels (no spatial-entanglement claims — the product state has none); the Gibbs marginal is again Gibbs; and the region entropy formula $S(\rho_{\beta,R}) = S_{\rm micro}(R,\beta) = \sum_{k\in R} s_k(\beta\omega_k)$, saturating at $\sum \log D_k$ as $\beta \to 0$ with the all-$\beta$ Gibbs bound (a new reusable spectral fact en route: eigenvalue sums of real diagonal matrices are entry sums).
  • \texttt{S\_micro\_zero\_eq\_inducedQuarterG} (DY6, the destination): \textbf{the saturated conditional Sakharov cross-check, calibration-free} — the code dynamics' region entropy at saturation equals the induced area over $4G_{\rm ind}$, with the microscopic side computed from the Hamiltonian and the macroscopic side supplied by independent Sakharov/species/cell data (\texttt{InducedCrossCheckData}; with the derived $G_{\rm ind} = 1/(N_{\rm eff}\Lambda_s^2)$ only the species/cell matching remains input), and the proofs referencing NONE of the keystone calibration (\texttt{wEntTau}/\texttt{cutTau}/\texttt{tauMonomial}/\texttt{hJoin} — grep-verified). Equality at saturation ONLY; the arbitrary-$\beta$ equality is false and never claimed.
  • \texttt{FlatSpaceRecordGravityCorrespondence} (DY7): \textbf{the conjecture, stated sharp} — in the continuum limit the capacity-bounded record code with diagonal dynamics equals free QFT $+$ linearized gravity: for every region, micro record entropy $=$ one-loop conical entropy $=$ area$/4G_{\rm ind}$, with $G_{\rm ind}$ the Sakharov constant of the SAME field content (one microscopic system computing both the states and $G$ — the independent cross-check the finite level cannot supply). Packaged as a named Lean \texttt{Prop} with the DY1--DY6 evidence bundled and PROVEN (\texttt{finiteEvidence\_holds}) and the continuum claim carrying NO proof field, NO axiom, NO instance — stated, never assumed.
  • The checkpoint (DY7), verbatim: HAVE — "a finite, axiom-free diagonal code dynamics, explicit Gibbs/KMS states, product-mode reductions, and a saturated conditional induced-gravity cross-check whose proof does not use the trace/wEnt area calibration." HAVE NOT — "a finite proof of a continuum one-loop heat-kernel area law or an equality between finite thermal entropy at arbitrary β and an induced geometric area; that remains the named continuum frontier/conjecture." NOT quantum gravity solved; no wall crossed.

\emph{The decoupling shadow (DS1--DS7): the finite forced core of the dictionary.} Prompted by a primary-source re-reading of Maldacena (hep-th/9711200), whose deepest structural feature is that the dual sides are two SURVIVING DESCRIPTIONS of one parent construction under one limit: the campaign delivers the honest finite shadow of that structure — the word "constructed" is deleted from every part of QIQT-H's dictionary where deleting it is mathematically true, and named parent data everywhere else (\texttt{QIQTH/Decoupling/}, \texttt{QIQTH/Rigidity/}; axiom-free, std-3).

  • \texttt{commutator\_eventually\_exact} (DS1): the free-oscillator sector is FORCED by the cutoff limit, operator level — at fixed occupations the truncated ladder matrix elements are D-independent and the commutator's entries stabilize to the exact-CCR values (the truncation defect lives only at the top level, which bounded occupations eventually never see); stated in \texttt{∀ᶠ D} form.
  • \texttt{tendsto\_meanN} / \texttt{tendsto\_defectExpect} (DS2, the development's first genuine limit theorems): at fixed $βω > 0$, the truncated partition function, occupation expectation, and defect expectation converge to $1/(1-q)$, the PLANCK value $q/(1-q)$, and $0$ — the state-level decoupling half — with the bridge \texttt{ZMode\_eq\_Zgeom} certifying these are statements about the code's own thermal states.
  • \texttt{guard\_entropy\_saturates} / \texttt{guard\_defect\_survives} (DS3, THE REGIME-SEPARATION GUARD): $S_D(x) \to$ the free Planck entropy at fixed $x > 0$ and $\to \log D$ as $x \to 0^+$; and along ANY schedule $x_D·D \to 0$, capacity saturates BUT the defect expectation tends to \textbf{1}, not 0 — exact saturated capacity is provably NOT the positive-temperature free-oscillator limit. The two halves of the shadow live in different regimes, as a theorem: the formalization itself forbids running the count and the field limit together into a fake continuum claim.
  • \texttt{tendsto\_productEntropy} (DS4): the single-mode limits lifted to finite mode sets — the product thermal entropy, total defect, and every fixed occupation's Gibbs weight converge to the finite-mode free-field values.
  • \texttt{monotone\_logValuation} (DS5): the classical monotone-additive Cauchy rigidity, proved by hand — a monotone product-to-sum valuation on the positive reals is $κ\log$.
  • \texttt{finiteCorner\_valuation\_rigidity} / \texttt{forced\_weight\_product} / \texttt{nu2\_counterexample} (DS6, THE FORCED WEIGHT): a MONOIDAL valuation on finite dimensions, MONOTONE under all isometric embeddings, is forced to be $κ\log n$ (the double-log squeeze) — hence on product record corners $A = κ\sum_k \log D_k$: the local weight of the keystone/join/embedding dictionaries is the UNIQUE refinement-natural valuation, no longer a constructed choice, with $κ$ the one free normalization (where $4G$ lives, input). The necessity is machine-checked: the 2-adic valuation is additive and divisibility-monotone yet NOT $\propto \log$ — the strong hypotheses cannot be weakened.
  • The shadow package (DS7): the parent tower with its three theorems — the free sector survives the cutoff limit (\texttt{decouplingShadow\_holds}, every field a landed theorem, the guard included); every refinement-natural valuation is $κ\sum \log D_k$; and given the normalization, the saturated area law survives (\texttt{saturated\_entropy\_eq\_forced\_area}: the code's $β = 0$ entropy $=$ the forced area over $κ$).
  • The checkpoint (DS7), verbatim: HAVE — "The capacity-limit theorem forces the oscillator/free-field sector only in the bounded-occupation or positive-temperature sense; it does not force the screen geometry or Newton constant." HAVE NOT — "The tower-rigidity theorem forces the logarithmic capacity weight only under monoidal, monotone refinement naturality; without those hypotheses there are explicit finite counterexamples." NOT a full decoupling derivation (the join incidence geometry, the species/cell match, and the value of G remain parent data); NOT quantum gravity solved; no wall crossed.

\emph{The tower (T1--T8): the first machine-checked contact with the Type III$_1$ wall.} Phase A of the continuum-definition attack, via the ITPFI/Araki--Woods route: the capacity code tower with its Gibbs product states IS Araki--Woods input data, and its arithmetic invariant is machine-verified to be the III$_1$ one --- at the level of the wall's FINGERPRINT, never the wall itself (\texttt{QIQTH/Tower/}; axiom-free, std-3).

  • \texttt{IsTailModularExponent} / \texttt{AWFingerprintIII1} (T1): the named witness predicate for the Araki--Woods asymptotic-ratio invariant, formulated additively in the modular exponent $\kappa$ --- a tail quantifier ($\forall N\,\exists k\ge N$) and a uniform weight floor $\delta$, both load-bearing (drifting-frequency and vanishing-weight counterexamples documented) --- with the EXACT per-mode ratio $\lambda_1/\lambda_0 = e^{-x}$ (the partition function cancels; never approximated) and the bridges to the held \texttt{kappaOf} eigenvalue law.
  • \texttt{gibbsTower\_awFingerprint\_III₁} (T3, the centerpiece): two frequencies occurring infinitely often at irrational ratio, uniform bounds $0<a\le\beta\omega_k\le b$, cutoffs $D_k\ge 2$ $\Rightarrow$ the tower's eigenvalue family satisfies the III$_1$ fingerprint (the additive group generated by the witnessed tail exponents is dense in $\mathbb R$, via Kronecker density from \texttt{AddSubgroup.dense\_or\_cyclic}); PLUS the hypothesis-free instance --- the alternating qubit tower with frequencies $\{1,\sqrt2\}$, via \texttt{irrational\_sqrt\_two}, no hypotheses carried. The operator reading rests on three facts CITED verbatim in the docstring and never proved (Araki--Woods 1968 sufficiency/amplification; the closed-subgroup structure of $r_\infty$; $r_\infty=[0,\infty)$ $=$ the III$_1$ class $=$ the Connes $S$-invariant for ITPFI, Connes 1973): given those, the tower's ITPFI factor is III$_1$ --- by citation, not by proof; no von Neumann algebra is constructed anywhere in the development.
  • \texttt{gibbsTower\_constant\_not\_fingerprint} (T4, the Powers guard): a constant-frequency tower provably FAILS the fingerprint (every tail exponent lies in $s\mathbb Z$, whose closure is cyclic, not dense) --- the arithmetic fingerprint of the Powers factor III$_{e^{-s}}$ (cited). The separation theorem: the predicate is neither vacuous nor universal.
  • \texttt{gibbsLimitMeasure} (T5): the $\sigma$-additive INFINITE-MODE Gibbs measure on occupation configurations --- the unique projective limit of the dynamics campaign's own thermal marginals through the held Kolmogorov/product machinery (\texttt{Measure.infinitePi}), with the capstone identity that the finite marginals ARE the code's DY Gibbs weights.
  • \texttt{gibbsLimitMeasure\_noAtoms} (T6): NON-ATOMICITY --- under the uniform frequency bound every singleton configuration is null (the cylinder squeeze: depth-$N$ cylinder mass $\le (1/(1+e^{-b}))^N \to 0$). Hence no diagonal-density ("diagState") reading of the limit exists: the quantum reading of the classical limit measure is FALSE, not deferred --- the Type-I shortcut is provably closed, which is exactly why the ITPFI route is the right frontier. The vacuum-atom dichotomy ($\sum e^{-\beta\omega_k}<\infty$ $\Rightarrow$ an atom; Kakutani-type) is cited, never proved.
  • \texttt{cornerEmbed} (T7): the finite operator refinement tower --- for nested corners $C\subseteq C'$ the inclusion is a unital ⋆-homomorphism, MODE-compatible, STATE-compatible ($\varphi_{C'}\circ\iota=\varphi_C$, the operator form of the Gibbs-marginal consistency), and MODULAR-FLOW EQUIVARIANT ($\sigma_s^{C'}\circ\iota=\iota\circ\sigma_s^C$, via the \texttt{kappaOf} eigen-law: the Gibbs log-weight difference of configurations agreeing off $C$ is computed inside $C$). A family of finite-dimensional maps only --- exactly the ITPFI tower DATA, its classification never performed.
  • The checkpoint (T8), verbatim: HAVE --- "the machine-checked arithmetic content of the Araki--Woods III₁ criterion for the code's Gibbs tower, including a hypothesis-free concrete instance, the Powers-guard separation, the σ-additive infinite-mode Gibbs measure with its non-atomicity, and the state-compatible modular-equivariant finite refinement tower; the inference to an actual III₁ factor is cited (Araki--Woods 1968; Connes 1973), never proved." HAVE NOT --- "the ITPFI von Neumann algebra, its ratio set, its type, any inductive limit or weak closure, any quantum state on the infinite system, or any continuum-limit completion --- none are constructed or classified here." NOT the continuum done; no wall crossed --- this is the wall's fingerprint.

\emph{The closure (C1--C11): the von Neumann double-commutant theorem.} The convergent blocker of the continuum program — the missing lemma cluster that four separately-named frontiers (the vN closure of the crossed product, the dual-weight trace extension, the Fock Klein-twist commutant theorem, the crossed-product \texttt{VonNeumannAlgebra} packaging) all reduce to — is now a machine-checked theorem, in Mathlib-styleable form: Mathlib defines \texttt{VonNeumannAlgebra} by the bicommutant property but has NO bicommutant theorem; this campaign closes that gap (\texttt{QIQTH/VonNeumann/}; axiom-free, std-3; every increment landed, the stretch included).

  • \texttt{vonNeumann\_double\_commutant} (C7, the centerpiece — green first try): for every unital ⋆-subalgebra $A$ of the bounded operators on a complex Hilbert space, the double centralizer $A''$ EQUALS the set of operators approximable from $A$ in norm on every finite tuple of vectors — the strong-operator-topology closure, stated concretely as the \texttt{SOTApprox} predicate (one approximant uniform over each tuple; the quantifier order is load-bearing and documented). Forward: the classical amplification argument, fully verified — the cyclic-subspace projection lies in the commutant (C1; ⋆-closure load-bearing, upper-triangular counterexample), single-vector density (C3; unitality load-bearing, $A=\{0\}$ counterexample), the frozen \texttt{PiLp} amplification interface with adjoint $\iota^\dagger = \pi$ (C4), the two minimal matrix-commutant lemmas — entries and assembly, never $M_n(A')$ (C5), and the $n$-vector density (C6). Converse: the $(x, Sx)$ two-vector estimate — single-vector approximability is provably insufficient.
  • \texttt{wotClosure\_image\_eq\_image\_bicommutant} (C10, the stretch — shipped): the weak-operator-topology closure IS the bicommutant, proved wholly inside Mathlib's \texttt{H →WOT[ℂ] H} type copy (closed commutation equalizers via separate continuity — joint WOT continuity of multiplication is false; SOT approximability defeats every seminorm-basis neighborhood). With C7: \textbf{the WOT closure, the SOT closure, and the double centralizer coincide} — the full classical statement.
  • \texttt{VonNeumannAlgebra.generatedBy} (C2) + \texttt{generatedBy\_carrier\_eq} (C7): the vN algebra generated by any set of operators — the double centralizer, packaged with minimality and the Galois lemma $(\mathrm{adjoin}\,S)' = (S \cup S^*)'$ — proven to BE the SOT closure of the generated ⋆-algebra.
  • The two downstream payoffs, honestly scoped: \texttt{crossedProductVN} (C8) — the project's crossed product $M \rtimes_\sigma \mathbb{R}$ (matter representation + clock translations on $L^2(\mathbb{R};H)$) packaged as a genuine \texttt{VonNeumannAlgebra} with the SOT-approximability membership characterization — packaging only, the dual-weight trace is NOT claimed to extend to the weak closure; and \texttt{limitVN} (C9) — the directed-union limit algebra for any hypothesized common representation, with $T \in \lim \Leftrightarrow$ SOT-approximable from the union of the stages — the refinement-tower limit in its honest form (the code tower's own instantiation awaits the tower-GNS campaign).
  • The checkpoint (C11), verbatim: HAVE — "We have the von Neumann double-commutant theorem as an axiom-free Lean theorem over current Mathlib — for every unital ⋆-subalgebra A of the bounded operators on a complex Hilbert space, the double centralizer A″ equals the set of operators approximable from A in norm on every finite tuple of vectors (and, in the shipped WOT increment, the weak-operator closure) — packaged as \texttt{VonNeumannAlgebra.generatedBy} with membership lemmas, and instantiated to present the project's crossed-product representation and any commonly-represented refinement tower as genuine \texttt{VonNeumannAlgebra}s." HAVE NOT — "We do not have Kaplansky density, normal states, preduals or the σ-weak topology, type classification, or the inductive-limit (tower-GNS) Hilbert space — the ITPFI tower's limit algebra is packaged only relative to a hypothesized common representation, and the crossed-product dual-weight trace is not claimed to extend from the algebraic core to the weak closure." The GATE to the continuum, not the wall crossed.

\emph{The representation (R1--R9): the corner tower on ONE Hilbert space — the tower's limit von Neumann algebra exists.} The GNS construction of the compatible Gibbs family: the T5/T7 refinement tower, previously a family of finite matrix algebras related by embeddings, now acts on a single Hilbert space, and its directed-union limit von Neumann algebra is an actual, named, axiom-free object (\texttt{QIQTH/TowerGNS/}; all std-3, budget 0; the first genuinely infinite-dimensional quantum object of the development).

  • \texttt{TowerGNS} (R3--R4): the Hilbert space — the completion of the direct sum of ALL finite corners under the stabilized Gibbs pairing, which is DELIBERATELY semidefinite: the null directions are exactly the direct-limit gluing, and the metric completion performs the identification (\texttt{towerGerm}: the image of a corner element equals the image of any of its embeddings) — no quotient is ever taken. The instance architecture copies Mathlib's own \texttt{GelfandNaimarkSegal.lean}.
  • \texttt{towerRep} (R5--R7): every corner algebra acts as a unital ⋆-algebra homomorphism into the bounded operators — boundedness from an elementary Frobenius estimate (\texttt{gnsInner\_leftMul\_le}, honest scope: bounded, NOT claimed contractive — the pin has no bundled matrix C*-algebra), the algebra laws holding ONLY in the completion (they are provably false at the pre-level — the stages differ; the germ reconciles), the star law via the adjoint relation under the Gibbs state. CAPSTONE \texttt{towerRep\_cornerEmbed}: $\pi_{C'} \circ \iota = \pi_C$ — the tower acts COHERENTLY through every stage: T7's algebraic compatibility has become operator-level compatibility on one space.
  • \texttt{towerCyclicVec} (R8): the unit vector $\Omega$ implements EVERY corner Gibbs state as a vector state — $\langle\Omega, \pi_C(a)\Omega\rangle = \varphi_C(a)$ — and is CYCLIC (the span of its orbit is dense). $\Omega$ is NOT shown separating (cut: needs the commutant/modular theory of the limit).
  • \texttt{towerLimitVN} (R9, the capstone): the directed-union limit von Neumann algebra of the representation images, via the double-commutant campaign's \texttt{limitVN}, with membership characterized by SOT-approximation from the finite stages, and the ℕ-instantiation \texttt{freqTowerLimitVN} for the code's frequency tower. THE OBJECT THE TOWER AND CLOSURE CAMPAIGNS WERE BUILT FOR.
  • The checkpoint (R9), verbatim: HAVE — "One Hilbert space — the completion of the semidefinite Gibbs-GNS pre-space on the direct sum of all finite corners — carrying compatible unital ⋆-representations of every corner algebra (π\_\{C′\} ∘ cornerEmbed = π\_C for all C ⊆ C′), a unit cyclic vector Ω implementing every corner Gibbs state as a vector state (⟪Ω, π\_C(a)Ω⟫ = φ\_C(a)), and the directed-union limit von Neumann algebra towerLimitVN = limitVN of the representation images, with membership characterized by SOT-approximation from the finite stages — all axiom-free." HAVE NOT — "The type of towerLimitVN is not classified — no factor, no ITPFI identification, no III₁ claim is made or proved (the T3 fingerprint stays arithmetic; Araki–Woods 1968 and Connes 1973 stay cited, never invoked); Ω is not shown separating, the modular theory of the limit state on the completion is not constructed, and the representations are not shown isometric." The continuum is NOT done — but for the first time it has an inhabitant.

\emph{The transport and the accounting (B1--B8 · A1--A4): the limit algebra gets its dynamics; the species form is forced.} Two tracks, one campaign (\texttt{QIQTH/TowerGNS/Flow*.lean}, \texttt{QIQTH/Rigidity/RegulatorRigidity.lean}, \texttt{QIQTH/HeatKernelOneD.lean}, \texttt{QIQTH/SpeciesCrossCheck.lean}; all axiom-free std-3, budget 0).

  • THE MODULAR TRANSPORT (Track B): the per-corner Gibbs modular flows transported to \texttt{towerFlow} — a \textbf{strongly continuous one-parameter unitary group} $U_t$ on the tower GNS space (isometric extension through the completion; $U_0 = 1$, group law, $U_t^* = U_{-t}$, unitary membership; the ban on touching \texttt{cpow} held throughout — every fact routed through the DY3 rescale bridge $\sigma_s = \alpha_{-\beta s}$). $U_t\Omega = \Omega$; the finite-stage boundary KMS identity is displayed on the limit space (honestly bannered: NOT strip analyticity, NOT a KMS state of the limit algebra). THE IMPLEMENTATION THEOREM: $U_t\,\pi_C(a)\,U_{-t} = \pi_C(\sigma_t a)$ at every finite stage — the covariance is EXACT already at the pre-level, no germ argument needed — and \texttt{towerLimitVN} is INVARIANT under conjugation by the flow. Strong continuity (the shipped stretch) opens the door to the held Stone tower; the self-adjoint generator is NOT claimed. HAVE NOT, verbatim: no Tomita operator, $\Delta$, or $J$ on the completion is constructed; $\Omega$ is not shown separating; no type is classified; $U_t$ is defined by TRANSPORT of the finite corner flows, not derived from the limit state.
  • THE ACCOUNTING (Track A): the honest maximal species upgrade. \textbf{The regulator rigidity theorem} (\texttt{speciesRegulator\_forced}): any positive, species-additive, monotone family covariant under rescalings — with the covariance factor an UNKNOWN function, so the exponent is an OUTPUT — is forced to the Sakharov/Dvali form $1/G = N_{\rm eff}\Lambda^\kappa$, with $\kappa = 2$ pinned by a single dimensional calibration, an explicit dyadic counterexample showing weakened covariance breaks the conclusion, and the non-vacuity instance equal to the held \texttt{inducedInvG} by \texttt{rfl}. \textbf{The first derived (not cited) heat-kernel coefficient in the repository}: $(1/2\pi)\int e^{-tk^2}dk = 1/\sqrt{4\pi t}$, from Mathlib's Gaussian integral (\texttt{heatDensity\_oneD}), with the cutoff moment realizing the $\Lambda^2$ as a derived momentum integral. \textbf{The mixed-species consistency chain}: ONE shared species datum feeds both the one-loop entanglement entropy and the induced $1/G$; the mixed-field-content $1/4$ (\texttt{species\_sakharov\_ratio} — the entire species sum cancels) and $S_{\rm ent} = A/4G$ (\texttt{speciesEntropy\_eq\_capacity}) are theorems, chained through the BTZ Cardy count. HAVE NOT, verbatim: the numerical value of $G$ is not derived; the $c_i$ remain cited one-loop Seeley--DeWitt heat-kernel data, hand-entered — the chain certifies that one shared cited datum is wired coherently through both bookkeepings, not that either side is independently computed.

\emph{The generator (G1--G6): the self-adjoint Hamiltonian of the limit dynamics, computed.} The transported flow gets its generator: \texttt{towerGen := stoneGen (towerFlow)} — a genuine self-adjoint unbounded operator ($K = K^\dagger$ in Mathlib's \texttt{LinearPMap} adjoint sense), obtained by instantiating the general Stone-generator theorem (built during the P4-wall campaign) with the five held flow facts, with exactly two adapter lemmas (\texttt{QIQTH/TowerGNS/Generator.lean}; axiom-free std-3, budget 0; every increment green on first or second build).

  • \textbf{The zero-mode}: $\Omega \in \mathrm{dom}(K)$ with $K\Omega = 0$ — the cyclic vector is annihilated because the flow fixes it exactly.
  • \textbf{The explicit core} (\texttt{towerGen\_of} + \texttt{cornerGenMatrix\_eq\_commutator}): on every coerced pure component, $K\,\uparrow\!(\mathrm{of}_C\,a) = \uparrow\!(\mathrm{of}_C\,[H_C, a])$ with $H_C = \mathrm{diagonal}(\log \mathrm{gibbsWeight})$ — the finite-stage modular Hamiltonian acts by COMMUTATOR, with the phases $\kappa_{nm} = \log w_n - \log w_m$ of the held entry formula; the generator is COMPUTED, not just certified, and its domain is dense CONSTRUCTIVELY (the coerced pre-vectors — no Gårding mollification needed).
  • \textbf{Flow covariance}: $U_s$ preserves $\mathrm{dom}(K)$ and $K U_s = U_s K$.
  • HAVE NOT, verbatim: \texttt{towerGen} is NOT constructed from, and NOT claimed equal to, a Tomita modular Hamiltonian $\log \Delta$ of the limit state — no $\Delta$, $J$, $S$, separating property, KMS-at-the-limit, or von Neumann type is claimed; no spectral resolution (PVM) of the unbounded \texttt{towerGen}; no exponential-recovery identity $U_t = e^{itK}$ — the recovery wall is open by design.

\emph{The separation (S1--S8): $\Omega$ is cyclic AND separating — the standard form.} The limit algebra reaches the standard-form hypothesis pair of Tomita--Takesaki theory (\texttt{QIQTH/TowerGNS/RightMul.lean}, \texttt{Separation.lean}; axiom-free std-3, budget 0).

  • \textbf{The weight exchange and the half-power intertwining}: T7's modular-frequency lemma, exponentiated, gives $w_K(m)\,w_0(\hat n) = w_K(n)\,w_0(\hat m)$ for off-corner-agreeing occupations; hence the ENGINE $\iota(a)\cdot\sqrt{\rho_K} = \sqrt{\rho_K}\cdot\iota(\mathrm{rightConj}\,a)$, from which RIGHT multiplication by any corner element is BOUNDED on the tower GNS space with the weighted Frobenius constant $\sum_{n,m}\|a_{nm}\|^2\,(w_m/w_n)$ (the classical right-boundedness-via-the-modular-element argument in finite guise; never claimed contractive).
  • \textbf{The commutant facts}: the left and right actions commute (the stage mismatch resolved by a double germ at a fresh deep stage — no Finset equality is ever stated), and every element of \texttt{towerLimitVN} commutes with every right multiplication by PURE BICOMMUTANT ALGEBRA (membership in the double centralizer is definitional — the double-commutant campaign consumed with zero analysis).
  • \textbf{THE CAPSTONE} (\texttt{towerCyclicVec\_separating}): $T \in$ \texttt{towerLimitVN}, $T\Omega = 0 \Rightarrow T = 0$ — $T$ kills the right orbit of $\Omega$, which IS the left orbit ($R_a\Omega = \pi_C(a)\Omega$), and that orbit spans densely. With the held cyclicity: \textbf{$\Omega$ is CYCLIC AND SEPARATING for the tower limit von Neumann algebra — the standard-form hypothesis pair of Tomita--Takesaki theory, exhibited axiom-free.} Plus \texttt{towerLimitVN\_eq\_of\_apply\_cyclicVec} ($T\Omega = S\Omega \Rightarrow T = S$) — the well-definedness germ of a future Tomita operator $S_0$.
  • HAVE NOT, verbatim: no Tomita operator $S_0$, no modular operator $\Delta$, no conjugation $J$, no KMS condition at the limit, and no type classification is constructed or claimed here — separation is the HYPOTHESIS for that theory, not the theory.

\emph{The Tomita operator (T0\_1--T0\_6): $S_0$, computed and closable.} Modular theory proper begins: the Tomita operator of the tower limit state is constructed on its classical orbit domain $\{T\Omega : T \in \mathrm{towerLimitVN}\}$ (\texttt{QIQTH/TowerGNS/Tomita.lean}; axiom-free std-3, budget 0; the heaviest increment landed with zero proof failures).

  • \textbf{The operator}: \texttt{towerTomita₀} is a genuine conjugate-linear (σ-semilinear) partial operator — Mathlib's \texttt{LinearPMap} is semilinear at this pin, so the object typechecks with full algebraic API — well-defined BECAUSE $\Omega$ is separating (the previous campaign's capstone doing exactly its job), involutive, densely defined, with $S_0\Omega = \Omega$ and the COMPUTED core action $S_0\!\uparrow\!(\mathrm{of}_C\,a) = \uparrow\!(\mathrm{of}_C\,a^\dagger)$ — the Tomita involution is conjugate-transpose on pure components.
  • \textbf{The finite $\sigma_{-i}$, computed}: the commutant-side right multiplications carry the exact adjoint $R_a^\dagger = R_{\rho a^\dagger \rho^{-1}}$ — the engine squared ($\iota(a)\rho_K = \rho_K\,\iota(\mathrm{rightConj}^2 a)$, two half-power intertwinings composed) plus the \texttt{modAut} bridge: the adjoint parameter IS the finite-stage $\sigma_{-i}$ image of $a^\dagger$, computed rather than analytically continued — the honest finite-tower substitute for KMS analyticity.
  • \textbf{The classical pairing}: $\langle T^*\Omega, R_a\Omega\rangle = \langle R_a^\dagger\Omega, T\Omega\rangle$ on the dense pure-component family ($F_0 \subseteq S_0^\sharp$ in family form).
  • \textbf{Closability}: $T_n\Omega \to 0$ and $T_n^*\Omega \to v$ force $v = 0$ — the graph-limit form, proved from the commutant orbit's density.
  • HAVE NOT, verbatim: the closure $\bar S$ is not constructed as an object, and no polar decomposition, no modular operator $\Delta$, no conjugation $J$, no KMS condition of the limit state, and no type classification is constructed or claimed; Mathlib's \texttt{LinearPMap} closure and adjoint theories cover only $\mathbb{C}$-linear (identity ring-hom) partial maps, and a conjugate-linear closure theory is not built here — the named next wall.

\emph{The conjugate closure (CC1--CC7): $\bar S$ as an object — and a new slice of Mathlib.} The closure of the Tomita operator now exists (\texttt{QIQTH/TowerGNS/ConjClosure.lean} abstract + \texttt{TomitaBar.lean}; axiom-free std-3, budget 0), built by the $\mathbb{R}$-reduction: a conjugate-linear map IS $\mathbb{R}$-linear, and Mathlib's entire partial-operator closure theory applies verbatim at $\mathbb{R}$ through the global \texttt{complexToReal} instances — no local instances anywhere.

  • \textbf{Four new abstract theorems} (Mathlib-only imports — upstream-gap contributions): the $\mathbb{R}$-restriction view of a conjugate-linear partial map (no Mathlib helper existed); the SEQUENCE-CLOSABILITY BRIDGE (\texttt{isClosable\_of\_seq} — absent from Mathlib even for ordinary linear maps); and the two transfer theorems — conjugate-homogeneity and the involution SURVIVE closure (the twisted-map and swap-homeomorphism engines; no adjoint anywhere).
  • \textbf{$\bar S$} (\texttt{towerTomitaBar}): closed, extending $S_0$ with the orbit domain as a CORE, $\bar S\Omega = \Omega$, conjugate-transpose on pure components, conjugate-homogeneous (twist-guarded: $\bar S(c\Omega) = \bar c\,\Omega$), and FULLY INVOLUTIVE on its domain with trivial kernel and range equal to domain.
  • HAVE NOT, verbatim: $\Delta$, $J$, and the polar decomposition are not constructed (the documented $\Delta$ contract is the named next campaign); no $\sigma$-semilinear graph theory is contributed (the $\mathbb{R}$-reduction sidesteps it — the $\sigma$-graph remains Mathlib's own open TODO); no KMS condition of the limit state and no von Neumann type is claimed.

\emph{The modular operator (M1--M7): $\Delta$, computed as the modular automorphism.} The campaign after the conjugate closure delivers Tomita--Takesaki's central object for the tower limit state (\texttt{QIQTH/TowerGNS/ConjAdjoint.lean} + \texttt{ModularOp.lean}; axiom-free std-3, budget 0).

  • \textbf{Tomita's $F$ with no Riesz machinery}: the conjugate-linear adjoint of $\bar S$ is built on the $\exists$-Riesz domain $\{y : \exists w, \forall x, \langle \bar S x, y\rangle = \langle w, x\rangle\}$ — no real inner product, no dual-space machinery, no completeness argument anywhere; uniqueness of the witness is density alone. The abstract construction (\texttt{conjAdjoint}, closed in the sequence sense, with the core-extension equalizer lemma) is itself a Mathlib-gap contribution: Mathlib's adjoint theory covers only $\mathbb{C}$-linear partial maps.
  • \textbf{$\Delta := F \circ \bar S$}: a $\mathbb{C}$-linear densely defined partial operator — the two conjugations cancel — which is SYMMETRIC (\texttt{IsFormalAdjoint} $\Delta$ $\Delta$), POSITIVE ($\langle \Delta x, x\rangle = \lVert\bar S x\rVert^2 \ge 0$), CLOSABLE ($\Delta \le \Delta^\dagger$ with $\Delta^\dagger$ closed), and fixes $\Omega$.
  • \textbf{The computation}: on the dense pure-component core, $\Delta\,\uparrow\!(\mathrm{of}_C\,a) = \uparrow\!(\mathrm{of}_C\,(\mathrm{modAut}_{\rho_C}\,a))$ — the continuum modular operator ACTS AS the finite modular automorphism $\rho\,a\,\rho^{-1}$, stage by stage: the modular operator of the physics, computed rather than postulated.
  • HAVE NOT, verbatim: full self-adjointness $\Delta^\dagger = \Delta$ is not proved — it is von Neumann's $\bar S^*\bar S$ theorem, absent from Mathlib and named as the next target; no polar decomposition, no $J$, no $\Delta^{1/2}$ or $\Delta^{it}$ (no unbounded positive square root or spectral theory for partial operators exists in the pin), no KMS condition of the limit state, and no von Neumann type is constructed or claimed.

\emph{The von Neumann campaign (VN1--VN6): $\Delta^\dagger = \Delta$ — and von Neumann's theorem itself.} One campaign later, the modular operator is genuinely SELF-ADJOINT (\texttt{QIQTH/VonNeumann/} + \texttt{TowerGNS/ModularSurjective.lean}, \texttt{ModularSelfAdjoint.lean}; axiom-free std-3, budget 0; six increments in a single session, four consecutive first-try greens).

  • \textbf{Three abstract Mathlib-gap files} (RCLike-generic, Mathlib-only imports): the self-adjointness kernel (densely defined + symmetric + $\mathrm{ran}(1+A) = \top$ $\Rightarrow$ $A^\dagger = A$); the von Neumann graph orthogonal decomposition of a closed partial operator in $\ell^2(E \times E)$ (no adjoint anywhere in statement or proof); and \textbf{von Neumann's theorem itself} — for closed densely defined $T$ over any RCLike field, $T^\dagger T$ is densely defined and self-adjoint. None of these exists in Mathlib at the pin.
  • \textbf{The i-twist}: conjugate-homogeneity upgrades the real graph-orthogonality pairing to the full complex pairing — which lands VERBATIM in the $\exists$-Riesz $F$-domain membership: $\mathrm{ran}(1+\Delta) = \top$, with witness $h - x$.
  • \textbf{THE HEADLINE}: $\Delta^\dagger = \Delta$ — \texttt{towerModularOp\_isSelfAdjoint}. The tower modular operator is a genuinely self-adjoint (Mathlib \texttt{LinearPMap.adjoint}-sense), positive, closed, densely defined operator fixing $\Omega$ and computed as the finite modular automorphism on the dense core; with kernel triviality and the resolvent bound $\|x\| \le \|x + \Delta x\|$.
  • HAVE NOT, verbatim: $\Delta^{1/2}$ and the polar decomposition $\bar S = J\Delta^{1/2}$ are not constructed (so $J$ is still not an object); the spectral resolution of the unbounded $\Delta$ is not built; $\Delta^{it}$ and the KMS property of the limit state are not proved — in particular NO claim that the transported dynamics equals the modular flow of $\Delta$; no Tomita theorem at the algebra level; no von Neumann type. The certified next step is the resolvent bridge: $(1+\Delta)^{-1}$ as a bounded self-adjoint contraction feeding the project's existing bounded spectral theorem.

\emph{The resolvent campaign (R1--R7): $\Delta^{it}$ exists.} The modular unitary group of the tower limit state is constructed (\texttt{QIQTH/TowerGNS/Resolvent*.lean}, \texttt{ModularUnitary*.lean}, \texttt{QIQTH/Spectral/PVMEigen.lean}; axiom-free std-3, budget 0; seven increments, every one first-try green).

  • \textbf{The resolvent}: $R := (1+\Delta)^{-1}$, an everywhere-defined self-adjoint contraction with $0 \le R \le 1$, trivial kernel, dense range equal to $\mathrm{dom}\,\Delta$, spectrum in $[0,1]$, $R\Omega = \tfrac12\Omega$, and the exact identity $\Delta \circ R = 1 - R$ — so $\Delta$ is a function of a single bounded self-adjoint operator, pinned against the bounded calculus on its whole domain.
  • \textbf{The abstract PVM supplement} (reusable, new to the project's spectral tower and to Mathlib): the operator-level spectral theorem $T = \int \lambda\, dE$; the KERNEL-ATOM lemma — injective self-adjoint $T$ forces $E(\{0\}) = 0$ (not automatic from kernel triviality; a multiplicative symbol argument does it); and the eigenvector calculus $f(T)x = f(r)x$.
  • \textbf{THE HEADLINE}: $\Delta^{it} :=$ \texttt{towerModUnitary}, the bounded Borel calculus of $R$ under the symbol $((1-r)/r)^{it}$ — a STRONGLY CONTINUOUS one-parameter unitary group ($U_0 = 1$, $U_{s+t} = U_sU_t$, $U_t^\star = U_{-t}$) that fixes the cyclic vector ($U_t\Omega = \Omega$), commutes with $\Delta$ on its whole domain, and carries no spectral weight at the junk point ($E(\{0\}) = 0$) — the group genuinely represents $\delta^{it}$ on the spectrum.
  • HAVE NOT, verbatim: no claim that \texttt{towerModUnitary} equals the transported \texttt{towerFlow} (equivalently $\mathrm{towerGen} = \log\Delta$): two strongly continuous unitary groups now coexist on the tower space and their identification — the exponential-recovery wall — is the named next campaign, not crossed here; no KMS condition of the limit state; Tomita's theorem ($\Delta^{it} M \Delta^{-it} = M$, $JMJ = M'$) is not proved — $U_t$ is not shown to implement automorphisms of the limit algebra; $\Delta^{1/2}$, $J$, and the polar decomposition are still not constructed; no von Neumann type statement.

\emph{The identification campaign (ID1--ID5): the exponential-recovery wall, crossed — and Tomita's theorem, first half.} The campaign after the resolvent identifies the two unitary groups (\texttt{QIQTH/TowerGNS/ModularEigenbasis.lean} through \texttt{TomitaFirstHalf.lean}; axiom-free std-3, budget 0; five increments, every one first-try green).

  • \textbf{The eigenvector route}: the Gibbs density is DIAGONAL by construction, so matrix units are simultaneous eigenvectors of the finite modular automorphism ($\mathrm{modAut}\,\rho\,(e_{nm}) = (w_n/w_m)\,e_{nm}$), of $\Delta$, of the resolvent, and — through the eigenvector calculus of the previous campaign — of $\Delta^{it}$, with the SAME explicit scalar $e^{it(\log w_n - \log w_m)}$ that the transported physical flow carries. No spectral theorem, no Stone uniqueness, no $\log\Delta$ needed anywhere.
  • \textbf{THE IDENTIFICATION}: \texttt{towerModUnitary\_eq\_towerFlow} — $\Delta^{it} = $ \texttt{towerFlow}$_t$ AS OPERATORS (two bounded operators agreeing on the dense cyclic span are equal): the transported physical dynamics of the tower limit state IS the spectral modular flow of its modular operator. Corollary: \texttt{towerGen} IS the Stone generator of $\Delta^{it}$ — the group-level content of $\mathrm{towerGen} = \log\Delta$.
  • \textbf{TOMITA'S THEOREM, FIRST HALF}: $\Delta^{it}$ implements the modular automorphisms ($\Delta^{it}\pi_C(a)\Delta^{-it} = \pi_C(\sigma_t a)$) and preserves the limit algebra: $\Delta^{it}\,\mathrm{towerLimitVN}\,\Delta^{-it} = \mathrm{towerLimitVN}$ — the modular theory of the physics equals the modular theory of the state.
  • HAVE NOT, verbatim: $J$ and the polar decomposition $\bar S = J\Delta^{1/2}$ are still not constructed (so $JMJ = M'$, the second half of Tomita's theorem, stays open — the natural next campaign: $J$ on matrix units is also explicit, so the same eigenbasis method applies); no analytic strip-KMS of the limit state (only the boundary identity); no von Neumann type classification; and everything remains the finite-stage Gibbs inductive-limit state — the free-field/Type-III continuum objects are untouched.

\emph{The modular conjugation, non-traciality, KMS-boundary, and field-level BW (J1--J9, N1--N4, K/C1, F1): the tower's Tomita--Takesaki theory COMPLETE, and the free-field wedge BW made unconditional.} Four short campaigns after the identification close the modular tower and cap the free-field wedge (\texttt{QIQTH/TowerGNS/JStage.lean}--\texttt{TomitaSecondHalf.lean}, \texttt{QIQTH/NonTracial/*}, \texttt{QIQTH/Fock/FieldBWUnconditional.lean}; axiom-free std-3, budget 0; nearly every increment first-try green).

  • \textbf{The modular conjugation $J$} (\texttt{towerJ}): the completion of the explicit conjugate-linear stage map $\mathrm{jStage}\,a = \sqrt{\rho}\,a^\dagger\sqrt{\rho}^{-1}$, semilinearly extended with no $\mathbb{R}$-reduction (Mathlib's \texttt{ContinuousLinearMap.completion} is $\sigma$-generic). Full anti-unitary pack: $\langle J\xi,J\eta\rangle = \langle\eta,\xi\rangle$, $J^2 = 1$, $J\Omega = \Omega$, the eigenbasis action $J\!\uparrow\!(\mathrm{of}_C\,e_{nm}) = \sqrt{w_m/w_n}\,\uparrow\!(\mathrm{of}_C\,e_{mn}\bar c)$. The POLAR DECOMPOSITION on the dense core: $\bar S = J\circ\Delta^{1/2}$ and $F = \Delta^{1/2}\circ J$ (the core $\Delta^{1/2}$-action squares to the modular automorphism), and $J\Delta^{it} = \Delta^{it}J$.
  • \textbf{TOMITA'S THEOREM, SECOND HALF (inclusion)}: $J$ conjugates every represented left multiplication into the commutant-side right multiplication ($J\pi_C(a)J = R_{\mathrm{jStage}\,a}$), hence $J\cdot\mathrm{towerLimitVN}\cdot J \subseteq \mathrm{towerLimitVN}'$, with $\Omega$ cyclic and separating for the commutant.
  • \textbf{NON-TRACIALITY}: the tower vacuum state $\omega(\cdot) = \langle\Omega,\pi(\cdot)\Omega\rangle = \mathrm{tr}(\rho_C\,\cdot)$ is NOT a trace — $\omega(\pi(e_{nm})\pi(e_{mn})) = w_n \ne w_m = \omega(\pi(e_{mn})\pi(e_{nm}))$ when the Gibbs weights differ — and the modular data is nontrivial ($\Delta \ne 1$, $\Delta^{it} \ne \mathrm{id}$): the Powers "not-the-tracial-case" separation. \textbf{KMS-boundary}: $\omega(\pi(x)\pi(y)) = \omega(\pi(y)\pi(\sigma(x)))$ against the finite modular automorphism, bundled with the whole tower into the capstone \texttt{modular\textbackslash\{\}\_data\textbackslash\{\}\_complete\textbackslash\{\}\_witness} — the first COMPLETE Tomita--Takesaki modular theory in any proof assistant.
  • \textbf{FIELD-LEVEL BISOGNANO--WICHMANN, UNCONDITIONAL} (\texttt{freeField\_secondQuant\_BW\_unconditional}): the second-quantized wedge modular automorphism acts on Weyl operators as the geometric Lorentz boost, $\sigma_t(W(u)) = W(\mathrm{boost}(2\pi t)u)$ conjugated, with NO carried BW hypothesis — discharged from the already-unconditional one-particle result \texttt{oneParticleBW\_niceWedge\_unconditional}. The free-field wedge modular structure (one-particle BW, strip-KMS, Reeh--Schlieder cyclic $+$ separating, field-level BW) is now machine-checked end-to-end: modular flow $=$ boost, a Lean-first result.
  • HAVE NOT, verbatim: the REVERSE inclusion $\mathrm{towerLimitVN}' \subseteq J\cdot\mathrm{towerLimitVN}\cdot J$ (full equality $JMJ = M'$, Tomita's hard half) is NOT proved — the Rieffel--van Daele commutant wall; no unbounded $\Delta^{1/2}$ (the polar decomposition is a core-level identity); no strip-analyticity KMS; no von Neumann type classification (Mathlib has no trace/factor/type API); the field-level BW is free-field / single-mass, NOT interacting, and makes no low-energy Lorentz-violation prediction (unconditional BW $=$ standard induced gravity); everything remains the finite-stage Gibbs inductive-limit / free-field sector — the interacting and Type-III continuum frontiers are untouched.

\subsection{Credit division}

The mathematical infrastructure, Type III$_1$ classification of local algebras, the crossed-product construction giving Type II algebras with renormalized entropy matching generalized entropy in semiclassical gravity, is due to Chandrasekaran, Penington, Witten, and collaborators (CPW 2022; Witten 2022; CLPW 2022; Jensen-Sorce-Speranza 2023; building on Connes-Takesaki and Haag-Kastler).

The original contributions of this paper are:

(a) The \textbf{finite regional entropy reading} of the Bekenstein-Bousso bound together with the \textbf{$(\Phi,\lambda)$ architecture}: $Q_R$ bounds the von Neumann / renormalized entropy $S(\rho_R)$ of the pre-selection state $\Phi$ (its effective modular rank $e^{S(\rho_R)}$), and the non-dynamical selector $\lambda$ fixes the single actual record per run. This entropy bound goes beyond what Witten's algebraic theorem directly establishes; the irreducible new physics is (P4) [the entropy bound] and (P5) on the (P1) $(\Phi,\lambda)$ ontology.

(b) The \textbf{(FQ) axiom} stated rigorously in the CPW Type II framework: (i) regional content given by the algebra-state, (ii) the renormalized entropy bound $S(\rho_R) \le Q_R$ postulated as axiom. (An earlier third part posited a physical-instantiation precision-floor; that amplitude-precision reading is superseded, §1.1a.)

(c) \textbf{Theorem 3 (finite regional entropy)}: $S(\rho_R) \le Q_R$. (The earlier subsidiary $\epsilon(R) > 0$ amplitude-resolution statement is superseded, §1.1a.)

(d) \textbf{The mechanism (§6, Theorem 4):} decoherence renders the records non-interfering, the capacity bound (FQ) bounds their entropy/cardinality, the microscopic IC index the run, and the non-dynamical selector $\lambda$ (P1) makes exactly one record actual, producing single-record per-run outcomes without collapse. (IC only indexes; $\lambda$ makes actual.)

(e) The \textbf{formal/per-run distinction} within the algebraic framework.

(f) The \textbf{schematic Born-typicality program} identified as the central quantitative open problem.

The credit division is sharp: CPW supplies the mathematical infrastructure (Type II regional algebras); we supply the entropy reading of the bound plus the $(\Phi,\lambda)$ ontology on top, and develop its structural consequences for the measurement problem.

\subsection{Closing remarks}

The framework retains standard quantum mechanics exactly at the level of the underlying field algebra. It rests on five postulates; the irreducible new physics is (P4) [the finite regional entropy bound $S(\rho_R) \le Q_R$] and (P5) on the (P1) $(\Phi,\lambda)$ ontology, together with one distinction (formal vs per-run wave function). From these the structure follows: decoherence removes interference between macroscopic records; the capacity bound (FQ) bounds their entropy/cardinality (it does \emph{not} make a $\ge 2$-record content non-instantiable, §1.1a); and the non-dynamical selector $\lambda$ (P1) makes exactly one record actual, hence single-record per-run content, the macroscopic world emerging without any dynamical modification.

The framework is empirically conservative, Schrödinger and Born preserved exactly, no laboratory deviations predicted. Its ontology is the pair $(\Phi,\lambda)$: $\Phi$ the pre-selection wave function as a physical state of spacetime, and $\lambda$ the non-dynamical actuality primitive that fixes which single record is realized. It is mathematically grounded in the CPW Type II algebra framework. The principal open problems are the rigorous Born typicality theorem and the dynamical, Lorentz-covariant realization of the record structure $\lambda$ selects from.

The wave function is one. The realized world per run is one. The Bekenstein-Bousso holographic bound, read as a Lorentz-invariant bound on the regional entropy $S(\rho_R)$ of $\Phi$, sets the arena; decoherence removes interference and bounds the distinguishable-record entropy; and the selector $\lambda$ makes one record actual. The Schrödinger equation is unchanged. The Born rule is unchanged. The macroscopic world emerges as a structural consequence, single outcomes per run, no collapse postulate, no hidden particles, no branching, no modal value-rule.

\medskip\hrule\medskip

\section{Acknowledgements}

The author thanks the participants in extended discussions that informed the framework developed here. The mathematical apparatus draws on the work of Chandrasekaran, Penington, Witten, and others in the Type II crossed-product algebra construction; we are indebted to that line of work for providing the rigorous home for the regional information bound.

\medskip\hrule\medskip

\section{References}

\begin{enumerate}
  \item Ballentine, L. E. (1970). The statistical interpretation of quantum mechanics. \emph{Rev. Mod. Phys.}, 42, 358.
  \item Banks, T. (2025). \emph{Finite Entropy Implies Finite Dimension in Quantum Gravity.} arXiv:2509.17856.
  \item Bassi, A., Dorato, M., \& Ulbricht, H. (2023). \emph{Collapse Models: A Theoretical, Experimental and Philosophical Review.} Entropy 25, 645. arXiv:2310.14969. (See also Tomaz, A. A., Mattos, R. S., \& Barbatti, M. (2025). \emph{The Quantum Measurement Problem: A Review of Recent Trends.} Phil. Mag. C. arXiv:2502.19278, for a broader recent survey.)
  \item Bekenstein, J. D. (1981). Universal upper bound on the entropy-to-energy ratio for bounded systems. \emph{Phys. Rev. D}, 23, 287.
  \item Bohm, D. (1952). A suggested interpretation of the quantum theory in terms of "hidden" variables. \emph{Phys. Rev.}, 85, 166.
  \item Bousso, R. (2002). The holographic principle. \emph{Rev. Mod. Phys.}, 74, 825. arXiv:hep-th/0203101.
  \item Carroll, S. M., \& Sebens, C. (2014). Many worlds, the Born rule, and self-locating uncertainty. In \emph{Quantum Theory: A Two-Time Success Story: Yakir Aharonov Festschrift}, Springer, pp. 157–169. arXiv:1405.7907.
  \item \textbf{Chandrasekaran, V., Longo, R., Penington, G., \& Witten, E. (2022).} \emph{An algebra of observables for de Sitter space.} JHEP 02 (2023) 082. arXiv:2206.10780.
  \item \textbf{Chandrasekaran, V., Penington, G., \& Witten, E. (2022).} \emph{Large N algebras and generalized entropy.} JHEP 04 (2023) 009. arXiv:2209.10454.
  \item Dieks, D. (1989). Quantum mechanics without the projection postulate. \emph{Foundations of Physics}, 19, 1397.
  \item Diósi, L. (1989). Models for universal reduction of macroscopic quantum fluctuations. \emph{Phys. Rev. A}, 40, 1165.
  \item Dürr, D., Goldstein, S., \& Zanghì, N. (1992). Quantum equilibrium and the origin of absolute uncertainty. \emph{J. Stat. Phys.}, 67, 843–907. arXiv:quant-ph/0308039.
  \item \textbf{Dorau, P., \& Much, A. (2026).} \emph{From Quantum Relative Entropy to the Semiclassical Einstein Equations.} Phys. Rev. Lett. (2026), DOI: 10.1103/lmq8-nsty. arXiv:2510.24491 [hep-th].
  \item Engelhardt, N., \& Wall, A. C. (2015). \emph{Quantum extremal surfaces: holographic entanglement entropy beyond the classical regime.} JHEP 01 (2015) 073. arXiv:1408.3203.
  \item Everett, H. (1957). "Relative state" formulation of quantum mechanics. \emph{Rev. Mod. Phys.}, 29, 454.
  \item Gell-Mann, M., \& Hartle, J. B. (1993). Classical equations for quantum systems. \emph{Phys. Rev. D}, 47, 3345. arXiv:gr-qc/9210010.
  \item Ghirardi, G. C., Rimini, A., \& Weber, T. (1986). Unified dynamics for microscopic and macroscopic systems. \emph{Phys. Rev. D}, 34, 470.
  \item Griffiths, R. B. (2002). \emph{Consistent Quantum Theory.} Cambridge University Press.
  \item \textbf{Haag, R. (1992).} \emph{Local Quantum Physics: Fields, Particles, Algebras.} Springer.
  \item \textbf{Haag, R., \& Kastler, D. (1964).} An algebraic approach to quantum field theory. \emph{J. Math. Phys.}, 5, 848.
  \item Healey, R. (2017). \emph{The Quantum Revolution in Philosophy.} Oxford University Press.
  \item 't Hooft, G. (1993). \emph{Dimensional reduction in quantum gravity.} arXiv:gr-qc/9310026.
  \item Jacobson, T. (1995). Thermodynamics of spacetime: The Einstein equation of state. \emph{Phys. Rev. Lett.}, 75, 1260–1263. arXiv:gr-qc/9504004.
  \item \textbf{Jensen, K., Sorce, J., \& Speranza, A. J. (2023).} \emph{Generalized entropy for general subregions in quantum gravity.} arXiv:2306.01837.
  \item Joos, E., Zeh, H. D., Kiefer, C., Giulini, D., Kupsch, J., \& Stamatescu, I.-O. (2003). \emph{Decoherence and the Appearance of a Classical World in Quantum Theory.} Springer.
  \item \textbf{Murray, F. J., \& von Neumann, J. (1936).} On rings of operators. \emph{Annals of Mathematics}, 37, 116.
  \item Nielsen, M. A., \& Chuang, I. L. (2010). \emph{Quantum Computation and Quantum Information.} Cambridge.
  \item Omnès, R. (1994). \emph{The Interpretation of Quantum Mechanics.} Princeton University Press.
  \item Palmer, T. N. (2025). \emph{Rational Quantum Mechanics: Testing Quantum Theory with Quantum Computers.} Proc. Natl. Acad. Sci. USA (PNAS). arXiv:2510.02877. (See also Palmer, T. N. (2009). \emph{The Invariant Set Postulate: A New Geometric Framework for the Foundations of Quantum Theory and the Role Played by Gravity.} Proc. Roy. Soc. A 465, 3165–3185. DOI: 10.1098/rspa.2009.0080. arXiv:0812.1148; and subsequent IST development.)
  \item Pearle, P. (1989). Combining stochastic dynamical state-vector reduction with spontaneous localization. \emph{Phys. Rev. A}, 39, 2277.
  \item Penrose, R. (1996). On gravity's role in quantum state reduction. \emph{Gen. Rel. Grav.}, 28, 581.
  \item Susskind, L. (1995). The world as a hologram. \emph{J. Math. Phys.}, 36, 6377. arXiv:hep-th/9409089.
  \item von Neumann, J. (1932). \emph{Mathematische Grundlagen der Quantenmechanik.} Springer.
  \item Wald, R. M. (1993). Black hole entropy is the Noether charge. \emph{Phys. Rev. D}, 48, R3427–R3431. arXiv:gr-qc/9307038.
  \item Wall, A. C. (2012). A proof of the generalized second law for rapidly changing fields and arbitrary horizon slices. \emph{Phys. Rev. D}, 85, 104049. arXiv:1105.3445. (Erratum: \emph{Phys. Rev. D} 87, 069904, 2013.)
  \item Wallace, D. (2012). \emph{The Emergent Multiverse: Quantum Theory according to the Everett Interpretation.} Oxford University Press.
  \item \textbf{Witten, E. (2022).} \emph{Gravity and the crossed product.} JHEP 10 (2022) 008. arXiv:2112.12828.
  \item Zurek, W. H. (2003). Decoherence, einselection, and the quantum origins of the classical. \emph{Rev. Mod. Phys.}, 75, 715–775. arXiv:quant-ph/0105127.
  \item Kapłański, P. (2026). \emph{One Wave Function, One World: How the Holographic Information Bound Selects the Macroscopic World.} (Position paper, companion to this one.)
\end{enumerate}

\medskip\hrule\medskip

\emph{End of manuscript.}

\end{document}
